% Generated mechanically from the frozen sites-modules.tex target.
% Do not edit this generated file; edit the bound locale source instead.

% OK, start here.
%

\setcounter{chapter}{17}
\chapter{站点上的模}



\phantomsection
\label{sites-modules-section-phantom}


\section{引言}
\label{sites-modules-section-introduction}

\noindent
本章阐述环拓扑斯或环站点上的模层之基本概念。我们先给出阿贝尔层的若干
基本事实，随后引入环站点与环拓扑斯。我们讨论 $\mathcal{O}$-模（预）层的
一些最基本概念；这些内容类似于空间上的层一章中关于 $\mathcal{O}$-模
（预）层的材料。完成这些准备后，我们重述《模层》中相当一部分讨论
（见《模层》第 \ref{modules-section-introduction} 节）。基本参考文献为
\cite{FAC}、\cite{EGA} 和 \cite{SGA4}。






\section{阿贝尔预层}
\label{sites-modules-section-abelian-pre-sheaves}

\noindent
设 $\mathcal{C}$ 为范畴。阿贝尔预层已在《站点》第
\ref{sites-section-presheaves}、\ref{sites-section-sheaves} 节中引入，
并在《站点》第 \ref{sites-section-sheaves-algebraic-structures} 节中作了
进一步讨论。我们沿用后一处的约定：将阿贝尔预层看作集合预层，其每个截面
集合上都带有与限制映射相容的加法法则。回顾 $\mathcal{C}$ 上阿贝尔预层的
范畴记为 $\textit{PAb}(\mathcal{C})$。

\medskip\noindent
依《同调》定义 \ref{homology-definition-abelian-category}，范畴
$\textit{PAb}(\mathcal{C})$ 是阿贝尔范畴。给定预层映射
$\varphi : \mathcal{G}_1 \to \mathcal{G}_2$，$\varphi$ 的核是阿贝尔预层
$U \mapsto \Ker(\mathcal{G}_1(U) \to \mathcal{G}_2(U))$，
而 $\varphi$ 的余核是预层
$U \mapsto \Coker(\mathcal{G}_1(U) \to \mathcal{G}_2(U))$。
由于阿贝尔群范畴是阿贝尔范畴，并且这一点对每个 $U$ 都成立，故
$\Coim = \Im$。阿贝尔预层列
$$
\mathcal{G}_1 \longrightarrow
\mathcal{G}_2 \longrightarrow
\mathcal{G}_3
$$
正合，当且仅当对所有 $U \in \Ob(\mathcal{C})$，阿贝尔群列
$\mathcal{G}_1(U) \to \mathcal{G}_2(U) \to \mathcal{G}_3(U)$
正合。验证留给读者。

\begin{lemma}
\label{sites-modules-lemma-limits-colimits-abelian-presheaves}
设 $\mathcal{C}$ 为范畴。
\begin{enumerate}
\item $\textit{PAb}(\mathcal{C})$ 中一切极限与余极限均存在。
\item 一切极限与余极限均与在 $\mathcal{C}$ 的对象上取截面可交换。
\end{enumerate}
\end{lemma}

\begin{proof}
设 $\mathcal{I} \to \textit{PAb}(\mathcal{C})$，$i \mapsto \mathcal{F}_i$
为图。我们可直接按下列法则定义阿贝尔预层 $L$ 与 $C$：
$$
L : U \longmapsto \lim_i \mathcal{F}_i(U)
$$
以及
$$
C : U \longmapsto \colim_i \mathcal{F}_i(U).
$$
在每个 $U$ 上的截面群中使用相应映射，显然得到阿贝尔预层映射
$L \to \mathcal{F}_i$ 与 $\mathcal{F}_i \to C$。容易验证，配备这些映射的
$L$ 与 $C$ 分别是该图在 $\textit{PAb}(\mathcal{C})$ 中的极限与余极限。
这证明了 (1) 与 (2)。细节从略。
\end{proof}


\section{阿贝尔层}
\label{sites-modules-section-abelian-sheaves}

\noindent
设 $\mathcal{C}$ 为站点。$\mathcal{C}$ 上阿贝尔层的范畴记为
$\textit{Ab}(\mathcal{C})$。它是 $\textit{PAb}(\mathcal{C})$ 的全子范畴，
由底层集合预层为层的那些阿贝尔预层组成。《站点》第
\ref{sites-section-sheaves-algebraic-structures} 节的性质
($\alpha$)--($\zeta$) 成立；见《站点》命题
\ref{sites-proposition-functoriality-algebraic-structures-topoi}。
特别地，包含函子
$\textit{Ab}(\mathcal{C}) \to \textit{PAb}(\mathcal{C})$
具有左伴随，即层化函子 $\mathcal{G} \mapsto \mathcal{G}^\#$。

\medskip\noindent
我们建议读者在草稿纸上自行证明下述引理，而不是阅读其证明。

\begin{lemma}
\label{sites-modules-lemma-abelian-abelian}
设 $\mathcal{C}$ 为站点，$\varphi : \mathcal{F} \to \mathcal{G}$ 为
$\mathcal{C}$ 上阿贝尔层的态射。
\begin{enumerate}
\item 范畴 $\textit{Ab}(\mathcal{C})$ 是阿贝尔范畴。
\item $\varphi$ 的核 $\Ker(\varphi)$ 与将 $\varphi$ 视为预层态射时的核相同。
\item 态射 $\varphi$ 是单射（《同调》定义
\ref{homology-definition-injective-surjective}），当且仅当它作为预层映射是单射
（《站点》定义 \ref{sites-definition-presheaves-injective-surjective}），又当且仅当
它作为层映射是单射（《站点》定义
\ref{sites-definition-sheaves-injective-surjective}）。
\item $\varphi$ 的余核 $\Coker(\varphi)$ 是将 $\varphi$ 视为预层态射时所得
余核的层化。
\item 态射 $\varphi$ 是满射（《同调》定义
\ref{homology-definition-injective-surjective}），当且仅当它作为层映射是满射
（《站点》定义 \ref{sites-definition-sheaves-injective-surjective}）。
\item 阿贝尔层复形
$$
\mathcal{F} \to \mathcal{G} \to \mathcal{H}
$$
在 $\mathcal{G}$ 处正合，当且仅当对每个 $U \in \Ob(\mathcal{C})$ 及每个
在 $\mathcal{H}(U)$ 中映为零的 $s \in \mathcal{G}(U)$，都存在
$\mathcal{C}$ 中的覆盖 $\{U_i \to U\}_{i \in I}$，使每个 $s|_{U_i}$
都属于 $\mathcal{F}(U_i) \to \mathcal{G}(U_i)$ 的像。
\end{enumerate}
\end{lemma}

\begin{proof}
我们断言，《同调》引理 \ref{homology-lemma-adjoint-get-abelian} 可用于范畴
$\mathcal{A} = \textit{Ab}(\mathcal{C})$、
$\mathcal{B} = \textit{PAb}(\mathcal{C})$ 以及函子
$a : \mathcal{A} \to \mathcal{B}$（包含）与
$b : \mathcal{B} \to \mathcal{A}$（层化）。下面核验《同调》引理
\ref{homology-lemma-adjoint-get-abelian} 的假设。
假设 (1) 要求 $\mathcal{A}$、$\mathcal{B}$ 为加性范畴，$a$、$b$ 为加性
函子，且 $a$ 是 $b$ 的右伴随。前两点显然，而伴随性就是《站点》第
\ref{sites-section-sheaves-algebraic-structures} 节的 ($\epsilon$)。假设 (2)
要求 $\textit{PAb}(\mathcal{C})$ 为阿贝尔范畴——这已在第
\ref{sites-modules-section-abelian-pre-sheaves} 节中看到——并要求层化为左正合函子；后者
即《站点》第 \ref{sites-section-sheaves-algebraic-structures} 节的 ($\zeta$)。
最后一个假设是 $ba \cong \text{id}_\mathcal{A}$，即《站点》第
\ref{sites-section-sheaves-algebraic-structures} 节的 ($\delta$)。
故《同调》引理 \ref{homology-lemma-adjoint-get-abelian} 适用，从而
$\textit{Ab}(\mathcal{C})$ 是阿贝尔范畴。

\medskip\noindent
《同调》引理 \ref{homology-lemma-adjoint-get-abelian} 的证明表明，
$\Ker(\varphi)$ 与 $\Coker(\varphi)$ 分别等于将 $\varphi$ 视为阿贝尔预层
态射时所得核与余核的层化。这证明了 (4)。由于核是一个等化子（因而是极限），
且层化与有限极限可交换，遂得 (2)。

\medskip\noindent
断言 (2) 蕴含 (3)。由我们对层化的描述，断言 (4) 蕴含 (5)。断言 (6) 中
对正合性的刻画来自 (2)、(5)，以及该列正合当且仅当
$\Im(\mathcal{F} \to \mathcal{G}) =
\Ker(\mathcal{G} \to \mathcal{H})$.
\end{proof}

\noindent
该引理第 (6) 项还可表述为：阿贝尔层列
$$
\mathcal{F}_1 \longrightarrow
\mathcal{F}_2 \longrightarrow
\mathcal{F}_3
$$
正合，当且仅当预层
$U \mapsto \Im(\mathcal{F}_1(U) \to \mathcal{F}_2(U))$
的层化等于 $\mathcal{F}_2 \to \mathcal{F}_3$ 的核。

\begin{lemma}
\label{sites-modules-lemma-limits-colimits-abelian-sheaves}
设 $\mathcal{C}$ 为站点。
\begin{enumerate}
\item $\textit{Ab}(\mathcal{C})$ 中一切极限与余极限均存在。
\item 极限与 $\mathcal{C}$ 上阿贝尔预层的相应极限相同（即与在
$\mathcal{C}$ 的对象上取截面可交换）。
\item 有限直和与 $\mathcal{C}$ 上阿贝尔预层范畴中的相应有限直和相同。
\item 余极限是阿贝尔预层范畴中相应余极限的层化。
\item 滤过余极限正合。
\end{enumerate}
\end{lemma}

\begin{proof}
由引理 \ref{sites-modules-lemma-limits-colimits-abelian-presheaves}，阿贝尔预层的极限与
余极限存在，并可描述为在对象上的截面层面取极限与余极限。

\medskip\noindent
设 $\mathcal{I} \to \textit{Ab}(\mathcal{C})$，$i \mapsto \mathcal{F}_i$
为图，并令 $\lim_i \mathcal{F}_i$ 为该图在阿贝尔预层范畴中的极限。由
《站点》引理 \ref{sites-lemma-limit-sheaf}，它是阿贝尔层。此时容易看出，
$\lim_i \mathcal{F}_i$ 也是该图在 $\textit{Ab}(\mathcal{C})$ 中的极限。
这证明了极限存在，且 (2) 成立。

\medskip\noindent
由《范畴论》引理 \ref{categories-lemma-adjoint-exact}，并利用层化是包含函子的
左伴随，可知 $\colim_i \mathcal{F}$ 存在，且它是
$\textit{PAb}(\mathcal{C})$ 中余极限的层化。这证明了余极限存在，且 (4) 成立。

\medskip\noindent
在任何阿贝尔范畴中，有限直和与有限直积相同。因此 (3) 来自 (2)。

\medskip\noindent
证明 (5)。该断言意指：给定由有向集 $I$ 参数化的阿贝尔层正合列系统
$0 \to \mathcal{F}_i \to \mathcal{G}_i \to \mathcal{H}_i \to 0$，列
$0 \to \colim \mathcal{F}_i \to \colim \mathcal{G}_i \to
\colim \mathcal{H}_i \to 0$ 也正合。利用《同调》引理
\ref{homology-lemma-check-exactness} 与余极限定义作形式论证，可知列
$\colim \mathcal{F}_i \to \colim \mathcal{G}_i \to \colim \mathcal{H}_i \to 0$
正合。注意，$\colim \mathcal{F}_i \to \colim \mathcal{G}_i$ 是预层余极限
映射的层化；由于每个 $\mathcal{F}_i \to \mathcal{G}_i$ 都是单射，后一映射
也是单射。层化正合，故得结论。
\end{proof}







\section{自由阿贝尔预层}
\label{sites-modules-section-free-abelian-presheaf}

\noindent
为下述定义准备符号，我们约定将集合 $S$ 上的自由阿贝尔群记为\footnote{在其他
章节中，符号 $\mathbf{Z}[S]$ 有时表示以 $S$ 为变量集的 $\mathbf{Z}$ 上多项式环。}
$\mathbf{Z}[S] = \bigoplus_{s \in S} \mathbf{Z}$。其刻画性质为
$$
\Mor_{\textit{Ab}}(\mathbf{Z}[S], A)
=
\Mor_{\textit{Sets}}(S, A)
$$
换言之，构造 $S \mapsto \mathbf{Z}[S]$ 是遗忘函子
$\textit{Ab} \to \textit{Sets}$ 的左伴随。

\begin{definition}
\label{sites-modules-definition-free-abelian-presheaf-on}
设 $\mathcal{C}$ 为范畴，$\mathcal{G}$ 为集合预层。$\mathcal{G}$ 上的
{\it 自由阿贝尔预层} $\mathbf{Z}_\mathcal{G}$ 是由下列法则定义的阿贝尔预层：
$$
U \longmapsto \mathbf{Z}[\mathcal{G}(U)].
$$
在 $\mathcal{G} = h_X$ 是与 $\mathcal{C}$ 的对象 $X$ 相伴的可表预层这一
特殊情形，我们采用符号 $\mathbf{Z}_X = \mathbf{Z}_{h_X}$。换言之，
$$
\mathbf{Z}_X(U) = \mathbf{Z}[\Mor_\mathcal{C}(U, X)].
$$
\end{definition}

\noindent
此构造关于预层 $\mathcal{G}$ 显然具有函子性。事实上，它是遗忘函子
$\textit{PAb}(\mathcal{C}) \to \textit{PSh}(\mathcal{C})$ 的左伴随。
准确陈述如下。

\begin{lemma}
\label{sites-modules-lemma-obvious-adjointness}
设 $\mathcal{C}$ 为范畴，$\mathcal{G}$、$\mathcal{F}$ 为集合预层，
$\mathcal{A}$ 为阿贝尔预层，$U$ 为 $\mathcal{C}$ 的对象。则
\begin{align*}
\Mor_{\textit{PSh}(\mathcal{C})}(h_U, \mathcal{F})
& =
\mathcal{F}(U), \\
\Mor_{\textit{PAb}(\mathcal{C})}(\mathbf{Z}_\mathcal{G}, \mathcal{A})
& =
\Mor_{\textit{PSh}(\mathcal{C})}(\mathcal{G}, \mathcal{A}), \\
\Mor_{\textit{PAb}(\mathcal{C})}(\mathbf{Z}_U, \mathcal{A})
& =
\mathcal{A}(U).
\end{align*}
所有这些等式均具有函子性。
\end{lemma}

\begin{proof}
从略。
\end{proof}

\begin{lemma}
\label{sites-modules-lemma-coproduct-sum-free-abelian-presheaf}
设 $\mathcal{C}$ 为范畴，$I$ 为集合。对每个 $i \in I$，设
$\mathcal{G}_i$ 为集合预层。则
$$
\mathbf{Z}_{\coprod_i \mathcal{G}_i}
=
\bigoplus\nolimits_{i \in I} \mathbf{Z}_{\mathcal{G}_i}
$$
在 $\textit{PAb}(\mathcal{C})$ 中成立。
\end{lemma}

\begin{proof}
从略。
\end{proof}



\section{自由阿贝尔层}
\label{sites-modules-section-free-abelian-sheaf}

\noindent
下面给出集合层上的自由阿贝尔层之概念。

\begin{definition}
\label{sites-modules-definition-free-abelian-sheaf-on}
设 $\mathcal{C}$ 为站点，$\mathcal{G}$ 为集合预层。$\mathcal{G}$ 上的
{\it 自由阿贝尔层} $\mathbf{Z}_\mathcal{G}^\#$ 是阿贝尔层
$\mathbf{Z}_\mathcal{G}^\#$，即 $\mathcal{G}$ 上自由阿贝尔预层的层化。
在 $\mathcal{G} = h_X$ 是与 $\mathcal{C}$ 的对象 $X$ 相伴的可表预层这一
特殊情形，我们采用符号 $\mathbf{Z}_X^\#$。
\end{definition}

\noindent
此构造关于预层 $\mathcal{G}$ 显然具有函子性。事实上，它给出遗忘函子
$\textit{Ab}(\mathcal{C}) \to \Sh(\mathcal{C})$ 的左伴随。准确陈述如下。

\begin{lemma}
\label{sites-modules-lemma-obvious-adjointness-sheaves}
设 $\mathcal{C}$ 为站点，$\mathcal{G}$、$\mathcal{F}$ 为集合层，
$\mathcal{A}$ 为阿贝尔层，$U$ 为 $\mathcal{C}$ 的对象。则
\begin{align*}
\Mor_{\Sh(\mathcal{C})}(h_U^\#, \mathcal{F})
& =
\mathcal{F}(U), \\
\Mor_{\textit{Ab}(\mathcal{C})}(\mathbf{Z}_\mathcal{G}^\#,
\mathcal{A})
& =
\Mor_{\Sh(\mathcal{C})}(\mathcal{G}, \mathcal{A}), \\
\Mor_{\textit{Ab}(\mathcal{C})}(\mathbf{Z}_U^\#, \mathcal{A})
& =
\mathcal{A}(U).
\end{align*}
所有这些等式均具有函子性。
\end{lemma}

\begin{proof}
从略。
\end{proof}

\begin{lemma}
\label{sites-modules-lemma-may-sheafify-before-abelianize}
设 $\mathcal{C}$ 为站点，$\mathcal{G}$ 为集合预层。则
$\mathbf{Z}_\mathcal{G}^\# = (\mathbf{Z}_{\mathcal{G}^\#})^\#$。
\end{lemma}

\begin{proof}
从略。
\end{proof}








\section{环站点}
\label{sites-modules-section-ringed-sites}

\noindent
本章主要研究环站点上的模层，因此需要定义这一概念。

\begin{definition}
\label{sites-modules-definition-ringed-site}
环站点。
\begin{enumerate}
\item {\it 环站点}是偶 $(\mathcal{C}, \mathcal{O})$，其中
$\mathcal{C}$ 为站点，$\mathcal{O}$ 为 $\mathcal{C}$ 上的环层。层
$\mathcal{O}$ 称为该环站点的{\it 结构层}。
\item 设 $(\mathcal{C}, \mathcal{O})$、$(\mathcal{C}', \mathcal{O}')$
为环站点。一个{\it 环站点态射}
$$
(f, f^\sharp) :
(\mathcal{C}, \mathcal{O})
\longrightarrow
(\mathcal{C}', \mathcal{O}')
$$
由站点态射 $f : \mathcal{C} \to \mathcal{C}'$（见《站点》定义
\ref{sites-definition-morphism-sites}）以及环层映射
$f^\sharp : f^{-1}\mathcal{O}' \to \mathcal{O}$ 给出；由伴随性，后者等价于
环层映射 $f^\sharp : \mathcal{O}' \to f_*\mathcal{O}$。
\item 设
$(f, f^\sharp) :
(\mathcal{C}_1, \mathcal{O}_1) \to (\mathcal{C}_2, \mathcal{O}_2)$ 以及
$(g, g^\sharp) :
(\mathcal{C}_2, \mathcal{O}_2) \to (\mathcal{C}_3, \mathcal{O}_3)$
为环站点态射。按下列法则定义{\it 环站点态射的复合}：
$$
(g, g^\sharp) \circ (f, f^\sharp) = (g \circ f, f^\sharp \circ g^\sharp).
$$
这里使用《站点》定义 \ref{sites-definition-composition-morphisms-sites}
中定义的站点态射之复合，而 $f^\sharp \circ g^\sharp$ 表示环层态射
$$
\mathcal{O}_3 \xrightarrow{g^\sharp} g_*\mathcal{O}_2
\xrightarrow{g_*f^\sharp} g_*f_*\mathcal{O}_1 = (g \circ f)_*\mathcal{O}_1
$$
\end{enumerate}
\end{definition}






\section{环拓扑斯}
\label{sites-modules-section-ringed-topoi}

\noindent
环拓扑斯本身就是环站点；区别只在于，环拓扑斯态射的概念不同于环站点态射。

\begin{definition}
\label{sites-modules-definition-ringed-topos}
环拓扑斯。
\begin{enumerate}
\item {\it 环拓扑斯}是偶
$(\Sh(\mathcal{C}), \mathcal{O})$
，其中 $\mathcal{C}$ 为站点，$\mathcal{O}$ 为 $\mathcal{C}$ 上的环层。
层 $\mathcal{O}$ 称为该环拓扑斯的{\it 结构层}。
\item 设 $(\Sh(\mathcal{C}), \mathcal{O})$、
$(\Sh(\mathcal{C}'), \mathcal{O}')$ 为环拓扑斯。一个{\it 环拓扑斯态射}
$$
(f, f^\sharp) :
(\Sh(\mathcal{C}), \mathcal{O})
\longrightarrow
(\Sh(\mathcal{C}'), \mathcal{O}')
$$
由拓扑斯态射 $f : \Sh(\mathcal{C}) \to \Sh(\mathcal{C}')$（见《站点》定义
\ref{sites-definition-topos}）以及环层映射
$f^\sharp : f^{-1}\mathcal{O}' \to \mathcal{O}$ 给出；由伴随性，后者等价于
环层映射 $f^\sharp : \mathcal{O}' \to f_*\mathcal{O}$。
\item 设
$(f, f^\sharp) :
(\Sh(\mathcal{C}_1), \mathcal{O}_1)
\to (\Sh(\mathcal{C}_2), \mathcal{O}_2)$ 以及
$(g, g^\sharp) :
(\Sh(\mathcal{C}_2), \mathcal{O}_2) \to
(\Sh(\mathcal{C}_3), \mathcal{O}_3)$
为环拓扑斯态射。按下列法则定义{\it 环拓扑斯态射的复合}：
$$
(g, g^\sharp) \circ (f, f^\sharp) = (g \circ f, f^\sharp \circ g^\sharp).
$$
这里使用《站点》定义 \ref{sites-definition-topos} 中定义的拓扑斯态射之复合，
而 $f^\sharp \circ g^\sharp$ 表示环层态射
$$
\mathcal{O}_3 \xrightarrow{g^\sharp} g_*\mathcal{O}_2
\xrightarrow{g_*f^\sharp} g_*f_*\mathcal{O}_1 = (g \circ f)_*\mathcal{O}_1
$$
\end{enumerate}
\end{definition}

\noindent
每个环拓扑斯态射都是一个环拓扑斯等价与一个由环站点态射诱导的环拓扑斯态射
之复合。准确陈述如下。

\begin{lemma}
\label{sites-modules-lemma-morphism-ringed-topoi-comes-from-morphism-ringed-sites}
设 $(f, f^\sharp) :
(\Sh(\mathcal{C}), \mathcal{O}_\mathcal{C})
\to (\Sh(\mathcal{D}), \mathcal{O}_\mathcal{D})$
为环拓扑斯态射。存在分解
$$
\xymatrix{
(\Sh(\mathcal{C}), \mathcal{O}_\mathcal{C})
\ar[rr]_{(f, f^\sharp)}
\ar[d]_{(g, g^\sharp)}
& &
(\Sh(\mathcal{D}), \mathcal{O}_\mathcal{D}) \ar[d]^{(e, e^\sharp)}
\\
(\Sh(\mathcal{C}'), \mathcal{O}_{\mathcal{C}'})
\ar[rr]^{(h, h^\sharp)} & &
(\Sh(\mathcal{D}'), \mathcal{O}_{\mathcal{D}'})
}
$$
其中
\begin{enumerate}
\item $g : \Sh(\mathcal{C}) \to \Sh(\mathcal{C}')$
是由特殊余连续函子 $\mathcal{C} \to \mathcal{C}'$ 诱导的拓扑斯等价
（见《站点》定义 \ref{sites-definition-special-cocontinuous-functor}），
\item $e : \Sh(\mathcal{D}) \to \Sh(\mathcal{D}')$
是由特殊余连续函子 $\mathcal{D} \to \mathcal{D}'$ 诱导的拓扑斯等价
（见《站点》定义 \ref{sites-definition-special-cocontinuous-functor}），
\item $\mathcal{O}_{\mathcal{C}'} = g_*\mathcal{O}_\mathcal{C}$
且 $g^\sharp$ 为显然映射，
\item $\mathcal{O}_{\mathcal{D}'} = e_*\mathcal{O}_\mathcal{D}$
且 $e^\sharp$ 为显然映射，
\item 站点 $\mathcal{C}'$ 与 $\mathcal{D}'$ 具有终对象与纤维积（即一切有限极限），
\item $h$ 是由与一切有限极限可交换的连续函子
$u : \mathcal{D}' \to \mathcal{C}'$ 诱导的站点态射（即它满足《站点》命题
\ref{sites-proposition-get-morphism} 的假设），以及
\item 给定 $\mathcal{C}$ 上任意一族层 $\mathcal{F}_i$（相应地，
$\mathcal{D}$ 上任意一族层 $\mathcal{G}_j$），可假设它们分别是
$\mathcal{C}'$（相应地，$\mathcal{D}'$）上的可表层。
\end{enumerate}
此外，若 $(f, f^\sharp)$ 是环拓扑斯等价，则可选择该图使得
$\mathcal{C}' = \mathcal{D}'$,
$\mathcal{O}_{\mathcal{C}'} = \mathcal{O}_{\mathcal{D}'}$
且 $(h, h^\sharp)$ 为恒等态射。
\end{lemma}

\begin{proof}
这由《站点》引理 \ref{sites-lemma-morphism-topoi-comes-from-morphism-sites}
以及《站点》注
\ref{sites-remark-morphism-topoi-comes-from-morphism-sites}、
\ref{sites-remark-equivalence-topoi-comes-from-morphism-sites} 得出；只需同时
携带相应环层。部分细节从略。
\end{proof}







\section{环拓扑斯的 2-态射}
\label{sites-modules-section-2-category}

\noindent
本节简要讨论环拓扑斯的 $2$-态射这一概念。

\begin{definition}
\label{sites-modules-definition-2-morphism-ringed-topoi}
设
$f, g :
(\Sh(\mathcal{C}), \mathcal{O}_\mathcal{C})
\to
(\Sh(\mathcal{D}), \mathcal{O}_\mathcal{D})$
为两个环拓扑斯态射。一个{\it 从 $f$ 到 $g$ 的 2-态射}由函子变换
$t : f_* \to g_*$ 给出，使得
$$
\xymatrix{
& \mathcal{O}_\mathcal{D}
\ar[ld]_{f^\sharp}
\ar[rd]^{g^\sharp} \\
f_*\mathcal{O}_\mathcal{C} \ar[rr]^t & &
g_*\mathcal{O}_\mathcal{C}
}
$$
可交换。
\end{definition}

\noindent
在图示中，我们有时将 $t$ 表示如下：
$$
\xymatrix{
(\Sh(\mathcal{C}), \mathcal{O}_\mathcal{C})
\rrtwocell^f_g{t}
&
&
(\Sh(\mathcal{D}), \mathcal{O}_\mathcal{D})
}
$$
如《站点》第 \ref{sites-section-2-category} 节所述，给出 2-态射
$t : f_* \to g_*$ 等价于给出 $t : g^{-1} \to f^{-1}$（通常仍以同一符号
表示），使得图
$$
\xymatrix{
f^{-1}\mathcal{O}_\mathcal{D}
\ar[rd]_{f^\sharp}  & &
g^{-1}\mathcal{O}_\mathcal{D} \ar[ll]^t \ar[ld]^{g^\sharp} \\
& \mathcal{O}_\mathcal{C}
}
$$
可交换。同样依《站点》第 \ref{sites-section-2-category} 节，按该处所述定义
水平复合与垂直复合后，严格 2-范畴的公理成立。










\section{模预层}
\label{sites-modules-section-presheaves-modules}

\noindent
设 $\mathcal{C}$ 为范畴，$\mathcal{O}$ 为 $\mathcal{C}$ 上的环预层。我们尚未
定义 $\mathcal{O}$-模预层，现立即给出定义。

\begin{definition}
\label{sites-modules-definition-presheaf-modules}
设 $\mathcal{C}$ 为范畴，$\mathcal{O}$ 为 $\mathcal{C}$ 上的环预层。
\begin{enumerate}
\item 一个{\it $\mathcal{O}$-模预层}由阿贝尔预层 $\mathcal{F}$ 以及集合
预层映射
$$
\mathcal{O} \times \mathcal{F} \longrightarrow \mathcal{F}
$$
给出，并要求对 $\mathcal{C}$ 的每个对象 $U$，映射
$\mathcal{O}(U) \times \mathcal{F}(U) \to \mathcal{F}(U)$
在阿贝尔群 $\mathcal{F}(U)$ 上定义 $\mathcal{O}(U)$-模结构。
\item 一个{\it $\mathcal{O}$-模预层态射
$\varphi : \mathcal{F} \to \mathcal{G}$}是阿贝尔预层态射
$\varphi : \mathcal{F} \to \mathcal{G}$，使图
$$
\xymatrix{
\mathcal{O} \times \mathcal{F} \ar[r] \ar[d]_{\text{id} \times \varphi} &
\mathcal{F} \ar[d]^{\varphi} \\
\mathcal{O} \times \mathcal{G} \ar[r] &
\mathcal{G}
}
$$
可交换。
\item 上述 $\mathcal{O}$-模态射组成的集合记为
$\Hom_\mathcal{O}(\mathcal{F}, \mathcal{G})$。
\item $\mathcal{O}$-模预层的范畴记为 $\textit{PMod}(\mathcal{O})$。
\end{enumerate}
\end{definition}

\noindent
设 $\mathcal{O}_1 \to \mathcal{O}_2$ 是范畴 $\mathcal{C}$ 上环预层的态射。
若 $\mathcal{F}$ 是 $\mathcal{O}_2$-模预层，则借助复合
$$
\mathcal{O}_1 \times \mathcal{F}
\to
\mathcal{O}_2 \times \mathcal{F}
\to
\mathcal{F}.
$$
可将 $\mathcal{F}$ 视为 $\mathcal{O}_1$-模预层。为表明标量限制，我们有时将
它记为 $\mathcal{F}_{\mathcal{O}_1}$，并称之为 $\mathcal{F}$ 的
{\it 标量限制}。由此得到标量限制函子
$$
\textit{PMod}(\mathcal{O}_2)
\longrightarrow
\textit{PMod}(\mathcal{O}_1)
$$

\medskip\noindent
另一方面，给定 $\mathcal{O}_1$-模预层 $\mathcal{G}$，可按下列法则构造
$\mathcal{O}_2$-模预层
$\mathcal{O}_2 \otimes_{p, \mathcal{O}_1} \mathcal{G}$：
$$
U \longmapsto
\left(\mathcal{O}_2 \otimes_{p, \mathcal{O}_1} \mathcal{G}\right)(U)
=
\mathcal{O}_2(U) \otimes_{\mathcal{O}_1(U)} \mathcal{G}(U)
$$
其中 $U \in \Ob(\mathcal{C})$，限制映射取显然的映射。下标 $p$ 表示
“预层”，而非“点”。该预层称为张量积预层。由此得到{\it 换环}函子
$$
\textit{PMod}(\mathcal{O}_1)
\longrightarrow
\textit{PMod}(\mathcal{O}_2)
$$

\begin{lemma}
\label{sites-modules-lemma-adjointness-tensor-restrict-presheaves}
沿用上述 $\mathcal{C}$、$\mathcal{O}_1 \to \mathcal{O}_2$、$\mathcal{F}$ 与
$\mathcal{G}$，存在典范双射
$$
\Hom_{\mathcal{O}_1}(\mathcal{G}, \mathcal{F}_{\mathcal{O}_1})
=
\Hom_{\mathcal{O}_2}(
\mathcal{O}_2 \otimes_{p, \mathcal{O}_1} \mathcal{G},
\mathcal{F}
)
$$
换言之，上述标量限制函子与换环函子彼此伴随。
\end{lemma}

\begin{proof}
这由下述事实得出：对环同态 $A \to B$，标量限制函子与换环函子彼此伴随。
\end{proof}


\section{模层}
\label{sites-modules-section-sheaves-modules}

\begin{definition}
\label{sites-modules-definition-sheaf-modules}
设 $\mathcal{C}$ 为站点，$\mathcal{O}$ 为 $\mathcal{C}$ 上的环层。
\begin{enumerate}
\item 一个{\it $\mathcal{O}$-模层}是 $\mathcal{O}$-模预层 $\mathcal{F}$
（见定义 \ref{sites-modules-definition-presheaf-modules}），并且其底层阿贝尔群预层
$\mathcal{F}$ 是层。
\item 一个{\it $\mathcal{O}$-模层态射}就是 $\mathcal{O}$-模预层态射。
\item 给定 $\mathcal{O}$-模层 $\mathcal{F}$ 与 $\mathcal{G}$，以
$\Hom_\mathcal{O}(\mathcal{F}, \mathcal{G})$ 表示 $\mathcal{O}$-模层态射
组成的集合。
\item $\mathcal{O}$-模层的范畴记为 $\textit{Mod}(\mathcal{O})$。
\end{enumerate}
\end{definition}

\noindent
即使 $\mathcal{O}$ 只是环预层，这一定义也勉强说得通；不过我们不知道它在
何种例子中有用，并且当 $\mathcal{O}$ 不是环层时，将避免使用
“$\mathcal{O}$-模层”这一术语。



\section{模预层的层化}
\label{sites-modules-section-sheafification-presheaves-modules}

\begin{lemma}
\label{sites-modules-lemma-sheafification-presheaf-modules}
设 $\mathcal{C}$ 为站点，$\mathcal{O}$ 为 $\mathcal{C}$ 上的环预层，
$\mathcal{F}$ 为 $\mathcal{O}$-模预层。设 $\mathcal{O}^\#$ 为将
$\mathcal{O}$ 视为环预层时的层化；见《站点》第
\ref{sites-section-sheaves-algebraic-structures} 节。设 $\mathcal{F}^\#$ 为将
$\mathcal{F}$ 视为阿贝尔群预层时的层化。存在唯一的集合层映射
$$
\mathcal{O}^\# \times \mathcal{F}^\#
\longrightarrow
\mathcal{F}^\#
$$
使得图
$$
\xymatrix{
\mathcal{O} \times \mathcal{F} \ar[r] \ar[d] &
\mathcal{F} \ar[d] \\
\mathcal{O}^\# \times \mathcal{F}^\# \ar[r] &
\mathcal{F}^\#
}
$$
可交换，并使 $\mathcal{F}^\#$ 成为 $\mathcal{O}^\#$-模层。此外，若
$\mathcal{G}$ 是 $\mathcal{O}^\#$-模层，则任意 $\mathcal{O}$-模预层态射
$\mathcal{F} \to \mathcal{G}$（其靶是 $\mathcal{G}$ 限制为
$\mathcal{O}$-模后的预层）都唯一分解为
$\mathcal{F} \to \mathcal{F}^\# \to \mathcal{G}$，其中
$\mathcal{F}^\# \to \mathcal{G}$ 是 $\mathcal{O}^\#$-模态射。
\end{lemma}

\begin{proof}
从略。
\end{proof}

\noindent
这实际上意味着函子
$i : \textit{Mod}(\mathcal{O}^\#) \to \textit{PMod}(\mathcal{O})$
（结合标量限制与将层包含进预层）和该引理中的层化函子
${}^\# : \textit{PMod}(\mathcal{O}) \to \textit{Mod}(\mathcal{O}^\#)$
彼此伴随。公式为
$$
\Mor_{\textit{PMod}(\mathcal{O})}(\mathcal{F}, i\mathcal{G})
=
\Mor_{\textit{Mod}(\mathcal{O}^\#)}(\mathcal{F}^\#, \mathcal{G})
$$
一个重要情形是 $\mathcal{O}$ 已经为环层。此时公式成为
$$
\Mor_{\textit{PMod}(\mathcal{O})}(\mathcal{F}, i\mathcal{G})
=
\Mor_{\textit{Mod}(\mathcal{O})}(\mathcal{F}^\#, \mathcal{G})
$$
因为此时 $\mathcal{O} = \mathcal{O}^\#$。

\begin{lemma}
\label{sites-modules-lemma-sheafification-exact}
设 $\mathcal{C}$ 为站点，$\mathcal{O}$ 为 $\mathcal{C}$ 上的环预层。层化函子
$$
\textit{PMod}(\mathcal{O}) \longrightarrow \textit{Mod}(\mathcal{O}^\#), \quad
\mathcal{F} \longmapsto \mathcal{F}^\#
$$
正合。
\end{lemma}

\begin{proof}
这是因为层化 $\textit{PAb}(\mathcal{C}) \to \textit{Ab}(\mathcal{C})$
具有这一性质；见第 \ref{sites-modules-section-abelian-sheaves} 节的讨论。
\end{proof}

\noindent
设 $\mathcal{C}$ 为站点，$\mathcal{O}_1 \to \mathcal{O}_2$ 为
$\mathcal{C}$ 上环层的态射。在第 \ref{sites-modules-section-presheaves-modules} 节中，
我们定义了此情形所伴随的模预层标量限制函子与换环函子。

\medskip\noindent
若 $\mathcal{F}$ 是 $\mathcal{O}_2$-模层，则 $\mathcal{F}$ 的标量限制
$\mathcal{F}_{\mathcal{O}_1}$ 显然为 $\mathcal{O}_1$-模层。由此得到标量
限制函子
$$
\textit{Mod}(\mathcal{O}_2)
\longrightarrow
\textit{Mod}(\mathcal{O}_1)
$$

\medskip\noindent
另一方面，给定 $\mathcal{O}_1$-模层 $\mathcal{G}$，$\mathcal{O}_2$-模预层
$\mathcal{O}_2 \otimes_{p, \mathcal{O}_1} \mathcal{G}$ 一般不是层。因此按公式
定义{\it 张量积层} $\mathcal{O}_2 \otimes_{\mathcal{O}_1} \mathcal{G}$：
$$
\mathcal{O}_2 \otimes_{\mathcal{O}_1} \mathcal{G}
=
(\mathcal{O}_2 \otimes_{p, \mathcal{O}_1} \mathcal{G})^\#
$$
即取上述预层构造的层化。由此得到{\it 换环}函子
$$
\textit{Mod}(\mathcal{O}_1)
\longrightarrow
\textit{Mod}(\mathcal{O}_2)
$$

\begin{lemma}
\label{sites-modules-lemma-adjointness-tensor-restrict}
沿用上述 $X$、$\mathcal{O}_1$、$\mathcal{O}_2$、$\mathcal{F}$ 与
$\mathcal{G}$，存在典范双射
$$
\Hom_{\mathcal{O}_1}(\mathcal{G}, \mathcal{F}_{\mathcal{O}_1})
=
\Hom_{\mathcal{O}_2}(
\mathcal{O}_2 \otimes_{\mathcal{O}_1} \mathcal{G},
\mathcal{F}
)
$$
换言之，标量限制函子与换环函子彼此伴随。
\end{lemma}

\begin{proof}
这由引理 \ref{sites-modules-lemma-adjointness-tensor-restrict-presheaves} 以及下述事实得出：
$\Hom_{\mathcal{O}_2}(
\mathcal{O}_2 \otimes_{\mathcal{O}_1} \mathcal{G},
\mathcal{F}
)
=
\Hom_{\mathcal{O}_2}(
\mathcal{O}_2 \otimes_{p, \mathcal{O}_1} \mathcal{G},
\mathcal{F}
)$
因为 $\mathcal{F}$ 是层。
\end{proof}

\begin{lemma}
\label{sites-modules-lemma-epimorphism-modules}
设 $\mathcal{C}$ 为站点，$\mathcal{O} \to \mathcal{O}'$ 为环层的满态射，
$\mathcal{G}_1, \mathcal{G}_2$ 为 $\mathcal{O}'$-模层。则
$$
\Hom_{\mathcal{O}'}(\mathcal{G}_1, \mathcal{G}_2) =
\Hom_\mathcal{O}(\mathcal{G}_1, \mathcal{G}_2).
$$
换言之，标量限制函子
$\textit{Mod}(\mathcal{O}') \to \textit{Mod}(\mathcal{O})$ 充分忠实。
\end{lemma}

\begin{proof}
这是《代数》引理 \ref{algebra-lemma-epimorphism-modules} 的层版本，证明完全相同。
\end{proof}




\section{拓扑斯态射与模层}
\label{sites-modules-section-sheaves-modules-functorial}

\noindent
这里的一切材料都完全直接。我们以拓扑斯态射的情形陈述各项结果；当然，
它们对站点态射也成立。

\begin{lemma}
\label{sites-modules-lemma-pushforward-module}
设 $\mathcal{C}$、$\mathcal{D}$ 为站点，
$f : \Sh(\mathcal{C}) \to \Sh(\mathcal{D})$ 为拓扑斯态射，
$\mathcal{O}$ 为 $\mathcal{C}$ 上的环层，$\mathcal{F}$ 为
$\mathcal{O}$-模层。存在自然的集合层映射
$$
f_*\mathcal{O} \times f_*\mathcal{F}
\longrightarrow
f_*\mathcal{F}
$$
使 $f_*\mathcal{F}$ 成为 $f_*\mathcal{O}$-模层。此构造关于
$\mathcal{F}$ 具有函子性。
\end{lemma}

\begin{proof}
以 $\mu : \mathcal{O} \times \mathcal{F} \to \mathcal{F}$ 表示乘法映射。
回顾 $f_*$（作用于集合层时）左正合，因而与直积可交换。因此
$f_*\mu$ 正是所示映射。这证明了引理。
\end{proof}

\begin{lemma}
\label{sites-modules-lemma-pullback-module}
设 $\mathcal{C}$、$\mathcal{D}$ 为站点，
$f : \Sh(\mathcal{C}) \to \Sh(\mathcal{D})$ 为拓扑斯态射，
$\mathcal{O}$ 为 $\mathcal{D}$ 上的环层，$\mathcal{G}$ 为
$\mathcal{O}$-模层。存在自然的集合层映射
$$
f^{-1}\mathcal{O} \times f^{-1}\mathcal{G}
\longrightarrow
f^{-1}\mathcal{G}
$$
使 $f^{-1}\mathcal{G}$ 成为 $f^{-1}\mathcal{O}$-模层。此构造关于
$\mathcal{G}$ 具有函子性。
\end{lemma}

\begin{proof}
以 $\mu : \mathcal{O} \times \mathcal{G} \to \mathcal{G}$ 表示乘法映射。
回顾 $f^{-1}$（作用于集合层时）正合，因而与直积可交换。因此
$f^{-1}\mu$ 正是所示映射。这证明了引理。
\end{proof}

\begin{lemma}
\label{sites-modules-lemma-adjoint-push-pull-modules}
设 $\mathcal{C}$、$\mathcal{D}$ 为站点，
$f : \Sh(\mathcal{C}) \to \Sh(\mathcal{D})$ 为拓扑斯态射，
$\mathcal{O}$ 为 $\mathcal{D}$ 上的环层，$\mathcal{G}$ 为
$\mathcal{O}$-模层，$\mathcal{F}$ 为 $f^{-1}\mathcal{O}$-模层。则
$$
\Mor_{\textit{Mod}(f^{-1}\mathcal{O})}(f^{-1}\mathcal{G}, \mathcal{F})
=
\Mor_{\textit{Mod}(\mathcal{O})}(\mathcal{G}, f_*\mathcal{F}).
$$
这里使用引理 \ref{sites-modules-lemma-pullback-module} 与
\ref{sites-modules-lemma-pushforward-module}，并通过
$\mathcal{O} \to f_*f^{-1}\mathcal{O}$ 作标量限制，将
$f_*\mathcal{F}$ 视为 $\mathcal{O}$-模层。
\end{lemma}

\begin{proof}
首先注意到
$$
\Mor_{\textit{Ab}(\mathcal{C})}(f^{-1}\mathcal{G}, \mathcal{F})
=
\Mor_{\textit{Ab}(\mathcal{D})}(\mathcal{G}, f_*\mathcal{F}).
$$
，这是《站点》命题
\ref{sites-proposition-functoriality-algebraic-structures-topoi}。
设 $\alpha : f^{-1}\mathcal{G} \to \mathcal{F}$ 与
$\beta : \mathcal{G} \to f_*\mathcal{F}$ 是在上述公式下对应的阿贝尔层
态射。我们须证明，$\alpha$ 是 $f^{-1}\mathcal{O}$-线性的，当且仅当
$\beta$ 是 $\mathcal{O}$-线性的。例如，设 $\alpha$ 是
$f^{-1}\mathcal{O}$-线性的，则 $f_*\alpha$ 显然是
$f_*f^{-1}\mathcal{O}$-线性的，因而（由于标量限制是函子）也是
$\mathcal{O}$-线性的。故只需证明伴随映射
$\mathcal{G} \to f_*f^{-1}\mathcal{G}$ 是 $\mathcal{O}$-线性的。利用
$f_*$ 与 $f^{-1}$（作用于集合层时）都与直积可交换，这归结为证明图
$$
\xymatrix{
\mathcal{O} \times \mathcal{G} \ar[r] \ar[d] &
f_*f^{-1}(\mathcal{O} \times \mathcal{G}) \ar[d] \\
\mathcal{G} \ar[r] & f_*f^{-1}\mathcal{G}
}
$$
可交换。这是因为伴随映射
$\text{id}_{\Sh(\mathcal{D})} \to f_*f^{-1}$ 是函子变换。蕴含关系
$\beta$ 线性 $\Rightarrow$ $\alpha$ 线性的证明从略。
\end{proof}

\begin{lemma}
\label{sites-modules-lemma-adjoint-pull-push-modules}
设 $\mathcal{C}$、$\mathcal{D}$ 为站点，
$f : \Sh(\mathcal{C}) \to \Sh(\mathcal{D})$ 为拓扑斯态射，
$\mathcal{O}$ 为 $\mathcal{C}$ 上的环层，$\mathcal{F}$ 为
$\mathcal{O}$-模层，$\mathcal{G}$ 为 $f_*\mathcal{O}$-模层。则
$$
\Mor_{\textit{Mod}(\mathcal{O})}(
\mathcal{O} \otimes_{f^{-1}f_*\mathcal{O}} f^{-1}\mathcal{G}, \mathcal{F})
=
\Mor_{\textit{Mod}(f_*\mathcal{O})}(\mathcal{G}, f_*\mathcal{F}).
$$
这里使用引理 \ref{sites-modules-lemma-pullback-module} 与
\ref{sites-modules-lemma-pushforward-module}，并在张量积定义中使用典范映射
$f^{-1}f_*\mathcal{O} \to \mathcal{O}$。
\end{lemma}

\begin{proof}
注意到
$$
\Mor_{\textit{Mod}(\mathcal{O})}(
\mathcal{O} \otimes_{f^{-1}f_*\mathcal{O}} f^{-1}\mathcal{G}, \mathcal{F})
=
\Mor_{\textit{Mod}(f^{-1}f_*\mathcal{O})}(
f^{-1}\mathcal{G}, \mathcal{F}_{f^{-1}f_*\mathcal{O}})
$$
，这是引理 \ref{sites-modules-lemma-adjointness-tensor-restrict}。于是结论由引理
\ref{sites-modules-lemma-adjoint-push-pull-modules} 得出。
\end{proof}






\section{环拓扑斯态射与模}
\label{sites-modules-section-functoriality-modules}

\noindent
现在已经引入足够多的符号，可以定义模沿环拓扑斯态射的拉回与直像。

\begin{definition}
\label{sites-modules-definition-pushforward}
设
$(f, f^\sharp) :
(\Sh(\mathcal{C}), \mathcal{O}_\mathcal{C})
\to
(\Sh(\mathcal{D}), \mathcal{O}_\mathcal{D})$
为环拓扑斯态射或环站点态射。
\begin{enumerate}
\item 设 $\mathcal{F}$ 为 $\mathcal{O}_\mathcal{C}$-模层。定义
$\mathcal{F}$ 的{\it 直像}为如下 $\mathcal{O}_\mathcal{D}$-模层：作为阿贝尔
群层，它等于 $f_*\mathcal{F}$；其模结构由下述模结构经
$f^\sharp : \mathcal{O}_\mathcal{D} \to f_*\mathcal{O}_\mathcal{C}$
作标量限制而得：
$$
f_*\mathcal{O}_\mathcal{C} \times f_*\mathcal{F}
\longrightarrow
f_*\mathcal{F}
$$
这里使用引理 \ref{sites-modules-lemma-pushforward-module}。
\item 设 $\mathcal{G}$ 为 $\mathcal{O}_\mathcal{D}$-模层。定义其
{\it 拉回} $f^*\mathcal{G}$ 为由下式定义的
$\mathcal{O}_\mathcal{C}$-模层：
$$
f^*\mathcal{G}
=
\mathcal{O}_\mathcal{C} \otimes_{f^{-1}\mathcal{O}_\mathcal{D}}
f^{-1}\mathcal{G}
$$
其中环同态 $f^{-1}\mathcal{O}_\mathcal{D} \to \mathcal{O}_\mathcal{C}$
为 $f^\sharp$，模结构由引理 \ref{sites-modules-lemma-pullback-module} 给出。
\end{enumerate}
\end{definition}

\noindent
这样便定义了函子
\begin{eqnarray*}
f_* : \textit{Mod}(\mathcal{O}_\mathcal{C})
& \longrightarrow &
\textit{Mod}(\mathcal{O}_\mathcal{D}) \\
f^* : \textit{Mod}(\mathcal{O}_\mathcal{D})
& \longrightarrow &
\textit{Mod}(\mathcal{O}_\mathcal{C})
\end{eqnarray*}
关于这两个函子的最后一个结果是：它们确如预期彼此伴随。

\begin{lemma}
\label{sites-modules-lemma-adjoint-pullback-pushforward-modules}
设
$(f, f^\sharp) :
(\Sh(\mathcal{C}), \mathcal{O}_\mathcal{C})
\to
(\Sh(\mathcal{D}), \mathcal{O}_\mathcal{D})$
为环拓扑斯态射或环站点态射。设 $\mathcal{F}$ 为
$\mathcal{O}_\mathcal{C}$-模层，$\mathcal{G}$ 为
$\mathcal{O}_\mathcal{D}$-模层。存在典范双射
$$
\Hom_{\mathcal{O}_\mathcal{C}}(f^*\mathcal{G}, \mathcal{F})
=
\Hom_{\mathcal{O}_\mathcal{D}}(\mathcal{G}, f_*\mathcal{F}).
$$
换言之，函子 $f^*$ 是 $f_*$ 的左伴随。
\end{lemma}

\begin{proof}
这由前面的工作得出：
\begin{eqnarray*}
\Hom_{\mathcal{O}_\mathcal{C}}(f^*\mathcal{G}, \mathcal{F})
& = &
\Mor_{\textit{Mod}(\mathcal{O}_\mathcal{C})}(
\mathcal{O}_\mathcal{C}
\otimes_{f^{-1}\mathcal{O}_\mathcal{D}} f^{-1}\mathcal{G},
\mathcal{F}) \\
& = &
\Mor_{\textit{Mod}(f^{-1}\mathcal{O}_\mathcal{D})}(
f^{-1}\mathcal{G}, \mathcal{F}_{f^{-1}\mathcal{O}_\mathcal{D}}) \\
& = &
\Hom_{\mathcal{O}_\mathcal{D}}(\mathcal{G}, f_*\mathcal{F}).
\end{eqnarray*}
这里使用引理 \ref{sites-modules-lemma-adjointness-tensor-restrict} 与
\ref{sites-modules-lemma-adjoint-push-pull-modules}。
\end{proof}

\begin{lemma}
\label{sites-modules-lemma-push-pull-composition-modules}
设 $(f, f^\sharp) :
(\Sh(\mathcal{C}_1), \mathcal{O}_1)
\to (\Sh(\mathcal{C}_2), \mathcal{O}_2)$ 以及
$(g, g^\sharp) :
(\Sh(\mathcal{C}_2), \mathcal{O}_2) \to
(\Sh(\mathcal{C}_3), \mathcal{O}_3)$
为环拓扑斯态射。存在函子的典范同构
$(g \circ f)_* \cong g_* \circ f_*$ 以及
$(g \circ f)^* \cong f^* \circ g^*$.
\end{lemma}

\begin{proof}
这由定义即得。
\end{proof}





\section{模层的阿贝尔范畴}
\label{sites-modules-section-kernels}

\noindent
设 $(\Sh(\mathcal{C}), \mathcal{O})$ 为环拓扑斯，$\mathcal{F}$、
$\mathcal{G}$ 为 $\mathcal{O}$-模层；见《层》定义
\ref{sheaves-definition-sheaf-modules}。设
$\varphi, \psi : \mathcal{F} \to \mathcal{G}$ 为 $\mathcal{O}$-模层态射。
定义 $\varphi + \psi : \mathcal{F} \to \mathcal{G}$ 为将 $\varphi$ 与
$\psi$ 视作阿贝尔层态射时的和。这显然仍为 $\mathcal{O}$-模映射。也显然，
$\mathcal{O}$-模映射的复合关于此加法是双线性的。因此
$\textit{Mod}(\mathcal{O})$ 为预加性范畴；见《同调》定义
\ref{homology-definition-preadditive}。

\medskip\noindent
以 $0$ 表示这样的 $\mathcal{O}$-模层：对 $\mathcal{C}$ 的每个对象 $U$，
其值恒为 $\{0\}$。它显然既是 $\textit{Mod}(\mathcal{O})$ 的终对象，
也是始对象。给定 $\mathcal{O}$-模态射
$\varphi : \mathcal{F} \to \mathcal{G}$，下列条件等价：
(a) $\varphi$ 为零；(b) $\varphi$ 经过 $0$ 分解；
(c) $\varphi$ 在每个对象 $U$ 上的截面中为零。

\medskip\noindent
此外，给定一对 $\mathcal{O}$-模层 $\mathcal{F}$、$\mathcal{G}$，可定义直和为
$$
\mathcal{F} \oplus \mathcal{G} = \mathcal{F} \times \mathcal{G}
$$
，并取《同调》定义 \ref{homology-definition-direct-sum} 中显然的映射
$(i, j, p, q)$。因此 $\textit{Mod}(\mathcal{O})$ 是加性范畴；见《同调》定义
\ref{homology-definition-additive-category}。

\medskip\noindent
设 $\varphi : \mathcal{F} \to \mathcal{G}$ 为 $\mathcal{O}$-模态射。
可定义 $\Ker(\varphi)$ 为将 $\varphi$ 视作阿贝尔层映射时的核。由第
\ref{sites-modules-section-abelian-sheaves} 节，这是 $\mathcal{F}$ 的子层，其截面为
$$
\Ker(\varphi)(U) =
\{ s \in \mathcal{F}(U) \mid \varphi(s) = 0 \text{ 于 } \mathcal{G}(U)\}
$$
，其中 $U$ 遍历 $\mathcal{C}$ 的所有对象。容易看出，这确为
$\mathcal{O}$-模范畴中的核。换言之，态射
$\alpha : \mathcal{H} \to \mathcal{F}$ 经过 $\Ker(\varphi)$ 分解，当且仅当
$\varphi \circ \alpha = 0$。

\medskip\noindent
类似地，定义 $\Coker(\varphi)$ 为将 $\varphi$ 视作阿贝尔层映射时的余核。
存在唯一的乘法映射
$$
\mathcal{O} \times \Coker(\varphi) \longrightarrow \Coker(\varphi)
$$
使映射 $\mathcal{G} \to \Coker(\varphi)$ 成为 $\mathcal{O}$-模态射（验证从略）。
映射 $\mathcal{G} \to \Coker(\varphi)$ 是满射（将其视为集合层映射；见第
\ref{sites-modules-section-abelian-sheaves} 节）。为证明 $\Coker(\varphi)$ 是
$\textit{Mod}(\mathcal{O})$ 中的余核，注意：若
$\beta : \mathcal{G} \to \mathcal{H}$ 是满足
$\beta \circ \varphi = 0$ 的 $\mathcal{O}$-模态射，则对 $\mathcal{C}$ 的每个
对象 $U$，$\beta$ 诱导映射
$\mathcal{G}(U)/\varphi(\mathcal{F}(U)) \to \mathcal{H}(U)$。由层化的泛性质
（见《层》引理 \ref{sheaves-lemma-sheafification-presheaf-modules}），得到典范
映射 $\Coker(\varphi) \to \mathcal{H}$，使原来的 $\beta$ 等于复合
$\mathcal{G} \to \Coker(\varphi) \to \mathcal{H}$。由上述满射性，态射
$\Coker(\varphi) \to \mathcal{H}$ 唯一。

\begin{lemma}
\label{sites-modules-lemma-abelian}
设 $(\Sh(\mathcal{C}), \mathcal{O})$ 为环拓扑斯。范畴
$\textit{Mod}(\mathcal{O})$ 是阿贝尔范畴。遗忘函子
$\textit{Mod}(\mathcal{O}) \to \textit{Ab}(\mathcal{C})$
正合。因此，$\mathcal{O}$-模的核、余核及正合性分别对应于阿贝尔层的相应概念。
\end{lemma}

\begin{proof}
上面已经看到，$\textit{Mod}(\mathcal{O})$ 是具有核与余核的加性范畴，且
$\textit{Mod}(\mathcal{O}) \to \textit{Ab}(\mathcal{C})$ 保持核与余核。
依《同调》定义 \ref{homology-definition-abelian-category}，只须证明像与余像
一致。这显然成立，因为它在 $\textit{Ab}(\mathcal{C})$ 中成立。引理得证。
\end{proof}

\begin{lemma}
\label{sites-modules-lemma-limits-colimits}
设 $(\Sh(\mathcal{C}), \mathcal{O})$ 为环拓扑斯。
$\textit{Mod}(\mathcal{O})$ 中一切极限与余极限均存在，且遗忘函子
$\textit{Mod}(\mathcal{O}) \to \textit{Ab}(\mathcal{C})$
与它们可交换。此外，滤过余极限正合。
\end{lemma}

\begin{proof}
最后一个断言来自第一个断言，因为由引理
\ref{sites-modules-lemma-limits-colimits-abelian-sheaves}，滤过余极限在
$\textit{Ab}(\mathcal{C})$ 中正合。设
$\mathcal{I} \to \textit{Mod}(\mathcal{C})$，$i \mapsto \mathcal{F}_i$
为图，并设 $\lim_i \mathcal{F}_i$ 为该图在
$\textit{Ab}(\mathcal{C})$ 中的极限。由引理
\ref{sites-modules-lemma-limits-colimits-abelian-sheaves} 对此极限的描述，立即可见存在乘法
$$
\mathcal{O} \times \lim_i \mathcal{F}_i
\longrightarrow
\lim_i \mathcal{F}_i
$$
使 $\lim_i \mathcal{F}_i$ 成为 $\mathcal{O}$-模层。容易看出，这是该图在
$\textit{Mod}(\mathcal{C})$ 中的极限。设 $\colim_i \mathcal{F}_i$ 为该图在
$\textit{PAb}(\mathcal{C})$ 中的余极限。由引理
\ref{sites-modules-lemma-limits-colimits-abelian-presheaves} 的证明中对此余极限的描述，
立即可见存在乘法
$$
\mathcal{O} \times \colim_i \mathcal{F}_i
\longrightarrow
\colim_i \mathcal{F}_i
$$
使 $\colim_i \mathcal{F}_i$ 成为 $\mathcal{O}$-模预层。施行层化，得到
$\mathcal{O}$-模层 $(\colim_i \mathcal{F}_i)^\#$；见引理
\ref{sites-modules-lemma-sheafification-presheaf-modules}。容易看出，
$(\colim_i \mathcal{F}_i)^\#$ 是该图在
$\textit{Mod}(\mathcal{O})$ 中的余极限；又由引理
\ref{sites-modules-lemma-limits-colimits-abelian-sheaves}，遗忘 $\mathcal{O}$-模结构后，
它就是 $\textit{Ab}(\mathcal{C})$ 中的余极限。
\end{proof}

\noindent
极限与余极限的存在，使我们可以像《范畴论》第
\ref{categories-section-exact-functor} 节那样，用极限与余极限讨论定义在
$\mathcal{O}$-模范畴上的函子的正合性。关于用短正合列描述正合性，见《同调》
引理 \ref{homology-lemma-exact-functor}。

\begin{lemma}
\label{sites-modules-lemma-exactness-pushforward-pullback}
设 $f : (\Sh(\mathcal{C}), \mathcal{O}_\mathcal{C})
\to (\Sh(\mathcal{D}), \mathcal{O}_\mathcal{D})$
为环拓扑斯态射。
\begin{enumerate}
\item 函子 $f_*$ 左正合；事实上，它与一切极限可交换。
\item 函子 $f^*$ 右正合；事实上，它与一切余极限可交换。
\end{enumerate}
\end{lemma}

\begin{proof}
这是因为 $(f^*, f_*)$ 是一对伴随函子；见引理
\ref{sites-modules-lemma-adjoint-pullback-pushforward-modules}。另见《范畴论》第
\ref{categories-section-adjoint} 节。
\end{proof}

\begin{lemma}
\label{sites-modules-lemma-check-exactness-stalks}
设 $\mathcal{C}$ 为站点。若 $\{p_i\}_{i \in I}$ 是保守点族，则可在各点
$p_i$（$i \in I$）处的茎上检验阿贝尔层列的正合性。若 $\mathcal{C}$ 有
足够多的点，则阿贝尔层列的正合性可在茎上检验。
\end{lemma}

\begin{proof}
这由《站点》引理 \ref{sites-lemma-exactness-stalks} 立即得出。
\end{proof}





\section{直像的正合性}
\label{sites-modules-section-pushforward}

\noindent
下面给出若干关于直像之正合性质的技术性引理。

\begin{lemma}
\label{sites-modules-lemma-reflect-surjections}
设 $f : \Sh(\mathcal{C}) \to \Sh(\mathcal{D})$ 为拓扑斯态射。下列条件等价：
\begin{enumerate}
\item 对 $\textit{Ab}(\mathcal{C})$ 中所有 $\mathcal{F}$，
$f^{-1}f_*\mathcal{F} \to \mathcal{F}$ 均为满射；以及
\item $f_* : \textit{Ab}(\mathcal{C}) \to \textit{Ab}(\mathcal{D})$
反映满射。
\end{enumerate}
在此情形下，函子
$f_* : \textit{Ab}(\mathcal{C}) \to \textit{Ab}(\mathcal{D})$
忠实。
\end{lemma}

\begin{proof}
假设 (1)。设 $a : \mathcal{F} \to \mathcal{F}'$ 为 $\mathcal{C}$ 上阿贝尔层
的映射，且 $f_*a$ 为满射。由于 $f^{-1}$ 正合，可知
$f^{-1}f_*a : f^{-1}f_*\mathcal{F} \to f^{-1}f_*\mathcal{F}'$
为满射。结合 (1)，可知 $a$ 为满射。这表明 (2) 成立。

\medskip\noindent
假设 (2)。设 $\mathcal{F}$ 为 $\mathcal{C}$ 上阿贝尔层。我们须证明映射
$f^{-1}f_*\mathcal{F} \to \mathcal{F}$ 为满射。由 (2)，只需证明
$f_*f^{-1}f_*\mathcal{F} \to f_*\mathcal{F}$ 为满射。这确实成立，因为存在
作为单侧逆的典范映射 $f_*\mathcal{F} \to f_*f^{-1}f_*\mathcal{F}$。

\medskip\noindent
最后一个断言的证明从略。
\end{proof}

\begin{lemma}
\label{sites-modules-lemma-exactness}
设 $f : \Sh(\mathcal{C}) \to \Sh(\mathcal{D})$ 为拓扑斯态射。假设下列性质
中至少一个成立：
\begin{enumerate}
\item $f_*$ 将集合层的满射变为满射；
\item $f_*$ 将阿贝尔层的满射变为满射；
\item $f_*$ 与集合层上的余等化子可交换；
\item $f_*$ 与集合层上的推出可交换。
\end{enumerate}
则 $f_* : \textit{Ab}(\mathcal{C}) \to \textit{Ab}(\mathcal{D})$ 正合。
\end{lemma}

\begin{proof}
由于 $f_* : \textit{Ab}(\mathcal{C}) \to \textit{Ab}(\mathcal{D})$
是右伴随，我们已经知道，它将 $\mathcal{C}$ 上阿贝尔层的短正合列
$0 \to \mathcal{F}_1 \to \mathcal{F}_2 \to \mathcal{F}_3 \to 0$
变为正合列
$$
0 \to f_*\mathcal{F}_1 \to f_*\mathcal{F}_2 \to f_*\mathcal{F}_3
$$
；见《范畴论》第 \ref{categories-section-exact-functor}、
\ref{categories-section-adjoint} 节以及《同调》第
\ref{homology-section-functors} 节。因此只需证明映射
$f_*\mathcal{F}_2 \to f_*\mathcal{F}_3$ 为满射。若 (1) 或 (2) 成立，
这由定义即得。由《站点》引理 \ref{sites-lemma-exactness-properties}，(3) 或
(4) 都形式地蕴含 (1)，故这些情形也完成了证明。
\end{proof}

\begin{lemma}
\label{sites-modules-lemma-morphism-ringed-sites-almost-cocontinuous}
设 $f : \mathcal{D} \to \mathcal{C}$ 为连续函子
$u : \mathcal{C} \to \mathcal{D}$ 所伴随的站点态射。假设 $u$ 几乎余连续。则
\begin{enumerate}
\item $f_* : \textit{Ab}(\mathcal{D}) \to \textit{Ab}(\mathcal{C})$ 正合。
\item 若给定 $f^\sharp : f^{-1}\mathcal{O}_\mathcal{C} \to
\mathcal{O}_\mathcal{D}$，使 $f$ 成为环站点态射，则
$f_* : \textit{Mod}(\mathcal{O}_\mathcal{D}) \to
\textit{Mod}(\mathcal{O}_\mathcal{C})$ 正合。
\end{enumerate}
\end{lemma}

\begin{proof}
由引理 \ref{sites-modules-lemma-limits-colimits}，第 (2) 项来自第 (1) 项。第 (1) 项来自
《站点》引理 \ref{sites-lemma-morphism-of-sites-almost-cocontinuous} 与
\ref{sites-lemma-exactness-properties}。
\end{proof}





\section{下叹号函子的正合性}
\label{sites-modules-section-exactness-lower-shriek}

\noindent
设 $u : \mathcal{C} \to \mathcal{D}$ 为站点之间的函子。假设
\begin{enumerate}
\item[(a)] $u$ 余连续；且
\item[(b)] $u$ 连续。
\end{enumerate}
设 $g : \Sh(\mathcal{C}) \to \Sh(\mathcal{D})$ 为 $u$ 所伴随的拓扑斯态射；
见《站点》引理 \ref{sites-lemma-cocontinuous-morphism-topoi}。回顾
$g^{-1} = u^p$，即 $g^{-1}$ 由简单公式
$(g^{-1}\mathcal{G})(U) = \mathcal{G}(u(U))$ 给出；见《站点》引理
\ref{sites-lemma-when-shriek}。我们要证明
$g^{-1} : \textit{Ab}(\mathcal{D}) \to \textit{Ab}(\mathcal{C})$
具有左伴随 $g_!$。由《站点》引理 \ref{sites-lemma-when-shriek}，函子
$g^{Sh}_! = (u_p\ )^\#$ 是集合层上的左伴随。此外，我们知道
$g^{Sh}_!\mathcal{F}$ 是与预层
$$
V \longmapsto \colim_{V \to u(U)} \mathcal{F}(U)
$$
相伴的层，其中余极限取遍 $(\mathcal{I}_V^u)^{opp}$，并在集合范畴中计算。
因此下述定义很自然。

\begin{definition}
\label{sites-modules-definition-g-shriek}
设 $u : \mathcal{C} \to \mathcal{D}$ 满足上述 (a)、(b)。对
$\mathcal{F} \in \textit{PAb}(\mathcal{C})$，定义
{\it $g_{p!}\mathcal{F}$} 为预层
$$
V \longmapsto \colim_{V \to u(U)} \mathcal{F}(U)
$$
，其中取遍 $(\mathcal{I}_V^u)^{opp}$ 的余极限在 $\textit{Ab}$ 中计算。
对 $\mathcal{F} \in \textit{PAb}(\mathcal{C})$，令
{\it $g_!\mathcal{F} = (g_{p!}\mathcal{F})^\#$}。
\end{definition}

\noindent
之所以如此明确，是因为函子 $g^{Sh}_!$ 与 $g_!$ 不同。每当同时使用二者时，
都必须谨慎地区分它们。

\begin{lemma}
\label{sites-modules-lemma-g-shriek-adjoint}
函子 $g_{p!}$ 是 $u^p$ 的左伴随；函子 $g_!$ 是 $g^{-1}$ 的左伴随。换言之，公式
\begin{align*}
\Mor_{\textit{PAb}(\mathcal{C})}(\mathcal{F}, u^p\mathcal{G})
& =
\Mor_{\textit{PAb}(\mathcal{D})}(g_{p!}\mathcal{F}, \mathcal{G}), \\
\Mor_{\textit{Ab}(\mathcal{C})}(\mathcal{F}, g^{-1}\mathcal{G})
& =
\Mor_{\textit{Ab}(\mathcal{D})}(g_!\mathcal{F}, \mathcal{G})
\end{align*}
关于 $\mathcal{F}$ 与 $\mathcal{G}$ 双函子性地成立。
\end{lemma}

\begin{proof}
第二个公式形式地来自第一个公式，因为若 $\mathcal{F}$ 与 $\mathcal{G}$
为阿贝尔层，则
\begin{align*}
\Mor_{\textit{Ab}(\mathcal{C})}(\mathcal{F}, g^{-1}\mathcal{G})
& =
\Mor_{\textit{PAb}(\mathcal{D})}(g_{p!}\mathcal{F}, \mathcal{G}) \\
& =
\Mor_{\textit{Ab}(\mathcal{D})}(g_!\mathcal{F}, \mathcal{G})
\end{align*}
，这是层化的泛性质。

\medskip\noindent
为证明第一个公式，设 $\mathcal{F}$、$\mathcal{G}$ 为阿贝尔预层。我们将从
左端群到右端群构造映射，并略去验证这些映射互为逆映射。

\medskip\noindent
注意，存在阿贝尔预层的典范映射
$\mathcal{F} \to u^pg_{p!}\mathcal{F}$；它在 $U$ 上的截面中是自然映射
$\mathcal{F}(U) \to \colim_{u(U) \to u(U')} \mathcal{F}(U')$；见《站点》
引理 \ref{sites-lemma-recover}。给定映射
$\alpha : g_{p!}\mathcal{F} \to \mathcal{G}$，得到
$u^p\alpha : u^pg_{p!}\mathcal{F} \to u^p\mathcal{G}$，再将它与映射
$\mathcal{F} \to u^pg_{p!}\mathcal{F}$ 前复合。

\medskip\noindent
注意，存在阿贝尔预层的典范映射
$g_{p!}u^p\mathcal{G} \to \mathcal{G}$；它在 $V$ 上的截面中是自然映射
$\colim_{V \to u(U)} \mathcal{G}(u(U)) \to \mathcal{G}(V)$。它把与
$t : V \to u(U)$ 对应的直和项中的截面 $s \in u(U)$ 映为
$t^*s \in \mathcal{G}(V)$。因此，给定映射
$\beta : \mathcal{F} \to u^p\mathcal{G}$，得到映射
$g_{p!}\beta : g_{p!}\mathcal{F} \to g_{p!}u^p\mathcal{G}$，再将它与上述映射
$g_{p!}u^p\mathcal{G} \to \mathcal{G}$ 后复合。
\end{proof}

\begin{lemma}
\label{sites-modules-lemma-exactness-lower-shriek}
设 $\mathcal{C}$ 与 $\mathcal{D}$ 为站点，$u : \mathcal{C} \to \mathcal{D}$
为函子。假设
\begin{enumerate}
\item[(a)] $u$ 余连续；
\item[(b)] $u$ 连续；且
\item[(c)] $\mathcal{C}$ 中存在纤维积与等化子，并且 $u$ 与它们可交换。
\end{enumerate}
在此情形下，函子
$g_! : \textit{Ab}(\mathcal{C}) \to \textit{Ab}(\mathcal{D})$
正合。
\end{lemma}

\begin{proof}
比较《站点》引理 \ref{sites-lemma-preserve-equalizers}。假设 (a)、(b)、(c)。
我们已经知道 $g_!$ 右正合，因为它是左伴随；见《范畴论》引理
\ref{categories-lemma-exact-adjoint} 与《同调》第
\ref{homology-section-functors} 节。我们有 $g_! = (g_{p!}\ )^\#$。须证明
$g_!$ 将阿贝尔层的单射映为阿贝尔预层的单射。回顾阿贝尔预层的层化正合；
见引理 \ref{sites-modules-lemma-limits-colimits-abelian-sheaves}。因此，只需证明
$g_{p!}$ 将阿贝尔预层的单射映为阿贝尔预层的单射。为此，只需《站点》第
\ref{sites-section-functoriality-PSh} 节中取遍范畴
$(\mathcal{I}_V^u)^{opp}$ 的余极限把图之间的单射映为单射。这由《站点》引理
\ref{sites-lemma-almost-directed} 与《代数》引理
\ref{algebra-lemma-almost-directed-colimit-exact} 得出。
\end{proof}

\begin{lemma}
\label{sites-modules-lemma-back-and-forth}
设 $\mathcal{C}$ 与 $\mathcal{D}$ 为站点，$u : \mathcal{C} \to \mathcal{D}$
为函子。假设
\begin{enumerate}
\item[(a)] $u$ 余连续；
\item[(b)] $u$ 连续；且
\item[(c)] $u$ 充分忠实。
\end{enumerate}
对上述 $g_!, g^{-1}, g_*$，典范映射
$\mathcal{F} \to g^{-1}g_!\mathcal{F}$ 与
$g^{-1}g_*\mathcal{F} \to \mathcal{F}$
对 $\mathcal{C}$ 上一切阿贝尔层 $\mathcal{F}$ 均为同构。
\end{lemma}

\begin{proof}
由《站点》引理 \ref{sites-lemma-back-and-forth}，以及阿贝尔层的逆像与直像
同其底层集合层的逆像与直像一致这一事实，映射
$g^{-1}g_*\mathcal{F} \to \mathcal{F}$ 为同构。

\medskip\noindent
取 $U \in \Ob(\mathcal{C})$。我们将证明
$g^{-1}g_!\mathcal{F}(U) = \mathcal{F}(U)$。先注意到
$g^{-1}g_!\mathcal{F}(U) = g_!\mathcal{F}(u(U))$。因此只需证明
$g_!\mathcal{F}(u(U)) = \mathcal{F}(U)$。我们知道，$g_!\mathcal{F}$ 是与
预层 $g_{p!}\mathcal{F}$ 相伴的阿贝尔层，而后者由法则
$$
V \longmapsto \colim_{V \to u(U')} \mathcal{F}(U')
$$
定义，其中余极限在 $\textit{Ab}$ 中计算。若 $V = u(U)$，则由于 $u$ 充分
忠实，此余极限取遍 $U \to U'$。故得到
$g_{p!}\mathcal{F}(u(U) = \mathcal{F}(U)$。由于 $u$ 既余连续又连续，
$\mathcal{D}$ 中 $u(U)$ 的任意覆盖都可由 $\mathcal{D}$ 的覆盖（！）
$\{u(U_i) \to u(U)\}$ 加细，其中 $\{U_i \to U\}$ 是 $\mathcal{C}$ 中的覆盖。
这还蕴含 $(g_{p!}\mathcal{F})^+(u(U)) = \mathcal{F}(U)$，因为在定义
$(g_{p!}\mathcal{F})^+$ 于 $u(U)$ 的值的余极限中，可限制到上述覆盖
$\{u(U_i) \to u(U)\}$ 所成的共终系统。于是对 $\mathcal{C}$ 的所有对象 $U$，
同样有 $(g_{p!}\mathcal{F})^+(u(U)) = \mathcal{F}(U)$。再重复一次此论证，
对 $\mathcal{C}$ 的所有对象 $U$ 得到
$(g_{p!}\mathcal{F})^\#(u(U)) = \mathcal{F}(U)$。这给出所需等式
$g^{-1}g_!\mathcal{F} = \mathcal{F}$。
\end{proof}

\begin{remark}
\label{sites-modules-remark-no-extension}
一般而言，当 $g$ 是环拓扑斯态射（的一部分）时，函子 $g_!$ 不能延伸到模
范畴。事实上，给定任意环同态 $A \to B$，函子
$M \mapsto B \otimes_A M$ 具有右伴随（标量限制），但一般没有左伴随
（因为左伴随的存在将蕴含 $A \to B$ 平坦）。下面将在第
\ref{sites-modules-section-localize} 节看到：对站点的局部化，可以在模层上定义 $j_!$。
我们将在第 \ref{sites-modules-section-lower-shriek-modules} 节以更大的一般性讨论此事。
\end{remark}

\begin{lemma}
\label{sites-modules-lemma-have-left-adjoint}
设 $\mathcal{C}$ 与 $\mathcal{D}$ 为站点，
$g : \Sh(\mathcal{C}) \to \Sh(\mathcal{D})$ 为连续且余连续的函子
$u : \mathcal{C} \to \mathcal{D}$ 所伴随的拓扑斯态射。
\begin{enumerate}
\item 若 $u$ 具有左伴随 $w$，则 $g_!$ 与 $g_!^{\Sh}$ 在底层集合层上
一致，且 $g_!$ 正合。
\item 若此外 $w$ 余连续，则 $g_! = h^{-1}$ 且 $g^{-1} = h_*$，其中
$h : \Sh(\mathcal{D}) \to \Sh(\mathcal{C})$ 是 $w$ 所伴随的拓扑斯态射。
\end{enumerate}
\end{lemma}

\begin{proof}
本引理对应于《站点》引理 \ref{sites-lemma-have-left-adjoint}。由《站点》引理
\ref{sites-lemma-adjoint-functors}，范畴 $\mathcal{I}_V^u$ 具有始对象。因此，
取遍 $(\mathcal{I}_V^u)^{opp}$ 的阿贝尔群系统之余极限的底层集合，就是底层
集合的余极限。结合定义 \ref{sites-modules-definition-g-shriek} 中对 $g_!$ 的构造，遂得
$g_!^{\Sh}$ 与 $g_!$ 一致。正合性及 (2) 立即来自《站点》引理
\ref{sites-lemma-have-left-adjoint} 中的相应断言。
\end{proof}





\section{模的整体类型}
\label{sites-modules-section-global}

\begin{definition}
\label{sites-modules-definition-global}
设 $(\Sh(\mathcal{C}), \mathcal{O})$ 为环拓扑斯，$\mathcal{F}$ 为
$\mathcal{O}$-模层。
\begin{enumerate}
\item 若 $\mathcal{F}$ 作为 $\mathcal{O}$-模层同构于形如
$\bigoplus_{i \in I} \mathcal{O}$ 的层，则称 $\mathcal{F}$ 为
{\it 自由 $\mathcal{O}$-模层}。
\item 若 $\mathcal{F}$ 作为 $\mathcal{O}$-模层同构于形如
$\bigoplus_{i \in I} \mathcal{O}$ 的层，其中指标集 $I$ 有限，则称
$\mathcal{F}$ {\it 有限自由}。
\item 若存在从自由 $\mathcal{O}$-模层到 $\mathcal{F}$ 的满射
$$
\bigoplus\nolimits_{i \in I} \mathcal{O} \longrightarrow \mathcal{F}
$$
，则称 $\mathcal{F}$ {\it 由整体截面生成}。
\item 给定 $r \geq 0$，若存在满射
$\mathcal{O}^{\oplus r} \to \mathcal{F}$，则称 $\mathcal{F}$
{\it 由 $r$ 个整体截面生成}。
\item 若对某个 $r \geq 0$，$\mathcal{F}$ 由 $r$ 个整体截面生成，则称
$\mathcal{F}$ {\it 由有限多个整体截面生成}。
\item 若存在 $\mathcal{O}$-模层正合列
$$
\bigoplus\nolimits_{j \in J} \mathcal{O} \longrightarrow
\bigoplus\nolimits_{i \in I} \mathcal{O} \longrightarrow
\mathcal{F} \longrightarrow 0
$$
，则称 $\mathcal{F}$ 具有{\it 整体表示}。
\item 若存在 $\mathcal{O}$-模层正合列
$$
\bigoplus\nolimits_{j \in J} \mathcal{O} \longrightarrow
\bigoplus\nolimits_{i \in I} \mathcal{O} \longrightarrow
\mathcal{F} \longrightarrow 0
$$
，其中 $I$ 与 $J$ 为有限集，则称 $\mathcal{F}$ 具有{\it 整体有限表示}。
\end{enumerate}
\end{definition}

\noindent
注意，对任意集合 $I$，直和 $\bigoplus_{i \in I} \mathcal{O}$ 存在
（引理 \ref{sites-modules-lemma-limits-colimits}），并且它是预层
$U \mapsto \bigoplus_{i \in I} \mathcal{O}(U)$ 的层化。该模层称为
{\it 集合 $I$ 上的自由 $\mathcal{O}$-模层}。

\begin{lemma}
\label{sites-modules-lemma-global-pullback}
设
$(f, f^\sharp) :
(\Sh(\mathcal{C}), \mathcal{O}_\mathcal{C})
\to
(\Sh(\mathcal{D}), \mathcal{O}_\mathcal{D})$
为环拓扑斯态射，$\mathcal{F}$ 为 $\mathcal{O}_\mathcal{D}$-模层。
\begin{enumerate}
\item 若 $\mathcal{F}$ 自由，则 $f^*\mathcal{F}$ 自由。
\item 若 $\mathcal{F}$ 有限自由，则 $f^*\mathcal{F}$ 有限自由。
\item 若 $\mathcal{F}$ 由整体截面生成，则 $f^*\mathcal{F}$ 由整体截面生成。
\item 给定 $r \geq 0$，若 $\mathcal{F}$ 由 $r$ 个整体截面生成，则
$f^*\mathcal{F}$ 由 $r$ 个整体截面生成。
\item 若 $\mathcal{F}$ 由有限多个整体截面生成，则 $f^*\mathcal{F}$ 亦然。
\item 若 $\mathcal{F}$ 具有整体表示，则 $f^*\mathcal{F}$ 亦然。
\item 若 $\mathcal{F}$ 具有整体有限表示，则 $f^*\mathcal{F}$ 亦然。
\end{enumerate}
\end{lemma}

\begin{proof}
这是因为 $f^*$ 与任意余极限可交换（引理
\ref{sites-modules-lemma-exactness-pushforward-pullback}），且
$f^*\mathcal{O}_\mathcal{D} = \mathcal{O}_\mathcal{C}$。
\end{proof}






\section{模的内蕴性质}
\label{sites-modules-section-intrinsic}

\noindent
设 $\mathcal{P}$ 为环拓扑斯上的模层之性质。若每当
$(f, f^\sharp) :
(\Sh(\mathcal{C}'), \mathcal{O}')
\to
(\Sh(\mathcal{C}), \mathcal{O})$
是环拓扑斯等价时，都有
$\mathcal{P}(\mathcal{F}) \Leftrightarrow \mathcal{P}(f^*\mathcal{F})$
，则称 $\mathcal{P}$ 为{\it 内蕴性质}。例如，自由性是内蕴的。事实上，
集合 $I$ 上的自由 $\mathcal{O}$-模层由下述性质刻画：
$$
\Mor_{\textit{Mod}(\mathcal{O})}(
\bigoplus\nolimits_{i \in I} \mathcal{O},
\mathcal{F})
=
\prod\nolimits_{i \in I} \Mor_{\Sh(\mathcal{C})}(\{*\},
\mathcal{F})
$$
其中 $\mathcal{F}$ 遍历 $\textit{Mod}(\mathcal{O})$。也可以利用引理
\ref{sites-modules-lemma-global-pullback} 看出自由性是内蕴的。事实上，出于同一理由，
定义 \ref{sites-modules-definition-global} 中的每个性质都是内蕴的。接下来如何定义
$\mathcal{O}$-模的其他内蕴性质？

\medskip\noindent
引理 \ref{sites-modules-lemma-morphism-ringed-topoi-comes-from-morphism-ringed-sites}
的要点如下：假设要定义拓扑斯上 $\mathcal{O}$-模的内蕴性质 $\mathcal{P}$，
可按以下步骤进行：
\begin{enumerate}
\item 给定任意站点 $\mathcal{C}$、$\mathcal{C}$ 上任意环层
$\mathcal{O}$ 及任意 $\mathcal{O}$-模层 $\mathcal{F}$，定义相应性质
$\mathcal{P}(\mathcal{C}, \mathcal{O}, \mathcal{F})$。
\item 对任意一对站点 $\mathcal{C}$、$\mathcal{C}'$，任意特殊余连续函子
$u : \mathcal{C} \to \mathcal{C}'$，$\mathcal{C}$ 上任意环层
$\mathcal{O}$ 及任意 $\mathcal{O}$-模层 $\mathcal{F}$，证明
$$
\mathcal{P}(\mathcal{C}, \mathcal{O}, \mathcal{F})
\Leftrightarrow
\mathcal{P}(\mathcal{C}', g_*\mathcal{O}, g_*\mathcal{F})
$$
，其中 $g : \Sh(\mathcal{C}) \to \Sh(\mathcal{C}')$ 是 $u$ 所伴随的
拓扑斯等价。
\end{enumerate}
在此情形下，给定任意环拓扑斯 $(\Sh(\mathcal{C}), \mathcal{O})$ 与任意
$\mathcal{O}$-模层 $\mathcal{F}$，若
$\mathcal{P}(\mathcal{C}, \mathcal{O}, \mathcal{F})$ 成立，就简称
$\mathcal{F}$ 具有性质 $\mathcal{P}$。引理
\ref{sites-modules-lemma-morphism-ringed-topoi-comes-from-morphism-ringed-sites} 与上述 (2)
共同保证此定义良定。

\medskip\noindent
此外，同一引理 \ref{sites-modules-lemma-morphism-ringed-topoi-comes-from-morphism-ringed-sites}
还保证，若进一步有
\begin{enumerate}
\item[(3)] 对任意环站点态射
$(f, f^\sharp) :
(\mathcal{C}, \mathcal{O}_\mathcal{C})
\to
(\mathcal{D}, \mathcal{O}_\mathcal{D})$
，其中 $f$ 由满足《站点》命题 \ref{sites-proposition-get-morphism} 之假设的
函子 $u : \mathcal{D} \to \mathcal{C}$ 给出，并且对任意
$\mathcal{O}_\mathcal{D}$-模层 $\mathcal{G}$，都有
$$
\mathcal{P}(\mathcal{D}, \mathcal{O}_\mathcal{D}, \mathcal{F})
\Rightarrow
\mathcal{P}(\mathcal{C}, \mathcal{O}_\mathcal{C}, f^*\mathcal{F})
$$
\end{enumerate}
，则 $\mathcal{P}$ 在模沿任意环拓扑斯态射拉回时得以保持。

\medskip\noindent
以下各节将用此方法看出：局部自由、局部由截面生成、局部由 $r$ 个截面生成、
有限型、有限表示、拟凝聚与凝聚，都是模的内蕴性质。

\medskip\noindent
与其采用此处概述的繁琐过程，也许为这些概念寻找内蕴定义更令人满意。另一方面，
在许多要应用这些定义的几何情形中，给定的是一个具体环站点与一个具体模层；
因此，已有适合这套语言的定义也很方便。




\section{环站点的局部化}
\label{sites-modules-section-localize}

\noindent
设 $(\mathcal{C}, \mathcal{O})$ 为环站点，$U \in \Ob(\mathcal{C})$。
我们说明《站点》第 \ref{sites-section-localize} 节中各结果在此情形下的对应物。

\medskip\noindent
以 $\mathcal{O}_U = j_U^{-1}\mathcal{O}$ 表示 $\mathcal{O}$ 到站点
$\mathcal{C}/U$ 的限制。它由简单法则
$\mathcal{O}_U(V/U) = \mathcal{O}(V)$ 描述。采用此符号后，局部化态射
$j_U$ 成为环拓扑斯态射
$$
(j_U, j_U^\sharp) :
(\Sh(\mathcal{C}/U), \mathcal{O}_U)
\longrightarrow
(\Sh(\mathcal{C}), \mathcal{O})
$$
；也就是说，取 $j_U^\sharp : j_U^{-1}\mathcal{O} \to \mathcal{O}_U$
为恒等映射。此外，模的直像与拉回具有下述描述。

\begin{definition}
\label{sites-modules-definition-localize-ringed-site}
设 $(\mathcal{C}, \mathcal{O})$ 为环站点，$U \in \Ob(\mathcal{C})$。
\begin{enumerate}
\item 环站点 $(\mathcal{C}/U, \mathcal{O}_U)$ 称为
{\it 环站点 $(\mathcal{C}, \mathcal{O})$ 在对象 $U$ 处的局部化}。
\item 环拓扑斯态射
$(j_U, j_U^\sharp) :
(\Sh(\mathcal{C}/U), \mathcal{O}_U)
\to
(\Sh(\mathcal{C}), \mathcal{O})$
称为{\it 局部化态射}。
\item 函子
$j_{U*} : \textit{Mod}(\mathcal{O}_U) \to \textit{Mod}(\mathcal{O})$
称为{\it 直像函子}。
\item 对 $\mathcal{C}$ 上的 $\mathcal{O}$-模层 $\mathcal{F}$，层
$j_U^*\mathcal{F}$ 称为{\it $\mathcal{F}$ 到 $\mathcal{C}/U$ 的限制}。
有时将它记为 $\mathcal{F}|_{\mathcal{C}/U}$，甚至简记为
$\mathcal{F}|_U$。它由简单法则
$j_U^*(\mathcal{F})(X/U) = \mathcal{F}(X)$ 描述。
\item 限制函子的左伴随
$j_{U!} : \textit{Mod}(\mathcal{O}_U) \to \textit{Mod}(\mathcal{O})$
称为{\it 以零延拓}。由引理 \ref{sites-modules-lemma-extension-by-zero} 与
\ref{sites-modules-lemma-extension-by-zero-exact}，它存在且正合。
\end{enumerate}
\end{definition}

\noindent
与拓扑情形相同（见《层》第 \ref{sheaves-section-open-immersions} 节），
以零延拓函子 $j_{U!}$ 不同于在集合层上定义的以空集延拓 $j_{U!}$。
下面的引理定义以零延拓。

\begin{lemma}
\label{sites-modules-lemma-extension-by-zero}
设 $(\mathcal{C}, \mathcal{O})$ 为环站点，$U \in \Ob(\mathcal{C})$。限制函子
$j_U^* : \textit{Mod}(\mathcal{O}) \to \textit{Mod}(\mathcal{O}_U)$
具有左伴随
$j_{U!} : \textit{Mod}(\mathcal{O}_U) \to \textit{Mod}(\mathcal{O})$.
因此
$$
\Mor_{\textit{Mod}(\mathcal{O}_U)}(\mathcal{G}, j_U^*\mathcal{F})
=
\Mor_{\textit{Mod}(\mathcal{O})}(j_{U!}\mathcal{G}, \mathcal{F})
$$
对 $\mathcal{F} \in \Ob(\textit{Mod}(\mathcal{O}))$ 与
$\mathcal{G} \in \Ob(\textit{Mod}(\mathcal{O}_U))$ 成立。此外，
$\mathcal{G}$ 的以零延拓 $j_{U!}\mathcal{G}$ 是与预层
$$
V
\longmapsto
\bigoplus\nolimits_{\varphi \in \Mor_\mathcal{C}(V, U)}
\mathcal{G}(V \xrightarrow{\varphi} U)
$$
相伴的层，其中限制映射与 $\mathcal{O}$-模结构均取显然者。
\end{lemma}

\begin{proof}
预层上的 $\mathcal{O}$-模结构定义如下。若 $f \in \mathcal{O}(V)$ 且
$s \in \mathcal{G}(V \xrightarrow{\varphi} U)$，则定义
$f \cdot s = fs$，其中
$f \in \mathcal{O}_U(\varphi : V \to U) = \mathcal{O}(V)$
（因为 $\mathcal{O}_U$ 是 $\mathcal{O}$ 到 $\mathcal{C}/U$ 的限制）。

\medskip\noindent
类似地，设 $\alpha : \mathcal{G} \to \mathcal{F}|_U$ 为
$\mathcal{O}_U$-模态射。此时可定义从引理中的预层到 $\mathcal{F}$ 的映射
$$
\bigoplus\nolimits_{\varphi \in \Mor_\mathcal{C}(V, U)}
\mathcal{G}(V \xrightarrow{\varphi} U)
\longrightarrow
\mathcal{F}(V)
$$
，法则为：将 $s \in \mathcal{G}(V \xrightarrow{\varphi} U)$ 映为
$\alpha(s) \in \mathcal{F}(V)$。这显然为 $\mathcal{O}$-线性的，因而由模层化
的性质（引理 \ref{sites-modules-lemma-sheafification-presheaf-modules}）诱导
$\mathcal{O}$-模态射 $\alpha' : j_{U!}\mathcal{G} \to \mathcal{F}$。

\medskip\noindent
反之，设 $\beta : j_{U!}\mathcal{G} \to \mathcal{F}$ 为
$\mathcal{O}$-模映射。回顾《站点》第 \ref{sites-section-localize} 节：在集合
层上存在以空集延拓
$j^{Sh}_{U!} : \Sh(\mathcal{C}/U) \to \Sh(\mathcal{C})$
，它是 $j_U^{-1}$ 的左伴随。此外，$j^{Sh}_{U!}\mathcal{G}$ 是与预层
$$
V
\longmapsto
\coprod\nolimits_{\varphi \in \Mor_\mathcal{C}(V, U)}
\mathcal{G}(V \xrightarrow{\varphi} U)
$$
相伴的层。因此存在集合层的自然映射
$j^{Sh}_{U!}\mathcal{G} \to j_{U!}\mathcal{G}$。将 $\beta$ 与此映射前复合，
得到集合层映射 $j^{Sh}_{U!}\mathcal{G} \to \mathcal{F}$；由伴随性，它对应于
集合层映射 $\beta' : \mathcal{G} \to \mathcal{F}|_U$。我们断言
$\beta'$ 是 $\mathcal{O}_U$-线性的。事实上，设
$\varphi : V \to U$ 为 $\mathcal{C}/U$ 的对象，
$s, s' \in \mathcal{G}(\varphi : V \to U)$，且
$f \in \mathcal{O}(V) = \mathcal{O}_U(\varphi : V \to U)$。由上述讨论，
$\mathcal{F}|_U(\varphi : V \to U)$ 中的 $\beta'(s + s')$、相应地
$\beta'(fs)$，分别对应于 $\mathcal{F}(V)$ 中的 $\beta(s + s')$、
相应地 $\beta(fs)$。由于 $\beta$ 是同态，结论成立。

\medskip\noindent
为完成引理的证明，须证明构造 $\alpha \mapsto \alpha'$ 与
$\beta \mapsto \beta'$ 互为逆构造。验证从略。
\end{proof}

\noindent
注意，在定义 \ref{sites-modules-definition-localize-ringed-site} 的情形中，对每个
$\mathcal{O}$-模 $\mathcal{F}$，都有
\begin{equation}
\label{sites-modules-equation-map-lower-shriek-OU-into-module}
\Hom_\mathcal{O}(j_{U!}\mathcal{O}_U, \mathcal{F}) =
\Hom_{\mathcal{O}_U}(\mathcal{O}_U, j_U^*\mathcal{F}) =
\mathcal{F}(U)
\end{equation}
。事实上，第一个等号来自 $j_{U!}$ 与 $j_U^*$ 的伴随性，第二个等号则来自
{\small $\Hom_{\mathcal{O}_U}(\mathcal{O}_U, j_U^*\mathcal{F}) =
j_U^*\mathcal{F}(U/U)$}；又有
{\small $j_U^*\mathcal{F}(U/U) = \mathcal{F}|_U(U/U) = \mathcal{F}(U)$}。

\begin{lemma}
\label{sites-modules-lemma-extension-by-zero-exact}
设 $(\mathcal{C}, \mathcal{O})$ 是环站点，且 $U \in \Ob(\mathcal{C})$。函子
$j_{U!} : \textit{Mod}(\mathcal{O}_U) \to \textit{Mod}(\mathcal{O})$
是正合的。
\end{lemma}

\begin{proof}
由于 $j_{U!}$ 是 $j_U^*$ 的左伴随，它是右正合的（见
《范畴》，引理 \ref{categories-lemma-exact-adjoint}，以及
《同调代数》，第 \ref{homology-section-functors} 节）。
因此，只需证明：若 $\mathcal{G}_1 \to \mathcal{G}_2$ 是
$\mathcal{O}_U$-模的单射，则
$j_{U!}\mathcal{G}_1 \to j_{U!}\mathcal{G}_2$ 也是单射。
在对象 $V$ 上的预层截面映射（如引理 \ref{sites-modules-lemma-extension-by-zero} 中所述）为
$$
\bigoplus\nolimits_{\varphi \in \Mor_\mathcal{C}(V, U)}
\mathcal{G}_1(V \xrightarrow{\varphi} U)
\longrightarrow
\bigoplus\nolimits_{\varphi \in \Mor_\mathcal{C}(V, U)}
\mathcal{G}_2(V \xrightarrow{\varphi} U)
$$
，由假设它是单射。根据引理 \ref{sites-modules-lemma-sheafification-exact}，层化是正合的，
故 $j_{U!}\mathcal{G}_1 \to j_{U!}\mathcal{G}_2$ 是单射，结论成立。
\end{proof}

\begin{lemma}
\label{sites-modules-lemma-j-shriek-reflects-exactness}
设 $(\mathcal{C}, \mathcal{O})$ 是环站点，且 $U \in \Ob(\mathcal{C})$。
$\mathcal{O}_U$-模的复形
$\mathcal{G}_1 \to \mathcal{G}_2 \to \mathcal{G}_3$ 正合，当且仅当
$j_{U!}\mathcal{G}_1 \to j_{U!}\mathcal{G}_2 \to j_{U!}\mathcal{G}_3$
作为 $\mathcal{O}$-模列正合。
\end{lemma}

\begin{proof}
由引理 \ref{sites-modules-lemma-extension-by-zero-exact}，我们已经知道 $j_{U!}$ 正合。
因此，只需证明
$j_{U!} :  \textit{Mod}(\mathcal{O}_U) \to \textit{Mod}(\mathcal{O})$
反映单射和满射。

\medskip\noindent
对 $\textit{Mod}(\mathcal{O}_U)$ 中的每个 $\mathcal{G}$，有该伴随的单位
$\mathcal{G} \to j_U^*j_{U!}\mathcal{G}$。我们断言此映射是层的单射。
事实上，考察引理 \ref{sites-modules-lemma-extension-by-zero} 的构造可知，此映射是如下规则的层化：
它将 $\mathcal{C}/U$ 的对象 $V/U$ 送到单射
$$
\mathcal{G}(V/U) \longrightarrow
\bigoplus\nolimits_{\varphi \in \Mor_\mathcal{C}(V, U)}
\mathcal{G}(V \xrightarrow{\varphi} U)
$$
，即嵌入到对应于结构态射 $V \to U$ 的直和项中。由于层化正合，断言成立。
略去若干细节。

\medskip\noindent
若 $\mathcal{G} \to \mathcal{G}'$ 是 $\mathcal{O}_U$-模映射，且
$j_{U!}\mathcal{G} \to j_{U!}\mathcal{G}'$ 是单射，则
$j_U^*j_{U!}\mathcal{G} \to j_U^*j_{U!}\mathcal{G}'$ 是单射
（限制是正合的），从而 $\mathcal{G} \to j_U^*j_{U!}\mathcal{G}'$
是单射，进而 $\mathcal{G} \to \mathcal{G}'$ 是单射。
所以 $j_{U!}$ 反映单射。

\medskip\noindent
设 $a : \mathcal{G} \to \mathcal{G}'$ 是 $\mathcal{O}_U$-模映射，且
$j_{U!}\mathcal{G} \to j_{U!}\mathcal{G}'$ 是满射。令 $\mathcal{H}$ 为
$a$ 的余核。由于 $j_{U!}$ 正合，有 $j_{U!}\mathcal{H} = 0$。
由上可知，映射 $\mathcal{H} \to j^*_U j_{U!}\mathcal{H}$ 是单射，
故 $\mathcal{H} = 0$，正合所需。
\end{proof}

\begin{lemma}
\label{sites-modules-lemma-relocalize}
设 $(\mathcal{C}, \mathcal{O})$ 是环站点，且 $f : V \to U$ 是
$\mathcal{C}$ 的态射。则存在环拓扑斯的交换图
$$
\xymatrix{
(\Sh(\mathcal{C}/V), \mathcal{O}_V)
\ar[rd]_{(j_V, j_V^\sharp)} \ar[rr]_{(j, j^\sharp)} & &
(\Sh(\mathcal{C}/U), \mathcal{O}_U)
\ar[ld]^{(j_U, j_U^\sharp)} \\
& (\Sh(\mathcal{C}), \mathcal{O}) &
}
$$
。这里 $(j, j^\sharp)$ 是与环站点
$(\mathcal{C}/V, \mathcal{O}_V)$ 的对象 $V/U$ 相伴的局部化态射。
\end{lemma}

\begin{proof}
只需检验 $j_V^\sharp = j^\sharp \circ j^{-1}(j_U^\sharp)$，
其余结论直接来自《站点》，引理 \ref{sites-lemma-relocalize} 和
公式 (\ref{sites-equation-relocalize})。略去该等式的验证。
\end{proof}

\begin{remark}
\label{sites-modules-remark-localize-shriek-equal}
在引理 \ref{sites-modules-lemma-extension-by-zero} 的情形中，图
$$
\xymatrix{
\textit{Mod}(\mathcal{O}_U) \ar[r]_{j_{U!}} \ar[d]_{forget} &
\textit{Mod}(\mathcal{O}_\mathcal{C}) \ar[d]^{forget} \\
\textit{Ab}(\mathcal{C}/U) \ar[r]^{j^{Ab}_{U!}} &
\textit{Ab}(\mathcal{C})
}
$$
交换。这直接来自该引理中函子 $j_{U!}$ 的显式描述。
\end{remark}

\begin{remark}
\label{sites-modules-remark-localize-presheaves}
局部化与模预层；见《站点》，注记 \ref{sites-remark-localize-presheaves}。
设 $\mathcal{C}$ 是范畴，$\mathcal{O}$ 是环预层，且 $U$ 是
$\mathcal{C}$ 的对象。严格地说，我们尚未对 $\mathcal{O}$-模预层定义
函子 $j_U^*$、$j_{U*}$ 和 $j_{U!}$。不过，可以把预层看成
$\mathcal{C}$ 上混沌拓扑的层（见《站点》，例 \ref{sites-example-indiscrete}）。
因此也得到函子
$$
j_U^* :
\textit{PMod}(\mathcal{O})
\longrightarrow
\textit{PMod}(\mathcal{O}_U)
$$
以及函子
$$
j_{U*}, j_{U!} :
\textit{PMod}(\mathcal{O}_U)
\longrightarrow
\textit{PMod}(\mathcal{O})
$$
，它们分别是 $j_U^*$ 的右伴随和左伴随。考察引理
\ref{sites-modules-lemma-extension-by-zero} 的证明可知，$j_{U!}\mathcal{G}$ 是预层
$$
V \longmapsto
\bigoplus\nolimits_{\varphi \in \Mor_\mathcal{C}(V, U)}
\mathcal{G}(V \xrightarrow{\varphi} U)
$$
。此外，函子 $j_{U!}$ 正合（对离散拓扑应用引理
\ref{sites-modules-lemma-extension-by-zero-exact}）。再者，若 $\mathcal{C}$ 实际上是站点，
且 $\mathcal{O}$ 实际上是环层，则图
$$
\xymatrix{
\textit{Mod}(\mathcal{O}_U) \ar[r]_{j_{U!}} \ar[d]_{forget} &
\textit{Mod}(\mathcal{O}) \\
\textit{PMod}(\mathcal{O}_U) \ar[r]^{j_{U!}} &
\textit{PMod}(\mathcal{O}) \ar[u]_{(\ )^\#}
}
$$
交换。
\end{remark}

\begin{lemma}
\label{sites-modules-lemma-restrict-back}
设 $\mathcal{C}$ 是站点，$U \in \Ob(\mathcal{C})$。假设对
$\mathcal{C}$ 中每个 $X$，至多有一个从 $X$ 到 $U$ 的态射。
设 $\mathcal{F}$ 是 $\mathcal{C}/U$ 上的阿贝尔层。则典范映射
$\mathcal{F} \to j_U^{-1}j_{U!}\mathcal{F}$ 和
$j_U^{-1}j_{U*}\mathcal{F} \to \mathcal{F}$ 都是同构。
\end{lemma}

\begin{proof}
这是引理 \ref{sites-modules-lemma-back-and-forth} 的特例，因为关于 $U$ 的假设等价于
局部化函子 $\mathcal{C}/U \to \mathcal{C}$ 是全忠实的。
\end{proof}









\section{环站点态射的局部化}
\label{sites-modules-section-localize-morphisms}

\noindent
本节是《站点》，第 \ref{sites-section-localize-morphisms} 节的类似版本。

\begin{lemma}
\label{sites-modules-lemma-localize-morphism-ringed-sites}
设
$(f, f^\sharp) :
(\mathcal{C}, \mathcal{O})
\longrightarrow
(\mathcal{D}, \mathcal{O}')$
是环站点态射，其中 $f$ 由连续函子
$u : \mathcal{D} \to \mathcal{C}$ 给出。设 $V$ 是 $\mathcal{D}$ 的对象，
并令 $U = u(V)$。则存在典范的环层映射 $(f')^\sharp$，使《站点》，
引理 \ref{sites-lemma-localize-morphism} 的图成为环拓扑斯的交换图
$$
\xymatrix{
(\Sh(\mathcal{C}/U), \mathcal{O}_U)
\ar[rr]_{(j_U, j_U^\sharp)} \ar[d]_{(f', (f')^\sharp)} & &
(\Sh(\mathcal{C}), \mathcal{O})
\ar[d]^{(f, f^\sharp)} \\
(\Sh(\mathcal{D}/V), \mathcal{O}'_V)
\ar[rr]^{(j_V, j_V^\sharp)} & &
(\Sh(\mathcal{D}), \mathcal{O}').
}
$$
。此外，在此情形中有 $f'_*j_U^{-1} = j_V^{-1}f_*$ 和
$f'_*j_U^* = j_V^*f_*$。
\end{lemma}

\begin{proof}
只需取 $(f')^\sharp$ 为
$$
(f')^{-1}\mathcal{O}'_V =
(f')^{-1}j_V^{-1}\mathcal{O}' =
j_U^{-1}f^{-1}\mathcal{O}' \xrightarrow{j_U^{-1}f^\sharp}
j_U^{-1}\mathcal{O} = \mathcal{O}_U
$$
，其余结论都来自《站点》，引理 \ref{sites-lemma-localize-morphism}。
（注意，若 $j$ 是局部化态射，则在模层上 $j^{-1} = j^*$，
因而第一个函子等式蕴含第二个。）
\end{proof}

\begin{lemma}
\label{sites-modules-lemma-relocalize-morphism-ringed-sites}
设
$(f, f^\sharp) :
(\mathcal{C}, \mathcal{O})
\longrightarrow
(\mathcal{D}, \mathcal{O}')$
是环站点态射，其中 $f$ 由连续函子
$u : \mathcal{D} \to \mathcal{C}$ 给出。设
$V \in \Ob(\mathcal{D})$、$U \in \Ob(\mathcal{C})$，且
$c : U \to u(V)$ 是 $\mathcal{C}$ 的态射。则存在环拓扑斯的交换图
$$
\xymatrix{
(\Sh(\mathcal{C}/U), \mathcal{O}_U)
\ar[rr]_{(j_U, j_U^\sharp)} \ar[d]_{(f_c, f_c^\sharp)} & &
(\Sh(\mathcal{C}), \mathcal{O}) \ar[d]^{(f, f^\sharp)} \\
(\Sh(\mathcal{D}/V), \mathcal{O}'_V)
\ar[rr]^{(j_V, j_V^\sharp)} & &
(\Sh(\mathcal{D}), \mathcal{O}').
}
$$
。态射 $(f_c, f_c^\sharp)$ 等于如下态射
$$
(f', (f')^\sharp) :
(\Sh(\mathcal{C}/u(V)), \mathcal{O}_{u(V)})
\longrightarrow
(\Sh(\mathcal{D}/V), \mathcal{O}'_V)
$$
（见引理 \ref{sites-modules-lemma-localize-morphism-ringed-sites}）与如下态射
$$
(j, j^\sharp) :
(\Sh(\mathcal{C}/U), \mathcal{O}_U)
\to
(\Sh(\mathcal{C}/u(V)), \mathcal{O}_{u(V)})
$$
（见引理 \ref{sites-modules-lemma-relocalize}）的复合。给定任意态射
$b : V' \to V$、$a : U' \to U$ 和 $c' : U' \to u(V')$，使得
$$
\xymatrix{
U' \ar[r]_-{c'} \ar[d]_a & u(V') \ar[d]^{u(b)} \\
U \ar[r]^-c & u(V)
}
$$
交换，则下列环拓扑斯之图
$$
\xymatrix{
(\Sh(\mathcal{C}/U'), \mathcal{O}_{U'})
\ar[rr]_{(j_{U'/U}, j_{U'/U}^\sharp)} \ar[d]_{(f_{c'}, f_{c'}^\sharp)} & &
(\Sh(\mathcal{C}/U), \mathcal{O}_U)
\ar[d]^{(f_c, f_c^\sharp)} \\
(\Sh(\mathcal{D}/V'), \mathcal{O}'_{V'})
\ar[rr]^{(j_{V'/V}, j_{V'/V}^\sharp)} & &
(\Sh(\mathcal{D}/V), \mathcal{O}'_{V'})
}
$$
交换。
\end{lemma}

\begin{proof}
在拓扑斯态射的层面，这就是《站点》，引理
\ref{sites-lemma-relocalize-morphism}。为检验这些图作为环拓扑斯态射之图交换，
使用引理 \ref{sites-modules-lemma-relocalize} 和
\ref{sites-modules-lemma-localize-morphism-ringed-sites}，完全仿照《站点》，引理
\ref{sites-lemma-relocalize-morphism} 的证明即可。
\end{proof}
















\section{环拓扑斯的局部化}
\label{sites-modules-section-localize-ringed-topoi}

\noindent
本节是在环拓扑斯的设定下，对《站点》，第
\ref{sites-section-localize-topoi} 节所作的类似论述。

\begin{lemma}
\label{sites-modules-lemma-localize-ringed-topos}
设 $(\Sh(\mathcal{C}), \mathcal{O})$ 是环拓扑斯，且
$\mathcal{F} \in \Sh(\mathcal{C})$ 是层。对 $\mathcal{C}$ 上的层
$\mathcal{H}$，以 $\mathcal{H}_\mathcal{F}$ 表示层
$\mathcal{H} \times \mathcal{F}$，并将它视为范畴
$\Sh(\mathcal{C})/\mathcal{F}$ 的对象。则
$(\Sh(\mathcal{C})/\mathcal{F}, \mathcal{O}_\mathcal{F})$
是环拓扑斯，且存在环拓扑斯的典范态射
$$
(j_\mathcal{F}, j_\mathcal{F}^\sharp) :
(\Sh(\mathcal{C})/\mathcal{F}, \mathcal{O}_\mathcal{F})
\longrightarrow
(\Sh(\mathcal{C}), \mathcal{O})
$$
，它是第 \ref{sites-modules-section-localize} 节意义下的局部化，并且
\begin{enumerate}
\item 函子 $j_\mathcal{F}^{-1}$ 是函子
$\mathcal{H} \mapsto \mathcal{H}_\mathcal{F}$,
\item 函子 $j_\mathcal{F}^*$ 是函子
$\mathcal{H} \mapsto \mathcal{H}_\mathcal{F}$,
\item 集合层上的函子 $j_{\mathcal{F}!}$ 是遗忘函子
$\mathcal{G}/\mathcal{F} \mapsto \mathcal{G}$，
\item 模层上的函子 $j_{\mathcal{F}!}$ 将 $\mathcal{O}_\mathcal{F}$-模
$\varphi : \mathcal{G} \to \mathcal{F}$ 送到如下预层之层化所得的
$\mathcal{O}$-模：
$$
V \longmapsto
\bigoplus\nolimits_{s \in \mathcal{F}(V)}
\{\sigma \in \mathcal{G}(V) \mid \varphi(\sigma) = s \}
$$
其中 $V \in \Ob(\mathcal{C})$。
\end{enumerate}
\end{lemma}

\begin{proof}
由《站点》，引理 \ref{sites-lemma-localize-topos} 可知，
$\Sh(\mathcal{C})/\mathcal{F}$ 是拓扑斯，且 (1) 和 (3) 成立。
特别地，这说明
$j_\mathcal{F}^{-1}\mathcal{O} = \mathcal{O}_\mathcal{F}$
，并说明 $\mathcal{O}_\mathcal{F}$ 是环层。因此可取映射
$j_\mathcal{F}^\sharp$ 为恒等映射，特别地，(2) 成立。
此外，《站点》，引理 \ref{sites-lemma-localize-topos} 的证明说明，
可以假设 $\mathcal{C}$ 是具有所有有限极限和次典范拓扑的站点，且对
$\mathcal{C}$ 的某个对象 $U$ 有 $\mathcal{F} = h_U$。
此时 (4) 来自引理 \ref{sites-modules-lemma-extension-by-zero} 中对
$j_{\mathcal{F}!}$ 的描述。另一种方法是直接证明 (4) 中描述的函子是
$j_\mathcal{F}^*$ 的左伴随。
\end{proof}

\begin{definition}
\label{sites-modules-definition-localize-ringed-topos}
设 $(\Sh(\mathcal{C}), \mathcal{O})$ 是环拓扑斯，且
$\mathcal{F} \in \Sh(\mathcal{C})$。
\begin{enumerate}
\item 环拓扑斯
$(\Sh(\mathcal{C})/\mathcal{F}, \mathcal{O}_\mathcal{F})$
称为{\it 环拓扑斯 $(\Sh(\mathcal{C}), \mathcal{O})$ 在
$\mathcal{F}$ 处的局部化}。
\item 引理 \ref{sites-modules-lemma-localize-ringed-topos} 的环拓扑斯态射
$(j_\mathcal{F}, j_\mathcal{F}^\sharp) :
(\Sh(\mathcal{C})/\mathcal{F}, \mathcal{O}_\mathcal{F})
\to
(\Sh(\mathcal{C}), \mathcal{O})$
称为{\it 局部化态射}。
\end{enumerate}
\end{definition}

\noindent
我们延续《站点》一章中确立的做法：每当有意义时，检验拓扑斯上的局部化构造
与站点上的局部化构造相容。

\begin{lemma}
\label{sites-modules-lemma-localize-compare}
设 $(\Sh(\mathcal{C}), \mathcal{O})$ 和
$\mathcal{F} \in \Sh(\mathcal{C})$ 如引理
\ref{sites-modules-lemma-localize-ringed-topos} 中所述。若对 $\mathcal{C}$ 的某个对象
$U$ 有 $\mathcal{F} = h_U^\#$，则经由《站点》，引理
\ref{sites-lemma-essential-image-j-shriek} 的等同
$\Sh(\mathcal{C}/U) = \Sh(\mathcal{C})/h_U^\#$，有
\begin{enumerate}
\item 典范地有 $\mathcal{O}_U = \mathcal{O}_\mathcal{F}$；
\item 在这些等同下，
$(j_\mathcal{F}, j_\mathcal{F}^\sharp) = (j_U, j_U^\sharp)$。
\end{enumerate}
\end{lemma}

\begin{proof}
关于底层拓扑斯的断言就是《站点》，引理
\ref{sites-lemma-localize-compare}。注意，$\mathcal{O}_U$ 是
$\mathcal{O}$ 的限制；由《站点》，引理
\ref{sites-lemma-compute-j-shriek-restrict}，在《站点》，引理
\ref{sites-lemma-essential-image-j-shriek} 的等价下，它对应于
$\mathcal{O} \times h_U^\#$。由 $\mathcal{O}_\mathcal{F}$ 的定义即得 (1)。
余下只须证明在此等同下 $j_\mathcal{F}^\sharp = j_U^\sharp$。验证从略。
\end{proof}

\noindent
局部化在以下两种意义下具有函子性：可以对一个局部化再次“局部化”（见引理
\ref{sites-modules-lemma-relocalize-ringed-topos}）；也可以给定环拓扑斯态射，在上方拓扑斯中
对下方拓扑斯上某层的逆像作局部化，从而得到环拓扑斯的交换图（见引理
\ref{sites-modules-lemma-localize-morphism-ringed-topoi}）。

\begin{lemma}
\label{sites-modules-lemma-relocalize-ringed-topos}
设 $(\Sh(\mathcal{C}), \mathcal{O})$ 是环拓扑斯。若
$s : \mathcal{G} \to \mathcal{F}$ 是 $\mathcal{C}$ 上的层态射，
则存在环拓扑斯态射的自然交换图
$$
\xymatrix{
(\Sh(\mathcal{C})/\mathcal{G}, \mathcal{O}_\mathcal{G})
\ar[rd]_{(j_\mathcal{G}, j_\mathcal{G}^\sharp)} \ar[rr]_{(j, j^\sharp)} & &
(\Sh(\mathcal{C})/\mathcal{F}, \mathcal{O}_\mathcal{F})
\ar[ld]^{(j_\mathcal{F}, j_\mathcal{F}^\sharp)} \\
& (\Sh(\mathcal{C}), \mathcal{O}) &
}
$$
，其中 $(j, j^\sharp)$ 是环拓扑斯
$(\Sh(\mathcal{C})/\mathcal{F}, \mathcal{O}_\mathcal{F})$
在对象 $\mathcal{G}/\mathcal{F}$ 处的局部化态射。
\end{lemma}

\begin{proof}
除等式
$j_\mathcal{G}^\sharp = j^\sharp \circ j^{-1}(j_\mathcal{F}^\sharp)$
外，所有断言都来自《站点》，引理 \ref{sites-lemma-relocalize-topos}。
该等式的验证从略。
\end{proof}

\begin{lemma}
\label{sites-modules-lemma-relocalize-compare}
设 $(\Sh(\mathcal{C}), \mathcal{O})$ 和
$s : \mathcal{G} \to \mathcal{F}$ 如引理
\ref{sites-modules-lemma-relocalize-ringed-topos} 中所述。若存在 $\mathcal{C}$ 的态射
$f : V \to U$，使得 $\mathcal{G} = h_V^\#$、
$\mathcal{F} = h_U^\#$，且 $s$ 由 $f$ 诱导，则引理
\ref{sites-modules-lemma-relocalize} 与引理
\ref{sites-modules-lemma-relocalize-ringed-topos} 的图经由等同
$(j_\mathcal{F}, j_\mathcal{F}^\sharp) = (j_U, j_U^\sharp)$
和
$(j_\mathcal{G}, j_\mathcal{G}^\sharp) = (j_V, j_V^\sharp)$
（见引理 \ref{sites-modules-lemma-localize-compare}）彼此一致。
\end{lemma}

\begin{proof}
除两个映射 $j^\sharp$ 一致这一断言外，所有断言都来自《站点》，引理
\ref{sites-lemma-relocalize-compare}。在两种情形中，映射 $j^\sharp$
都只是恒等映射，故该断言成立。略去若干细节。
\end{proof}










\section{环拓扑斯态射的局部化}
\label{sites-modules-section-localize-morphisms-ringed-topoi}

\noindent
本节是《站点》，第 \ref{sites-section-localize-morphisms-topoi} 节的类似版本。

\begin{lemma}
\label{sites-modules-lemma-localize-morphism-ringed-topoi}
设
$$
f :
(\Sh(\mathcal{C}), \mathcal{O})
\longrightarrow
(\Sh(\mathcal{D}), \mathcal{O}')
$$
是环拓扑斯态射。设 $\mathcal{G}$ 是 $\mathcal{D}$ 上的层，并令
$\mathcal{F} = f^{-1}\mathcal{G}$。则存在环拓扑斯的交换图
$$
\xymatrix{
(\Sh(\mathcal{C})/\mathcal{F}, \mathcal{O}_\mathcal{F})
\ar[rr]_{(j_\mathcal{F}, j_\mathcal{F}^\sharp)}
\ar[d]_{(f', (f')^\sharp)} & &
(\Sh(\mathcal{C}), \mathcal{O}) \ar[d]^{(f, f^\sharp)} \\
(\Sh(\mathcal{D})/\mathcal{G}, \mathcal{O}'_\mathcal{G})
\ar[rr]^{(j_\mathcal{G}, j_\mathcal{G}^\sharp)} & &
(\Sh(\mathcal{D}), \mathcal{O}')
}
$$
。有 $f'_*j_\mathcal{F}^{-1} = j_\mathcal{G}^{-1}f_*$ 和
$f'_*j_\mathcal{F}^* = j_\mathcal{G}^*f_*$。此外，态射 $f'$ 由规则
$$
(f')^{-1}(\mathcal{H} \xrightarrow{\varphi} \mathcal{G})
=
(f^{-1}\mathcal{H} \xrightarrow{f^{-1}\varphi} \mathcal{F}).
$$
\end{lemma}

\begin{proof}
由《站点》，引理 \ref{sites-lemma-localize-morphism-topoi}，得到底层拓扑斯之图、
等式 $f'_*j_\mathcal{F}^{-1} = j_\mathcal{G}^{-1}f_*$，以及
对 $(f')^{-1}$ 的描述。为定义 $(f')^\sharp$，使用映射
$$
(f')^\sharp :
\mathcal{O}'_\mathcal{G} =
j_\mathcal{G}^{-1} \mathcal{O}'
\xrightarrow{j_\mathcal{G}^{-1}f^\sharp}
j_\mathcal{G}^{-1} f_*\mathcal{O} =
f'_* j_\mathcal{F}^{-1}\mathcal{O} =
f'_* \mathcal{O}_\mathcal{F}
$$
，或等价地使用映射
$$
(f')^\sharp :
(f')^{-1}\mathcal{O}'_\mathcal{G} =
(f')^{-1}j_\mathcal{G}^{-1} \mathcal{O}' =
j_\mathcal{F}^{-1}f^{-1}\mathcal{O}'
\xrightarrow{j_\mathcal{F}^{-1}f^\sharp}
j_\mathcal{F}^{-1} \mathcal{O} =
\mathcal{O}_\mathcal{F}.
$$
。略去这两个映射确实互为伴随映射的验证。$(f')^\sharp$ 的第二种构造说明，
该图在环拓扑斯的 $2$-范畴中交换（因为映射
$j_\mathcal{F}^\sharp$ 和 $j_\mathcal{G}^\sharp$ 都是恒等映射）。
最后，等式 $f'_*j_\mathcal{F}^* = j_\mathcal{G}^*f_*$ 来自等式
$f'_*j_\mathcal{F}^{-1} = j_\mathcal{G}^{-1}f_*$
以及模层与集合层的拉回一致这一事实；见引理
\ref{sites-modules-lemma-localize-ringed-topos}。
\end{proof}

\begin{lemma}
\label{sites-modules-lemma-localize-morphism-compare}
设
$$
f :
(\Sh(\mathcal{C}), \mathcal{O})
\longrightarrow
(\Sh(\mathcal{D}), \mathcal{O}')
$$
是环拓扑斯态射。设 $\mathcal{G}$ 是 $\mathcal{D}$ 上的层，并令
$\mathcal{F} = f^{-1}\mathcal{G}$。若 $f$ 由连续函子
$u : \mathcal{D} \to \mathcal{C}$ 给出，且 $\mathcal{G} = h_V^\#$，
则引理 \ref{sites-modules-lemma-localize-morphism-ringed-sites} 与引理
\ref{sites-modules-lemma-localize-morphism-ringed-topoi} 的交换图经由引理
\ref{sites-modules-lemma-localize-compare} 的等同彼此一致。
\end{lemma}

\begin{proof}
在拓扑斯态射的层面，这就是《站点》，引理
\ref{sites-lemma-localize-morphism-compare}。这在环拓扑斯态射的层面也成立，
因为引理 \ref{sites-modules-lemma-localize-morphism-ringed-sites} 与引理
\ref{sites-modules-lemma-localize-morphism-ringed-topoi} 的证明中定义
$(f')^\sharp$ 的公式一致。
\end{proof}

\begin{lemma}
\label{sites-modules-lemma-relocalize-morphism-ringed-topoi}
设
$(f, f^\sharp) :
(\Sh(\mathcal{C}), \mathcal{O})
\to
(\Sh(\mathcal{D}), \mathcal{O}')$
是环拓扑斯态射。设 $\mathcal{G}$ 是 $\mathcal{D}$ 上的层，
$\mathcal{F}$ 是 $\mathcal{C}$ 上的层，且
$s : \mathcal{F} \to f^{-1}\mathcal{G}$ 是层态射。则存在环拓扑斯的交换图
$$
\xymatrix{
(\Sh(\mathcal{C})/\mathcal{F}, \mathcal{O}_\mathcal{F})
\ar[rr]_{(j_\mathcal{F}, j_\mathcal{F}^\sharp)}
\ar[d]_{(f_c, f_c^\sharp)} & &
(\Sh(\mathcal{C}), \mathcal{O})
\ar[d]^{(f, f^\sharp)} \\
(\Sh(\mathcal{D})/\mathcal{G}, \mathcal{O}'_\mathcal{G})
\ar[rr]^{(j_\mathcal{G}, j_\mathcal{G}^\sharp)} & &
(\Sh(\mathcal{D}), \mathcal{O}').
}
$$
。态射 $(f_s, f_s^\sharp)$ 等于如下态射
$$
(f', (f')^\sharp) :
(\Sh(\mathcal{C})/f^{-1}\mathcal{G}, \mathcal{O}_{f^{-1}\mathcal{G}})
\longrightarrow
(\Sh(\mathcal{D})/{\mathcal{G}}, \mathcal{O}'_\mathcal{G})
$$
（见引理 \ref{sites-modules-lemma-localize-morphism-ringed-topoi}）与如下态射
$$
(j, j^\sharp) :
(\Sh(\mathcal{C})/\mathcal{F}, \mathcal{O}_\mathcal{F})
\to
(\Sh(\mathcal{C})/f^{-1}\mathcal{G}, \mathcal{O}_{f^{-1}\mathcal{G}})
$$
（见引理 \ref{sites-modules-lemma-relocalize-ringed-topos}）的复合。给定任意态射
$b : \mathcal{G}' \to \mathcal{G}$、$a : \mathcal{F}' \to \mathcal{F}$ 和
$s' : \mathcal{F}' \to f^{-1}\mathcal{G}'$，使得
$$
\xymatrix{
\mathcal{F}' \ar[r]_-{s'} \ar[d]_a &
f^{-1}\mathcal{G}' \ar[d]^{f^{-1}b} \\
\mathcal{F} \ar[r]^-s &
f^{-1}\mathcal{G}
}
$$
交换，则下列环拓扑斯之图
$$
\xymatrix{
(\Sh(\mathcal{C})/\mathcal{F}', \mathcal{O}_{\mathcal{F}'})
\ar[rr]_{(j_{\mathcal{F}'/\mathcal{F}}, j_{\mathcal{F}'/\mathcal{F}}^\sharp)}
\ar[d]_{(f_{s'}, f_{s'}^\sharp)} & &
(\Sh(\mathcal{C})/\mathcal{F}, \mathcal{O}_\mathcal{F})
\ar[d]^{(f_s, f_s^\sharp)} \\
(\Sh(\mathcal{D})/\mathcal{G}', \mathcal{O}'_{\mathcal{G}'})
\ar[rr]^{(j_{\mathcal{G}'/\mathcal{G}}, j_{\mathcal{G}'/\mathcal{G}}^\sharp)}
& &
(\Sh(\mathcal{D})/\mathcal{G}, \mathcal{O}'_{\mathcal{G}'})
}
$$
交换。
\end{lemma}

\begin{proof}
在拓扑斯态射的层面，这就是《站点》，引理
\ref{sites-lemma-relocalize-morphism-topoi}。为检验这些图作为环拓扑斯态射之图
交换，使用引理 \ref{sites-modules-lemma-relocalize-ringed-topos} 和
\ref{sites-modules-lemma-localize-morphism-ringed-topoi} 的交换图即可。
\end{proof}

\begin{lemma}
\label{sites-modules-lemma-relocalize-morphism-compare}
设
$(f, f^\sharp) :
(\Sh(\mathcal{C}), \mathcal{O})
\to
(\Sh(\mathcal{D}), \mathcal{O}')$,
$s : \mathcal{F} \to f^{-1}\mathcal{G}$ 如引理
\ref{sites-modules-lemma-relocalize-morphism-ringed-topoi} 中所述。若 $f$ 由连续函子
$u : \mathcal{D} \to \mathcal{C}$
给出，且 $\mathcal{G} = h_V^\#$、$\mathcal{F} = h_U^\#$，而 $s$
来自态射 $c : U \to u(V)$，则引理
\ref{sites-modules-lemma-relocalize-morphism-ringed-sites} 与引理
\ref{sites-modules-lemma-relocalize-morphism-ringed-topoi} 的交换图经由引理
\ref{sites-modules-lemma-localize-compare} 的等同彼此一致。
\end{lemma}

\begin{proof}
这是引理 \ref{sites-modules-lemma-relocalize-compare} 和
\ref{sites-modules-lemma-localize-morphism-compare} 的形式推论。
\end{proof}

















\section{模的局部类型}
\label{sites-modules-section-local}

\noindent
按照第 \ref{sites-modules-section-intrinsic} 节说明的一般策略，我们先定义环站点上模层的
局部类型，随后立即证明这些类型是内蕴的，因而对环拓扑斯上的模层也有意义。

\begin{definition}
\label{sites-modules-definition-site-local}
设 $(\mathcal{C}, \mathcal{O})$ 是环站点，且 $\mathcal{F}$ 是
$\mathcal{O}$-模层。我们将自由使用定义 \ref{sites-modules-definition-global} 中定义的概念。
\begin{enumerate}
\item 称 $\mathcal{F}$ 是{\it 局部自由的}，如果对 $\mathcal{C}$ 的每个对象
$U$，存在 $\mathcal{C}$ 的覆盖 $\{U_i \to U\}_{i \in I}$，使每个限制
$\mathcal{F}|_{\mathcal{C}/U_i}$ 都是自由 $\mathcal{O}_{U_i}$-模。
\item 称 $\mathcal{F}$ 是{\it 有限局部自由的}，如果对 $\mathcal{C}$ 的每个对象
$U$，存在 $\mathcal{C}$ 的覆盖 $\{U_i \to U\}_{i \in I}$，使每个限制
$\mathcal{F}|_{\mathcal{C}/U_i}$ 都是有限自由 $\mathcal{O}_{U_i}$-模。
\item 称 $\mathcal{F}$ 是{\it 局部由截面生成的}，如果对 $\mathcal{C}$ 的每个
对象 $U$，存在覆盖 $\{U_i \to U\}_{i \in I}$，使每个限制
$\mathcal{F}|_{\mathcal{C}/U_i}$ 都是由整体截面生成的
$\mathcal{O}_{U_i}$-模。
\item 给定 $r \geq 0$，称 $\mathcal{F}$ 是{\it 局部由 $r$ 个截面生成的}，
如果对 $\mathcal{C}$ 的每个对象 $U$，存在覆盖
$\{U_i \to U\}_{i \in I}$，使每个限制
$\mathcal{F}|_{\mathcal{C}/U_i}$ 都是由 $r$ 个整体截面生成的
$\mathcal{O}_{U_i}$-模。
\item 称 $\mathcal{F}$ 是{\it 有限型的}，如果对 $\mathcal{C}$ 的每个对象
$U$，存在覆盖 $\{U_i \to U\}_{i \in I}$，使每个限制
$\mathcal{F}|_{\mathcal{C}/U_i}$ 都是由有限多个整体截面生成的
$\mathcal{O}_{U_i}$-模。
\item 称 $\mathcal{F}$ 是{\it 拟凝聚的}，如果对 $\mathcal{C}$ 的每个对象
$U$，存在覆盖 $\{U_i \to U\}_{i \in I}$，使每个限制
$\mathcal{F}|_{\mathcal{C}/U_i}$ 都是具有整体表示的
$\mathcal{O}_{U_i}$-模。
\item 称 $\mathcal{F}$ 是{\it 有限表示的}，如果对 $\mathcal{C}$ 的每个对象
$U$，存在覆盖 $\{U_i \to U\}_{i \in I}$，使每个限制
$\mathcal{F}|_{\mathcal{C}/U_i}$ 都是具有有限整体表示的
$\mathcal{O}_{U_i}$-模。
\item 称 $\mathcal{F}$ 是{\it 凝聚的}，当且仅当 $\mathcal{F}$ 是有限型的，
并且对 $\mathcal{C}$ 的每个对象 $U$ 及任意
$s_1, \ldots, s_n \in \mathcal{F}(U)$，映射
$\bigoplus_{i = 1, \ldots, n} \mathcal{O}_U \to \mathcal{F}|_U$
的核在 $(\mathcal{C}/U, \mathcal{O}_U)$ 上是有限型的。
\end{enumerate}
\end{definition}

\begin{lemma}
\label{sites-modules-lemma-special-locally-free}
定义 \ref{sites-modules-definition-site-local} 的性质 (1)--(8) 中任意一个都是内蕴的
（见第 \ref{sites-modules-section-intrinsic} 节的讨论）。
\end{lemma}

\begin{proof}
设 $\mathcal{C}$、$\mathcal{D}$ 是站点，且
$u : \mathcal{C} \to \mathcal{D}$ 是特殊余连续函子。设
$\mathcal{O}$ 是 $\mathcal{C}$ 上的环层，$\mathcal{F}$ 是
$\mathcal{C}$ 上的 $\mathcal{O}$-模层。设
$g : \Sh(\mathcal{C}) \to \Sh(\mathcal{D})$ 是与 $u$ 相伴的拓扑斯等价。
令 $\mathcal{O}' = g_*\mathcal{O}$，并令
$g^\sharp : \mathcal{O}' \to g_*\mathcal{O}$ 为恒等映射。最后，令
$\mathcal{F}' = g_*\mathcal{F}$。设 $\mathcal{P}_l$ 是定义
\ref{sites-modules-definition-site-local} 所列性质 (1)--(7) 之一（凝聚情形将在证明末尾讨论）。
以 $\mathcal{P}_g$ 表示定义 \ref{sites-modules-definition-global} 中相应的性质。
我们已经看到 $\mathcal{P}_g$ 是内蕴的。须证明
$\mathcal{P}_l(\mathcal{C}, \mathcal{O}, \mathcal{F})$
成立，当且仅当
$\mathcal{P}_l(\mathcal{D}, \mathcal{O}', \mathcal{F}')$
成立。

\medskip\noindent
假设 $\mathcal{F}$ 具有性质 $\mathcal{P}_l$。设 $V$ 是 $\mathcal{D}$ 的对象。
特殊余连续函子的性质之一是：站点 $\mathcal{D}$ 中存在覆盖
$\{u(U_i) \to V\}_{i \in I}$。由假设，对每个 $i$，在 $\mathcal{C}$ 中
存在覆盖 $\{U_{ij} \to U_i\}_{j \in J_i}$，使每个限制
$\mathcal{F}|_{U_{ij}}$ 都具有性质 $\mathcal{P}_g$。由《站点》，引理
\ref{sites-lemma-localize-special-cocontinuous}，有环拓扑斯的交换图
$$
\xymatrix{
(\Sh(\mathcal{C}/U_{ij}), \mathcal{O}_{U_{ij}}) \ar[r] \ar[d] &
(\Sh(\mathcal{C}), \mathcal{O}) \ar[d] \\
(\Sh(\mathcal{D}/u(U_{ij})), \mathcal{O}'_{u(U_{ij})}) \ar[r] &
(\Sh(\mathcal{D}), \mathcal{O}')
}
$$
，其中竖直箭头都是等价。因此 $\mathcal{F}'|_{u(U_{ij})}$ 也具有性质
$\mathcal{P}_g$。此外，$\{u(U_{ij}) \to V\}_{i \in I, j \in J_i}$
是站点 $\mathcal{D}$ 的覆盖，故 $\mathcal{F}'$ 具有性质
$\mathcal{P}_l$。

\medskip\noindent
假设 $\mathcal{F}'$ 具有性质 $\mathcal{P}_l$。设 $U$ 是
$\mathcal{C}$ 的对象。由假设，存在覆盖
$\{V_i \to u(U)\}_{i \in I}$，使 $\mathcal{F}'|_{V_i}$ 具有性质
$\mathcal{P}_g$。由于 $u$ 余连续，可以用族
$\{u(U_j) \to u(U)\}_{j \in J}$ 加细此覆盖，其中
$\{U_j \to U\}_{j \in J}$ 是 $\mathcal{C}$ 中的覆盖。设此加细由
$\alpha : J \to I$ 和 $u(U_j) \to V_{\alpha(j)}$ 给出。
限制具有传递性，即
$(\mathcal{F}'|_{V_{\alpha(j)}})|_{u(U_j)} = \mathcal{F}'|_{u(U_j)}$.
因此由引理 \ref{sites-modules-lemma-global-pullback} 可知，
$\mathcal{F}'|_{u(U_j)}$ 具有性质 $\mathcal{P}_g$。于是图
$$
\xymatrix{
(\Sh(\mathcal{C}/U_j), \mathcal{O}_{U_j}) \ar[r] \ar[d] &
(\Sh(\mathcal{C}), \mathcal{O}) \ar[d] \\
(\Sh(\mathcal{D}/u(U_j)), \mathcal{O}'_{u(U_j)})
\ar[r] &
(\Sh(\mathcal{D}), \mathcal{O}')
}
$$
（其中竖直箭头都是等价）说明 $\mathcal{F}|_{U_j}$ 也具有性质
$\mathcal{P}_g$。故 $\mathcal{F}$ 具有性质 $\mathcal{P}_l$，正合所需。

\medskip\noindent
最后证明 $\mathcal{P}_l = coherent$ 的情形\footnote{这里的操作略显繁琐，
我们建议读者跳过证明的这一部分。}。假设 $\mathcal{F}$ 凝聚。
于是 $\mathcal{F}$ 是有限型的，由证明的第一部分，$\mathcal{F}'$ 也是有限型的。
设 $V$ 是 $\mathcal{D}$ 的对象，且
$s_1, \ldots, s_n \in \mathcal{F}'(V)$。须证明映射
$\bigoplus_{j = 1, \ldots, n} \mathcal{O}_V \to \mathcal{F}'|_V$ 的核
$\mathcal{K}'$
在 $\mathcal{D}/V$ 上是有限型的。这意味着须证明：对任意 $V'/V$，
存在覆盖 $\{V'_i \to V'\}$，使 $\mathcal{F}'|_{V'_i}$ 由有限多个截面生成。
以 $V'$ 替换 $V$（并把截面 $s_j$ 限制到 $V'$），可约化到
$V' = V$ 的情形。由于 $u$ 是特殊余连续函子，站点 $\mathcal{D}$ 中
存在覆盖 $\{u(U_i) \to V\}_{i \in I}$。利用拓扑斯同构
$\Sh(\mathcal{C}/U_i) = \Sh(\mathcal{D}/u(U_i))$
可知，$\mathcal{K}'|_{u(U_i)}$ 对应于映射
$\bigoplus_{j = 1, \ldots, n} \mathcal{O}_{U_i} \to \mathcal{F}|_{U_i}$.
的核 $\mathcal{K}_i$。由于 $\mathcal{F}$ 凝聚，$\mathcal{K}_i$ 是有限型的。
再次由证明的第一部分可知，$\mathcal{K}|_{u(U_i)}$ 是有限型的。因此存在覆盖
$\{V_{il} \to u(U_i)\}$，使 $\mathcal{K}|_{V_{il}}$ 由有限多个整体截面生成。
由于 $\{V_{il} \to V\}$ 是 $\mathcal{D}$ 的覆盖，得到
$\mathcal{K}$ 是有限型的，正合所需。

\medskip\noindent
假设 $\mathcal{F}'$ 凝聚。于是 $\mathcal{F}'$ 是有限型的，由证明的第一部分，
$\mathcal{F}$ 也是有限型的。设 $U$ 是 $\mathcal{C}$ 的对象，且
$s_1, \ldots, s_n \in \mathcal{F}(U)$。须证明映射
$\bigoplus_{j = 1, \ldots, n} \mathcal{O}_U \to \mathcal{F}|_U$ 的核
$\mathcal{K}$
在 $\mathcal{C}/U$ 上是有限型的。利用拓扑斯同构
$\Sh(\mathcal{C}/U) = \Sh(\mathcal{D}/u(U))$
可知，$\mathcal{K}|_U$ 对应于映射
$\bigoplus_{j = 1, \ldots, n} \mathcal{O}_{u(U)} \to \mathcal{F}'|_{u(U)}$
的核 $\mathcal{K}'$。
由于 $\mathcal{F}'$ 凝聚，$\mathcal{K}'$ 是有限型的。再次由证明的第一部分，
得到 $\mathcal{K}$ 是有限型的。
\end{proof}

\noindent
因此，从现在起可以不指定站点，直接使用定义 \ref{sites-modules-definition-site-local}
中定义的 $\mathcal{O}$-模性质。

\begin{lemma}
\label{sites-modules-lemma-local-final-object}
设 $(\Sh(\mathcal{C}), \mathcal{O})$ 是环拓扑斯，且 $\mathcal{F}$ 是
$\mathcal{O}$-模。假设站点 $\mathcal{C}$ 有终对象 $X$。则
\begin{enumerate}
\item 下列条件等价：
\begin{enumerate}
\item $\mathcal{F}$ 局部自由；
\item $\mathcal{C}$ 中存在覆盖 $\{X_i \to X\}$，使每个限制
$\mathcal{F}|_{\mathcal{C}/X_i}$ 都是局部自由 $\mathcal{O}_{X_i}$-模；
\item $\mathcal{C}$ 中存在覆盖 $\{X_i \to X\}$，使每个限制
$\mathcal{F}|_{\mathcal{C}/X_i}$ 都是自由 $\mathcal{O}_{X_i}$-模。
\end{enumerate}
\item 下列条件等价：
\begin{enumerate}
\item $\mathcal{F}$ 有限局部自由；
\item $\mathcal{C}$ 中存在覆盖 $\{X_i \to X\}$，使每个限制
$\mathcal{F}|_{\mathcal{C}/X_i}$ 都是有限局部自由
$\mathcal{O}_{X_i}$-模；
\item $\mathcal{C}$ 中存在覆盖 $\{X_i \to X\}$，使每个限制
$\mathcal{F}|_{\mathcal{C}/X_i}$ 都是有限自由 $\mathcal{O}_{X_i}$-模。
\end{enumerate}
\item 下列条件等价：
\begin{enumerate}
\item $\mathcal{F}$ 局部由截面生成；
\item $\mathcal{C}$ 中存在覆盖 $\{X_i \to X\}$，使每个限制
$\mathcal{F}|_{\mathcal{C}/X_i}$ 都是局部由截面生成的
$\mathcal{O}_{X_i}$-模；
\item $\mathcal{C}$ 中存在覆盖 $\{X_i \to X\}$，使每个限制
$\mathcal{F}|_{\mathcal{C}/X_i}$ 都是由整体截面生成的
$\mathcal{O}_{X_i}$-模。
\end{enumerate}
\item 给定 $r \geq 0$，下列条件等价：
\begin{enumerate}
\item $\mathcal{F}$ 局部由 $r$ 个截面生成；
\item $\mathcal{C}$ 中存在覆盖 $\{X_i \to X\}$，使每个限制
$\mathcal{F}|_{\mathcal{C}/X_i}$ 都是局部由 $r$ 个截面生成的
$\mathcal{O}_{X_i}$-模；
\item $\mathcal{C}$ 中存在覆盖 $\{X_i \to X\}$，使每个限制
$\mathcal{F}|_{\mathcal{C}/X_i}$ 都是由 $r$ 个整体截面生成的
$\mathcal{O}_{X_i}$-模。
\end{enumerate}
\item 下列条件等价：
\begin{enumerate}
\item $\mathcal{F}$ 是有限型的；
\item $\mathcal{C}$ 中存在覆盖 $\{X_i \to X\}$，使每个限制
$\mathcal{F}|_{\mathcal{C}/X_i}$ 都是有限型 $\mathcal{O}_{X_i}$-模；
\item $\mathcal{C}$ 中存在覆盖 $\{X_i \to X\}$，使每个限制
$\mathcal{F}|_{\mathcal{C}/X_i}$ 都是由有限多个整体截面生成的
$\mathcal{O}_{X_i}$-模。
\end{enumerate}
\item 下列条件等价：
\begin{enumerate}
\item $\mathcal{F}$ 拟凝聚；
\item $\mathcal{C}$ 中存在覆盖 $\{X_i \to X\}$，使每个限制
$\mathcal{F}|_{\mathcal{C}/X_i}$ 都是拟凝聚 $\mathcal{O}_{X_i}$-模；
\item $\mathcal{C}$ 中存在覆盖 $\{X_i \to X\}$，使每个限制
$\mathcal{F}|_{\mathcal{C}/X_i}$ 都是具有整体表示的
$\mathcal{O}_{X_i}$-模。
\end{enumerate}
\item 下列条件等价：
\begin{enumerate}
\item $\mathcal{F}$ 是有限表示的；
\item $\mathcal{C}$ 中存在覆盖 $\{X_i \to X\}$，使每个限制
$\mathcal{F}|_{\mathcal{C}/X_i}$ 都是有限表示 $\mathcal{O}_{X_i}$-模；
\item $\mathcal{C}$ 中存在覆盖 $\{X_i \to X\}$，使每个限制
$\mathcal{F}|_{\mathcal{C}/X_i}$ 都是具有有限整体表示的
$\mathcal{O}_{X_i}$-模。
\end{enumerate}
\item 下列条件等价：
\begin{enumerate}
\item $\mathcal{F}$ 凝聚；
\item $\mathcal{C}$ 中存在覆盖 $\{X_i \to X\}$，使每个限制
$\mathcal{F}|_{\mathcal{C}/X_i}$ 都是凝聚 $\mathcal{O}_{X_i}$-模。
\end{enumerate}
\end{enumerate}
\end{lemma}

\begin{proof}
在每种情形中都有 (a) $\Rightarrow (b)$。对情形 (1)--(6)，由站点的公理 (2)
（见《站点》，定义 \ref{sites-definition-site}）和模的局部类型定义，
条件 (b) 蕴含条件 (c)。设 $\{X_i \to X\}$ 是覆盖。则对
$\mathcal{C}$ 的每个对象 $U$，得到诱导覆盖
$\{X_i \times_X U \to U\}$。此外，由引理 \ref{sites-modules-lemma-global-pullback}，
(c) 中 $\mathcal{F}|_{\mathcal{C}/X_i}$ 的整体性质蕴含
$\mathcal{F}|_{\mathcal{C}/X_i \times_X U}$ 的相应整体性质，
因而按定义该层具有性质 (a)。略去情形 (7) 中 (b) $\Rightarrow$ (a) 的证明。
\end{proof}

\begin{lemma}
\label{sites-modules-lemma-local-pullback}
设
$(f, f^\sharp) :
(\Sh(\mathcal{C}), \mathcal{O}_\mathcal{C})
\to
(\Sh(\mathcal{D}), \mathcal{O}_\mathcal{D})$
是环拓扑斯态射，且 $\mathcal{F}$ 是 $\mathcal{O}_\mathcal{D}$-模。
\begin{enumerate}
\item 若 $\mathcal{F}$ 局部自由，则 $f^*\mathcal{F}$ 局部自由。
\item 若 $\mathcal{F}$ 有限局部自由，则 $f^*\mathcal{F}$ 有限局部自由。
\item 若 $\mathcal{F}$ 局部由截面生成，则 $f^*\mathcal{F}$ 局部由截面生成。
\item 若 $\mathcal{F}$ 局部由 $r$ 个截面生成，则
$f^*\mathcal{F}$ 局部由 $r$ 个截面生成。
\item 若 $\mathcal{F}$ 是有限型的，则 $f^*\mathcal{F}$ 是有限型的。
\item 若 $\mathcal{F}$ 拟凝聚，则 $f^*\mathcal{F}$ 拟凝聚。
\item 若 $\mathcal{F}$ 是有限表示的，则 $f^*\mathcal{F}$ 是有限表示的。
\end{enumerate}
\end{lemma}

\begin{proof}
按照第 \ref{sites-modules-section-intrinsic} 节的讨论，只须对环站点态射
$(f, f^\sharp) :
(\mathcal{C}, \mathcal{O}_\mathcal{C})
\to
(\mathcal{D}, \mathcal{O}_\mathcal{D})$
检验拉回保持这些性质，其中 $f$ 由具有所有有限极限的站点之间的左正合连续函子
$u : \mathcal{D} \to \mathcal{C}$ 给出。设 $\mathcal{G}$ 是
$\mathcal{O}_\mathcal{D}$-模层，并具有定义
\ref{sites-modules-definition-site-local} 的性质 (1)--(6) 之一。我们知道
$\mathcal{D}$ 有终对象 $Y$，且 $X = u(Y)$ 是 $\mathcal{C}$ 的终对象。
由假设，存在覆盖 $\{Y_i \to Y\}$，使
$\mathcal{G}|_{\mathcal{D}/Y_i}$ 具有相应的整体性质。令
$X_i = u(Y_i)$，于是 $\{X_i \to X\}$ 是 $\mathcal{C}$ 中的覆盖。
由《站点》，引理 \ref{sites-lemma-localize-morphism-strong}，得到环站点态射的交换图
$$
\xymatrix{
(\mathcal{C}/X_i, \mathcal{O}_\mathcal{C}|_{X_i}) \ar[r] \ar[d] &
(\mathcal{C}, \mathcal{O}_\mathcal{C}) \ar[d] \\
(\mathcal{D}/Y_i, \mathcal{O}_\mathcal{D}|_{Y_i}) \ar[r] &
(\mathcal{D}, \mathcal{O}_\mathcal{D})
}
$$
。因此由引理 \ref{sites-modules-lemma-global-pullback}，
$f^*\mathcal{G}|_{X_i}$ 具有相应的整体性质。故由引理
\ref{sites-modules-lemma-local-final-object}，得到 $\mathcal{G}$ 具有起初所考虑的局部性质。
\end{proof}







\section{模的局部类型的基本结果}
\label{sites-modules-section-basics}

\noindent
下面给出与上述定义有关的基本引理。

\begin{lemma}
\label{sites-modules-lemma-cokernel-finite-finite-presentation}
设 $(\mathcal{C},\mathcal{O})$ 是环站点，$\mathcal{F}$ 是有限表示
$\mathcal{O}$-模，且
$\varphi : \mathcal{G} \to \mathcal{F}$ 是 $\mathcal{O}$-模态射。
若 $\mathcal{G}$ 是有限型的，则 $\Coker(\varphi)$ 是有限表示的。
\end{lemma}

\begin{proof}
从略。提示：见《模层》，引理
\ref{modules-lemma-cokernel-finite-finite-presentation}。
\end{proof}

\begin{lemma}
\label{sites-modules-lemma-kernel-surjection-finite-onto-finite-presentation}
设 $(\mathcal{C}, \mathcal{O})$ 是环站点。设
$\theta : \mathcal{G} \to \mathcal{F}$ 是 $\mathcal{O}$-模的满射，
其中 $\mathcal{F}$ 是有限表示的，$\mathcal{G}$ 是有限型的。
则 $\Ker(\theta)$ 是有限型的。
\end{lemma}

\begin{proof}
从略。提示：见《模层》，引理
\ref{modules-lemma-kernel-surjection-finite-free-onto-finite-presentation}。
\end{proof}

\begin{lemma}
\label{sites-modules-lemma-chaotic-quasi-coherent}
设 $\mathcal{C}$ 是范畴，并赋予混沌拓扑而视为站点；见《站点》，例
\ref{sites-example-indiscrete}。设 $\mathcal{O}$ 是 $\mathcal{C}$ 上的环层，
$\mathcal{F}$ 是 $\mathcal{O}$-模层。则 $\mathcal{F}$ 拟凝聚，当且仅当
对 $\mathcal{C}$ 中所有 $U \to V$，典范映射
$$
\mathcal{F}(V) \otimes_{\mathcal{O}(V)} \mathcal{O}(U)
\longrightarrow
\mathcal{F}(U)
$$
是同构。
\end{lemma}

\begin{proof}
假设 $\mathcal{F}$ 拟凝聚，并设 $U \to V$ 是 $\mathcal{C}$ 的态射。
由于 $V$ 的每个覆盖都由一个同构给出，定义
\ref{sites-modules-definition-site-local} 说明存在表示
$$
\bigoplus\nolimits_{j \in J} \mathcal{O}_V \longrightarrow
\bigoplus\nolimits_{i \in I} \mathcal{O}_V \longrightarrow
\mathcal{F}|_{\mathcal{C}/V} \longrightarrow 0
$$
。由于 $\mathcal{C}$ 上的拓扑是混沌拓扑，在 $\mathcal{C}$ 的任意对象上
取截面是正合的。因此得到 $\mathcal{O}(V)$-模 $\mathcal{F}(V)$ 的表示
$$
\bigoplus\nolimits_{j \in J} \mathcal{O}(V) \longrightarrow
\bigoplus\nolimits_{i \in I} \mathcal{O}(V) \longrightarrow
\mathcal{F}(V) \longrightarrow 0
$$
，对 $\mathcal{F}(U)$ 亦然。这立即说明引理陈述中的映射是同构。

\medskip\noindent
反之，假设对 $\mathcal{C}$ 中每个态射 $U \to V$，引理陈述中的映射都是同构。
固定 $V$，并选取 $\mathcal{O}(V)$-模 $\mathcal{F}(V)$ 的表示
$$
\bigoplus\nolimits_{j \in J} \mathcal{O}(V) \longrightarrow
\bigoplus\nolimits_{i \in I} \mathcal{O}(V) \longrightarrow
\mathcal{F}(V) \longrightarrow 0
$$
。则关于 $\mathcal{F}$ 的假设恰好意味着相应的列
$$
\bigoplus\nolimits_{j \in J} \mathcal{O}_V \longrightarrow
\bigoplus\nolimits_{i \in I} \mathcal{O}_V \longrightarrow
\mathcal{F}|_{\mathcal{C}/V} \longrightarrow 0
$$
正合，故 $\mathcal{F}$ 拟凝聚。
\end{proof}

\begin{lemma}
\label{sites-modules-lemma-chaotic-flat-quasi-coherent}
设 $\mathcal{C}$ 是范畴，并赋予混沌拓扑而视为站点；见《站点》，例
\ref{sites-example-indiscrete}。设 $\mathcal{O}$ 是 $\mathcal{C}$ 上的环层。
假设对 $\mathcal{C}$ 中所有 $U \to V$，限制映射
$\mathcal{O}(V) \to \mathcal{O}(U)$ 都是平坦环同态。
则拟凝聚 $\mathcal{O}$-模构成 $\textit{Mod}(\mathcal{O})$
的弱 Serre 子范畴。
\end{lemma}

\begin{proof}
检验弱 Serre 子范畴的定义；见《同调代数》，定义
\ref{homology-definition-serre-subcategory}。为此使用引理
\ref{sites-modules-lemma-chaotic-quasi-coherent} 给出的拟凝聚模刻画。考虑正合列
$$
\mathcal{F}_0 \to \mathcal{F}_1 \to
\mathcal{F}_2 \to \mathcal{F}_3 \to
\mathcal{F}_4
$$
，它在 $\textit{Mod}(\mathcal{O})$ 中，且 $\mathcal{F}_0$、
$\mathcal{F}_1$、$\mathcal{F}_3$ 和 $\mathcal{F}_4$ 拟凝聚。
设 $U \to V$ 是 $\mathcal{C}$ 的态射，并考虑交换图
$$
\vcenter{\hbox{\resizebox{0.98\textwidth}{!}{$
\xymatrix{
\mathcal{F}_0(V) \otimes_{\mathcal{O}(V)} \mathcal{O}(U) \ar[r] \ar[d] &
\mathcal{F}_1(V) \otimes_{\mathcal{O}(V)} \mathcal{O}(U) \ar[r] \ar[d] &
\mathcal{F}_2(V) \otimes_{\mathcal{O}(V)} \mathcal{O}(U) \ar[r] \ar[d] &
\mathcal{F}_3(V) \otimes_{\mathcal{O}(V)} \mathcal{O}(U) \ar[r] \ar[d] &
\mathcal{F}_4(V) \otimes_{\mathcal{O}(V)} \mathcal{O}(U) \ar[d] \\
\mathcal{F}_0(U) \ar[r] &
\mathcal{F}_1(U) \ar[r] &
\mathcal{F}_2(U) \ar[r] &
\mathcal{F}_3(U) \ar[r] &
\mathcal{F}_4(U)
}
$}}}
$$
由假设，指标为 $0$、$1$、$3$、$4$ 的竖直箭头都是同构。
由于 $\mathcal{C}$ 上的拓扑是混沌拓扑，在 $\mathcal{C}$ 的对象上取截面正合，
故下行正合。又因 $\mathcal{O}(V) \to \mathcal{O}(U)$ 平坦，上行也正合。
因此由五引理（《同调代数》，引理 \ref{homology-lemma-five-lemma}），
中间箭头是同构。
\end{proof}






\section{环拓扑斯的闭浸入}
\label{sites-modules-section-closed-immersion}

\noindent
何时称环拓扑斯态射
$i : (\Sh(\mathcal{C}), \mathcal{O}) \to (\Sh(\mathcal{D}), \mathcal{O}')$
为闭浸入？类比《模层》，第 \ref{modules-section-closed-immersion} 节的讨论，
至少作如下假设似乎是自然的：
\begin{enumerate}
\item 函子 $i$ 是拓扑斯的闭浸入（《站点》，定义
\ref{sites-definition-immersion-topoi}）。
\item 相伴映射 $\mathcal{O}' \to i_*\mathcal{O}$ 是满射。
\end{enumerate}
这些条件已经蕴含若干令人满意的结果，本节将讨论它们。不过，由于存在多种可用的
定义，不实际定义环拓扑斯闭浸入这一概念似乎更为审慎。

\begin{lemma}
\label{sites-modules-lemma-i-star-equivalence}
设 $i : (\Sh(\mathcal{C}), \mathcal{O}) \to (\Sh(\mathcal{D}), \mathcal{O}')$
是环拓扑斯态射。假设 $i$ 是拓扑斯的闭浸入，且
$i^\sharp : \mathcal{O}' \to i_*\mathcal{O}$ 是满射。以
$\mathcal{I} \subset \mathcal{O}'$ 表示 $i^\sharp$ 的核。函子
$$
i_* :
\textit{Mod}(\mathcal{O})
\longrightarrow
\textit{Mod}(\mathcal{O}')
$$
正合且全忠实，其本质像恰为满足
$\mathcal{I}\mathcal{G} = 0$ 的 $\mathcal{O}'$-模 $\mathcal{G}$。
\end{lemma}

\begin{proof}
由引理 \ref{sites-modules-lemma-exactness} 和《站点》，引理
\ref{sites-lemma-closed-immersion} 可知 $i_*$ 正合。由 $i_*$ 在集合层上全忠实，
以及 $i^\sharp$ 是满射，得到 $i_*$ 作为函子
$\textit{Mod}(\mathcal{O}) \to \textit{Mod}(\mathcal{O}')$ 也是全忠实的。
事实上，设
$\alpha : i_*\mathcal{F}_1 \to i_*\mathcal{F}_2$ 是
$\mathcal{O}'$-模映射。由 $i_*$ 的全忠实性，得到集合层映射
$\beta : \mathcal{F}_1 \to \mathcal{F}_2$。为证明 $\beta$ 是模映射，
须证明
$$
\xymatrix{
\mathcal{O} \times \mathcal{F}_1 \ar[r] \ar[d] &
\mathcal{F}_1 \ar[d] \\
\mathcal{O} \times \mathcal{F}_2 \ar[r] &
\mathcal{F}_2
}
$$
交换。只须在应用 $i_*$ 后证明交换性。考虑
$$
\xymatrix{
\mathcal{O}' \times i_*\mathcal{F}_1 \ar[r] \ar[d] &
i_*\mathcal{O} \times i_*\mathcal{F}_1 \ar[r] \ar[d] &
i_*\mathcal{F}_1 \ar[d] \\
\mathcal{O}' \times i_*\mathcal{F}_2 \ar[r] &
i_*\mathcal{O} \times i_*\mathcal{F}_2 \ar[r] &
i_*\mathcal{F}_2
}
$$
外矩形交换。由于 $i^\sharp$ 是满射，结论成立。

\medskip\noindent
为完成证明，还须证明关于 $i_*$ 本质像的断言。显然，对任意
$\mathcal{O}$-模 $\mathcal{F}$，$i_*\mathcal{F}$ 被 $\mathcal{I}$ 零化。
反之，设 $\mathcal{G}$ 是满足 $\mathcal{I}\mathcal{G} = 0$ 的
$\mathcal{O}'$-模。由闭子拓扑斯的定义，$\mathcal{D}$ 的终对象存在一个子层
$\mathcal{U}$，使 $i_*$ 在集合层上的本质像恰为如下集合层
$\mathcal{H}$ 所成的类：映射
$\mathcal{H} \times \mathcal{U} \to \mathcal{U}$ 是同构。
特别地，$i_*\mathcal{O} \times \mathcal{U} = \mathcal{U}$。这蕴含
$\mathcal{I} \times \mathcal{U} = \mathcal{O} \times \mathcal{U}$.
因此模 $\mathcal{G}$ 满足
$\mathcal{G} \times \mathcal{U} = \{0\} \times \mathcal{U} = \mathcal{U}$
（因为零模同构于集合层的终对象）。从而存在 $\mathcal{C}$ 上的集合层
$\mathcal{F}$，使 $i_*\mathcal{F} = \mathcal{G}$。由于 $i_*$ 在集合层上
全忠实，为定义加法 $\mathcal{F} \times \mathcal{F} \to \mathcal{F}$ 和
乘法 $\mathcal{O} \times \mathcal{F} \to \mathcal{F}$，只须使用加法
$$
\mathcal{G} \times \mathcal{G} \longrightarrow \mathcal{G}
$$
（这是由 $\mathcal{G}$ 的 $\mathcal{O}'$-模结构给定的）以及乘法
$$
i_*\mathcal{O} \times \mathcal{G} \to \mathcal{G}
$$
。该乘法由 $\mathcal{O}'$ 的乘法给出，因为按假设它零化
$\mathcal{I}$，且 $i_*\mathcal{O} = \mathcal{O}'/\mathcal{I}$。
依构造，$\mathcal{G}$ 同构于由此构造的 $\mathcal{O}$-模
$\mathcal{F}$ 的直像。
\end{proof}












\section{张量积}
\label{sites-modules-section-tensor-product}

\noindent
在第 \ref{sites-modules-section-presheaves-modules} 节和第
\ref{sites-modules-section-sheafification-presheaves-modules} 节中，我们用张量积构造定义了
换环函子。这一构造也可以用来定义 $\mathcal{O}$-模预层的张量积。
确切地，设 $\mathcal{C}$ 是范畴，$\mathcal{O}$ 是环预层，且
$\mathcal{F}$、$\mathcal{G}$ 是 $\mathcal{O}$-模预层。此时定义
$\mathcal{F} \otimes_{p, \mathcal{O}} \mathcal{G}$ 为预层
$$
U
\longmapsto
(\mathcal{F} \otimes_{p, \mathcal{O}} \mathcal{G})(U)
=
\mathcal{F}(U) \otimes_{\mathcal{O}(U)} \mathcal{G}(U)
$$
。若 $\mathcal{C}$ 是站点，$\mathcal{O}$ 是环层，且
$\mathcal{F}$、$\mathcal{G}$ 是 $\mathcal{O}$-模层，则定义
$$
\mathcal{F} \otimes_\mathcal{O} \mathcal{G}
=
(\mathcal{F} \otimes_{p, \mathcal{O}} \mathcal{G})^\#
$$
为与预层 $\mathcal{F} \otimes_{p, \mathcal{O}} \mathcal{G}$ 相伴的
$\mathcal{O}$-模层。

\medskip\noindent
下面列出一些此后将直接使用而不再说明的公式：
$$
(\mathcal{F}
\otimes_{p, \mathcal{O}} \mathcal{G})
\otimes_{p, \mathcal{O}} \mathcal{H}
=
\mathcal{F}
\otimes_{p, \mathcal{O}} (\mathcal{G}
\otimes_{p, \mathcal{O}} \mathcal{H}),
$$
，对层也有类似公式。若 $\mathcal{O}_1 \to \mathcal{O}_2$ 是环预层映射，则
$$
(\mathcal{F} \otimes_{p, \mathcal{O}_1} \mathcal{G})
\otimes_{p, \mathcal{O}_1} \mathcal{O}_2 =
(\mathcal{F} \otimes_{p, \mathcal{O}_1} \mathcal{O}_2)
\otimes_{p, \mathcal{O}_2}
(\mathcal{G} \otimes_{p, \mathcal{O}_1} \mathcal{O}_2),
$$
，对层也有类似公式。这些公式来自相应的代数公式和层化。

\begin{lemma}
\label{sites-modules-lemma-sheafification-tensor}
设 $\mathcal{C}$ 是站点，$\mathcal{O}$ 是环预层，且
$\mathcal{F}$、$\mathcal{G}$ 是 $\mathcal{O}$-模预层。则
$\mathcal{F}^\# \otimes_{\mathcal{O}^\#} \mathcal{G}^\#$
等于
$(\mathcal{F} \otimes_{p, \mathcal{O}} \mathcal{G})^\#$.
\end{lemma}

\begin{proof}
从略。提示：使用下文由双线性映射给出的张量积刻画，以及层化的泛性质。
\end{proof}

\noindent
设 $\mathcal{C}$ 是站点，$\mathcal{O}$ 是环层，且
$\mathcal{F}$、$\mathcal{G}$、$\mathcal{H}$ 是 $\mathcal{O}$-模层。
此时定义
$$
\text{Bilin}_\mathcal{O}(\mathcal{F} \times \mathcal{G}, \mathcal{H})
=
\{\varphi \in
\Mor_{\Sh(\mathcal{C})}(
\mathcal{F} \times \mathcal{G}, \mathcal{H}) \mid
\varphi \text{ 是 }\mathcal{O}\text{-双线性的}\}.
$$
由此定义，有
$$
\Hom_\mathcal{O}
(\mathcal{F} \otimes_\mathcal{O} \mathcal{G}, \mathcal{H})
=
\text{Bilin}_\mathcal{O}(\mathcal{F} \times \mathcal{G}, \mathcal{H}).
$$
换言之，$\mathcal{F} \otimes_\mathcal{O} \mathcal{G}$ 表示如下函子：
它将 $\mathcal{H}$ 送到双线性映射
$\mathcal{F} \times \mathcal{G} \to \mathcal{H}$ 的集合。
特别地，由于在环拓扑斯上两个模之间的双线性映射概念有意义，模层的张量积
对拓扑斯而言是内蕴的（比较第 \ref{sites-modules-section-intrinsic} 节的讨论）。
事实上，有下述结果。

\begin{lemma}
\label{sites-modules-lemma-tensor-product-pullback}
设 $f : (\Sh(\mathcal{C}), \mathcal{O}_\mathcal{C})
\to (\Sh(\mathcal{D}), \mathcal{O}_\mathcal{D})$ 是环拓扑斯态射，且
$\mathcal{F}$、$\mathcal{G}$ 是 $\mathcal{O}_\mathcal{D}$-模。则
$f^*(\mathcal{F} \otimes_{\mathcal{O}_\mathcal{D}} \mathcal{G})
= f^*\mathcal{F} \otimes_{\mathcal{O}_\mathcal{C}} f^*\mathcal{G}$
，且关于 $\mathcal{F}$、$\mathcal{G}$ 具有函子性。
\end{lemma}

\begin{proof}
对 $\mathcal{O}_\mathcal{C}$-模层 $\mathcal{H}$，有
\begin{align*}
\Hom_{\mathcal{O}_\mathcal{C}}(
f^*(\mathcal{F} \otimes_\mathcal{O} \mathcal{G}), \mathcal{H})
& =
\Hom_{\mathcal{O}_\mathcal{D}}(
\mathcal{F} \otimes_\mathcal{O} \mathcal{G}, f_*\mathcal{H}) \\
& =
\text{Bilin}_{\mathcal{O}_\mathcal{D}}(
\mathcal{F} \times \mathcal{G}, f_*\mathcal{H}) \\
& =
\text{Bilin}_{f^{-1}\mathcal{O}_\mathcal{D}}(
f^{-1}\mathcal{F} \times f^{-1}\mathcal{G}, \mathcal{H}) \\
& =
\Hom_{f^{-1}\mathcal{O}_\mathcal{D}}(
f^{-1}\mathcal{F} \otimes_{f^{-1}\mathcal{O}_\mathcal{D}} f^{-1}\mathcal{G},
\mathcal{H}) \\
& =
\Hom_{\mathcal{O}_\mathcal{C}}(
f^*\mathcal{F} \otimes_{f^*\mathcal{O}_\mathcal{D}} f^*\mathcal{G},
\mathcal{H})
\end{align*}
这一串等式中有实质内容的是第三个“$=$”。它直接来自定义以及前面各节所讨论的
$f_*$ 与 $f^{-1}$ 的伴随性。
\end{proof}

\begin{lemma}
\label{sites-modules-lemma-tensor-product-permanence}
设 $(\mathcal{C}, \mathcal{O})$ 是环站点，且 $\mathcal{F}$、
$\mathcal{G}$ 是 $\mathcal{O}$-模层。
\begin{enumerate}
\item 若 $\mathcal{F}$、$\mathcal{G}$ 局部自由，则
$\mathcal{F} \otimes_\mathcal{O} \mathcal{G}$ 也局部自由。
\item 若 $\mathcal{F}$、$\mathcal{G}$ 有限局部自由，则
$\mathcal{F} \otimes_\mathcal{O} \mathcal{G}$ 也有限局部自由。
\item 若 $\mathcal{F}$、$\mathcal{G}$ 局部由截面生成，则
$\mathcal{F} \otimes_\mathcal{O} \mathcal{G}$ 也局部由截面生成。
\item 若 $\mathcal{F}$、$\mathcal{G}$ 是有限型的，则
$\mathcal{F} \otimes_\mathcal{O} \mathcal{G}$ 也是有限型的。
\item 若 $\mathcal{F}$、$\mathcal{G}$ 拟凝聚，则
$\mathcal{F} \otimes_\mathcal{O} \mathcal{G}$ 也拟凝聚。
\item 若 $\mathcal{F}$、$\mathcal{G}$ 是有限表示的，则
$\mathcal{F} \otimes_\mathcal{O} \mathcal{G}$ 也是有限表示的。
\item 若 $\mathcal{F}$ 是有限表示的，且 $\mathcal{G}$ 凝聚，则
$\mathcal{F} \otimes_\mathcal{O} \mathcal{G}$ 凝聚。
\item 若 $\mathcal{F}$、$\mathcal{G}$ 凝聚，则
$\mathcal{F} \otimes_\mathcal{O} \mathcal{G}$ 也凝聚。
\end{enumerate}
\end{lemma}

\begin{proof}
从略。提示：比较《模层》，引理 \ref{modules-lemma-tensor-product-permanence}。
\end{proof}



\section{内 Hom}
\label{sites-modules-section-internal-hom}

\noindent
设 $\mathcal{C}$ 是范畴，$\mathcal{O}$ 是环预层，且
$\mathcal{F}$、$\mathcal{G}$ 是 $\mathcal{O}$-模预层。考虑规则
$$
U \longmapsto \Hom_{\mathcal{O}_U}(\mathcal{F}|_U, \mathcal{G}|_U).
$$
。对 $\mathcal{C}$ 中的 $\varphi : V \to U$，定义限制映射
$$
\Hom_{\mathcal{O}_U}(\mathcal{F}|_U, \mathcal{G}|_U)
\longrightarrow
\Hom_{\mathcal{O}_V}(\mathcal{F}|_V, \mathcal{G}|_V)
$$
，即经由再局部化态射 $j : \mathcal{C}/V \to \mathcal{C}/U$ 作限制；
见《站点》，引理 \ref{sites-lemma-relocalize}。由此定义预层
$\SheafHom_\mathcal{O}(\mathcal{F}, \mathcal{G})$。此外，给定元素
$\varphi \in \Hom_{\mathcal{O}|_U}(\mathcal{F}|_U, \mathcal{G}|_U)$
和截面 $f \in \mathcal{O}(U)$，可以定义
$f\varphi \in \Hom_{\mathcal{O}|_U}(\mathcal{F}|_U, \mathcal{G}|_U)$
，方法是与 $\mathcal{F}|_U$ 上乘以 $f$ 的映射前复合，或与
$\mathcal{G}|_U$ 上乘以 $f$ 的映射后复合（两者结果相同）。
因此实际得到 $\mathcal{O}$-模预层。存在典范的“求值”态射
$$
\mathcal{F}
\otimes_{p, \mathcal{O}}
\SheafHom_\mathcal{O}(\mathcal{F}, \mathcal{G})
\longrightarrow
\mathcal{G}.
$$

\begin{lemma}
\label{sites-modules-lemma-internal-hom}
若 $\mathcal{C}$ 是站点，$\mathcal{O}$ 是环层，$\mathcal{F}$ 是
$\mathcal{O}$-模预层，且 $\mathcal{G}$ 是 $\mathcal{O}$-模层，则
$\SheafHom_\mathcal{O}(\mathcal{F}, \mathcal{G})$
是 $\mathcal{O}$-模层。
\end{lemma}

\begin{proof}
从略。提示：先注意
$\SheafHom_\mathcal{O}(\mathcal{F}, \mathcal{G})
= \SheafHom_\mathcal{O}(\mathcal{F}^\#, \mathcal{G})$，这把问题约化到
$\mathcal{F}$ 和 $\mathcal{G}$ 都是层的情形。关于集合层的结果见
《站点》，引理 \ref{sites-lemma-glue-maps}。
\end{proof}

\begin{lemma}
\label{sites-modules-lemma-internal-hom-restriction}
设 $(\mathcal{C}, \mathcal{O})$ 是环站点，且 $\mathcal{F}$、
$\mathcal{G}$ 是 $\mathcal{O}$-模层。则对 $U \in \Ob(\mathcal{C})$，
$\SheafHom_\mathcal{O}(\mathcal{F}, \mathcal{G})$ 的形成与到 $U$ 的限制交换。
\end{lemma}

\begin{proof}
由定义立即可得。
\end{proof}

\begin{remark}
\label{sites-modules-remark-pullback-internal-hom}
设 $f : (\mathcal{C}, \mathcal{O}_\mathcal{C}) \to
(\mathcal{D}, \mathcal{O}_\mathcal{D})$ 是环站点态射，且
$\mathcal{F}$、$\mathcal{G}$ 是 $\mathcal{O}_\mathcal{D}$-模层。
存在典范映射
$$
f^*\SheafHom_{\mathcal{O}_\mathcal{D}}(\mathcal{F}, \mathcal{G})
\longrightarrow
\SheafHom_{\mathcal{O}_\mathcal{C}}(f^*\mathcal{F}, f^*\mathcal{G})
$$
。事实上，此映射与如下映射伴随：
$$
\SheafHom_{\mathcal{O}_\mathcal{D}}(\mathcal{F}, \mathcal{G})
\longrightarrow
f_*\SheafHom_{\mathcal{O}_\mathcal{C}}(f^*\mathcal{F}, f^*\mathcal{G})
$$
该映射定义如下。设 $f$ 由连续函子 $u : \mathcal{D} \to \mathcal{C}$ 给出。
对 $V \in \Ob(\mathcal{D})$ 上的截面，使用映射
\begin{align*}
\Gamma(V, \SheafHom_{\mathcal{O}_\mathcal{D}}(\mathcal{F}, \mathcal{G}))
& =
\Hom_{\mathcal{O}_V}(\mathcal{F}|_V, \mathcal{G}|_V) \\
& \longrightarrow
\Hom_{\mathcal{O}_{u(V)}}(f^*\mathcal{F}|_{u(V)}, f^*\mathcal{G}|_{u(V)}) \\
& =
\Gamma(u(V), \SheafHom_{\mathcal{O}_\mathcal{C}}(f^*\mathcal{F},
f^*\mathcal{G})) \\
& =
\Gamma(V, f_*\SheafHom_{\mathcal{O}_\mathcal{C}}(f^*\mathcal{F},
f^*\mathcal{G}))
\end{align*}
，其中箭头使用由 $f$ 诱导的态射
$(\mathcal{C}/u(V), \mathcal{O}_{u(V)}) \to
(\mathcal{D}/V, \mathcal{O}_V)$ 的拉回。
\end{remark}

\noindent
在引理 \ref{sites-modules-lemma-internal-hom} 的情形中，“求值”态射通过模层的张量积分解：
$$
\mathcal{F}
\otimes_\mathcal{O}
\SheafHom_\mathcal{O}(\mathcal{F}, \mathcal{G})
\longrightarrow
\mathcal{G}.
$$

\begin{lemma}
\label{sites-modules-lemma-internal-hom-commute-limits}
内 Hom 与（余）极限。设 $\mathcal{C}$ 是范畴，且 $\mathcal{O}$ 是环预层。
\begin{enumerate}
\item 对任意 $\mathcal{O}$-模预层 $\mathcal{F}$，函子
$$
\textit{PMod}(\mathcal{O}) \longrightarrow \textit{PMod}(\mathcal{O})
, \quad
\mathcal{G} \longmapsto \SheafHom_\mathcal{O}(\mathcal{F}, \mathcal{G})
$$
与任意极限交换。
\item 对任意 $\mathcal{O}$-模预层 $\mathcal{G}$，函子
$$
\textit{PMod}(\mathcal{O}) \longrightarrow \textit{PMod}(\mathcal{O})^{opp}
, \quad
\mathcal{F} \longmapsto \SheafHom_\mathcal{O}(\mathcal{F}, \mathcal{G})
$$
与任意余极限交换；公式为
$$
\SheafHom_\mathcal{O}(\colim_i \mathcal{F}_i, \mathcal{G})
=
\lim_i \SheafHom_\mathcal{O}(\mathcal{F}_i, \mathcal{G}).
$$
\end{enumerate}
再假设 $\mathcal{C}$ 是站点，且 $\mathcal{O}$ 是环层。
\begin{enumerate}
\item[(3)] 对任意 $\mathcal{O}$-模层 $\mathcal{F}$，函子
$$
\textit{Mod}(\mathcal{O}) \longrightarrow \textit{Mod}(\mathcal{O})
, \quad
\mathcal{G} \longmapsto \SheafHom_\mathcal{O}(\mathcal{F}, \mathcal{G})
$$
与任意极限交换。
\item[(4)] 对任意 $\mathcal{O}$-模层 $\mathcal{G}$，函子
$$
\textit{Mod}(\mathcal{O}) \longrightarrow \textit{Mod}(\mathcal{O})^{opp}
, \quad
\mathcal{F} \longmapsto \SheafHom_\mathcal{O}(\mathcal{F}, \mathcal{G})
$$
与任意余极限交换；公式为
$$
\SheafHom_\mathcal{O}(\colim_i \mathcal{F}_i, \mathcal{G})
=
\lim_i \SheafHom_\mathcal{O}(\mathcal{F}_i, \mathcal{G}).
$$
\end{enumerate}
\end{lemma}

\begin{proof}
设 $\mathcal{I} \to \textit{PMod}(\mathcal{O})$、
$i \mapsto \mathcal{G}_i$ 是图，且 $U$ 是范畴 $\mathcal{C}$ 的对象。
由于 $j_U^*$ 既是左伴随又是右伴随，有
$\lim_i j_U^*\mathcal{G}_i = j_U^* \lim_i \mathcal{G}_i$.
因此
\begin{align*}
\SheafHom_\mathcal{O}(\mathcal{F}, \lim_i \mathcal{G}_i)(U)
& =
\Hom_{\mathcal{O}_U}(\mathcal{F}|_U, \lim_i \mathcal{G}_i|_U) \\
& =
\lim_i \Hom_{\mathcal{O}_U}(\mathcal{F}|_U, \mathcal{G}_i|_U) \\
& = \lim_i \SheafHom_\mathcal{O}(\mathcal{F}, \mathcal{G}_i)(U)
\end{align*}
，其中用到了极限的定义。这证明了 (1)。(2) 的证明完全相同。
(3) 来自 (1)，因为层图的极限与预层范畴中的极限相同。
最后，(4) 成立，因为公式中有
$$
\Mor_{\textit{Mod}(\mathcal{O})}(
\colim_i \mathcal{F}_i, \mathcal{G})
=
\Mor_{\textit{PMod}(\mathcal{O})}(
\colim^{PSh}_i \mathcal{F}_i, \mathcal{G})
$$
，这是因为余极限 $\colim_i \mathcal{F}_i$ 是
$\textit{PMod}(\mathcal{O})$ 中余极限
$\colim^{PSh}_i \mathcal{F}_i$ 的层化。因此 (4) 来自 (2)
（并再次使用上述关于极限的说明）。
\end{proof}

\begin{lemma}
\label{sites-modules-lemma-internal-hom-exact}
设 $(\mathcal{C}, \mathcal{O})$ 是环站点，且 $\mathcal{F}$、
$\mathcal{G}$ 是 $\mathcal{O}$-模。
\begin{enumerate}
\item 若 $\mathcal{F}_2 \to \mathcal{F}_1 \to \mathcal{F} \to 0$
是 $\mathcal{O}$-模的正合列，则
$$
0 \to
\SheafHom_\mathcal{O}(\mathcal{F}, \mathcal{G}) \to
\SheafHom_\mathcal{O}(\mathcal{F}_1, \mathcal{G}) \to
\SheafHom_\mathcal{O}(\mathcal{F}_2, \mathcal{G})
$$
正合。
\item 若 $0 \to \mathcal{G} \to \mathcal{G}_1 \to \mathcal{G}_2$
是 $\mathcal{O}$-模的正合列，则
$$
0 \to
\SheafHom_\mathcal{O}(\mathcal{F}, \mathcal{G}) \to
\SheafHom_\mathcal{O}(\mathcal{F}, \mathcal{G}_1) \to
\SheafHom_\mathcal{O}(\mathcal{F}, \mathcal{G}_2)
$$
正合。
\end{enumerate}
\end{lemma}

\begin{proof}
由引理 \ref{sites-modules-lemma-internal-hom-commute-limits} 和《同调代数》，引理
\ref{homology-lemma-exact-functor} 可得。
\end{proof}

\begin{lemma}
\label{sites-modules-lemma-internal-hom-adjoint-tensor}
设 $\mathcal{C}$ 是范畴，且 $\mathcal{O}$ 是环预层。
\begin{enumerate}
\item 设 $\mathcal{F}$、$\mathcal{G}$、$\mathcal{H}$ 是
$\mathcal{O}$-模预层。存在典范同构
$$
\SheafHom_\mathcal{O}
(\mathcal{F} \otimes_{p, \mathcal{O}} \mathcal{G}, \mathcal{H})
\longrightarrow
\SheafHom_\mathcal{O}
(\mathcal{F}, \SheafHom_\mathcal{O}(\mathcal{G}, \mathcal{H}))
$$
，它关于三个变元均具有函子性（三处都取层 Hom）。特别地，
$$
\Mor_{\textit{PMod}(\mathcal{O})}(
\mathcal{F} \otimes_{p, \mathcal{O}} \mathcal{G}, \mathcal{H})
=
\Mor_{\textit{PMod}(\mathcal{O})}(
\mathcal{F}, \SheafHom_\mathcal{O}(\mathcal{G}, \mathcal{H}))
$$
\item
假设 $\mathcal{C}$ 是站点，$\mathcal{O}$ 是环层，且
$\mathcal{F}$、$\mathcal{G}$、$\mathcal{H}$ 是 $\mathcal{O}$-模层。
存在典范同构
$$
\SheafHom_\mathcal{O}
(\mathcal{F} \otimes_\mathcal{O} \mathcal{G}, \mathcal{H})
\longrightarrow
\SheafHom_\mathcal{O}
(\mathcal{F}, \SheafHom_\mathcal{O}(\mathcal{G}, \mathcal{H}))
$$
，它关于三个变元均具有函子性（三处都取层 Hom）。特别地，
$$
\Mor_{\textit{Mod}(\mathcal{O})}(
\mathcal{F} \otimes_\mathcal{O} \mathcal{G}, \mathcal{H})
=
\Mor_{\textit{Mod}(\mathcal{O})}(
\mathcal{F}, \SheafHom_\mathcal{O}(\mathcal{G}, \mathcal{H}))
$$
\end{enumerate}
\end{lemma}

\begin{proof}
这是《代数》，引理 \ref{algebra-lemma-hom-from-tensor-product} 的类似版本。
证明相同，从略。
\end{proof}

\begin{lemma}
\label{sites-modules-lemma-tensor-commute-colimits}
张量积与余极限。设 $\mathcal{C}$ 是范畴，且 $\mathcal{O}$ 是环预层。
\begin{enumerate}
\item 对任意 $\mathcal{O}$-模预层 $\mathcal{F}$，函子
$$
\textit{PMod}(\mathcal{O}) \longrightarrow \textit{PMod}(\mathcal{O})
, \quad
\mathcal{G} \longmapsto \mathcal{F} \otimes_{p, \mathcal{O}} \mathcal{G}
$$
与任意余极限交换。
\item
假设 $\mathcal{C}$ 是站点，且 $\mathcal{O}$ 是环层。对任意
$\mathcal{O}$-模层 $\mathcal{F}$，函子
$$
\textit{Mod}(\mathcal{O}) \longrightarrow \textit{Mod}(\mathcal{O})
, \quad
\mathcal{G} \longmapsto \mathcal{F} \otimes_\mathcal{O} \mathcal{G}
$$
与任意余极限交换。
\end{enumerate}
\end{lemma}

\begin{proof}
这是因为由引理 \ref{sites-modules-lemma-internal-hom-adjoint-tensor}，张量积与内 Hom 伴随。
见《范畴》，引理 \ref{categories-lemma-adjoint-exact}。
\end{proof}

\begin{lemma}
\label{sites-modules-lemma-adjoint-hom-restrict}
设 $\mathcal{C}$ 是范畴，相应地是站点；设
$\mathcal{O} \to \mathcal{O}'$ 是环预层映射，相应地是环层映射。则
$$
\Hom_\mathcal{O}(\mathcal{G}, \mathcal{F}) =
\Hom_{\mathcal{O}'}(\mathcal{G},
\SheafHom_\mathcal{O}(\mathcal{O}', \mathcal{F}))
$$
，其中 $\mathcal{G}$ 是任意 $\mathcal{O}'$-模，$\mathcal{F}$ 是任意
$\mathcal{O}$-模。
\end{lemma}

\begin{proof}
这是《代数》，引理 \ref{algebra-lemma-adjoint-hom-restrict} 的类似版本。
证明相同，从略。
\end{proof}

\begin{lemma}
\label{sites-modules-lemma-j-shriek-and-tensor}
设 $(\mathcal{C}, \mathcal{O})$ 是环站点，且
$U \in \Ob(\mathcal{C})$。对
$\textit{Mod}(\mathcal{O}_U)$ 中的 $\mathcal{G}$ 和
$\textit{Mod}(\mathcal{O})$ 中的 $\mathcal{F}$，有
$j_{U!}\mathcal{G} \otimes_\mathcal{O} \mathcal{F} =
j_{U!}(\mathcal{G} \otimes_{\mathcal{O}_U} \mathcal{F}|_U)$.
\end{lemma}

\begin{proof}
设 $\mathcal{H}$ 是 $\textit{Mod}(\mathcal{O})$ 的对象。则
\begin{align*}
\Hom_\mathcal{O}(j_{U!}(\mathcal{G} \otimes_{\mathcal{O}_U} \mathcal{F}|_U),
\mathcal{H})
& =
\Hom_{\mathcal{O}_U}(\mathcal{G} \otimes_{\mathcal{O}_U} \mathcal{F}|_U,
\mathcal{H}|_U) \\
& =
\Hom_{\mathcal{O}_U}(\mathcal{G},
\SheafHom_{\mathcal{O}_U}(\mathcal{F}|_U, \mathcal{H}|_U)) \\
& =
\Hom_{\mathcal{O}_U}(\mathcal{G},
\SheafHom_\mathcal{O}(\mathcal{F}, \mathcal{H})|_U) \\
& =
\Hom_\mathcal{O}(j_{U!}\mathcal{G},
\SheafHom_\mathcal{O}(\mathcal{F}, \mathcal{H})) \\
& =
\Hom_\mathcal{O}(j_{U!}\mathcal{G} \otimes_\mathcal{O} \mathcal{F},
\mathcal{H})
\end{align*}
第一个等式成立是因为 $j_{U!}$ 是模限制的左伴随；第二个等式由引理
\ref{sites-modules-lemma-internal-hom-adjoint-tensor} 得到；第三个等式由引理
\ref{sites-modules-lemma-internal-hom-restriction} 得到；第四个等式成立是因为
$j_{U!}$ 是模限制的左伴随；第五个等式由引理
\ref{sites-modules-lemma-internal-hom-adjoint-tensor} 得到。再由 Yoneda 引理即得本引理。
\end{proof}

\begin{remark}
\label{sites-modules-remark-j-shriek-tensor}
设 $\mathcal{C}$ 是站点。设 $\mathcal{F}$ 是 $\mathcal{C}$ 上的集合层，并考虑局部化态射
$j : \Sh(\mathcal{C})/\mathcal{F} \to \Sh(\mathcal{C})$.
见《站点》，定义 \ref{sites-definition-localize-topos}。
我们断言：(a) $j_!\mathbf{Z} = \mathbf{Z}_\mathcal{F}^\#$；
(b) 对 $\mathcal{C}$ 上任意阿贝尔层 $\mathcal{H}$，都有
$j_!(j^{-1}\mathcal{H}) = j_!\mathbf{Z} \otimes_\mathbf{Z} \mathcal{H}$。
令 $\mathcal{G}$ 为 $\mathcal{C}$ 上的阿贝尔层。由 Yoneda 引理，(a) 成立，因为
$$
\Hom(j_!\mathbf{Z}, \mathcal{G}) =
\Hom(\mathbf{Z}, j^{-1}\mathcal{G}) =
\Hom(\mathbf{Z}_\mathcal{F}^\#, \mathcal{G})
$$
，其中第二个等式成立，是因为等式两边都给出从 $\mathcal{F}$ 到
$\mathcal{G}$ 的映射集，并将其视为阿贝尔群。对于 (b)，使用 Yoneda 引理以及
\begin{align*}
\Hom(j_!(j^{-1}\mathcal{H}), \mathcal{G})
& =
\Hom(j^{-1}\mathcal{H}, j^{-1}\mathcal{G}) \\
& =
\Hom(\mathbf{Z}, \SheafHom(j^{-1}\mathcal{H}, j^{-1}\mathcal{G})) \\
& =
\Hom(\mathbf{Z}, j^{-1}\SheafHom(\mathcal{H}, \mathcal{G})) \\
& =
\Hom(j_!\mathbf{Z}, \SheafHom(\mathcal{H}, \mathcal{G})) \\
& =
\Hom(j_!\mathbf{Z} \otimes_\mathbf{Z} \mathcal{H}, \mathcal{G})
\end{align*}
即可。这里用到了伴随性、取 $\SheafHom$ 与局部化交换这一事实，以及引理
\ref{sites-modules-lemma-internal-hom-adjoint-tensor}。
\end{remark}

\begin{lemma}
\label{sites-modules-lemma-hom-finite-presentation-colimit}
设 $(\mathcal{C}, \mathcal{O})$ 是环站点，且 $\mathcal{F}$ 是有限表示
$\mathcal{O}$-模。设
$\mathcal{G} = \colim_{\lambda \in \Lambda} \mathcal{G}_\lambda$
是 $\mathcal{O}$-模的滤过余极限。则典范映射
$$
\colim_\lambda \SheafHom_\mathcal{O}(\mathcal{F}, \mathcal{G}_\lambda)
\longrightarrow
\SheafHom_\mathcal{O}(\mathcal{F}, \mathcal{G})
$$
是同构。
\end{lemma}

\begin{proof}
只需证明对 $\mathcal{C}$ 中每个 $U$ 作限制后，该箭头是同构。模层的余极限和
内 Hom 都与到 $U$ 的限制交换；例如见引理
\ref{sites-modules-lemma-exactness-pushforward-pullback} 和
\ref{sites-modules-lemma-internal-hom-restriction}。固定 $U$。给定覆盖
$\{U_i \to U\}_{i \in I}$，只需证明到每个 $U_i$ 的限制是同构。因此可以假设
$\mathcal{F}$ 有整体表示
$$
\bigoplus\nolimits_{j = 1, \ldots, m}
\mathcal{O}
\longrightarrow
\bigoplus\nolimits_{i = 1, \ldots, n}
\mathcal{O}
\to
\mathcal{F}
\to
0
$$
。函子 $\SheafHom_\mathcal{O}(-, -)$ 对任一变元都与有限直和交换，且
$\SheafHom_\mathcal{O}(\mathcal{O}, -)$ 是恒等函子。由此以及引理
\ref{sites-modules-lemma-internal-hom-exact}，得到正合列
$$
0 \to
\SheafHom_\mathcal{O}(\mathcal{F}, \mathcal{G}) \to
\bigoplus\nolimits_{i = 1, \ldots, n} \mathcal{G} \to
\bigoplus\nolimits_{j = 1, \ldots, m} \mathcal{G}
$$
。由引理 \ref{sites-modules-lemma-limits-colimits}，滤过余极限在
$\textit{Mod}(\mathcal{O})$ 中正合，所以下列交换图的上行也正合：
{\small
$$
\xymatrix{
0 \ar[r] &
\colim_\lambda \SheafHom_\mathcal{O}(\mathcal{F}, \mathcal{G}_\lambda)
\ar[r] \ar[d] &
\colim_\lambda \bigoplus\nolimits_{i = 1, \ldots, n} \mathcal{G}_\lambda
\ar[r] \ar[d] &
\colim_\lambda \bigoplus\nolimits_{j = 1, \ldots, m} \mathcal{G}_\lambda
\ar[d] \\
0 \ar[r] &
\SheafHom_\mathcal{O}(\mathcal{F}, \mathcal{G}) \ar[r] &
\bigoplus\nolimits_{i = 1, \ldots, n} \mathcal{G} \ar[r] &
\bigoplus\nolimits_{j = 1, \ldots, m} \mathcal{G}
}
$$
}
右侧两个竖直箭头都是同构，故结论成立。
\end{proof}

\begin{lemma}
\label{sites-modules-lemma-finite-presentation-quasi-compact-colimit}
设 $(\mathcal{C}, \mathcal{O})$ 是环站点。设
$\mathcal{G} = \colim_{\lambda \in \Lambda} \mathcal{G}_\lambda$
是 $\mathcal{O}$-模的滤过余极限，且 $\mathcal{F}$ 是有限表示
$\mathcal{O}$-模。若站点 $\mathcal{C}$ 满足《站点》，引理
\ref{sites-lemma-directed-colimits-global-sections} 第 (4) 部分的假设，则有
$$
\colim_\lambda \Hom_\mathcal{O}(\mathcal{F}, \mathcal{G}_\lambda)
=
\Hom_\mathcal{O}(\mathcal{F}, \mathcal{G}).
$$
另见《站点》，注 \ref{sites-remark-stronger-conditions}。
\end{lemma}

\begin{proof}
置 $\mathcal{H} =
\SheafHom_\mathcal{O}(\mathcal{F}, \colim \mathcal{G}_\lambda)$
以及 $\mathcal{H}_\lambda =
\SheafHom_\mathcal{O}(\mathcal{F}, \mathcal{G}_\lambda)$。由构造可知
$$
\Hom_\mathcal{O}(\mathcal{F}, \mathcal{G}) = \Gamma(\mathcal{C}, \mathcal{H})
\quad\text{且}\quad
\Hom_\mathcal{O}(\mathcal{F}, \mathcal{G}_\lambda) =
\Gamma(\mathcal{C}, \mathcal{H}_\lambda)
$$
。由引理 \ref{sites-modules-lemma-hom-finite-presentation-colimit}，有
$\mathcal{H} = \colim \mathcal{H}_\lambda$。故本引理由《站点》，引理
\ref{sites-lemma-directed-colimits-global-sections} 得出。
\end{proof}









\section{平坦模}
\label{sites-modules-section-flat}

\noindent
可以与环上模的情形完全相同地定义平坦模。

\begin{definition}
\label{sites-modules-definition-flat}
设 $\mathcal{C}$ 是范畴，且 $\mathcal{O}$ 是环预层。
\begin{enumerate}
\item 如果函子
$$
\textit{PMod}(\mathcal{O})
\longrightarrow
\textit{PMod}(\mathcal{O}), \quad
\mathcal{G} \mapsto \mathcal{G} \otimes_{p, \mathcal{O}} \mathcal{F}
$$
正合，则称 $\mathcal{O}$-模预层 $\mathcal{F}$ 是{\it 平坦的}。
\item 如果 $\mathcal{O}'$ 作为 $\mathcal{O}$-模预层是平坦的，则称环预层映射
$\mathcal{O} \to \mathcal{O}'$ 是{\it 平坦的}。
\item 假设 $\mathcal{C}$ 是站点，$\mathcal{O}$ 是环层，且
$\mathcal{F}$ 是 $\mathcal{O}$-模层。如果函子
$$
\textit{Mod}(\mathcal{O})
\longrightarrow
\textit{Mod}(\mathcal{O}), \quad
\mathcal{G} \mapsto \mathcal{G} \otimes_\mathcal{O} \mathcal{F}
$$
正合，则称 $\mathcal{F}$ 是{\it 平坦的}。
\item 如果 $\mathcal{O}'$ 作为 $\mathcal{O}$-模层是平坦的，则称站点上的环层映射
$\mathcal{O} \to \mathcal{O}'$ 是{\it 平坦的}。
\end{enumerate}
\end{definition}

\noindent
平坦模和平坦环映射的概念是内蕴的（第 \ref{sites-modules-section-intrinsic} 节）。

\begin{lemma}
\label{sites-modules-lemma-flatness-presheaves}
设 $\mathcal{C}$ 是范畴，$\mathcal{O}$ 是环预层，且
$\mathcal{F}$ 是 $\mathcal{O}$-模预层。如果每个 $\mathcal{F}(U)$ 都是平坦
$\mathcal{O}(U)$-模，则 $\mathcal{F}$ 是平坦的。
\end{lemma}

\begin{proof}
由定义立即得到。
\end{proof}

\begin{lemma}
\label{sites-modules-lemma-flatness-sheafification}
设 $\mathcal{C}$ 是站点，$\mathcal{O}$ 是环预层，且
$\mathcal{F}$ 是 $\mathcal{O}$-模预层。如果 $\mathcal{F}$ 是平坦
$\mathcal{O}$-模，则 $\mathcal{F}^\#$ 是平坦 $\mathcal{O}^\#$-模。
\end{lemma}

\begin{proof}
略。（提示：层化是正合的。）
\end{proof}

\begin{lemma}
\label{sites-modules-lemma-flatness-sheafification-refined}
设 $\mathcal{C}$ 是站点，$\mathcal{O}$ 是环预层，且
$\mathcal{F}$ 是 $\mathcal{O}$-模预层。假设 $\mathcal{C}$ 的每个对象 $U$
都有覆盖 $\{U_i \to U\}_{i \in I}$，使得 $\mathcal{F}(U_i)$ 是平坦
$\mathcal{O}(U_i)$-模。则 $\mathcal{F}^\#$ 是平坦 $\mathcal{O}^\#$-模。
\end{lemma}

\begin{proof}
设 $\mathcal{G} \subset \mathcal{G}'$ 是 $\mathcal{O}^\#$-模的包含。需要证明
$$
\mathcal{G} \otimes_{\mathcal{O}^\#} \mathcal{F}^\#
\longrightarrow
\mathcal{G}' \otimes_{\mathcal{O}^\#} \mathcal{F}^\#
$$
是单射。由引理 \ref{sites-modules-lemma-sheafification-tensor}，该箭头的源是预层
$\mathcal{G} \otimes_{p, \mathcal{O}} \mathcal{F}$ 的层化，靶也类似。
如果 $U$ 是 $\mathcal{C}$ 的对象，且 $\mathcal{F}(U)$ 是平坦
$\mathcal{O}(U)$-模，则
$$
(\mathcal{G} \otimes_{p, \mathcal{O}} \mathcal{F})(U) =
\mathcal{G}(U) \otimes_{\mathcal{O}(U)} \mathcal{F}(U)
\longrightarrow
\mathcal{G}'(U) \otimes_{\mathcal{O}(U)} \mathcal{F}(U) =
(\mathcal{G}' \otimes_{p, \mathcal{O}} \mathcal{F})(U)
$$
是单射。因此归结为证明：给定 $\mathcal{C}$ 上预层映射
$f : \mathcal{H} \to \mathcal{H}'$，如果 $\mathcal{C}$ 中每个 $U$ 都有覆盖
$\{U_i \to U\}_{i \in I}$，使得
$\mathcal{H}(U_i) \to \mathcal{H}'(U_i)$ 是单射，则 $f^\#$ 是单射。
将此留给读者作为练习。
\end{proof}

\begin{lemma}
\label{sites-modules-lemma-colimits-flat}
余极限与张量积。
\begin{enumerate}
\item 平坦模预层的滤过余极限是平坦的。平坦模预层的直和是平坦的。
\item 平坦模层的滤过余极限是平坦的。平坦模层的直和是平坦的。
\end{enumerate}
\end{lemma}

\begin{proof}
考察对象上的截面，由引理 \ref{sites-modules-lemma-tensor-commute-colimits} 和《代数》，引理
\ref{algebra-lemma-directed-colimit-exact} 即得 (1)。对于 (2)，使用引理
\ref{sites-modules-lemma-tensor-commute-colimits} 以及正合复形的滤过余极限仍为正合复形这一事实
（这里用到层化正合且与余极限交换）。略去若干细节。
\end{proof}

\begin{lemma}
\label{sites-modules-lemma-restriction-flat}
设 $(\mathcal{C}, \mathcal{O})$ 是环站点，且 $U$ 是 $\mathcal{C}$ 的对象。
如果 $\mathcal{F}$ 是平坦 $\mathcal{O}$-模，则 $\mathcal{F}|_U$ 是平坦
$\mathcal{O}_U$-模。
\end{lemma}

\begin{proof}
设 $\mathcal{G}_1 \to \mathcal{G}_2 \to \mathcal{G}_3$ 是
$\mathcal{O}_U$-模的正合复形。由于 $j_{U!}$ 正合（引理
\ref{sites-modules-lemma-extension-by-zero-exact}），且 $\mathcal{F}$ 作为
$\mathcal{O}$-模是平坦的，故由下列模组成的复形正合：
$$
j_{U!}(\mathcal{G}_i \otimes_{\mathcal{O}_U} \mathcal{F}|_U) =
j_{U!}\mathcal{G}_i \otimes_\mathcal{O} \mathcal{F}
$$
（引理 \ref{sites-modules-lemma-j-shriek-and-tensor}）。由引理
\ref{sites-modules-lemma-j-shriek-reflects-exactness} 可知
$\mathcal{G}_1 \otimes_{\mathcal{O}_U} \mathcal{F}|_U \to
\mathcal{G}_2 \otimes_{\mathcal{O}_U} \mathcal{F}|_U \to
\mathcal{G}_3 \otimes_{\mathcal{O}_U} \mathcal{F}|_U$
正合。
\end{proof}

\begin{lemma}
\label{sites-modules-lemma-j-shriek-flat}
设 $\mathcal{C}$ 是范畴，$\mathcal{O}$ 是环预层，且 $U$ 是
$\mathcal{C}$ 的对象。考虑函子 $j_U : \mathcal{C}/U \to \mathcal{C}$。
\begin{enumerate}
\item $\mathcal{O}$-模预层 $j_{U!}\mathcal{O}_U$ 是平坦的（见注
\ref{sites-modules-remark-localize-presheaves}）。
\item 如果 $\mathcal{C}$ 是站点且 $\mathcal{O}$ 是环层，则
$j_{U!}\mathcal{O}_U$ 是平坦 $\mathcal{O}$-模层。
\end{enumerate}
\end{lemma}

\begin{proof}
证明 (1)。由注 \ref{sites-modules-remark-localize-presheaves} 中的讨论可知
$$
j_{U!}\mathcal{O}_U(V)
=
\bigoplus\nolimits_{\varphi \in \Mor_\mathcal{C}(V, U)}
\mathcal{O}(V)
$$
，它是平坦 $\mathcal{O}(V)$-模。因此 (1) 由引理
\ref{sites-modules-lemma-flatness-presheaves} 得到。又因为
$j_{U!}\mathcal{O}_U = (j_{U!}\mathcal{O}_U)^\#$（第一个 $j_{U!}$
作用于层，第二个作用于预层），故由引理
\ref{sites-modules-lemma-flatness-sheafification} 得到 (2)。
\end{proof}

\begin{lemma}
\label{sites-modules-lemma-module-quotient-flat}
设 $\mathcal{C}$ 是范畴，且 $\mathcal{O}$ 是环预层。
\begin{enumerate}
\item 任意 $\mathcal{O}$-模预层都是某个直和
$\bigoplus j_{U_i!}\mathcal{O}_{U_i}$ 的商。
\item 任意 $\mathcal{O}$-模预层都是某个平坦 $\mathcal{O}$-模预层的商。
\item 如果 $\mathcal{C}$ 是站点且 $\mathcal{O}$ 是环层，则任意
$\mathcal{O}$-模层都是某个直和
$\bigoplus j_{U_i!}\mathcal{O}_{U_i}$ 的商。
\item 如果 $\mathcal{C}$ 是站点且 $\mathcal{O}$ 是环层，则任意
$\mathcal{O}$-模层都是某个平坦 $\mathcal{O}$-模层的商。
\end{enumerate}
\end{lemma}

\begin{proof}
证明 (1)。对 $\mathcal{C}$ 的每个对象 $U$ 以及每个
$s \in \mathcal{F}(U)$，得到态射 $j_{U!}\mathcal{O}_U \to \mathcal{F}$，
即态射 $\mathcal{O}_U \to \mathcal{F}|_U$，$1 \mapsto s$ 的伴随态射。
显然，映射
$$
\bigoplus\nolimits_{(U, s)} j_{U!}\mathcal{O}_U
\longrightarrow
\mathcal{F}
$$
是满射。结合引理 \ref{sites-modules-lemma-colimits-flat} 和
\ref{sites-modules-lemma-j-shriek-flat} 可知其源是平坦的，这就证明了 (2)。层的情形可通过
层化得出，也可重复同一论证。
\end{proof}

\begin{lemma}
\label{sites-modules-lemma-flat-tor-zero}
设 $\mathcal{C}$ 是范畴，$\mathcal{O}$ 是环预层。设
$$
0 \to \mathcal{F}'' \to \mathcal{F}' \to \mathcal{F} \to 0
$$
是 $\mathcal{O}$-模预层的短正合列，且 $\mathcal{G}$ 是
$\mathcal{O}$-模预层。
\begin{enumerate}
\item 如果 $\mathcal{F}$ 是平坦模预层，则序列
$$
0 \to
\mathcal{F}'' \otimes_{p, \mathcal{O}} \mathcal{G} \to
\mathcal{F}' \otimes_{p, \mathcal{O}} \mathcal{G} \to
\mathcal{F} \otimes_{p, \mathcal{O}} \mathcal{G} \to 0
$$
正合。
\item 如果 $\mathcal{C}$ 是站点，$\mathcal{O}$、$\mathcal{F}$、
$\mathcal{F}'$、$\mathcal{F}''$ 和 $\mathcal{G}$ 都是层，并且
$\mathcal{F}$ 作为模层是平坦的，则序列
$$
0 \to
\mathcal{F}'' \otimes_\mathcal{O} \mathcal{G} \to
\mathcal{F}' \otimes_\mathcal{O} \mathcal{G} \to
\mathcal{F} \otimes_\mathcal{O} \mathcal{G} \to 0
$$
正合。
\end{enumerate}
\end{lemma}

\begin{proof}
选择满射到 $\mathcal{G}$ 的平坦 $\mathcal{O}$-模预层 $\mathcal{G}'$；
由引理 \ref{sites-modules-lemma-module-quotient-flat}，这是可能的。令
$\mathcal{G}'' = \Ker(\mathcal{G}' \to \mathcal{G})$。对下图应用蛇引理即可：
$$
\begin{matrix}
 & & 0 & & 0 & & 0 & & \\
 & & \uparrow & & \uparrow & & \uparrow & & \\
 & & \mathcal{F}'' \otimes_{p, \mathcal{O}} \mathcal{G} & \to &
     \mathcal{F}' \otimes_{p, \mathcal{O}} \mathcal{G} & \to &
     \mathcal{F} \otimes_{p, \mathcal{O}} \mathcal{G} & \to & 0 \\
 & & \uparrow & & \uparrow & & \uparrow & & \\
0 & \to & \mathcal{F}'' \otimes_{p, \mathcal{O}} \mathcal{G}' & \to &
          \mathcal{F}' \otimes_{p, \mathcal{O}} \mathcal{G}' & \to &
          \mathcal{F} \otimes_{p, \mathcal{O}} \mathcal{G}' & \to & 0 \\
 & & \uparrow & & \uparrow & & \uparrow & & \\
 & & \mathcal{F}'' \otimes_{p, \mathcal{O}} \mathcal{G}'' & \to &
     \mathcal{F}' \otimes_{p, \mathcal{O}} \mathcal{G}'' & \to &
     \mathcal{F} \otimes_{p, \mathcal{O}} \mathcal{G}'' & \to & 0 \\
 & & & & & & \uparrow & & \\
 & & & & & & 0 & &
\end{matrix}
$$
其中各行各列均正合。中间一行正合，是因为与平坦模 $\mathcal{G}'$
作张量积是正合的。层的情形证明完全相同。
\end{proof}

\begin{lemma}
\label{sites-modules-lemma-flat-ses}
设 $\mathcal{C}$ 是范畴，$\mathcal{O}$ 是环预层。设
$$
0 \to
\mathcal{F}_2 \to
\mathcal{F}_1 \to
\mathcal{F}_0 \to 0
$$
是 $\mathcal{O}$-模预层的短正合列。
\begin{enumerate}
\item 如果 $\mathcal{F}_2$ 和 $\mathcal{F}_0$ 平坦，则
$\mathcal{F}_1$ 也平坦。
\item 如果 $\mathcal{F}_1$ 和 $\mathcal{F}_0$ 平坦，则
$\mathcal{F}_2$ 也平坦。
\end{enumerate}
如果 $\mathcal{C}$ 是站点且 $\mathcal{O}$ 是环层，则同一结论在
$\textit{Mod}(\mathcal{O})$ 中成立。
\end{lemma}

\begin{proof}
设 $\mathcal{G}^\bullet$ 是 $\mathcal{O}$-模预层的任意正合复形。假设
$\mathcal{F}_0$ 平坦。由引理 \ref{sites-modules-lemma-flat-tor-zero} 可知
$$
0 \to
\mathcal{G}^\bullet \otimes_{p, \mathcal{O}} \mathcal{F}_2 \to
\mathcal{G}^\bullet \otimes_{p, \mathcal{O}} \mathcal{F}_1 \to
\mathcal{G}^\bullet \otimes_{p, \mathcal{O}} \mathcal{F}_0 \to 0
$$
是 $\mathcal{O}$-模预层复形的短正合列。因此 (1) 和 (2) 由蛇引理得出。
模层的情形以相同方式证明。
\end{proof}

\begin{lemma}
\label{sites-modules-lemma-flat-resolution-of-flat}
设 $\mathcal{C}$ 是范畴，$\mathcal{O}$ 是环预层。设
$$
\ldots \to
\mathcal{F}_2 \to
\mathcal{F}_1 \to
\mathcal{F}_0 \to
\mathcal{Q} \to 0
$$
是 $\mathcal{O}$-模预层的正合复形。如果 $\mathcal{Q}$ 和所有
$\mathcal{F}_i$ 都是平坦 $\mathcal{O}$-模，则对任意
$\mathcal{O}$-模预层 $\mathcal{G}$，复形
$$
\ldots \to
\mathcal{F}_2 \otimes_{p, \mathcal{O}} \mathcal{G} \to
\mathcal{F}_1 \otimes_{p, \mathcal{O}} \mathcal{G} \to
\mathcal{F}_0 \otimes_{p, \mathcal{O}} \mathcal{G} \to
\mathcal{Q} \otimes_{p, \mathcal{O}} \mathcal{G} \to 0
$$
也正合。如果 $\mathcal{C}$ 是站点且 $\mathcal{O}$ 是环层，则同一结论在
$\textit{Mod}(\mathcal{O})$ 中成立。
\end{lemma}

\begin{proof}
本结论由引理 \ref{sites-modules-lemma-flat-tor-zero} 得出：将复形拆成短正合列，并利用引理
\ref{sites-modules-lemma-flat-ses} 归纳证明
$\Im(\mathcal{F}_{i + 1} \to \mathcal{F}_i)$ 平坦。
\end{proof}

\begin{lemma}
\label{sites-modules-lemma-tensor-flats}
设 $(\mathcal{C}, \mathcal{O})$ 是环站点。如果 $\mathcal{G}$ 和
$\mathcal{F}$ 是平坦 $\mathcal{O}$-模，则
$\mathcal{G} \otimes_\mathcal{O} \mathcal{F}$ 是平坦 $\mathcal{O}$-模。
\end{lemma}

\begin{proof}
这是因为
$$
(\mathcal{G} \otimes_\mathcal{O} \mathcal{F}) \otimes_\mathcal{O} \mathcal{H}
=
\mathcal{G} \otimes_\mathcal{O} (\mathcal{F} \otimes_\mathcal{O} \mathcal{H})
$$
，而正合函子的复合仍然正合。
\end{proof}

\begin{lemma}
\label{sites-modules-lemma-flat-change-of-rings}
设 $\mathcal{O}_1 \to \mathcal{O}_2$ 是站点 $\mathcal{C}$ 上的环层映射。
如果 $\mathcal{G}$ 是平坦 $\mathcal{O}_1$-模，则
$\mathcal{G} \otimes_{\mathcal{O}_1} \mathcal{O}_2$
是平坦 $\mathcal{O}_2$-模。
\end{lemma}

\begin{proof}
这是因为
$$
(\mathcal{G} \otimes_{\mathcal{O}_1} \mathcal{O}_2)
\otimes_{\mathcal{O}_2} \mathcal{H}
=
\mathcal{G} \otimes_{\mathcal{O}_1} \mathcal{F}
$$
（例如作为阿贝尔群层）。
\end{proof}

\noindent
下述引理是平坦性的方程判别准则（《代数》，引理
\ref{algebra-lemma-flat-eq}）的类似版本。

\begin{lemma}
\label{sites-modules-lemma-flat-eq}
设 $(\mathcal{C}, \mathcal{O})$ 是环站点，且 $\mathcal{F}$ 是
$\mathcal{O}$-模。下列条件等价：
\begin{enumerate}
\item $\mathcal{F}$ 是平坦 $\mathcal{O}$-模。
\item 设 $U$ 是 $\mathcal{C}$ 的对象，并设
$$
\mathcal{O}_U \xrightarrow{(f_1, \ldots, f_n)}
\mathcal{O}_U^{\oplus n} \xrightarrow{(s_1, \ldots, s_n)}
\mathcal{F}|_U
$$
是 $\mathcal{O}_U$-模复形。则存在覆盖 $\{U_i \to U\}$，且对每个 $i$
存在分解
$$
\mathcal{O}_{U_i}^{\oplus n}
\xrightarrow{B_i}
\mathcal{O}_{U_i}^{\oplus l_i} \xrightarrow{(t_{i1}, \ldots, t_{il_i})}
\mathcal{F}|_{U_i}
$$
，它分解 $(s_1, \ldots, s_n)|_{U_i}$，并且
$B_i \circ (f_1, \ldots, f_n)|_{U_i} = 0$。
\item 设 $U$ 是 $\mathcal{C}$ 的对象，并设
$$
\mathcal{O}_U^{\oplus m} \xrightarrow{A}
\mathcal{O}_U^{\oplus n} \xrightarrow{(s_1, \ldots, s_n)}
\mathcal{F}|_U
$$
是 $\mathcal{O}_U$-模复形。则存在覆盖 $\{U_i \to U\}$，且对每个 $i$
存在分解
$$
\mathcal{O}_{U_i}^{\oplus n}
\xrightarrow{B_i}
\mathcal{O}_{U_i}^{\oplus l_i} \xrightarrow{(t_{i1}, \ldots, t_{il_i})}
\mathcal{F}|_{U_i}
$$
，它分解 $(s_1, \ldots, s_n)|_{U_i}$，并且
$B_i \circ A|_{U_i} = 0$。
\end{enumerate}
\end{lemma}

\begin{proof}
假设 (1)。设 $\mathcal{I} \subset \mathcal{O}_U$ 是由
$f_1, \ldots, f_n$ 生成的理想层。则 $\sum f_j \otimes s_j$ 是
$\mathcal{I} \otimes_{\mathcal{O}_U} \mathcal{F}|_U$ 的截面，并在
$\mathcal{F}|_U$ 中映为零。因为 $\mathcal{F}|_U$ 平坦（引理
\ref{sites-modules-lemma-restriction-flat}），映射
$\mathcal{I} \otimes_{\mathcal{O}_U} \mathcal{F}|_U \to \mathcal{F}|_U$
是单射。由于 $\mathcal{I} \otimes_{\mathcal{O}_U} \mathcal{F}|_U$
是预层张量积的关联层，故存在覆盖 $\{U_i \to U\}$，使得
$\sum f_j|_{U_i} \otimes s_j|_{U_i}$ 在
$\mathcal{I}(U_i) \otimes_{\mathcal{O}(U_i)} \mathcal{F}(U_i)$.
中为零。利用《代数》，引理 \ref{algebra-lemma-relations} 展开定义，可找到
$t_{i1}, \ldots, t_{i l_i} \in \mathcal{F}(U_i)$ 和
$a_{ijk} \in \mathcal{O}(U_i)$
，使得 $\sum_j a_{ijk}f_j|_{U_i} = 0$ 且
$s_j|_{U_i} = \sum_k a_{ijk}t_{ik}$。故 (2) 成立。

\medskip\noindent
假设 (2)。给定 (3) 中的 $U$、$n$、$m$、$A$ 以及
$s_1, \ldots, s_n$。注意 $A$ 有 $m$ 列。对 $m$ 归纳证明 (3) 的断言。
在初始情形 $m = 0$，可取覆盖 $\{U \to U\}$，并使用分解
$\mathcal{O}_U^n \to \mathcal{O}_U^n \to \mathcal{F}|_U$，其中第一个映射是
恒等映射，第二个是 $(s_1, \ldots, s_n)$。假设 $m > 0$。令
$(f_1, \ldots, f_n)$ 为 $A$ 的最后一列，并应用 (2)。这给出新图
$$
\mathcal{O}_{U_i}^{\oplus m} \xrightarrow{B_i \circ A|_{U_i}}
\mathcal{O}_{U_i}^{\oplus l_i} \xrightarrow{(t_{i1}, \ldots, t_{il_i})}
\mathcal{F}|_{U_i}
$$
，但 $A_i = B_i \circ A|_{U_i}$ 的最后一列为零。因此可将归纳假设用于
$U_i$、$l_i$、$m - 1$、由 $A_i$ 前 $m - 1$ 列组成的矩阵以及
$t_{i1}, \ldots, t_{il_i}$，得到覆盖 $\{U_{ij} \to U_j\}$ 和分解
$$
\mathcal{O}_{U_{ij}}^{\oplus l_i}
\xrightarrow{C_{ij}}
\mathcal{O}_{U_{ij}}^{\oplus k_{ij}}
\xrightarrow{(v_{ij1}, \ldots, v_{ij k_{ij}})}
\mathcal{F}|_{U_{ij}}
$$
，它分解 $(t_{i1}, \ldots, t_{il_i})|_{U_{ij}}$，并且
$C_i \circ B_i|_{U_{ij}} \circ A|_{U_{ij}} = 0$。于是
$\{U_{ij} \to U\}$ 是覆盖，使用
$B_{ij} = C_i \circ B_i|_{U_{ij}}$ 和 $v_{ija}$ 即得到所需分解。
由此可见 (2) 蕴含 (3)。

\medskip\noindent
假设 (3)。设 $\mathcal{G} \to \mathcal{H}$ 是 $\mathcal{O}$-模的单同态。
需要证明
$\mathcal{G} \otimes_\mathcal{O} \mathcal{F} \to
\mathcal{H} \otimes_\mathcal{O} \mathcal{F}$
是单射。设 $U$ 是 $\mathcal{C}$ 的对象，并设
$s \in (\mathcal{G} \otimes_\mathcal{O} \mathcal{F})(U)$ 是在
$\mathcal{H} \otimes_\mathcal{O} \mathcal{F}$ 中映为零的截面。需要证明
$s$ 为零。由于
$\mathcal{G} \otimes_\mathcal{O} \mathcal{F}$
是层，只需找到 $\mathcal{C}$ 的覆盖 $\{U_i \to U\}_{i \in I}$，使得对所有
$i \in I$，$s|_{U_i}$ 为零。因此总可用某个覆盖的成员替换 $U$。特别地，
由于 $\mathcal{G} \otimes_\mathcal{O} \mathcal{F}$ 是
$\mathcal{G} \otimes_{p, \mathcal{O}} \mathcal{F}$ 的层化，可以假设 $s$ 是某个
$s' \in 
\mathcal{G}(U) \otimes_{\mathcal{O}(U)} \mathcal{F}(U)$.
的像。对 $\mathcal{H} \otimes_\mathcal{O} \mathcal{F}$ 作类似论证，可以假设
$s'$ 在 $\mathcal{H}(U) \otimes_{\mathcal{O}(U)} \mathcal{F}(U)$ 中映为零。
将 $\mathcal{F}(U) = \colim M_\alpha$ 写成有限表示
$\mathcal{O}(U)$-模 $M_\alpha$ 的滤过余极限（《代数》，引理
\ref{algebra-lemma-module-colimit-fp}）。由于张量积与滤过余极限交换
（《代数》，引理 \ref{algebra-lemma-tensor-products-commute-with-limits}），
可以选择 $\alpha$，使得 $s'$ 来自某个
$s'' \in \mathcal{G}(U) \otimes_{\mathcal{O}(U)} M_\alpha$，且 $s''$ 在
$\mathcal{H}(U) \otimes_{\mathcal{O}(U)} M_\alpha$ 中映为零。
固定 $\alpha$ 和 $s''$，并选择一个表示
$$
\mathcal{O}(U)^{\oplus m} \xrightarrow{A} \mathcal{O}(U)^{\oplus n}
\to M_\alpha \to 0
$$
。对相应的 $\mathcal{O}_U$-模复形应用 (3)：
$$
\mathcal{O}_U^{\oplus m} \xrightarrow{A}
\mathcal{O}_U^{\oplus n} \xrightarrow{(s_1, \ldots, s_n)}
\mathcal{F}|_U
$$
。用覆盖 $U_i$ 的成员替换 $U$ 后，可知映射
$$
M_\alpha \to \mathcal{F}(U)
$$
分解通过某个自由模 $\mathcal{O}(U)^{\oplus l}$。由于
$\mathcal{G}(U) \to \mathcal{H}(U)$ 是单射，可知
$$
\mathcal{G}(U) \otimes_{\mathcal{O}(U)} \mathcal{O}(U)^{\oplus l}
\to
\mathcal{H}(U) \otimes_{\mathcal{O}(U)} \mathcal{O}(U)^{\oplus l}
$$
也是单射。因此既然 $s''$ 在右侧模中映为零，它在左侧模中也映为零，
即 $s$ 如所需为零。
\end{proof}

\begin{lemma}
\label{sites-modules-lemma-flat-over-thickening}
设 $\mathcal{C}$ 是站点。设 $\mathcal{O}' \to \mathcal{O}$ 是环层的满射，
其核 $\mathcal{I}$ 是平方为零的理想。设 $\mathcal{F}'$ 是
$\mathcal{O}'$-模，并置
$\mathcal{F} = \mathcal{F}'/\mathcal{I}\mathcal{F}'$。下列条件等价：
\begin{enumerate}
\item $\mathcal{F}'$ 是平坦 $\mathcal{O}'$-模；
\item $\mathcal{F}$ 是平坦 $\mathcal{O}$-模，并且
$\mathcal{I} \otimes_\mathcal{O} \mathcal{F} \to \mathcal{F}'$
是单射。
\end{enumerate}
\end{lemma}

\begin{proof}
如果 (1) 成立，则由引理 \ref{sites-modules-lemma-flat-change-of-rings}，
$\mathcal{F} = \mathcal{F}' \otimes_{\mathcal{O}'} \mathcal{O}$ 在
$\mathcal{O}$ 上平坦。对正合列
$0 \to \mathcal{I} \to \mathcal{O}' \to \mathcal{O} \to 0$
应用 $- \otimes_{\mathcal{O}'} \mathcal{F}'$，可知映射
$\mathcal{I} \otimes_\mathcal{O} \mathcal{F} \to \mathcal{F}'$
是单射；见引理 \ref{sites-modules-lemma-flat-tor-zero}。假设 (2)。在以下证明中将不再说明地使用：
对任意被 $\mathcal{I}$ 零化的 $\mathcal{O}'$-模 $\mathcal{K}$，有
$\mathcal{K} \otimes_{\mathcal{O}'} \mathcal{F}' =
\mathcal{K} \otimes_\mathcal{O} \mathcal{F}$
。设 $\alpha : \mathcal{G}' \to \mathcal{H}'$ 是
$\mathcal{O}'$-模的单射。分别令 $\mathcal{G} \subset \mathcal{G}'$、
$\mathcal{H} \subset \mathcal{H}'$ 为由被 $\mathcal{I}$ 零化的截面组成的子层。
考虑交换图
$$
\xymatrix{
\mathcal{G} \otimes_{\mathcal{O}'} \mathcal{F}' \ar[r] \ar[d] &
\mathcal{G}' \otimes_{\mathcal{O}'} \mathcal{F}' \ar[r] \ar[d] &
\mathcal{G}'/\mathcal{G} \otimes_{\mathcal{O}'} \mathcal{F}' \ar[r] \ar[d] &
0 \\
\mathcal{H} \otimes_{\mathcal{O}'} \mathcal{F}' \ar[r] &
\mathcal{H}' \otimes_{\mathcal{O}'} \mathcal{F}' \ar[r] &
\mathcal{H}'/\mathcal{H} \otimes_{\mathcal{O}'} \mathcal{F}' \ar[r] & 0
}
$$
注意，$\mathcal{G}'/\mathcal{G}$ 和 $\mathcal{H}'/\mathcal{H}$
都被 $\mathcal{I}$ 零化，且
$\mathcal{G}'/\mathcal{G} \to \mathcal{H}'/\mathcal{H}$ 是单射。
由于 $\mathcal{F}$ 在 $\mathcal{O}$ 上平坦，右侧竖直箭头是单射；左侧竖直箭头
也如此。因此中间的竖直箭头是单射，故 $\mathcal{F}'$ 平坦。
\end{proof}

\begin{lemma}
\label{sites-modules-lemma-neighbourhood-isomorphism}
设 $\mathcal{C}$ 是站点。设 $\mathcal{O} \to \mathcal{O}'$
是环层的平坦同态。设
$\mathcal{I} \subset \mathcal{O}$
是理想层，使得诱导映射
$\mathcal{O}/\mathcal{I} \to \mathcal{O}'/\mathcal{I}\mathcal{O}'$
是同构。对任意被某个 $\mathcal{I}^n$（$n \geq 0$）零化的
$\mathcal{O}$-模 $\mathcal{F}$，映射
$\text{id} \otimes 1 :
\mathcal{F} \to \mathcal{F} \otimes_\mathcal{O} \mathcal{O}'$
是同构。
\end{lemma}

\begin{proof}
略。提示：见《代数详论》，引理
\ref{more-algebra-lemma-neighbourhood-isomorphism}。
\end{proof}





\section{对偶}
\label{sites-modules-section-duals}

\noindent
设 $(\mathcal{C}, \mathcal{O})$ 是环站点。赋予第
\ref{sites-modules-section-tensor-product} 节构造的张量积后，$\mathcal{O}$-模范畴成为
对称幺半范畴。对 $\mathcal{O}$-模 $\mathcal{F}$，下列条件等价：
\begin{enumerate}
\item $\mathcal{F}$ 在 $\mathcal{O}$-模幺半范畴中有左对偶；
\item 对 $\mathcal{C}$ 的每个对象 $U$，存在覆盖 $\{U_i \to U\}$，使得
$\mathcal{F}|_{U_i}$ 是某个有限自由 $\mathcal{O}|_{U_i}$-模的直和项；
\item $\mathcal{F}$ 是有限表示且平坦的 $\mathcal{O}$-模。
\end{enumerate}
这将在本节的例 \ref{sites-modules-example-dual} 以及引理
\ref{sites-modules-lemma-left-dual-module} 和
\ref{sites-modules-lemma-flat-locally-finite-presentation} 中证明。

\begin{example}
\label{sites-modules-example-dual}
设 $(\mathcal{C}, \mathcal{O})$ 是环站点。设 $\mathcal{F}$ 是
$\mathcal{O}$-模，并且对 $\mathcal{C}$ 的每个对象 $U$，都存在覆盖
$\{U_i \to U\}$，使得 $\mathcal{F}|_{U_i}$ 是某个有限自由
$\mathcal{O}|_{U_i}$-模的直和项。则映射
$$
\mathcal{F} \otimes_\mathcal{O}
\SheafHom_\mathcal{O}(\mathcal{F}, \mathcal{O})
\longrightarrow
\SheafHom_\mathcal{O}(\mathcal{F}, \mathcal{F})
$$
是同构。事实上，这是局部问题；当 $\mathcal{F}$ 有限自由时成立，并且若对某个模
成立，则对它的任一直和项也成立（略去细节）。记
$$
\eta :
\mathcal{O}
\longrightarrow
\mathcal{F} \otimes_\mathcal{O}
\SheafHom_\mathcal{O}(\mathcal{F}, \mathcal{O})
$$
为把 $1$ 送到上述同构下对应于 $\text{id}_\mathcal{F}$ 的截面的映射。记
$$
\epsilon : 
\SheafHom_\mathcal{O}(\mathcal{F}, \mathcal{O})
\otimes_\mathcal{O} \mathcal{F}
\longrightarrow
\mathcal{O}
$$
为求值映射。于是可见
$\SheafHom_\mathcal{O}(\mathcal{F}, \mathcal{O}), \eta, \epsilon$
是《范畴》，定义 \ref{categories-definition-dual} 意义下
$\mathcal{F}$ 的左对偶。略去下列等式的验证：
$(1 \otimes \epsilon) \circ (\eta \otimes 1) = \text{id}_\mathcal{F}$
以及
$(\epsilon \otimes 1) \circ (1 \otimes \eta) =
\text{id}_{\SheafHom_\mathcal{O}(\mathcal{F}, \mathcal{O})}$.
\end{example}

\begin{lemma}
\label{sites-modules-lemma-left-dual-module}
设 $(\mathcal{C}, \mathcal{O})$ 是环站点，且 $\mathcal{F}$ 是
$\mathcal{O}$-模。设 $\mathcal{G}, \eta, \epsilon$ 是
$\mathcal{F}$ 在 $\mathcal{O}$-模幺半范畴中的左对偶；见《范畴》，定义
\ref{categories-definition-dual}。则
\begin{enumerate}
\item 对 $\mathcal{C}$ 的每个对象 $U$，存在覆盖 $\{U_i \to U\}$，使得
$\mathcal{F}|_{U_i}$ 是某个有限自由 $\mathcal{O}|_{U_i}$-模的直和项；
\item 映射
$e : \SheafHom_\mathcal{O}(\mathcal{F}, \mathcal{O}) \to \mathcal{G}$
把局部截面 $\lambda$ 送到 $(\lambda \otimes 1)(\eta)$，它是同构；
\item 对 $\mathcal{F}$ 和 $\mathcal{G}$ 的局部截面 $f$ 和 $g$，有
$\epsilon(f, g) = e^{-1}(g)(f)$。
\end{enumerate}
\end{lemma}

\begin{proof}
这些假设意味着
$$
\mathcal{F} \xrightarrow{\eta \otimes 1}
\mathcal{F} \otimes_\mathcal{O} \mathcal{G}
\otimes_\mathcal{O} \mathcal{F}
\xrightarrow{1 \otimes \epsilon} \mathcal{F}
\quad\text{且}\quad
\mathcal{G} \xrightarrow{1 \otimes \eta}
\mathcal{G} \otimes_\mathcal{O} \mathcal{F}
\otimes_\mathcal{O} \mathcal{G}
\xrightarrow{\epsilon \otimes 1} \mathcal{G}
$$
都是恒等映射。设 $U$ 是 $\mathcal{C}$ 的对象。用 $U$ 的某个覆盖的成员替换
$U$ 后，可以找到 $\mathcal{F}$ 和 $\mathcal{G}$ 在 $U$ 上的有限多个截面
$f_1, \ldots, f_n$ 和 $g_1, \ldots, g_n$，使得
$\eta(1) = \sum f_i g_i$。记
$$
\mathcal{O}_U^{\oplus n} \to \mathcal{F}|_U
$$
为把第 $i$ 个基向量送到 $f_i$ 的映射。于是映射 $\eta|_U$ 分解通过映射
$\tilde \eta : \mathcal{O}_U \to
\mathcal{O}_U^{\oplus n} \otimes_{\mathcal{O}_U} \mathcal{G}|_U$。
得到交换图
$$
\xymatrix{
\mathcal{F}|_U
\ar[rr]_-{\eta \otimes 1} \ar[rrd]_-{\tilde \eta \otimes 1} & &
\mathcal{F}|_U \otimes \mathcal{G}|_U \otimes \mathcal{F}|_U
\ar[r]_-{1 \otimes \epsilon} &
\mathcal{F}|_U \\
& &
\mathcal{O}_U^{\oplus n} \otimes \mathcal{G}|_U \otimes \mathcal{F}|_U
\ar[u] \ar[r]^-{1 \otimes \epsilon} &
\mathcal{O}_U^{\oplus n} \ar[u]
}
$$
这表明 $\mathcal{F}|_U$ 上的恒等映射分解通过有限自由
$\mathcal{O}_U$-模，从而证明了 (1)。(2) 由《范畴》，引理
\ref{categories-lemma-left-dual} 及其证明得到。(3) 由本证明中的第一个等式得到。
也可以从左对偶的唯一性（《范畴》，注
\ref{categories-remark-left-dual-adjoint}）以及例
\ref{sites-modules-example-dual} 中左对偶的构造推出 (2) 和 (3)。
\end{proof}

\begin{lemma}
\label{sites-modules-lemma-flat-locally-finite-presentation}
设 $(\mathcal{C}, \mathcal{O})$ 是环站点。设 $\mathcal{F}$ 局部有限表示且
平坦。则给定 $\mathcal{C}$ 的对象 $U$，存在覆盖 $\{U_i \to U\}$，使得
$\mathcal{F}|_{U_i}$ 是某个有限自由 $\mathcal{O}_{U_i}$-模的直和项。
\end{lemma}

\begin{proof}
选择 $\mathcal{C}$ 的对象 $U$。用某个覆盖的成员替换 $U$ 后，可以假设存在表示
$$
\mathcal{O}_U^{\oplus r} \to
\mathcal{O}_U^{\oplus n} \to \mathcal{F}|_U \to 0
$$
。由引理 \ref{sites-modules-lemma-flat-eq}，再次用某个覆盖的成员替换 $U$ 后，可以假设存在分解
$$
\mathcal{O}_U^{\oplus n} \to
\mathcal{O}_U^{\oplus n_1} \to \mathcal{F}|_U
$$
，使得复合
$\mathcal{O}_U^{\oplus r} \to \mathcal{O}_U^{\oplus n}
\to \mathcal{O}_U^{\oplus n_r}$ 为零。这意味着满射
$\mathcal{O}_U^{\oplus n_r} \to \mathcal{F}|_U$ 有截面，结论成立。
\end{proof}











\section{走向可构造模层}
\label{sites-modules-section-constructible}

\noindent
回顾一下，站点的拟紧对象粗略说来是其每个覆盖都可由有限覆盖加细的对象
（实际定义略为复杂；见《站点》，第 \ref{sites-section-quasi-compact} 节）。
事实表明，如果站点的每个对象都有由拟紧对象组成的覆盖，那么当 $U$ 拟紧时，
模层 $j_!\mathcal{O}_U$ 构成全体模层范畴的一组特别良好的生成元。

\begin{lemma}
\label{sites-modules-lemma-covering-gives-surjection}
设 $(\mathcal{C}, \mathcal{O})$ 是环站点，且 $\{U_i \to U\}$ 是
$\mathcal{C}$ 的覆盖。则序列
$$
\bigoplus j_{U_i \times_U U_j!}\mathcal{O}_{U_i \times_U U_j} \to
\bigoplus j_{U_i!}\mathcal{O}_{U_i} \to j_!\mathcal{O}_U \to 0
$$
正合。
\end{lemma}

\begin{proof}
对任意 $\mathcal{O}$-模 $\mathcal{F}$，函子
$\Hom_\mathcal{O}(-, \mathcal{F})$ 将上述序列变为正合列
$0 \to \mathcal{F}(U) \to \prod \mathcal{F}(U_i) \to
\prod \mathcal{F}(U_i \times_U U_j)$；见
(\ref{sites-modules-equation-map-lower-shriek-OU-into-module})。再由《同调代数》，引理
\ref{homology-lemma-check-exactness} 即得本引理。
\end{proof}

\begin{lemma}
\label{sites-modules-lemma-silly-quasi-compact}
设 $(\mathcal{C}, \mathcal{O})$ 是环站点。设
$\mathcal{U} = \{U_i \to U\}_{i \in I}$ 是 $\mathcal{C}$ 的覆盖。
如果 $U$ 拟紧，则存在有限子集 $I' \subset I$，使得序列
$$
\bigoplus\nolimits_{i, i' \in I'}
j_{U_i \times_U U_{i'}!}\mathcal{O}_{U_i \times_U U_{i'}} \to
\bigoplus\nolimits_{i \in I'}
j_{U_i!}\mathcal{O}_{U_i} \to
j_!\mathcal{O}_U \to 0
$$
正合。
\end{lemma}

\begin{proof}
本引理由引理 \ref{sites-modules-lemma-covering-gives-surjection} 立即得到，只要 $U$ 满足
《站点》，引理 \ref{sites-lemma-conclude-quasi-compact} 的条件 (3)。
我们建议读者跳过一般情形的证明。由定义，存在覆盖
$\mathcal{V} = \{V_j \to U\}_{j \in J}$ 以及具有固定靶的映射族的态射
$\mathcal{V} \to \mathcal{U}$，它由 $\text{id} : U \to U$、
$\alpha : J \to I$ 和
$f_j : V_j \to U_{\alpha(j)}$
给出（见《站点》，定义 \ref{sites-definition-morphism-coverings}），并且
$\alpha$ 的像 $I' \subset I$ 有限。由《同调代数》，引理
\ref{homology-lemma-check-exactness}，只需证明对任意 $\mathcal{O}$-模层
$\mathcal{F}$，函子 $\Hom_\mathcal{O}(-, \mathcal{F})$ 将本引理的序列变为
正合列。由
(\ref{sites-modules-equation-map-lower-shriek-OU-into-module})
得到通常的序列
$$
0 \to
\mathcal{F}(U) \to
\prod\nolimits_{i \in I'} \mathcal{F}(U_i) \to
\prod\nolimits_{i, i' \in I'} \mathcal{F}(U_i \times_U U_{i'})
$$
。对由覆盖 $\mathcal{V}$ 加细的映射族
$\{U_i \to U\}_{i \in I'}$ 应用《站点》，引理
\ref{sites-lemma-compare-sheaf-condition}，可知这是正合列。
\end{proof}

\begin{lemma}
\label{sites-modules-lemma-sections-over-quasi-compact}
设 $\mathcal{C}$ 是站点，且 $W$ 是 $\mathcal{C}$ 的拟紧对象。
\begin{enumerate}
\item 函子 $\Sh(\mathcal{C}) \to \textit{Sets}$，
$\mathcal{F} \mapsto \mathcal{F}(W)$ 与余积交换。
\item 设 $\mathcal{O}$ 是 $\mathcal{C}$ 上的环层。函子
$\textit{Mod}(\mathcal{O}) \to \textit{Ab}$,
$\mathcal{F} \mapsto \mathcal{F}(W)$
与直和交换。
\end{enumerate}
\end{lemma}

\begin{proof}
证明 (1)。由《站点》，引理 \ref{sites-lemma-directed-colimits-sections}，
取 $W$ 上截面与转移映射为单射的滤过余极限交换。设 $\mathcal{F}_i$
是由集合 $I$ 标号的集合层族。则 $\coprod \mathcal{F}_i$ 是余积
$\mathcal{F}_E = \coprod_{i \in E} \mathcal{F}_i$ 在有限子集
$E \subset I$ 的偏序集上的滤过余极限。由于转移映射为单射，结论成立。

\medskip\noindent
证明 (2)。设 $\mathcal{F}_i$ 是由集合 $I$ 标号的 $\mathcal{O}$-模层族。
则 $\bigoplus \mathcal{F}_i$ 是直和
$\mathcal{F}_E = \bigoplus_{i \in E} \mathcal{F}_i$ 在有限子集
$E \subset I$ 的偏序集上的滤过余极限。阿贝尔层的滤过余极限可在集合层范畴中
计算。此外，对 $E \subset E'$，转移映射
$\mathcal{F}_E \to \mathcal{F}_{E'}$ 是单射（因为层化正合，且在底层预层上
单射性显然）。因此由《站点》，引理
\ref{sites-lemma-directed-colimits-sections}，只需证明有限指标集的情形。
有限情形已由引理 \ref{sites-modules-lemma-limits-colimits-abelian-sheaves} 处理
（它在 $\mathcal{C}$ 的任意对象上都成立）。
\end{proof}

\begin{lemma}
\label{sites-modules-lemma-quasi-compact-hom-from}
设 $(\mathcal{C}, \mathcal{O})$ 是环站点，且 $U$ 是 $\mathcal{C}$ 的拟紧对象。
则函子 $\Hom_\mathcal{O}(j_!\mathcal{O}_U, -)$ 与直和交换。
\end{lemma}

\begin{proof}
这是因为由 (\ref{sites-modules-equation-map-lower-shriek-OU-into-module})，有
$\Hom_\mathcal{O}(j_!\mathcal{O}_U, \mathcal{F}) = \mathcal{F}(U)$
，而由引理 \ref{sites-modules-lemma-sections-over-quasi-compact}，函子
$\mathcal{F} \mapsto \mathcal{F}(U)$ 与直和交换。
\end{proof}

\noindent
为陈述下文尽可能精确的结果，引入一些记号。

\begin{situation}
\label{sites-modules-situation-quasi-compact-objects}
设 $\mathcal{C}$ 是站点，并设
$\mathcal{B} \subset \text{Ob}(\mathcal{C})$ 是对象集。考虑下列条件：
\begin{enumerate}
\item
\label{sites-modules-item-enough}
$\mathcal{C}$ 的每个对象都有由 $\mathcal{B}$ 中元素组成的覆盖。
\item
\label{sites-modules-item-enough-qc}
每个 $U \in \mathcal{B}$ 都是拟紧的
（《站点》，第 \ref{sites-section-quasi-compact} 节）。
\item
\label{sites-modules-item-enough-qc-qs}
对满足 $U_i, U \in \mathcal{B}$ 的覆盖 $\{U_i \to U\}$，纤维积
$U_i \times_U U_j$ 是拟紧的。
\end{enumerate}
\end{situation}

\begin{lemma}
\label{sites-modules-lemma-module-quotient-direct-sum}
在情形 \ref{sites-modules-situation-quasi-compact-objects} 中，假设
(\ref{sites-modules-item-enough}) 成立。
\begin{enumerate}
\item 每个集合层都是某个满射的靶，该满射的源是余积
$\coprod h_{U_i}^\#$，其中 $U_i$ 属于 $\mathcal{B}$。
\item 如果 $\mathcal{O}$ 是环层，则每个 $\mathcal{O}$-模都是某个直和
$\bigoplus\nolimits j_{U_i!}\mathcal{O}_{U_i}$ 的商，其中 $U_i$
属于 $\mathcal{B}$。
\end{enumerate}
\end{lemma}

\begin{proof}
第 (1) 部分由《站点》，引理
\ref{sites-lemma-sheaf-coequalizer-representable} 和
\ref{sites-lemma-covering-surjective-after-sheafification} 得出。
第 (2) 部分由引理 \ref{sites-modules-lemma-module-quotient-flat} 和
\ref{sites-modules-lemma-covering-gives-surjection} 得出。
\end{proof}

\begin{lemma}
\label{sites-modules-lemma-module-filtered-colimit-constructibles}
在情形 \ref{sites-modules-situation-quasi-compact-objects} 中，假设
(\ref{sites-modules-item-enough}) 和 (\ref{sites-modules-item-enough-qc}) 成立。
\begin{enumerate}
\item 每个集合层都是下列形式的层的滤过余极限：
\begin{equation}
\label{sites-modules-equation-towards-constructible-sets}
\text{余等化子}\left(
\xymatrix{
\coprod\nolimits_{j = 1, \ldots, m} h_{V_j}^\#
\ar@<1ex>[r] \ar@<-1ex>[r] &
\coprod\nolimits_{i = 1, \ldots, n} h_{U_i}^\#
}
\right)
\end{equation}
，其中 $U_i$ 和 $V_j$ 属于 $\mathcal{B}$。
\item 如果 $\mathcal{O}$ 是环层，则每个 $\mathcal{O}$-模都是下列形式的层的
滤过余极限：
\begin{equation}
\label{sites-modules-equation-towards-constructible}
\Coker\left(
\bigoplus\nolimits_{j = 1, \ldots, m} j_{V_j!}\mathcal{O}_{V_j}
\longrightarrow
\bigoplus\nolimits_{i = 1, \ldots, n} j_{U_i!}\mathcal{O}_{U_i}
\right)
\end{equation}
，其中 $U_i$ 和 $V_j$ 属于 $\mathcal{B}$。
\end{enumerate}
\end{lemma}

\begin{proof}
证明 (1)。由引理 \ref{sites-modules-lemma-module-quotient-direct-sum}，每个集合层
$\mathcal{F}$ 都是某个满射的靶，其源 $\mathcal{F}_0$ 是形如
$h_U^\#$（$U \in \mathcal{B}$）的层的余积。把这一点应用于
$\mathcal{F}_0 \times_\mathcal{F} \mathcal{F}_0$，可知 $\mathcal{F}$ 是一对映射
$$
\xymatrix{
\coprod\nolimits_{j \in J} h_{V_j}^\#
\ar@<1ex>[r] \ar@<-1ex>[r] &
\coprod\nolimits_{i \in I} h_{U_i}^\#
}
$$
的余等化子，其中 $I$、$J$ 是某些指标集，且 $V_j$ 和 $U_i$ 属于
$\mathcal{B}$。对每个有限子集 $J' \subset J$，存在有限子集
$I' \subset I$，使得 $j \in J'$ 上的余积经两个映射都映入
$i \in I'$ 上的余积；见《站点》，引理
\ref{sites-lemma-directed-colimits-sections}。（略去细节；提示：无限余积是有限子余积的
滤过余极限。）因此我们的层是这些有限余积之间映射的余核的余极限。

\medskip\noindent
证明 (2)。由引理 \ref{sites-modules-lemma-module-quotient-direct-sum}，每个模都是形如
$j_{U!}\mathcal{O}_U$（$U \in \mathcal{B}$）的模的某个直和的商。因此每个模
都是一个余核
$$
\Coker\left(
\bigoplus\nolimits_{j \in J} j_{V_j!}\mathcal{O}_{V_j}
\longrightarrow
\bigoplus\nolimits_{i \in I} j_{U_i!}\mathcal{O}_{U_i}
\right)
$$
，其中 $I$、$J$ 是某些指标集，且 $V_j$ 和 $U_i$ 属于 $\mathcal{B}$。
对每个有限子集 $J' \subset J$，存在有限子集 $I' \subset I$，使得
$j \in J'$ 上的直和映入 $i \in I'$ 上的直和；见引理
\ref{sites-modules-lemma-quasi-compact-hom-from}。因此我们的模是这些有限直和之间映射的余核的
余极限。
\end{proof}

\begin{lemma}
\label{sites-modules-lemma-cokernel-map-towards-constructibles}
在情形 \ref{sites-modules-situation-quasi-compact-objects} 中，假设
(\ref{sites-modules-item-enough}) 和 (\ref{sites-modules-item-enough-qc}) 成立。设 $\mathcal{O}$ 是环层。
则形如 (\ref{sites-modules-equation-towards-constructible}) 的模之间映射的余核，仍是形如
(\ref{sites-modules-equation-towards-constructible}) 的模。
\end{lemma}

\begin{proof}
设 $\mathcal{F} = \Coker(\bigoplus j_{V_j!}\mathcal{O}_{V_j} \to
\bigoplus j_{U_i!}\mathcal{O}_{U_i})$
，如 (\ref{sites-modules-equation-towards-constructible}) 所示。只需证明，当
$W \in \mathcal{B}$ 时，映射
$\varphi : j_{W!}\mathcal{O}_W \to \mathcal{F}$ 的余核仍是同类型的模。
映射 $\varphi$ 对应于 $s \in \mathcal{F}(W)$。由于
$\bigoplus j_{U_i!}\mathcal{O}_{U_i} \to \mathcal{F}$ 是满射，由
(\ref{sites-modules-item-enough}) 可选择满足 $W_k \in \mathcal{B}$ 的覆盖
$\{W_k \to W\}_{k \in K}$，使得 $s|_{W_k}$ 是
$\bigoplus j_{U_i!}\mathcal{O}_{U_i})$ 的某个截面 $s_k$ 的像。
由 (\ref{sites-modules-item-enough-qc})，对象 $W$ 拟紧。由引理
\ref{sites-modules-lemma-silly-quasi-compact}，存在有限子集 $K' \subset K$，使得
$\bigoplus_{k \in K'} j_{W_k!}\mathcal{O}_{W_k} \to j_{W!}\mathcal{O}_W$
是满射。由此可知 $\Coker(\varphi)$ 等于
$$
\Coker\left(
\bigoplus\nolimits_{k \in K'} j_{W_k!}\mathcal{O}_{W_k} \oplus
\bigoplus j_{V_j!}\mathcal{O}_{V_j}
\longrightarrow
\bigoplus j_{U_i!}\mathcal{O}_{U_i}
\right)
$$
，其中映射 $\bigoplus_{k \in K'} j_{W_k!}\mathcal{O}_{W_k}
\to \bigoplus j_{U_i!}\mathcal{O}_{U_i}$ 对应于
$\sum_{k \in K'} s_k$。证明完毕。
\end{proof}

\begin{lemma}
\label{sites-modules-lemma-change-presentation-towards-constructibles}
在情形 \ref{sites-modules-situation-quasi-compact-objects} 中，假设
(\ref{sites-modules-item-enough})、(\ref{sites-modules-item-enough-qc}) 和
(\ref{sites-modules-item-enough-qc-qs}) 成立。设 $\mathcal{O}$ 是环层。假设给定映射
$$
\bigoplus\nolimits_{j = 1, \ldots, m} j_{V_j!}\mathcal{O}_{V_j}
\longrightarrow
\bigoplus\nolimits_{i = 1, \ldots, n} j_{U_i!}\mathcal{O}_{U_i}
$$
，其中 $U_i$ 和 $V_j$ 属于 $\mathcal{B}$；并给定满足
$U_{ik} \in \mathcal{B}$ 的覆盖
$\{U_{ik} \to U_i\}_{k \in K_i}$。则存在有限子集
$K'_i \subset K_i$、由 $W_l \in \mathcal{B}$ 组成的有限集 $L$，以及交换图
$$
\xymatrix{
\bigoplus_{l \in L} j_{W_l!}\mathcal{O}_{W_l} \ar[d] \ar[r] &
\bigoplus_{i = 1, \ldots, n} \bigoplus_{k \in K'_i}
j_{U_{ik}!}\mathcal{O}_{U_{ik}} \ar[d] \\
\bigoplus_{j = 1, \ldots, m} j_{V_j!}\mathcal{O}_{V_j} \ar[r] &
\bigoplus_{i = 1, \ldots, n} j_{U_i!}\mathcal{O}_{U_i}
}
$$
，它在横向映射的余核上诱导同构。
\end{lemma}

\begin{proof}
由于 $U_i$ 拟紧，可以按引理 \ref{sites-modules-lemma-silly-quasi-compact} 选择有限子集
$K'_i \subset K_i$。又因为
$\bigoplus_{i = 1, \ldots, n}
\bigoplus_{k \in K'_i} j_{U_{ik}!}\mathcal{O}_{U_{ik}} \to
\bigoplus_{i = 1, \ldots, n} j_{U_i!}\mathcal{O}_{U_i}$ 是满射，
可以找到满足 $V_{jm} \in \mathcal{B}$ 的覆盖
$\{V_{jm} \to V_j\}_{m \in M_j}$，使得存在交换图
$$
\xymatrix{
\bigoplus_{j = 1, \ldots, m} \bigoplus_{m \in M_j}
j_{V_{jm}!}\mathcal{O}_{V_{jm}} \ar[d] \ar[r] &
\bigoplus_{i = 1, \ldots n} \bigoplus_{k \in K'_i}
j_{U_{ik}!}\mathcal{O}_{U_{ik}} \ar[d] \\
\bigoplus_{j = 1, \ldots, m} j_{V_j!}\mathcal{O}_{V_j} \ar[r] &
\bigoplus_{i = 1, \ldots, n} j_{U_i!}\mathcal{O}_{U_i}
}
$$
。由于 $V_j$ 拟紧，可以按引理 \ref{sites-modules-lemma-silly-quasi-compact}
选择有限子集 $M'_j \subset M_j$。置
$$
L = \left(\coprod\nolimits_{i = 1, \ldots, n} K'_i \times K'_i \right)
\coprod
\left(\coprod\nolimits_{j = 1, \ldots, m} M'_j\right)
$$
；对 $l = (k, k') \in K'_i \times K'_i \subset L$，置
$W_l = U_{ik} \times_{U_i} U_{ik'}$；对
$l = m \in M'_j \subset L$，置 $W_l = V_{jm}$。对映射族
$\{U_{ik} \to U_i\}_{k \in K'_i}$，已有引理
\ref{sites-modules-lemma-silly-quasi-compact} 的正合列，因此得到本引理陈述中的图
（略去细节），但尚不清楚 $W_l \in \mathcal{B}$。不过，对所有
$l \in L$，$W_l$ 都拟紧，所以再次应用引理
\ref{sites-modules-lemma-silly-quasi-compact}，找到满足 $W_{lt} \in \mathcal{B}$ 的有限映射族
$\{W_{lt} \to W_l\}_{t \in T_l}$，使得
$\bigoplus j_{W_{lt}!}\mathcal{O}_{W_{lt}} \to j_{W_l!}\mathcal{O}_{W_l}$
是满射。然后用 $\coprod_{l \in L} T_l$ 替换 $L$，结论即清楚。
\end{proof}

\begin{lemma}
\label{sites-modules-lemma-extension-towards-constructibles}
在情形 \ref{sites-modules-situation-quasi-compact-objects} 中，假设
(\ref{sites-modules-item-enough})、(\ref{sites-modules-item-enough-qc}) 和
(\ref{sites-modules-item-enough-qc-qs}) 成立。设 $\mathcal{O}$ 是环层。
则形如 (\ref{sites-modules-equation-towards-constructible}) 的模的扩张仍是形如
(\ref{sites-modules-equation-towards-constructible}) 的模。
\end{lemma}

\begin{proof}
设 $0 \to \mathcal{F}_1 \to \mathcal{F}_2 \to \mathcal{F}_3 \to 0$
是 $\mathcal{O}$-模的短正合列，其中 $\mathcal{F}_1$ 和 $\mathcal{F}_3$
形如 (\ref{sites-modules-equation-towards-constructible})。选择表示
$$
\bigoplus A_{V_j} \to \bigoplus A_{U_i} \to \mathcal{F}_1 \to 0
\quad\text{且}\quad
\bigoplus A_{T_j} \to \bigoplus A_{W_i} \to \mathcal{F}_3 \to 0
$$
。在此证明中直和总是有限的，并且对 $U \in \mathcal{B}$，记
$A_U = j_{U!}\mathcal{O}_U$。由于 $\mathcal{F}_2 \to \mathcal{F}_3$
是满射，可选择满足 $W_{ik} \in \mathcal{B}$ 的覆盖
$\{W_{ik} \to W_i\}$，使得 $A_{W_{ik}} \to \mathcal{F}_3$ 可提升为映射
$A_{W_{ik}} \to \mathcal{F}_2$。由引理
\ref{sites-modules-lemma-change-presentation-towards-constructibles}，可以用集合
$\{W_{ik}\}$ 的某个有限子集替换集合 $\{W_i\}$，并假设映射
$\bigoplus A_{W_i} \to \mathcal{F}_3$ 可提升为到 $\mathcal{F}_2$ 的映射。
考虑核
$$
\mathcal{K}_2 = \Ker(\bigoplus A_{U_i} \oplus \bigoplus A_{W_i}
\longrightarrow \mathcal{F}_2)
$$
。由蛇引理，该核满射到
$\mathcal{K}_3 = \Ker(\bigoplus A_{W_i} \to \mathcal{F}_3)$.
。因此按上述论证，在用 $\mathcal{B}$ 中元素的有限族替换每个 $T_j$ 后
（由引理 \ref{sites-modules-lemma-silly-quasi-compact}，这是允许的），可以假设存在映射
$\bigoplus A_{T_j} \to \mathcal{K}_2$，提升给定映射
$\bigoplus A_{T_j} \to \mathcal{K}_3$。于是
$\bigoplus A_{V_j} \oplus \bigoplus A_{T_j} \to \mathcal{K}_2$
是满射，证明完毕。
\end{proof}

\begin{lemma}
\label{sites-modules-lemma-towards-constructible-when-serre-subcategory}
在情形 \ref{sites-modules-situation-quasi-compact-objects} 中，假设
(\ref{sites-modules-item-enough})、(\ref{sites-modules-item-enough-qc}) 和
(\ref{sites-modules-item-enough-qc-qs}) 成立。设 $\mathcal{O}$ 是环层，并设
$\mathcal{A} \subset \textit{Mod}(\mathcal{O})$ 是由同构于形如
(\ref{sites-modules-equation-towards-constructible}) 的余核的模组成的全子范畴。
如果每个形如
$$
\bigoplus\nolimits_{j = 1, \ldots, m} j_{V_j!}\mathcal{O}_{V_j}
\longrightarrow
\bigoplus\nolimits_{i = 1, \ldots, n} j_{U_i!}\mathcal{O}_{U_i}
$$
的 $\mathcal{O}$-模映射的核都属于 $\mathcal{A}$，其中 $U_i$ 和 $V_j$
属于 $\mathcal{B}$，则 $\mathcal{A}$ 是
$\textit{Mod}(\mathcal{O})$ 的弱 Serre 子范畴。
\end{lemma}

\begin{proof}
使用《同调代数》，引理
\ref{homology-lemma-characterize-weak-serre-subcategory} 的判别准则。
由引理 \ref{sites-modules-lemma-cokernel-map-towards-constructibles} 和
\ref{sites-modules-lemma-extension-towards-constructibles} 的结果，只需证明
$\mathcal{A}$ 中对象之间的映射 $\mathcal{F} \to \mathcal{G}$ 的核属于
$\mathcal{A}$。为此选择表示
$$
\bigoplus A_{V_j} \to \bigoplus A_{U_i} \to \mathcal{F} \to 0
\quad\text{且}\quad
\bigoplus A_{T_j} \to \bigoplus A_{W_i} \to \mathcal{G} \to 0
$$
。在此证明中直和总是有限的，并且对 $U \in \mathcal{B}$，记
$A_U = j_{U!}\mathcal{O}_U$。使用引理
\ref{sites-modules-lemma-covering-gives-surjection} 和
\ref{sites-modules-lemma-change-presentation-towards-constructibles}，并像引理
\ref{sites-modules-lemma-extension-towards-constructibles} 的证明那样论证，可以假设映射
$\mathcal{F} \to \mathcal{G}$ 提升为表示之间的映射
$$
\xymatrix{
\bigoplus A_{V_j} \ar[r] \ar[d] &
\bigoplus A_{U_i} \ar[r] \ar[d] &
\mathcal{F} \ar[r] \ar[d] & 0 \\
\bigoplus A_{T_j} \ar[r] &
\bigoplus A_{W_i} \ar[r] &
\mathcal{G} \ar[r] & 0
}
$$
。于是可见
$$
\Ker(\mathcal{F} \to \mathcal{G}) =
\Coker\left(\bigoplus A_{V_j} \to
\Ker\left(
\bigoplus A_{T_j} \oplus \bigoplus A_{U_i} \to \bigoplus A_{W_i}\right)\right)
$$
，本引理由假设以及引理
\ref{sites-modules-lemma-cokernel-map-towards-constructibles} 得出。
\end{proof}








\section{平坦态射}
\label{sites-modules-section-flat-morphisms}

\begin{definition}
\label{sites-modules-definition-flat-morphism}
设
$(f, f^\sharp) :
(\Sh(\mathcal{C}), \mathcal{O})
\longrightarrow
(\Sh(\mathcal{C}'), \mathcal{O}')$
是环拓扑斯的态射。如果环映射
$f^\sharp : f^{-1}\mathcal{O}' \to \mathcal{O}$ 平坦，则称
$(f, f^\sharp)$ 是{\it 平坦的}。如果环站点的态射所关联的环拓扑斯态射平坦，
则称该环站点态射是{\it 平坦的}。
\end{definition}

\begin{lemma}
\label{sites-modules-lemma-flat-pullback-exact}
设 $f : \Sh(\mathcal{C}) \to \Sh(\mathcal{C}')$ 是环拓扑斯的态射。则
$$
f^{-1} : \textit{Ab}(\mathcal{C}') \longrightarrow \textit{Ab}(\mathcal{C}),
\quad
\mathcal{F} \longmapsto f^{-1}\mathcal{F}
$$
正合。如果
$(f, f^\sharp) :
(\Sh(\mathcal{C}), \mathcal{O})
\to
(\Sh(\mathcal{C}'), \mathcal{O}')$
是环拓扑斯的平坦态射，则
$$
f^* : \textit{Mod}(\mathcal{O}') \longrightarrow \textit{Mod}(\mathcal{O}),
\quad
\mathcal{F} \longmapsto f^*\mathcal{F}
$$
正合。
\end{lemma}

\begin{proof}
给定 $\mathcal{C}'$ 上的阿贝尔层 $\mathcal{G}$，$f^{-1}\mathcal{G}$ 的底层
集合层与 $\mathcal{G}$ 的底层集合层的 $f^{-1}$ 相同；见《站点》，第
\ref{sites-section-sheaves-algebraic-structures} 节。因此 $f^{-1}$ 对集合层的
正合性（拓扑斯态射定义中的要求；见《站点》，定义
\ref{sites-definition-topos}）蕴含 $f^{-1}$ 作为阿贝尔层上的函子正合。

\medskip\noindent
对于模层的陈述，回顾 $f^*\mathcal{F}$ 定义为张量积
$f^{-1}\mathcal{F} \otimes_{f^{-1}\mathcal{O}', f^\sharp} \mathcal{O}$。
因此 $f^*$ 是两个正合函子的复合。
\end{proof}

\begin{definition}
\label{sites-modules-definition-flat-module}
设 $f : (\Sh(\mathcal{C}), \mathcal{O}) \to
(\Sh(\mathcal{D}), \mathcal{O}')$ 是环拓扑斯的态射，且
$\mathcal{F}$ 是 $\mathcal{O}$-模层。如果 $\mathcal{F}$ 作为
$f^{-1}\mathcal{O}'$-模平坦，则称 $\mathcal{F}$
{\it 在 $(\Sh(\mathcal{D}), \mathcal{O}')$ 上平坦}。
\end{definition}

\noindent
这与对环空间态射定义的概念相容；见《模层》，定义
\ref{modules-definition-flat-module} 及其后的讨论。

\begin{lemma}
\label{sites-modules-lemma-pullback-internal-hom}
设 $f : (\mathcal{C}, \mathcal{O}_\mathcal{C}) \to
(\mathcal{D}, \mathcal{O}_\mathcal{D})$ 是环站点的态射，且
$\mathcal{F}$、$\mathcal{G}$ 是 $\mathcal{O}_\mathcal{D}$-模。
如果 $\mathcal{F}$ 有限表示且 $f$ 平坦，则注
\ref{sites-modules-remark-pullback-internal-hom} 中的典范映射
$$
f^*\SheafHom_{\mathcal{O}_\mathcal{D}}(\mathcal{F}, \mathcal{G})
\longrightarrow
\SheafHom_{\mathcal{O}_\mathcal{C}}(f^*\mathcal{F}, f^*\mathcal{G})
$$
是同构。
\end{lemma}

\begin{proof}
设 $f$ 由连续函子 $u : \mathcal{D} \to \mathcal{C}$ 给出。需要证明对任意
$U \in \Ob(\mathcal{C})$，该映射到 $\mathcal{C}/U$ 的限制是同构。
可以用 $U$ 的某个覆盖的成员替换 $U$。因此由《站点》，引理
\ref{sites-lemma-morphism-of-sites-covering}，可以假设对某个
$V \in \Ob(\mathcal{C})$ 存在态射 $U \to u(V)$。当然，此时可用 $u(V)$
替换 $U$。又因为 $u$ 连续，可以用某个覆盖替换 $V$，并假设存在表示
$\mathcal{O}_V^{\oplus m} \to \mathcal{O}_V^{\oplus n} \to
\mathcal{F}|_V \to 0$
于 $\mathcal{D}/V$ 上。由于 $\SheafHom$ 的形成与局部化交换（引理
\ref{sites-modules-lemma-internal-hom-restriction}），可以用由 $f$ 诱导的态射
$(\mathcal{C}/u(V), \mathcal{O}_{u(V)}) \to
(\mathcal{D}/V, \mathcal{O}_V)$ 替换 $f$。因此归结到 $\mathcal{F}$ 有整体表示
$\mathcal{O}_\mathcal{D}^{\oplus m} \to \mathcal{O}_\mathcal{D}^{\oplus n} \to
\mathcal{F} \to 0$ 的情形。因为 $f$ 平坦且
$f^*\mathcal{O}_\mathcal{D} = \mathcal{O}_\mathcal{C}$，得到相应的表示
$\mathcal{O}_\mathcal{C}^{\oplus m} \to \mathcal{O}_\mathcal{C}^{\oplus n} \to
f^*\mathcal{F} \to 0$；见引理 \ref{sites-modules-lemma-flat-pullback-exact}。
利用 $\SheafHom$ 对第一个变元与有限直和交换、
$\SheafHom_{\mathcal{O}_\mathcal{C}}(\mathcal{O}_\mathcal{C}, -)$ 和
$\SheafHom_{\mathcal{O}_\mathcal{D}}(\mathcal{O}_\mathcal{D}, -)$
都是恒等函子，以及注 \ref{sites-modules-remark-pullback-internal-hom} 中构造的函子性，
得到交换图
$$
\xymatrix{
0 \ar[r] &
f^*\SheafHom_{\mathcal{O}_\mathcal{D}}(\mathcal{F}, \mathcal{G})
\ar[d] \ar[r] &
f^*\mathcal{G}^{\oplus n} \ar[d] \ar[r] &
f^*\mathcal{G}^{\oplus n} \ar[d] \\
0 \ar[r] &
\SheafHom_{\mathcal{O}_\mathcal{C}}(f^*\mathcal{F}, f^*\mathcal{G})
\ar[r] &
f^*\mathcal{G}^{\oplus n} \ar[r] &
f^*\mathcal{G}^{\oplus n}
}
$$
，其中右侧两个竖直箭头是同构。由引理
\ref{sites-modules-lemma-internal-hom-exact}，两行都正合。再由五引理即得结论。
\end{proof}







\section{可逆模}
\label{sites-modules-section-invertible}

\noindent
定义如下。

\begin{definition}
\label{sites-modules-definition-invertible-sheaf}
设 $(\mathcal{C}, \mathcal{O})$ 为环站点。
\begin{enumerate}
\item 若对 $\mathcal{C}$ 的每个对象 $U$，都存在 $U$ 的一个覆盖
$\{U_i \to U\}$，使得作为 $\mathcal{O}_{U_i}$-模，
$\mathcal{F}|_{U_i}$ 同构于 $\mathcal{O}_{U_i}^{\oplus r}$，
则称有限局部自由 $\mathcal{O}$-模 $\mathcal{F}$ 的{\it 秩为 $r$}。
\item 若函子
$$
\textit{Mod}(\mathcal{O}) \longrightarrow \textit{Mod}(\mathcal{O}),\quad
\mathcal{F}  \longmapsto \mathcal{F} \otimes_\mathcal{O} \mathcal{L}
$$
是等价，则称 $\mathcal{O}$-模 $\mathcal{L}$ {\it 可逆}。
\item 层 {\it $\mathcal{O}^*$} 是 $\mathcal{O}$ 的子层，按如下规则定义：
$$
U \longmapsto \mathcal{O}^*(U) = \{f \in \mathcal{O}(U) \mid
\exists g \in \mathcal{O}(U)\text{ 使得 }fg = 1\}
$$
它是以乘法为群运算的阿贝尔群层。
\end{enumerate}
\end{definition}

\noindent
下面的引理 \ref{sites-modules-lemma-invertible-is-locally-free-rank-1}
说明了它与秩为 $1$ 的局部自由模之间的关系。

\begin{lemma}
\label{sites-modules-lemma-invertible}
设 $(\mathcal{C}, \mathcal{O})$ 为环站点，$\mathcal{L}$
为 $\mathcal{O}$-模。下列条件等价：
\begin{enumerate}
\item $\mathcal{L}$ 可逆；
\item 存在 $\mathcal{O}$-模 $\mathcal{N}$，使得
$\mathcal{L} \otimes_\mathcal{O} \mathcal{N} \cong \mathcal{O}$.
\end{enumerate}
在此情形下，还有：
\begin{enumerate}
\item[(a)] $\mathcal{L}$ 是有限表示的平坦 $\mathcal{O}$-模；
\item[(b)] 对 $\mathcal{C}$ 的每个对象 $U$，都存在一个覆盖
$\{U_i \to U\}$，使得 $\mathcal{L}|_{U_i}$
是某个有限自由模的直和因子；
\item[(c)] (2) 中的模 $\mathcal{N}$ 同构于
$\SheafHom_\mathcal{O}(\mathcal{L}, \mathcal{O})$.
\end{enumerate}
\end{lemma}

\begin{proof}
假设 (1) 成立。此时函子 $- \otimes_\mathcal{O} \mathcal{L}$
本质满，故存在 (2) 中所述的 $\mathcal{O}$-模
$\mathcal{N}$。反之，若 (2) 成立，则函子
$- \otimes_\mathcal{O} \mathcal{N}$ 是函子
$- \otimes_\mathcal{O} \mathcal{L}$ 的拟逆，因而 (1) 成立。

\medskip\noindent
假设 (1) 与 (2) 成立。由于 $- \otimes_\mathcal{O} \mathcal{L}$
是等价，它是正合函子，故 $\mathcal{L}$ 平坦。记给定的同构为
$\psi : \mathcal{L} \otimes_\mathcal{O} \mathcal{N} \to \mathcal{O}$。
设 $U$ 为 $\mathcal{C}$ 的一个对象。我们将证明：在 $U$
的某个覆盖的各成员上，$\mathcal{L}$ 的限制是自由模的直和因子；
这显然蕴含 $\mathcal{L}$ 是有限表示的。根据 $\otimes$ 的构造，
将 $U$ 替换为某个覆盖的各成员后，可以假设存在整数
$n \geq 1$ 以及截面 $x_i \in \mathcal{L}(U)$、
$y_i \in \mathcal{N}(U)$，使得 $\psi(\sum x_i \otimes y_i) = 1$。
考虑同构
$$
\mathcal{L}|_U \to
\mathcal{L}|_U \otimes_{\mathcal{O}_U}
\mathcal{L}|_U \otimes_{\mathcal{O}_U} \mathcal{N}|_U \to \mathcal{L}|_U
$$
其中第一个箭头将 $x$ 映到 $\sum x_i \otimes x \otimes y_i$，
第二个箭头将 $x \otimes x' \otimes y$ 映到 $\psi(x' \otimes y)x$。
由此可知，$x \mapsto \sum \psi(x \otimes y_i)x_i$
是 $\mathcal{L}|_U$ 的自同构。该自同构分解为
$$
\mathcal{L}|_U \to \mathcal{O}_U^{\oplus n} \to \mathcal{L}|_U
$$
其中第一个箭头由
$x \mapsto (\psi(x \otimes y_1), \ldots, \psi(x \otimes y_n))$
给出，第二个箭头由 $(a_1, \ldots, a_n) \mapsto \sum a_i x_i$
给出。因此 $\mathcal{L}|_U$ 是某个有限自由
$\mathcal{O}_U$-模的直和因子。

\medskip\noindent
仍假设 (1) 与 (2) 成立。考虑求值映射
$$
\mathcal{L} \otimes_\mathcal{O}
\SheafHom_\mathcal{O}(\mathcal{L}, \mathcal{O}_X)
\longrightarrow \mathcal{O}_X
$$
为完成引理的证明，我们将证明它是同构。由引理
\ref{sites-modules-lemma-internal-hom-adjoint-tensor}，有
$$
\Hom_\mathcal{O}(\mathcal{O}, \mathcal{O}) =
\Hom_\mathcal{O}
(\mathcal{N} \otimes_\mathcal{O} \mathcal{L}, \mathcal{O})
\longrightarrow
\Hom_\mathcal{O}
(\mathcal{N}, \SheafHom_\mathcal{O}(\mathcal{L}, \mathcal{O}))
$$
$1$ 的像给出态射
$\mathcal{N} \to \SheafHom_\mathcal{O}(\mathcal{L}, \mathcal{O})$。
与 $\mathcal{L}$ 作张量积，得到
$$
\mathcal{O} = \mathcal{L} \otimes_\mathcal{O} \mathcal{N}
\longrightarrow
\mathcal{L} \otimes_\mathcal{O} \SheafHom_\mathcal{O}(\mathcal{L}, \mathcal{O})
$$
此映射是求值映射的逆；计算从略。
\end{proof}

\begin{lemma}
\label{sites-modules-lemma-pullback-invertible}
设 $f : (\Sh(\mathcal{C}), \mathcal{O}_\mathcal{C}) \to
(\Sh(\mathcal{D}), \mathcal{O}_\mathcal{D})$ 为环拓扑斯态射。
可逆 $\mathcal{O}_\mathcal{D}$-模
$\mathcal{L}$ 的拉回 $f^*\mathcal{L}$ 是可逆的。
\end{lemma}

\begin{proof}
由引理 \ref{sites-modules-lemma-invertible}，存在
$\mathcal{O}_\mathcal{D}$-模 $\mathcal{N}$，使得
$\mathcal{L} \otimes_{\mathcal{O}_\mathcal{D}} \mathcal{N} \cong
\mathcal{O}_\mathcal{D}$。作拉回并应用引理
\ref{sites-modules-lemma-tensor-product-pullback}，得到
$f^*\mathcal{L} \otimes_{\mathcal{O}_\mathcal{C}} f^*\mathcal{N} \cong
\mathcal{O}_\mathcal{C}$。
因此由引理 \ref{sites-modules-lemma-invertible}，$f^*\mathcal{L}$ 可逆。
\end{proof}

\begin{lemma}
\label{sites-modules-lemma-constructions-invertible}
设 $(\mathcal{C}, \mathcal{O})$ 为环站点。
\begin{enumerate}
\item 若 $\mathcal{L}$、$\mathcal{N}$ 是可逆
$\mathcal{O}$-模，则
$\mathcal{L} \otimes_\mathcal{O} \mathcal{N}$.
也可逆。
\item 若 $\mathcal{L}$ 是可逆 $\mathcal{O}$-模，则
$\SheafHom_\mathcal{O}(\mathcal{L}, \mathcal{O})$ 也可逆，且求值映射
$\mathcal{L} \otimes_\mathcal{O}
\SheafHom_\mathcal{O}(\mathcal{L}, \mathcal{O}) \to \mathcal{O}$
是同构。
\end{enumerate}
\end{lemma}

\begin{proof}
(1) 由定义即得，(2) 由引理 \ref{sites-modules-lemma-invertible} 及其证明即得。
\end{proof}

\begin{lemma}
\label{sites-modules-lemma-pic-set}
设 $(\mathcal{C}, \mathcal{O})$ 为环站点。存在一族可逆模
$\{\mathcal{L}_i\}_{i \in I}$，使得
$(\mathcal{C}, \mathcal{O})$ 上的每个可逆模恰与其中一个
$\mathcal{L}_i$ 同构。
\end{lemma}

\begin{proof}
从略；见《模层》，引理 \ref{modules-lemma-pic-set}。
\end{proof}

\noindent
引理 \ref{sites-modules-lemma-pic-set} 表明，可逆层的同构类全体构成一个集合。
引理 \ref{sites-modules-lemma-constructions-invertible} 表明，张量积在此集合上
定义了一个阿贝尔群结构，其中 $\mathcal{L}$ 的逆元由
$\SheafHom_\mathcal{O}(\mathcal{L}, \mathcal{O})$ 给出。

\medskip\noindent
事实上，给定可逆 $\mathcal{O}$-模 $\mathcal{L}$
及 $n \in \mathbf{Z}$，将等价
$\mathcal{F} \mapsto \mathcal{F} \otimes_\mathcal{O} \mathcal{L}$
对 $\mathcal{O}$ 迭代应用恰好 $n$ 次所得的像，定义为
$\mathcal{L}$ 的第 $n$ 个{\it 张量幂}
$\mathcal{L}^{\otimes n}$。因为我们把可逆 $\mathcal{O}$-模
定义为与其作张量积给出等价的模，所以该定义对负整数 $n$
也有意义。更明确地，
$$
\mathcal{L}^{\otimes n} =
\left\{
\begin{matrix}
\mathcal{O} & \text{若} & n = 0 \\
\SheafHom_\mathcal{O}(\mathcal{L}, \mathcal{O}) & \text{若} & n = -1\\
\mathcal{L} \otimes_\mathcal{O} \ldots \otimes_\mathcal{O} \mathcal{L}
& \text{若} & n > 0 \\
\mathcal{L}^{\otimes -1} \otimes_\mathcal{O} \ldots
\otimes_\mathcal{O} \mathcal{L}^{\otimes -1}
& \text{若} & n < -1
\end{matrix}
\right.
$$
见引理 \ref{sites-modules-lemma-constructions-invertible}。
按此定义，我们有典范同构
$\mathcal{L}^{\otimes n} \otimes_\mathcal{O}
\mathcal{L}^{\otimes m} \to
\mathcal{L}^{\otimes n + m}$，并且这些同构满足交换性约束与结合性约束
（表述从略）。

\begin{definition}
\label{sites-modules-definition-pic}
设 $(\mathcal{C}, \mathcal{O})$ 为环站点。该环站点的
{\it Picard 群} $\Pic(\mathcal{O})$ 是以可逆
$\mathcal{O}$-模的同构类为元素、加法对应于张量积的阿贝尔群。
\end{definition}









\section{微分模}
\label{sites-modules-section-differentials}

\noindent
本节简要说明如何为环拓扑斯态射定义相对微分模。建议读者参阅
交换代数一章中的相应一节
（《代数》，第 \ref{algebra-section-differentials} 节）。

\begin{definition}
\label{sites-modules-definition-derivation}
设 $\mathcal{C}$ 为站点，$\varphi : \mathcal{O}_1 \to \mathcal{O}_2$
为环层同态，$\mathcal{F}$ 为 $\mathcal{O}_2$-模。取值于
$\mathcal{F}$ 的{\it $\mathcal{O}_1$-导子}，更准确地说是
{\it $\varphi$-导子}，是一个映射
$D : \mathcal{O}_2 \to \mathcal{F}$，它可加，零化
$\mathcal{O}_1 \to \mathcal{O}_2$ 的像，并满足{\it Leibniz 法则}
$$
D(ab) = aD(b) + D(a)b
$$
其中 $a,b$ 是 $\mathcal{O}_2$ 的任意两个局部截面
（在二者均有定义之处）。记
$\text{Der}_{\mathcal{O}_1}(\mathcal{O}_2, \mathcal{F})$
为取值于 $\mathcal{F}$ 的 $\varphi$-导子所成的集合。
\end{definition}

\noindent
这是《代数》，定义 \ref{sites-modules-definition-derivation} 的层论版本。
给定定义中所述的导子 $D : \mathcal{O}_2 \to \mathcal{F}$，
整体截面上的映射
$$
D : \Gamma(\mathcal{O}_2) \longrightarrow \Gamma(\mathcal{F})
$$
显然是代数定义意义下的 $\Gamma(\mathcal{O}_1)$-导子。注意，若
$\alpha : \mathcal{F} \to \mathcal{G}$ 是
$\mathcal{O}_2$-模的映射，则有诱导映射
$$
\text{Der}_{\mathcal{O}_1}(\mathcal{O}_2, \mathcal{F})
\longrightarrow
\text{Der}_{\mathcal{O}_1}(\mathcal{O}_2, \mathcal{G})
$$
由规则 $D \mapsto \alpha \circ D$ 给出。换言之，我们得到一个函子。

\begin{lemma}
\label{sites-modules-lemma-universal-module}
设 $\mathcal{C}$ 为站点，$\varphi : \mathcal{O}_1 \to \mathcal{O}_2$
为环层同态。函子
$$
\textit{Mod}(\mathcal{O}_2) \longrightarrow \textit{Ab}, \quad
\mathcal{F} \longmapsto \text{Der}_{\mathcal{O}_1}(\mathcal{O}_2, \mathcal{F})
$$
可表。
\end{lemma}

\begin{proof}
证明与代数中的相应命题完全相同。为叙述该证明，对
$\mathcal{C}$ 上任意集合层 $\mathcal{F}$，记
$\mathcal{O}_2[\mathcal{F}]$ 为预层
$U \mapsto \mathcal{O}_2(U)[\mathcal{F}(U)]$ 的层化，其中右端表示
由集合 $\mathcal{F}(U)$ 生成的自由 $\mathcal{O}_2(U)$-模。
对 $s \in \mathcal{F}(U)$，以 $[s]$ 表示
$\mathcal{O}_2[\mathcal{F}]$ 在 $U$ 上的相应截面。若
$\mathcal{F}$ 是 $\mathcal{O}_2$-模层，则有典范映射
$$
c : \mathcal{O}_2[\mathcal{F}] \longrightarrow \mathcal{F}
$$
它在预层层次上由规则
$\sum f_s[s] \mapsto \sum f_s s$ 给出。以下以简写
$[s] \mapsto s$ 描述此映射，其他映射亦同。考虑
$\mathcal{O}_2$-模的映射
\begin{equation}
\label{sites-modules-equation-define-module-differentials}
\begin{matrix}
\mathcal{O}_2[\mathcal{O}_2 \times \mathcal{O}_2] \oplus
\mathcal{O}_2[\mathcal{O}_2 \times \mathcal{O}_2] \oplus
\mathcal{O}_2[\mathcal{O}_1] &
\longrightarrow &
\mathcal{O}_2[\mathcal{O}_2] \\
[(a, b)] \oplus [(f, g)] \oplus [h] & \longmapsto & [a + b] - [a] - [b] + \\
& & [fg] - g[f] - f[g] + \\
& & [\varphi(h)]
\end{matrix}
\end{equation}
其中采用上述简写。令 $\Omega_{\mathcal{O}_2/\mathcal{O}_1}$
为此映射的余核。显然存在集合层的映射
$$
\text{d} : \mathcal{O}_2 \longrightarrow \Omega_{\mathcal{O}_2/\mathcal{O}_1}
$$
它将局部截面 $f$ 映到 $[f]$ 在
$\Omega_{\mathcal{O}_2/\mathcal{O}_1}$ 中的像。由构造，
$\text{d}$ 是 $\mathcal{O}_1$-导子。再设 $\mathcal{F}$
为 $\mathcal{O}_2$-模层，且
$D : \mathcal{O}_2 \to \mathcal{F}$ 为 $\mathcal{O}_1$-导子。
考虑将 $[g]$ 映到 $D(g)$ 的 $\mathcal{O}_2$-线性映射
$\mathcal{O}_2[\mathcal{O}_2] \to \mathcal{F}$。
由导子的定义，该映射零化映射
(\ref{sites-modules-equation-define-module-differentials}) 像中的截面，故定义了映射
$$
\alpha_D : \Omega_{\mathcal{O}_2/\mathcal{O}_1} \longrightarrow \mathcal{F}
$$
显然 $D = \alpha_D \circ \text{d}$，故引理得证。
\end{proof}

\begin{definition}
\label{sites-modules-definition-module-differentials}
设 $\mathcal{C}$ 为站点，$\varphi : \mathcal{O}_1 \to \mathcal{O}_2$
为环层同态。环映射 $\varphi$ 的{\it 微分模}是表示函子
$\mathcal{F} \mapsto \text{Der}_{\mathcal{O}_1}(\mathcal{O}_2, \mathcal{F})$
的对象；由引理 \ref{sites-modules-lemma-universal-module}，该对象存在。
记它为 $\Omega_{\mathcal{O}_2/\mathcal{O}_1}$，并将{\it 泛
$\varphi$-导子}记为
$\text{d} : \mathcal{O}_2 \to \Omega_{\mathcal{O}_2/\mathcal{O}_1}$.
\end{definition}

\noindent
由于此模及该导子构成表示一个函子的泛对象，这一概念显然是内蕴的
（即不依赖于环拓扑斯所基于的站点的选择；见第
\ref{sites-modules-section-intrinsic} 节）。注意，
$\Omega_{\mathcal{O}_2/\mathcal{O}_1}$ 是
$\mathcal{O}_2$-模映射
(\ref{sites-modules-equation-define-module-differentials}) 的余核。此外，映射
$\text{d}$ 由如下规则描述：$\text{d}f$ 是局部截面 $[f]$ 的像。

\begin{lemma}
\label{sites-modules-lemma-differentials-sheafify}
设 $\mathcal{C}$ 为站点，$\varphi : \mathcal{O}_1 \to \mathcal{O}_2$
为环预层同态。则
$\Omega_{\mathcal{O}_2^\#/\mathcal{O}_1^\#}$ 是预层
$U \mapsto \Omega_{\mathcal{O}_2(U)/\mathcal{O}_1(U)}$ 的关联层。
\end{lemma}

\begin{proof}
考虑映射 (\ref{sites-modules-equation-define-module-differentials})。
存在一个类似的预层映射，它在 $U \in \Ob(\mathcal{C})$ 上的取值为
$$
\mathcal{O}_2(U)[\mathcal{O}_2(U) \times \mathcal{O}_2(U)] \oplus
\mathcal{O}_2(U)[\mathcal{O}_2(U) \times \mathcal{O}_2(U)] \oplus
\mathcal{O}_2(U)[\mathcal{O}_1(U)]
\longrightarrow
\mathcal{O}_2(U)[\mathcal{O}_2(U)]
$$
由《代数》，定义 \ref{algebra-definition-differentials}
中微分模的构造，此映射的余核在 $U$ 上的取值为
$\Omega_{\mathcal{O}_2(U)/\mathcal{O}_1(U)}$。另一方面，
(\ref{sites-modules-equation-define-module-differentials}) 中的各层正是上述
各预层的层化。由于层化正合，结论随即成立。
\end{proof}

\begin{lemma}
\label{sites-modules-lemma-pullback-differentials}
设 $f : \Sh(\mathcal{D}) \to \Sh(\mathcal{C})$ 为拓扑斯态射，
$\varphi : \mathcal{O}_1 \to \mathcal{O}_2$
为 $\mathcal{C}$ 上的环层同态。则有与泛导子相容的典范等同
$f^{-1}\Omega_{\mathcal{O}_2/\mathcal{O}_1} =
\Omega_{f^{-1}\mathcal{O}_2/f^{-1}\mathcal{O}_1}$。
\end{lemma}

\begin{proof}
这是因为层 $\Omega_{\mathcal{O}_2/\mathcal{O}_1}$ 是映射
(\ref{sites-modules-equation-define-module-differentials}) 的余核，
$\Omega_{f^{-1}\mathcal{O}_2/f^{-1}\mathcal{O}_1}$ 也有类似描述，
函子 $f^{-1}$ 正合，并且
$f^{-1}(\mathcal{O}_2[\mathcal{O}_2]) =
f^{-1}\mathcal{O}_2[f^{-1}\mathcal{O}_2]$,
$f^{-1}(\mathcal{O}_2[\mathcal{O}_2 \times \mathcal{O}_2]) =
f^{-1}\mathcal{O}_2[f^{-1}\mathcal{O}_2 \times f^{-1}\mathcal{O}_2]$, and
$f^{-1}(\mathcal{O}_2[\mathcal{O}_1]) =
f^{-1}\mathcal{O}_2[f^{-1}\mathcal{O}_1]$.
\end{proof}

\begin{lemma}
\label{sites-modules-lemma-localize-differentials}
设 $\mathcal{C}$ 为站点，$\varphi : \mathcal{O}_1 \to \mathcal{O}_2$
为环层同态。对 $\mathcal{C}$ 的任意对象 $U$，有与泛导子相容的
典范同构
$$
\Omega_{\mathcal{O}_2/\mathcal{O}_1}|_U =
\Omega_{(\mathcal{O}_2|_U)/(\mathcal{O}_1|_U)}
$$
\end{lemma}

\begin{proof}
这是引理 \ref{sites-modules-lemma-pullback-differentials} 的特例。
\end{proof}

\begin{lemma}
\label{sites-modules-lemma-functoriality-differentials}
设 $\mathcal{C}$ 为站点。设
$$
\xymatrix{
\mathcal{O}_2 \ar[r]_\varphi & \mathcal{O}_2' \\
\mathcal{O}_1 \ar[r] \ar[u] & \mathcal{O}'_1 \ar[u]
}
$$
为 $\mathcal{C}$ 上环层的交换图。将映射
$\mathcal{O}_2 \to \mathcal{O}'_2$ 与映射
$\text{d} : \mathcal{O}'_2 \to \Omega_{\mathcal{O}'_2/\mathcal{O}'_1}$
复合，得到一个 $\mathcal{O}_1$-导子。因此得到典范
$\mathcal{O}_2$-模映射
$\Omega_{\mathcal{O}_2/\mathcal{O}_1} \to
\Omega_{\mathcal{O}'_2/\mathcal{O}'_1}$.
它由下述性质唯一刻画：对 $\mathcal{O}_2$ 的任意局部截面 $f$，
$\text{d}(f)$ 映到 $\text{d}(\varphi(f))$。
由此，$\Omega_{-/-}$ 成为环层箭头范畴上的函子。
\end{lemma}

\begin{proof}
此引理由其陈述即得。
\end{proof}

\begin{lemma}
\label{sites-modules-lemma-differential-seq}
在引理 \ref{sites-modules-lemma-functoriality-differentials} 中，假设
$\mathcal{O}_2 \to \mathcal{O}'_2$ 满射，其核为
$\mathcal{I} \subset \mathcal{O}_2$，并假设
$\mathcal{O}_1 = \mathcal{O}'_1$。则有典范的
$\mathcal{O}'_2$-模正合列
$$
\mathcal{I}/\mathcal{I}^2
\longrightarrow
\Omega_{\mathcal{O}_2/\mathcal{O}_1} \otimes_{\mathcal{O}_2} \mathcal{O}'_2
\longrightarrow
\Omega_{\mathcal{O}'_2/\mathcal{O}_1}
\longrightarrow
0
$$
最左侧映射由下述规则刻画：$\mathcal{I}$ 的局部截面
$f$ 映到 $\text{d}f \otimes 1$。
\end{lemma}

\begin{proof}
对 $\mathcal{I}$ 的局部截面 $f$，记 $\overline{f}$ 为
$f$ 在 $\mathcal{I}/\mathcal{I}^2$ 中的像。为证明映射
$\overline{f} \mapsto \text{d}f \otimes 1$ 良定义，只需验证：
若 $f_1,f_2$ 是 $\mathcal{I}$ 的局部截面，则
$\text{d} f_1f_2 \otimes 1 = 0$。这由 Leibniz 法则立即得到：
$\text{d} f_1f_2 \otimes 1 =
(f_1 \text{d}f_2 + f_2 \text{d} f_1 )\otimes 1 =
\text{d}f_2 \otimes f_1 + \text{d}f_2 \otimes f_1 = 0$.
类似的计算表明，此映射是
$\mathcal{O}'_2 = \mathcal{O}_2/\mathcal{I}$-线性的。右侧映射是
引理 \ref{sites-modules-lemma-functoriality-differentials} 中的映射。

\medskip\noindent
为证明该列正合，作如下论证。令
$\mathcal{O}''_2 \subset \mathcal{O}'_2$ 为
$\mathcal{O}_1$-代数预层，它在 $U$ 上的取值是
$\mathcal{O}_2(U) \to \mathcal{O}'_2(U)$ 的像。由
《代数》，引理 \ref{algebra-lemma-differential-seq}，对
$\mathcal{C}$ 的每个对象 $U$，下列序列正合：
$$
\mathcal{I}(U)/\mathcal{I}(U)^2
\longrightarrow
\Omega_{\mathcal{O}_2(U)/\mathcal{O}_1(U)}
\otimes_{\mathcal{O}_2(U)} \mathcal{O}''_2(U)
\longrightarrow
\Omega_{\mathcal{O}''_2(U)/\mathcal{O}_1(U)}
\longrightarrow
0
$$
由于层化正合，这给出一个
$(\mathcal{O}'_2)^\#$-模层的正合列。再由引理
\ref{sites-modules-lemma-differentials-sheafify} 及
$(\mathcal{O}''_2)^\# = \mathcal{O}'_2$，即得结论。
\end{proof}

\noindent
下面给出导子自然出现的一种特殊情形。

\begin{lemma}
\label{sites-modules-lemma-double-structure-gives-derivation}
设 $\mathcal{C}$ 为站点，$\varphi : \mathcal{O}_1 \to \mathcal{O}_2$
为环层同态。考虑短正合列
$$
0 \to \mathcal{F} \to \mathcal{A} \to \mathcal{O}_2 \to 0
$$
这里 $\mathcal{A}$ 是 $\mathcal{O}_1$-代数层，
$\pi : \mathcal{A} \to \mathcal{O}_2$ 是
$\mathcal{O}_1$-代数层的满射，且
$\mathcal{F} = \Ker(\pi)$ 是其核。假设 $\mathcal{F}$
是 $\mathcal{A}$ 中平方为零的理想层。因此
$\mathcal{F}$ 自然带有 $\mathcal{O}_2$-模结构。
$\pi$ 的截面 $s : \mathcal{O}_2 \to \mathcal{A}$
是满足 $\pi \circ s = \text{id}$ 的 $\mathcal{O}_1$-代数映射。
给定 $\pi$ 的任意截面 $s : \mathcal{O}_2 \to \mathcal{F}$
以及任意 $\varphi$-导子 $D : \mathcal{O}_1 \to \mathcal{F}$，映射
$$
s + D : \mathcal{O}_1 \to \mathcal{A}
$$
是 $\pi$ 的截面，并且每个截面 $s'$ 都唯一地写成
$s + D$，其中 $D$ 是 $\varphi$-导子。
\end{lemma}

\begin{proof}
回忆，$\mathcal{F}$ 上的 $\mathcal{O}_2$-模结构由
$h \tau = \tilde h \tau$ 给出（右侧是在 $\mathcal{A}$ 中的乘法）；
这里 $h$ 是 $\mathcal{O}_2$ 的局部截面，
$\tilde h$ 是 $h$ 到 $\mathcal{A}$ 的局部截面的一个局部提升，
而 $\tau$ 是 $\mathcal{F}$ 的局部截面。特别地，给定 $s$ 后，
可以取 $\tilde h = s(h)$。为验证 $s + D$ 是环层同态，计算
\begin{eqnarray*}
(s + D)(ab) & = & s(ab) + D(ab) \\
& = & s(a)s(b) + aD(b) + D(a)b \\
& = & s(a) s(b) + s(a)D(b) + D(a)s(b) \\
& = & (s(a) + D(a))(s(b) + D(b))
\end{eqnarray*}
其中使用了 Leibniz 法则。类似地，由于 $D$ 是
$\mathcal{O}_1$-导子，可证 $s + D$ 是
$\mathcal{O}_1$-代数映射。反之，给定 $s'$，令
$D = s' - s$ 即可。细节从略。
\end{proof}

\begin{definition}
\label{sites-modules-definition-sheaf-differentials}
设 $X = (\Sh(\mathcal{C}), \mathcal{O})$ 与
$Y = (\Sh(\mathcal{C}'), \mathcal{O}')$ 为环拓扑斯，并设
$(f, f^\sharp) : X \to Y$ 为环拓扑斯态射。在此情形下：
\begin{enumerate}
\item 对 $\mathcal{O}$-模层 $\mathcal{F}$，{\it $Y$-导子}
$D : \mathcal{O} \to \mathcal{F}$ 就是 $f^\sharp$-导子；
\item {\it $X$ 在 $Y$ 上的微分层 $\Omega_{X/Y}$}
是如下映射的微分模：
$f^\sharp : f^{-1}\mathcal{O}' \to \mathcal{O}$,
见定义 \ref{sites-modules-definition-module-differentials}。
\end{enumerate}
因此 $\Omega_{X/Y}$ 配备一个{\it 泛 $Y$-导子}
$\text{d}_{X/Y} : \mathcal{O} \longrightarrow \Omega_{X/Y}$。
有时也写作 $\Omega_{X/Y} = \Omega_f$。
\end{definition}

\noindent
回忆 $f^\sharp : f^{-1}\mathcal{O}' \to \mathcal{O}$，故此定义有意义。

\begin{lemma}
\label{sites-modules-lemma-functoriality-differentials-ringed-topoi}
设
$X = (\Sh(\mathcal{C}_X), \mathcal{O}_X)$,
$Y = (\Sh(\mathcal{C}_Y), \mathcal{O}_Y)$,
$X' = (\Sh(\mathcal{C}_{X'}), \mathcal{O}_{X'})$, and
$Y' = (\Sh(\mathcal{C}_{Y'}), \mathcal{O}_{Y'})$ 为环拓扑斯。设
$$
\xymatrix{
X' \ar[d] \ar[r]_f & X \ar[d] \\
Y' \ar[r] & Y
}
$$
为环拓扑斯态射的交换图。将映射
$f^\sharp : \mathcal{O}_X \to f_*\mathcal{O}_{X'}$ 与映射
$f_*\text{d}_{X'/Y'} : f_*\mathcal{O}_{X'} \to f_*\Omega_{X'/Y'}$
复合，得到一个 $Y$-导子。因此得到典范
$\mathcal{O}_X$-模映射
$\Omega_{X/Y} \to f_*\Omega_{X'/Y'}$；再由
$f_*$ 与 $f^*$ 的伴随性，得到典范的
$\mathcal{O}_{X'}$-模同态
$$
c_f : f^*\Omega_{X/Y} \longrightarrow \Omega_{X'/Y'}.
$$
它由下述性质唯一刻画：对 $\mathcal{O}_X$ 的任意局部截面 $t$，
$f^*\text{d}_{X/Y}(t)$ 映到 $\text{d}_{X'/Y'}(f^* t)$。
\end{lemma}

\begin{proof}
除最后一个断言外均显然。下面解释该断言的含义。设
$U \in \Ob(\mathcal{C}_X)$ 且 $t \in \mathcal{O}_X(U)$；
这正是说 $t$ 是 $\mathcal{O}_X$ 的局部截面。可以把 $t$
视为集合层的映射
$t : h_U^\# \to \mathcal{O}_X$。于是有
$f^{-1}t : f^{-1}h_U^\# \to f^{-1}\mathcal{O}_X$。
所谓 $f^*t$，是指复合
$$
\xymatrix{
f^{-1}h_U^\# \ar[rr]^{f^{-1}t} \ar@/^4ex/[rrrr]^{f^*t} & &
f^{-1}\mathcal{O}_X \ar[rr]^{f^\sharp} & &
\mathcal{O}_{X'}
}
$$
注意 $\text{d}_{X/Y}(t) \in \Omega_{X/Y}(U)$。因此可以把
$\text{d}_{X/Y}(t)$ 视为映射
$\text{d}_{X/Y}(t) : h_U^\# \to \Omega_{X/Y}$。于是
$f^{-1}\text{d}_{X/Y}(t) : f^{-1}h_U^\# \to f^{-1}\Omega_{X/Y}$。
所谓 $f^*\text{d}_{X/Y}(t)$，是指复合
$$
\xymatrix{
f^{-1}h_U^\#
\ar[rr]^{f^{-1}\text{d}_{X/Y}(t)}
\ar@/^4ex/[rrrr]^{f^*\text{d}_{X/Y}(t)} & &
f^{-1}\Omega_{X/Y} \ar[rr]^{1 \otimes \text{id}} & &
f^*\Omega_{X/Y}
}
$$
于是，引理的断言是指
$$
c_f \circ f^*t = f^*\text{d}_{X/Y}(t)
$$
作为从 $f^{-1}h_U^\#$ 到 $\Omega_{X'/Y'}$ 的映射相等。
略去验证引理中定义的 $c_f$ 满足此性质。（提示：第一个映射
$c'_f : \Omega_{X/Y} \to f_*\Omega_{X'/Y'}$ 满足
$c'_f(\text{d}_{X/Y}(t)) = f_*\text{d}_{X'/Y'}(f^\sharp(t))$，
这是 $f_*\Omega_{X'/Y'}$ 在 $U$ 上的截面等式；须利用伴随性
将其化为上述等式。）此性质唯一刻画 $c_f$，因为
$f^*\text{d}_{X/Y}(t)$ 的像生成 $\mathcal{O}_{X'}$-模
$f^*\Omega_{X/Y}$；这又是因为局部截面
$\text{d}_{X/Y}(t)$ 生成 $\mathcal{O}_X$-模 $\Omega_{X/Y}$。
\end{proof}





\section{有限阶微分算子}
\label{sites-modules-section-differential-operators}

\noindent
本节介绍有限阶微分算子。建议读者参阅交换代数一章中的相应一节
（《代数》，第 \ref{algebra-section-differential-operators} 节）。

\begin{definition}
\label{sites-modules-definition-differential-operators}
设 $\mathcal{C}$ 为站点，$\varphi : \mathcal{O}_1 \to \mathcal{O}_2$
为环层同态，$k \geq 0$ 为整数，并设 $\mathcal{F}$、$\mathcal{G}$
为 $\mathcal{O}_2$-模层。所谓{\it 阶数为 $k$ 的微分算子
$D : \mathcal{F} \to \mathcal{G}$}，是指一个
$\mathcal{O}_1$-线性映射，满足：对 $\mathcal{O}_2$
的任意局部截面 $g$，映射 $s \mapsto D(gs) - gD(s)$
是阶数为 $k - 1$ 的微分算子。作为起始情形，当 $k = 0$ 时，
将阶数为 $0$ 的微分算子定义为 $\mathcal{O}_2$-线性映射。
\end{definition}

\noindent
若 $D : \mathcal{F} \to \mathcal{G}$ 是阶数为 $k$ 的微分算子，
则对 $\mathcal{O}_2$ 的任意局部截面 $g$，映射 $gD$
也是阶数为 $k$ 的微分算子。两个阶数为 $k$ 的微分算子之和
仍是同阶微分算子。因此所有这些算子所成的集合
$$
\text{Diff}^k(\mathcal{F}, \mathcal{G}) =
\text{Diff}^k_{\mathcal{O}_2/\mathcal{O}_1}(\mathcal{F}, \mathcal{G})
$$
是 $\Gamma(\mathcal{C}, \mathcal{O}_2)$-模。并且
$$
\text{Diff}^0(\mathcal{F}, \mathcal{G}) \subset
\text{Diff}^1(\mathcal{F}, \mathcal{G}) \subset
\text{Diff}^2(\mathcal{F}, \mathcal{G}) \subset \ldots
$$
把 $U \in \Ob(\mathcal{C})$ 映到阶数为 $k$ 的微分算子
$D : \mathcal{F}|_U \to \mathcal{G}|_U$ 所成模的规则，
给出站点 $\mathcal{C}$ 上的一个 $\mathcal{O}_2$-模层。
由此得到微分算子层（若以后需要，我们将在此补充定义）。

\begin{lemma}
\label{sites-modules-lemma-composition-differential-operators}
设 $\mathcal{C}$ 为站点，
$\mathcal{O}_1 \to \mathcal{O}_2$ 为环层映射，并设
$\mathcal{E}, \mathcal{F}, \mathcal{G}$ 为
$\mathcal{O}_2$-模层。若
$D : \mathcal{E} \to \mathcal{F}$ 与
$D' : \mathcal{F} \to \mathcal{G}$ 分别是阶数为 $k$ 与 $k'$
的微分算子，则 $D' \circ D$ 是阶数为 $k + k'$ 的微分算子。
\end{lemma}

\begin{proof}
设 $g$ 为 $\mathcal{O}_2$ 的局部截面。将 $\mathcal{E}$
的局部截面 $x$ 映到
$$
D'(D(gx)) - gD'(D(x)) = D'(D(gx)) - D'(gD(x)) + D'(gD(x)) - gD'(D(x))
$$
的映射，是两个较低阶微分算子复合的和。因此对 $k + k'$
作归纳即可证明引理。
\end{proof}

\begin{lemma}
\label{sites-modules-lemma-module-principal-parts}
设 $\mathcal{C}$ 为站点，
$\mathcal{O}_1 \to \mathcal{O}_2$ 为环层映射，
$\mathcal{F}$ 为 $\mathcal{O}_2$-模层，且 $k \geq 0$。
则存在 $\mathcal{O}_2$-模层
$\mathcal{P}^k_{\mathcal{O}_2/\mathcal{O}_1}(\mathcal{F})$
及典范同构
$$
\text{Diff}^k_{\mathcal{O}_2/\mathcal{O}_1}(\mathcal{F}, \mathcal{G}) =
\Hom_{\mathcal{O}_2}(
\mathcal{P}^k_{\mathcal{O}_2/\mathcal{O}_1}(\mathcal{F}), \mathcal{G})
$$
它关于 $\mathcal{O}_2$-模 $\mathcal{G}$ 函子性成立。
\end{lemma}

\begin{proof}
存在性可由一般范畴论论证得到（此处将来插入参考文献），但我们还将
给出一个直接构造，因为该构造会用于后续证明。以下自由使用引理
\ref{sites-modules-lemma-universal-module} 的证明中引入的记号。
给定任意微分算子 $D : \mathcal{F} \to \mathcal{G}$，
得到 $\mathcal{O}_2$-线性映射
$L_D : \mathcal{O}_2[\mathcal{F}] \to \mathcal{G}$
将 $[m]$ 映到 $D(m)$。若 $D$ 的阶数为 $0$，则
$L_D$ 零化局部截面
$$
[m + m'] - [m] - [m'],\quad
g_0[m] - [g_0m]
$$
其中 $g_0$ 是 $\mathcal{O}_2$ 的局部截面，而 $m,m'$
是 $\mathcal{F}$ 的局部截面。若 $D$ 的阶数为 $1$，则
$L_D$ 零化局部截面
$$
[m + m'] - [m] - [m'],\quad
f[m] - [fm], \quad
g_0g_1[m] - g_0[g_1m] - g_1[g_0m] + [g_1g_0m]
$$
其中 $f$ 是 $\mathcal{O}_1$ 的局部截面，
$g_0,g_1$ 是 $\mathcal{O}_2$ 的局部截面，而
$m,m'$ 是 $\mathcal{F}$ 的局部截面。若 $D$ 的阶数为 $k$，
则 $L_D$ 零化局部截面
$[m + m'] - [m] - [m']$、$f[m] - [fm]$ 以及
$$
g_0g_1\ldots g_k[m] - \sum g_0 \ldots \hat g_i \ldots g_k[g_im] + \ldots
+(-1)^{k + 1}[g_0\ldots g_km]
$$
反之，若 $L : \mathcal{O}_2[\mathcal{F}] \to \mathcal{G}$
是零化上句所列全部局部截面的 $\mathcal{O}_2$-线性映射，
则 $m \mapsto L([m])$ 是阶数为 $k$ 的微分算子。因此
$\mathcal{P}^k_{\mathcal{O}_2/\mathcal{O}_1}(\mathcal{F})$
是 $\mathcal{O}_2[\mathcal{F}]$ 对由这些局部截面生成的
$\mathcal{O}_2$-子模的商。
\end{proof}

\begin{definition}
\label{sites-modules-definition-module-principal-parts}
设 $\mathcal{C}$ 为站点，
$\mathcal{O}_1 \to \mathcal{O}_2$ 为环层映射，
$\mathcal{F}$ 为 $\mathcal{O}_2$-模层。引理
\ref{sites-modules-lemma-module-principal-parts} 中构造的模
$\mathcal{P}^k_{\mathcal{O}_2/\mathcal{O}_1}(\mathcal{F})$
称为 $\mathcal{F}$ 的{\it $k$ 阶主部模}。
\end{definition}

\noindent
注意，包含关系
$$
\text{Diff}^0(\mathcal{F}, \mathcal{G}) \subset
\text{Diff}^1(\mathcal{F}, \mathcal{G}) \subset
\text{Diff}^2(\mathcal{F}, \mathcal{G}) \subset \ldots
$$
由 Yoneda 引理（《范畴》，引理 \ref{categories-lemma-yoneda}）
对应于满射
$$
\ldots \to \mathcal{P}^2_{\mathcal{O}_2/\mathcal{O}_1}(\mathcal{F})
\to \mathcal{P}^1_{\mathcal{O}_2/\mathcal{O}_1}(\mathcal{F})
\to \mathcal{P}^0_{\mathcal{O}_2/\mathcal{O}_1}(\mathcal{F}) = \mathcal{F}
$$

\begin{lemma}
\label{sites-modules-lemma-differential-operators-sheafify}
设 $\mathcal{C}$ 为站点，$\mathcal{O}_1 \to \mathcal{O}_2$
为环预层同态，并设 $\mathcal{F}$ 为
$\mathcal{O}_2$-模预层。则
$\mathcal{P}^k_{\mathcal{O}_2^\#/\mathcal{O}_1^\#}(\mathcal{F}^\#)$
是预层
$U \mapsto P^k_{\mathcal{O}_2(U)/\mathcal{O}_1(U)}(\mathcal{F}(U))$
的关联层。
\end{lemma}

\begin{proof}
这可以用与引理 \ref{sites-modules-lemma-differentials-sheafify}
中微分层完全相同的方法证明。或许更自然的方法是直接利用引理
\ref{sites-modules-lemma-module-principal-parts} 的泛性质看出该等式。细节从略。
\end{proof}

\begin{lemma}
\label{sites-modules-lemma-sequence-of-principal-parts}
设 $\mathcal{C}$ 为站点，$\mathcal{O}_1 \to \mathcal{O}_2$
为环层同态，并设 $\mathcal{F}$ 为 $\mathcal{O}_2$-模层。
有典范短正合列
$$
0 \to
\Omega_{\mathcal{O}_2/\mathcal{O}_1} \otimes_{\mathcal{O}_2} \mathcal{F} \to
\mathcal{P}^1_{\mathcal{O}_2/\mathcal{O}_1}(\mathcal{F}) \to
\mathcal{F} \to 0
$$
它关于 $\mathcal{F}$ 函子性成立，称为{\it 主部序列}。
\end{lemma}

\begin{proof}
由交换代数版本（《代数》，引理
\ref{algebra-lemma-sequence-of-principal-parts}）及引理
\ref{sites-modules-lemma-differentials-sheafify}、\ref{sites-modules-lemma-differential-operators-sheafify}
即得。
\end{proof}

\begin{remark}
\label{sites-modules-remark-functoriality-principal-parts}
设 $\mathcal{C}$ 为站点，并给定环层的交换图
$$
\xymatrix{
\mathcal{B} \ar[r] & \mathcal{B}' \\
\mathcal{A} \ar[u] \ar[r] & \mathcal{A}' \ar[u]
}
$$
以及 $\mathcal{B}$-模 $\mathcal{F}$、
$\mathcal{B}'$-模 $\mathcal{F}'$ 和
$\mathcal{B}$-线性映射 $\mathcal{F} \to \mathcal{F}'$。
则得到相容的模映射系
$$
\xymatrix{
\ldots \ar[r] &
\mathcal{P}^2_{\mathcal{B}'/\mathcal{A}'}(\mathcal{F}') \ar[r] &
\mathcal{P}^1_{\mathcal{B}'/\mathcal{A}'}(\mathcal{F}') \ar[r] &
\mathcal{P}^0_{\mathcal{B}'/\mathcal{A}'}(\mathcal{F}') \\
\ldots \ar[r] &
\mathcal{P}^2_{\mathcal{B}/\mathcal{A}}(\mathcal{F}) \ar[r] \ar[u] &
\mathcal{P}^1_{\mathcal{B}/\mathcal{A}}(\mathcal{F}) \ar[r] \ar[u] &
\mathcal{P}^0_{\mathcal{B}/\mathcal{A}}(\mathcal{F}) \ar[u]
}
$$
这些映射与此类映射的进一步复合相容。最容易的验证方式，是使用
引理 \ref{sites-modules-lemma-module-principal-parts} 的证明中对模
$\mathcal{P}^k_{\mathcal{B}/\mathcal{A}}(\mathcal{M})$
用（局部）生成元与关系给出的描述；也可以直接从这些模的泛性质看出。
此外，这些映射与引理
\ref{sites-modules-lemma-sequence-of-principal-parts} 的短正合列相容。
\end{remark}








\section{朴素余切复形}
\label{sites-modules-section-netherlander}

\noindent
本节是《代数》，第 \ref{algebra-section-netherlander} 节
及《模》，第 \ref{modules-section-netherlander} 节的对应版本。
建议读者先阅读这些章节。

\medskip\noindent
设 $\mathcal{C}$ 为站点，$\mathcal{A} \to \mathcal{B}$
为 $\mathcal{C}$ 上的环层同态。本节中，对 $\mathcal{C}$
上的任意集合层 $\mathcal{E}$，记
$\mathcal{A}[\mathcal{E}]$ 为预层
$U \mapsto \mathcal{A}(U)[\mathcal{E}(U)]$ 的层化。这里
$\mathcal{A}(U)[\mathcal{E}(U)]$
表示 $\mathcal{A}(U)$ 上的多项式代数，其变量与
$\mathcal{E}(U)$ 的元素一一对应。以
$[e] \in \mathcal{A}(U)[\mathcal{E}(U)]$ 表示对应于
$e \in \mathcal{E}(U)$ 的变量。有典范的
$\mathcal{A}$-代数满射
\begin{equation}
\label{sites-modules-equation-canonical-presentation}
\mathcal{A}[\mathcal{B}] \longrightarrow \mathcal{B},\quad [b] \longmapsto b
\end{equation}
记其核为 $\mathcal{I} \subset \mathcal{A}[\mathcal{B}]$。
容易看出，$\mathcal{I}$ 由局部截面
$[b][b'] - [bb']$ 与 $[a] - a$ 生成。由引理
\ref{sites-modules-lemma-differential-seq}，有典范映射
\begin{equation}
\label{sites-modules-equation-naive-cotangent-complex}
\mathcal{I}/\mathcal{I}^2
\longrightarrow
\Omega_{\mathcal{A}[\mathcal{B}]/\mathcal{A}}
\otimes_{\mathcal{A}[\mathcal{B}]} \mathcal{B}
\end{equation}
其余核典范同构于 $\Omega_{\mathcal{B}/\mathcal{A}}$。

\begin{definition}
\label{sites-modules-definition-naive-cotangent-complex}
设 $\mathcal{C}$ 为站点，$\mathcal{A} \to \mathcal{B}$
为 $\mathcal{C}$ 上的环层同态。{\it 朴素余切复形}
$\NL_{\mathcal{B}/\mathcal{A}}$ 是链复形
(\ref{sites-modules-equation-naive-cotangent-complex})
$$
\NL_{\mathcal{B}/\mathcal{A}} =
\left(\mathcal{I}/\mathcal{I}^2
\longrightarrow
\Omega_{\mathcal{A}[\mathcal{B}]/\mathcal{A}}
\otimes_{\mathcal{A}[\mathcal{B}]} \mathcal{B}\right)
$$
其中 $\mathcal{I}/\mathcal{I}^2$ 置于次数 $-1$，而
$\Omega_{\mathcal{A}[\mathcal{B}]/\mathcal{A}}
\otimes_{\mathcal{A}[\mathcal{B}]} \mathcal{B}$
置于次数 $0$。
\end{definition}

\noindent
此构造满足与引理 \ref{sites-modules-lemma-functoriality-differentials}
中微分模的函子性类似的函子性。具体地，给定交换图
\begin{equation}
\label{sites-modules-equation-commutative-square-sheaves}
\vcenter{
\xymatrix{
\mathcal{B} \ar[r] & \mathcal{B}' \\
\mathcal{A} \ar[u] \ar[r] & \mathcal{A}' \ar[u]
}
}
\end{equation}
它是 $\mathcal{C}$ 上环层的图，则有典范的
$\mathcal{B}$-线性复形映射
$$
\NL_{\mathcal{B}/\mathcal{A}} \longrightarrow \NL_{\mathcal{B}'/\mathcal{A}'}
$$
事实上，交换图中的映射诱导典范映射
$\mathcal{A}[\mathcal{B}] \to \mathcal{A}'[\mathcal{B}']$
它把 $\mathcal{I}$ 映入
$\mathcal{I}' = \Ker(\mathcal{A}'[\mathcal{B}'] \to \mathcal{B}')$。
于是得到映射
$\mathcal{I}/\mathcal{I}^2 \to \mathcal{I}'/(\mathcal{I}')^2$
及微分模之间的映射；二者共同给出朴素余切复形之间所需的映射。

\medskip\noindent
也可以选取另一个把 $\mathcal{B}$ 表示为
$\mathcal{A}$ 上多项式代数之商的表示，而仍在 $D(\mathcal{B})$
中得到同一对象。为说明这一点，设 $\mathcal{E}$ 是
$\mathcal{C}$ 上的集合层，且
$\alpha : \mathcal{E} \to \mathcal{B}$ 是集合层的映射。
于是得到 $\mathcal{A}$-代数同态
$\mathcal{A}[\mathcal{E}] \to \mathcal{B}$。假设该映射满射，
并令 $\mathcal{J} \subset \mathcal{A}[\mathcal{E}]$ 为其核。置
$$
\NL(\alpha) = \left(
\mathcal{J}/\mathcal{J}^2
\longrightarrow
\Omega_{\mathcal{A}[\mathcal{E}]/\mathcal{A}}
\otimes_{\mathcal{A}[\mathcal{E}]} \mathcal{B}\right)
$$
结论如下。

\begin{lemma}
\label{sites-modules-lemma-NL-up-to-qis}
在上述情形下，$D(\mathcal{B})$ 中有典范同构
$\NL(\alpha) = \NL_{\mathcal{B}/\mathcal{A}}$。
\end{lemma}

\begin{proof}
注意 $\NL_{\mathcal{B}/\mathcal{A}} = \NL(\text{id}_\mathcal{B})$。
因此只需证明：给定上述两个映射
$\alpha_i : \mathcal{E}_i \to \mathcal{B}$，在 $D(\mathcal{B})$
中存在典范拟同构
$\NL(\alpha_1) = \NL(\alpha_2)$。为此，置
$\mathcal{E} = \mathcal{E}_1 \amalg \mathcal{E}_2$ 且
$\alpha = \alpha_1 \amalg \alpha_2 : \mathcal{E} \to \mathcal{B}$.
再置
$\mathcal{J}_i = \Ker(\mathcal{A}[\mathcal{E}_i] \to \mathcal{B})$
以及
$\mathcal{J} = \Ker(\mathcal{A}[\mathcal{E}] \to \mathcal{B})$.
得到把 $\mathcal{J}_i$ 映入 $\mathcal{J}$ 的映射
$\mathcal{A}[\mathcal{E}_i] \to \mathcal{A}[\mathcal{E}]$。
从而得到典范复形映射
$$
\NL(\alpha_i) \longrightarrow \NL(\alpha)
$$
只需证明这些映射是拟同构。为此作如下论证。首先，由引理
\ref{sites-modules-lemma-differential-seq}，
$H^0(\NL(\alpha_i)) = \Omega_{\mathcal{B}/\mathcal{A}}$ 且
$H^0(\NL(\alpha)) = \Omega_{\mathcal{B}/\mathcal{A}}$；
故该映射在次数 $0$ 的上同调层上是同构。类似地，我们断言
$H^{-1}(\NL(\alpha_i))$ 与 $H^{-1}(\NL(\alpha))$ 均为预层
$U \mapsto H_1(L_{\mathcal{B}(U)/\mathcal{A}(U)})$ 的关联层，其中
$H_1(L_{-/-})$ 如《代数》，定义
\ref{algebra-definition-naive-cotangent-complex}。若此断言成立，
证明便完成。

\medskip\noindent
证明该断言。设 $\alpha : \mathcal{E} \to \mathcal{B}$
如上。令 $\mathcal{B}' \subset \mathcal{B}$ 为
$\mathcal{A}$-代数子预层，它在 $U$ 上的取值是
$\mathcal{A}(U)[\mathcal{E}(U)] \to \mathcal{B}(U)$ 的像。
令 $\mathcal{I}'$ 为预层，它在 $U$ 上的取值是
$\mathcal{A}(U)[\mathcal{E}(U)] \to \mathcal{B}(U)$ 的核。
于是 $\mathcal{I}$ 是 $\mathcal{I}'$ 的层化，而
$\mathcal{B}$ 是 $\mathcal{B}'$ 的层化。类似地，
$H^{-1}(\NL(\alpha))$ 是下列预层的层化：
$$
U \longmapsto
\Ker(\mathcal{I}'(U)/\mathcal{I}'(U)^2 \to
\Omega_{\mathcal{A}(U)[\mathcal{E}(U)]/\mathcal{A}(U)}
\otimes_{\mathcal{A}(U)[\mathcal{E}(U)]} \mathcal{B}'(U))
$$
这是引理 \ref{sites-modules-lemma-differentials-sheafify} 的结论。
由《代数》，引理 \ref{algebra-lemma-NL-homotopy}，可知
$H^{-1}(\NL(\alpha))$ 是预层
$U \mapsto H_1(L_{\mathcal{B}'(U)/\mathcal{A}(U)})$ 的关联层。
因此只需证明映射
$H_1(L_{\mathcal{B}'(U)/\mathcal{A}(U)}) \to
H_1(L_{\mathcal{B}(U)/\mathcal{A}(U)})$ 在层化后诱导同构
$\mathcal{H}'_1 \to \mathcal{H}_1$。

\medskip\noindent
证明 $\mathcal{H}'_1 \to \mathcal{H}_1$ 的单射性。设
$f \in H_1(L_{\mathcal{B}'(U)/\mathcal{A}(U)})$ 在
$\mathcal{H}_1(U)$ 中映到零。须证：$f$ 在
$\mathcal{H}'_1(U)$ 中映到零。该假设意味着存在覆盖
$\{U_i \to U\}$，使得对所有 $i$，$f$ 在
$H_1(L_{\mathcal{B}(U_i)/\mathcal{A}(U_i)})$ 中映到零。
以 $U_i$ 替换 $U$，即可归约到 $f$ 在
$H_1(L_{\mathcal{B}(U)/\mathcal{A}(U)})$ 中映到零的情形。
由《代数》，引理 \ref{algebra-lemma-colimits-NL}，可找到有限生成子代数
$\mathcal{B}'(U) \subset B \subset \mathcal{B}(U)$，使得
$f$ 在 $H_1(L_{B/\mathcal{A}(U)})$ 中映到零。由于
$\mathcal{B} = (\mathcal{B}')^\#$，可找到覆盖
$\{U_i \to U\}$，使得 $B \to \mathcal{B}(U_i)$ 通过
$\mathcal{B}'(U_i)$ 分解。因此 $f$ 在
$H_1(L_{\mathcal{B}'(U_i)/\mathcal{A}(U_i)})$ 中映到零，如所需。

\medskip\noindent
$\mathcal{H}'_1 \to \mathcal{H}_1$ 的满射性完全类似地证明。
\end{proof}

\begin{lemma}
\label{sites-modules-lemma-pullback-NL}
设 $f : \Sh(\mathcal{C}) \to \Sh(\mathcal{D})$ 为拓扑斯态射，
$\mathcal{A} \to \mathcal{B}$ 为 $\mathcal{D}$ 上的环层同态。
则 $f^{-1}\NL_{\mathcal{B}/\mathcal{A}} =
\NL_{f^{-1}\mathcal{B}/f^{-1}\mathcal{A}}$.
\end{lemma}

\begin{proof}
从略。提示：使用引理 \ref{sites-modules-lemma-pullback-differentials}。
\end{proof}

\noindent
环拓扑斯态射的余切复形用上面定义的余切复形来定义。

\begin{definition}
\label{sites-modules-definition-cotangent-complex-morphism-ringed-topoi}
设 $X = (\Sh(\mathcal{C}), \mathcal{O})$ 与
$Y = (\Sh(\mathcal{C}'), \mathcal{O}')$ 为环拓扑斯，并设
$(f, f^\sharp) : X \to Y$ 为环拓扑斯态射。给定环拓扑斯态射的
{\it 朴素余切复形} $\NL_f = \NL_{X/Y}$ 定义为
$\NL_{\mathcal{O}/f^{-1}\mathcal{O}'}$。
有时也写成 $\NL_{X/Y} = \NL_{\mathcal{O}/\mathcal{O}'}$。
\end{definition}









\section{模的茎}
\label{sites-modules-section-stalks}

\noindent
在点处取茎时必须稍加小心，因为定义茎的余极限（见《站点》，公式
\ref{sites-equation-stalk}）未必是滤过的\footnote{当然，在几乎所有
自然出现的情形中，该余极限都是滤过的，因而本节的部分讨论可以简化。}。
另一方面，由站点之点的定义，茎函子正合且与任意余极限交换。
换言之，它的行为恰如该余极限是滤过的一样。

\begin{lemma}
\label{sites-modules-lemma-stalk-exact}
设 $\mathcal{C}$ 为站点，$p$ 为 $\mathcal{C}$ 的一个点。
\begin{enumerate}
\item 对 $\mathcal{C}$ 上任意集合预层，有
$(\mathcal{F}^\#)_p = \mathcal{F}_p$。
\item 茎函子
$\Sh(\mathcal{C}) \to \textit{Sets}$,
$\mathcal{F} \mapsto \mathcal{F}_p$ 正合（见《范畴》，定义
\ref{categories-definition-exact}），并与任意余极限交换。
\item 茎函子
$\textit{PSh}(\mathcal{C}) \to \textit{Sets}$,
$\mathcal{F} \mapsto \mathcal{F}_p$ 正合（见《范畴》，定义
\ref{categories-definition-exact}），并与任意余极限交换。
\end{enumerate}
\end{lemma}

\begin{proof}
由《站点》，引理 \ref{sites-lemma-point-pushforward-sheaf}
得到 (1)。由《站点》，引理
\ref{sites-lemma-adjoint-point-push-stalk} 可知
$\textit{PSh}(\mathcal{C}) \to \textit{Sets}$,
$\mathcal{F} \mapsto \mathcal{F}_p$ 是左伴随；由《站点》，引理
\ref{sites-lemma-point-pushforward-sheaf}，对
$\Sh(\mathcal{C}) \to \textit{Sets}$,
$\mathcal{F} \mapsto \mathcal{F}_p$ 也有相同结论。
故茎函子与任意余极限交换（见《范畴》，引理
\ref{categories-lemma-adjoint-exact}）。由站点之点的定义
（见《站点》，定义 \ref{sites-definition-point}），
$\Sh(\mathcal{C}) \to \textit{Sets}$,
$\mathcal{F} \mapsto \mathcal{F}_p$
正合。由于层化正合（《站点》，引理
\ref{sites-lemma-sheafification-exact}），可知
$\textit{PSh}(\mathcal{C}) \to \textit{Sets}$,
$\mathcal{F} \mapsto \mathcal{F}_p$
也正合。
\end{proof}

\noindent
特别地，由于预层上的茎函子
$\mathcal{F} \mapsto \mathcal{F}_p$ 与所有有限极限和余极限交换，
可以应用《站点》，命题
\ref{sites-proposition-functoriality-algebraic-structures-topoi}
证明中的论证。其结论是：若 $\mathcal{F}$ 是《站点》命题
\ref{sites-proposition-functoriality-algebraic-structures-topoi} 所列代数结构的
（预）层，则茎 $\mathcal{F}_p$ 自然成为同类代数结构。下面在
$\mathcal{F}$ 为阿贝尔预层时详细说明。此时加法映射
$+ : \mathcal{F} \times \mathcal{F} \to \mathcal{F}$ 诱导映射
$$
+ :
\mathcal{F}_p \times \mathcal{F}_p
=
(\mathcal{F} \times \mathcal{F})_p
\longrightarrow
\mathcal{F}_p
$$
其中等号用到了集合预层上的茎函子与有限极限交换。这在茎
$\mathcal{F}_p$ 上定义群结构。由此得到茎函子
$$
\textit{PAb}(\mathcal{C}) \longrightarrow \textit{Ab}, \quad
\mathcal{F} \longmapsto \mathcal{F}_p
$$
按构造，$\mathcal{F}_p$ 的底集合就是其底集合预层的茎。
再与包含
$\textit{Ab}(\mathcal{C}) \subset \textit{PAb}(\mathcal{C})$
预复合，也就定义了阿贝尔群层的茎函子。

\begin{lemma}
\label{sites-modules-lemma-stalk-exact-abelian}
设 $\mathcal{C}$ 为站点，$p$ 为 $\mathcal{C}$ 的一个点。
\begin{enumerate}
\item 函子
$\textit{Ab}(\mathcal{C}) \to \textit{Ab}$,
$\mathcal{F} \mapsto \mathcal{F}_p$ 正合。
\item 茎函子
$\textit{PAb}(\mathcal{C}) \to \textit{Ab}$,
$\mathcal{F}  \mapsto  \mathcal{F}_p$
正合。
\item 对 $\mathcal{F} \in \Ob(\textit{PAb}(\mathcal{C}))$，
有 $\mathcal{F}_p = \mathcal{F}^\#_p$。
\end{enumerate}
\end{lemma}

\begin{proof}
由引理 \ref{sites-modules-lemma-stalk-exact} 的结论和上述茎函子的构造形式地得到。
\end{proof}

\noindent
下面转向模层的情形。设 $(\mathcal{C}, \mathcal{O})$ 为环站点
（对以下讨论，只需 $\mathcal{O}$ 是环预层），
$\mathcal{F}$ 为 $\mathcal{O}$-模预层，$p$ 为
$\mathcal{C}$ 的一个点。此时得到映射
$$
\cdot :
\mathcal{O}_p \times \mathcal{O}_p
=
(\mathcal{O} \times \mathcal{O})_p
\longrightarrow
\mathcal{O}_p
$$
它是乘法映射的茎；以及
$$
\cdot :
\mathcal{O}_p \times \mathcal{F}_p
=
(\mathcal{O} \times \mathcal{F})_p
\longrightarrow
\mathcal{F}_p
$$
它也是乘法映射的茎。略去验证这在 $\mathcal{O}_p$
上定义环结构，并在 $\mathcal{F}_p$ 上定义
$\mathcal{O}_p$-模结构。由此得到函子
$$
\textit{PMod}(\mathcal{O}) \longrightarrow \textit{Mod}(\mathcal{O}_p), \quad
\mathcal{F} \longmapsto \mathcal{F}_p
$$
按构造，$\mathcal{F}_p$ 的底集合就是其底集合预层的茎。
再与包含
$\textit{Mod}(\mathcal{O}) \subset \textit{PMod}(\mathcal{O})$
预复合，也就定义了 $\mathcal{O}$-模层的茎函子。

\begin{lemma}
\label{sites-modules-lemma-stalk-exact-modules}
设 $(\mathcal{C}, \mathcal{O})$ 为环站点，
$p$ 为 $\mathcal{C}$ 的一个点。
\begin{enumerate}
\item 函子
$\textit{Mod}(\mathcal{O}) \to \textit{Mod}(\mathcal{O}_p)$,
$\mathcal{F} \mapsto \mathcal{F}_p$ 正合。
\item 茎函子
$\textit{PMod}(\mathcal{O}) \to \textit{Mod}(\mathcal{O}_p)$,
$\mathcal{F} \mapsto \mathcal{F}_p$
正合。
\item 对 $\mathcal{F} \in \Ob(\textit{PMod}(\mathcal{O}))$，
有 $\mathcal{F}_p = \mathcal{F}^\#_p$。
\end{enumerate}
\end{lemma}

\begin{proof}
由引理 \ref{sites-modules-lemma-stalk-exact-abelian} 的结论、
上述茎函子的构造及引理 \ref{sites-modules-lemma-abelian} 形式地得到。
\end{proof}

\begin{lemma}
\label{sites-modules-lemma-pullback-stalk}
设
$(f, f^\sharp) :
(\Sh(\mathcal{C}), \mathcal{O}_\mathcal{C})
\to
(\Sh(\mathcal{D}), \mathcal{O}_\mathcal{D})$
为环拓扑斯或环站点的态射。设 $p$ 为 $\mathcal{C}$
或 $\Sh(\mathcal{C})$ 的一个点，并置 $q = f \circ p$。则对任意
$\mathcal{O}_\mathcal{D}$-模 $\mathcal{F}$，有
$$
(f^*\mathcal{F})_p =
\mathcal{F}_q \otimes_{\mathcal{O}_{\mathcal{D}, q}}
\mathcal{O}_{\mathcal{C}, p}
$$
\end{lemma}

\begin{proof}
由定义，
$$
f^*\mathcal{F} =
f^{-1}\mathcal{F} \otimes_{f^{-1}\mathcal{O}_\mathcal{D}}
\mathcal{O}_\mathcal{C}
$$
由于在 $p$ 处取茎（即应用 $p^{-1}$）由引理
\ref{sites-modules-lemma-tensor-product-pullback} 与 $\otimes$ 交换，
再用《站点》，引理 \ref{sites-lemma-point-morphism-sites}
或引理 \ref{sites-lemma-point-morphism-topoi}
中所述的 $p$ 处拉回之茎与 $q$ 处茎的关系，即得结论。
\end{proof}






\section{摩天楼层}
\label{sites-modules-section-skyscraper}

\noindent
设 $p$ 为站点 $\mathcal{C}$ 或拓扑斯
$\Sh(\mathcal{C})$ 的一个点。本节研究把阿贝尔群 $A$
映到摩天楼层 $p_*A$ 的函子的正合性。首先回忆，
$p_* : \textit{Sets} \to \Sh(\mathcal{C})$ 具有许多正合性；
见《站点》，引理 \ref{sites-lemma-stalk-skyscraper} 及
\ref{sites-lemma-skyscraper-functor-exact}。

\begin{lemma}
\label{sites-modules-lemma-skyscraper-exact}
设 $\mathcal{C}$ 为站点，$p$ 为 $\mathcal{C}$
或其关联拓扑斯的一个点。
\begin{enumerate}
\item 函子 $p_* : \textit{Ab} \to \textit{Ab}(\mathcal{C})$,
$A \mapsto p_*A$ 正合。
\item 有函子性的直和分解
$$
p^{-1}p_*A = A \oplus I(A)
$$
其中 $A \in \Ob(\textit{Ab})$。
\end{enumerate}
\end{lemma}

\begin{proof}
由《站点》，引理 \ref{sites-lemma-stalk-skyscraper}，有函子性映射
$A \to p^{-1}p_*A \to A$，其复合等于 $\text{id}_A$。
故得到 (2) 中的函子性直和分解，其中 $I(A)$ 是伴随映射
$p^{-1}p_*A \to A$ 的核。由引理
\ref{sites-modules-lemma-exactness-pushforward-pullback}，函子 $p_*$ 左正合。
由《站点》，引理 \ref{sites-lemma-skyscraper-functor-exact}，
函子 $p_*$ 把满射变为满射。因此 (1) 成立。
\end{proof}

\noindent
为对模层作相同处理，设 $p$ 为环拓扑斯
$(\Sh(\mathcal{C}), \mathcal{O})$ 的一个点。回忆
$p^{-1}$ 就是茎函子。因此可以把 $p$ 看成环拓扑斯态射
$$
(p, \text{id}_{\mathcal{O}_p}) :
(\Sh(pt), \mathcal{O}_p)
\longrightarrow
(\Sh(\mathcal{C}), \mathcal{O}).
$$
于是得到拉回函子
$p^* : \textit{Mod}(\mathcal{O}) \to \textit{Mod}(\mathcal{O}_p)$
它等于茎函子，已在引理 \ref{sites-modules-lemma-stalk-exact-modules} 中讨论。
本节考虑函子
$p_* : \textit{Mod}(\mathcal{O}_p) \to \textit{Mod}(\mathcal{O})$.

\begin{lemma}
\label{sites-modules-lemma-skyscraper-modules-exact}
设 $(\Sh(\mathcal{C}), \mathcal{O})$ 为环拓扑斯，
$p$ 为拓扑斯 $\Sh(\mathcal{C})$ 的一个点。
\begin{enumerate}
\item 函子
$p_* : \textit{Mod}(\mathcal{O}_p) \to \textit{Mod}(\mathcal{O})$,
$M \mapsto p_*M$ 正合。
\item 典范满射 $p^{-1}p_*M \to M$ 是 $\mathcal{O}_p$-线性的。
\item 引理 \ref{sites-modules-lemma-skyscraper-exact} 的函子性直和分解
$p^{-1}p_*M = M \oplus I(M)$
一般而言{\bf 不是} $\mathcal{O}_p$-线性的。
\end{enumerate}
\end{lemma}

\begin{proof}
(1) 及 (2) 中的满射性立即由引理 \ref{sites-modules-lemma-skyscraper-exact}
中阿贝尔层的相应结论得到。由于
$p^{-1}\mathcal{O} = \mathcal{O}_p$，有 $p^{-1} = p^*$；
故 $p^{-1}p_*M \to M$ 就是模的伴随中的余单位
$p^*p_*M \to M$，因而是线性的。

\medskip\noindent
证明 (3)。设 $G$ 为群。考虑拓扑斯
$G\textit{-Sets} = \Sh(\mathcal{T}_G)$
及点 $p : \textit{Sets} \to G\textit{-Sets}$。见《站点》，第
\ref{sites-section-example-sheaf-G-sets} 节及例
\ref{sites-example-point-G-sets}。这里 $p^{-1}$
是遗忘 $G$-作用的函子，而 $p_*$ 是该遗忘函子的右伴随，
将 $M$ 映到 $\text{Map}(G, M)$。直和分解中的映射为
$$
M \to \text{Map}(G, M) \to M
$$
其中第一个映射把 $m \in M$ 映到值恒为 $m$ 的常值映射，
第二个映射是在 $G$ 的单位元 $1$ 处取值。再设 $R$
为配备 $G$-作用的环。这在 $\mathcal{T}_G$ 上确定环层
$\mathcal{O}$。$\mathcal{O}$-模范畴就是配备与 $R$
上的作用相容的 $G$-作用的 $R$-模 $M$ 所成范畴。
$\text{Map}(G, M)$ 上的 $R$-模结构由
$$
( r f ) (\sigma) = \sigma(r) f(\sigma)
$$
给出，其中 $r \in R$ 且 $f \in \text{Map}(G, M)$。这是因为它是
与在 $1$ 处取值相容的唯一 $G$-不变 $R$-模结构。容易看出，
一般而言 $M \to \text{Map}(G, M)$ 的像并不是 $R$-子模
（例如取 $M = R$ 并假设 $G$-作用非平凡），证明完成。
\end{proof}

\begin{example}
\label{sites-modules-example-ring-with-group-action}
设 $G$ 为群。考虑站点 $\mathcal{T}_G$ 及其点 $p$；
见《站点》，例 \ref{sites-example-point-G-sets}。设
$R$ 为配备 $G$-作用的环，它对应于 $\mathcal{T}_G$
上的环层 $\mathcal{O}$。则在遗忘 $G$-作用后，
$\mathcal{O}_p = R$。在此情形下
$p^{-1}p_*M = \text{Map}(G, M)$，且
$I(M) = \{f : G \to M \mid f(1_G) = 0\}$，而
$M \to \text{Map}(G, M)$ 把 $m \in M$ 映到值恒为 $m$ 的常值函数。
\end{example}




\section{局部化与点}
\label{sites-modules-section-localize-points}

\begin{lemma}
\label{sites-modules-lemma-stalk-j-shriek}
设 $(\mathcal{C}, \mathcal{O})$ 为环站点，
$p$ 为 $\mathcal{C}$ 的一个点，$U$ 为 $\mathcal{C}$ 的对象。
对 $\textit{Mod}(\mathcal{O}_U)$ 中的 $\mathcal{G}$，有
$$
(j_{U!}\mathcal{G})_p =
\bigoplus\nolimits_q \mathcal{G}_q
$$
其中余积遍历位于 $p$ 上方的 $\mathcal{C}/U$ 的各点 $q$；
见《站点》，引理 \ref{sites-lemma-points-above-point}。
\end{lemma}

\begin{proof}
使用引理 \ref{sites-modules-lemma-extension-by-zero} 中的描述：
$j_{U!}\mathcal{G}$ 是预层
$V \mapsto
\bigoplus\nolimits_{\varphi \in \Mor_\mathcal{C}(V, U)}
\mathcal{G}(V/_\varphi U)$
的关联层。由引理 \ref{sites-modules-lemma-stalk-exact-modules}，
$j_{U!}\mathcal{G}$ 在 $p$ 处的茎等于此预层的茎。
设 $u : \mathcal{C} \to \textit{Sets}$ 为对应于 $p$ 的函子
（见《站点》，第 \ref{sites-section-points} 节）。于是
$$
(j_{U!}\mathcal{G})_p = \colim_{(V, y)}
\bigoplus\nolimits_{\varphi : V \to U} \mathcal{G}(V/_\varphi U)
$$
其中余极限在阿贝尔群范畴中取。对该余极限中出现的四元组
$(V, y, \varphi, s)$，可以赋予
$x = u(\varphi)(y) \in u(U)$。因此得到
$$
(j_{U!}\mathcal{G})_p =
\bigoplus\nolimits_{x \in u(U)}
\colim_{(\varphi : V \to U, y), \ u(\varphi)(y) = x} \mathcal{G}(V/_\varphi U).
$$
由《站点》，引理 \ref{sites-lemma-points-above-point}
中对位于 $x$ 上方的点 $q$ 的描述，这等于引理中的表达式。
\end{proof}

\begin{remark}
\label{sites-modules-remark-not-pushforward}
警告：引理 \ref{sites-modules-lemma-stalk-j-shriek} 的结论对
$j_{U, *}$ 没有类似物。
\end{remark}












\section{平坦模的拉回}
\label{sites-modules-section-pullback-flat}

\noindent
平坦模沿环拓扑斯态射的拉回仍然平坦。证明这一点略有困难。

\begin{lemma}
\label{sites-modules-lemma-pullback-flat}
\begin{reference}
\cite[Expos\'e V, Corollary 1.7.1]{SGA4}
\end{reference}
设
$(f, f^\sharp) :
(\Sh(\mathcal{C}), \mathcal{O}_\mathcal{C})
\to
(\Sh(\mathcal{D}), \mathcal{O}_\mathcal{D})$
为环拓扑斯或环站点的态射。若 $\mathcal{F}$ 是平坦
$\mathcal{O}_\mathcal{D}$-模，则 $f^*\mathcal{F}$
是平坦 $\mathcal{O}_\mathcal{C}$-模。
\end{lemma}

\begin{proof}
选取引理
\ref{sites-modules-lemma-morphism-ringed-topoi-comes-from-morphism-ringed-sites}
中所述的图。回忆，成为平坦模是内蕴性质（见第
\ref{sites-modules-section-intrinsic} 节及定义 \ref{sites-modules-definition-flat}）。
故只需对态射 $(h, h^\sharp) :
(\Sh(\mathcal{C}'), \mathcal{O}_{\mathcal{C}'})
\to
(\Sh(\mathcal{D}'), \mathcal{O}_{\mathcal{D}'})$ 证明引理。
换言之，可以假设站点 $\mathcal{C}$ 与 $\mathcal{D}$
具有所有有限极限，并且 $f$ 是由与有限极限交换的连续函子
$u : \mathcal{D} \to \mathcal{C}$ 诱导的站点态射。

\medskip\noindent
回忆 $f^*\mathcal{F} =
\mathcal{O}_\mathcal{C} \otimes_{f^{-1}\mathcal{O}_\mathcal{D}}
f^{-1}\mathcal{F}$（定义 \ref{sites-modules-definition-pushforward}）。
由引理 \ref{sites-modules-lemma-flat-change-of-rings}，只需证明
$f^{-1}\mathcal{F}$ 是平坦
$f^{-1}\mathcal{O}_\mathcal{D}$-模。结合上一段，问题归约到下一段的情形。

\medskip\noindent
假设 $\mathcal{C}$ 与 $\mathcal{D}$ 是具有所有有限极限的站点，
且 $u : \mathcal{D} \to \mathcal{C}$ 是与有限极限交换的连续函子。
设 $\mathcal{O}$ 为 $\mathcal{D}$ 上的环层，
$\mathcal{F}$ 为平坦 $\mathcal{O}$-模。则 $u$
定义站点态射 $f : \mathcal{C} \to \mathcal{D}$
（《站点》，命题 \ref{sites-proposition-get-morphism}）。
须证：$f^{-1}\mathcal{F}$ 是平坦
$f^{-1}\mathcal{O}$-模。设 $U$ 为 $\mathcal{C}$ 的对象，并设
$$
f^{-1}\mathcal{O}|_U \xrightarrow{(f_1, \ldots, f_n)}
f^{-1}\mathcal{O}|_U^{\oplus n} \xrightarrow{(s_1, \ldots, s_n)}
f^{-1}\mathcal{F}|_U
$$
为 $f^{-1}\mathcal{O}|_U$-模的复形。目标是在 $U$
的某个覆盖的各成员上，构造 $(s_1, \ldots, s_n)$ 的分解，
如引理 \ref{sites-modules-lemma-flat-eq} 的 (2) 所述。考虑元素
$s_a \in f^{-1}\mathcal{F}(U)$ 与
$f_a \in f^{-1}\mathcal{O}(U)$。由于
$f^{-1}\mathcal{F}$（相应地 $f^{-1}\mathcal{O}$）是
$u_p\mathcal{F}$ 的层化，将 $U$ 替换为某个覆盖的各成员后，
可以假设 $s_a$ 是某个 $s'_a \in u_p\mathcal{F}(U)$ 的像，
而 $f_a$ 是某个 $f'_a \in u_p\mathcal{O}(U)$ 的像。再次将
$U$ 替换为某个覆盖的各成员后，可以假设
$\sum f'_as'_a$ 在 $u_p\mathcal{F}(U)$ 中为零。回忆范畴
$(\mathcal{I}_U^u)^{opp}$ 是有向的（《站点》，引理
\ref{sites-lemma-directed}），并且
$u_p\mathcal{F}(U) = \colim_{(\mathcal{I}_U^u)^{opp}} \mathcal{F}(V)$，
$u_p\mathcal{O}(U) = \colim_{(\mathcal{I}_U^u)^{opp}} \mathcal{O}(V)$。
因此可以假设存在
$(V, \phi) \in \Ob(\mathcal{I}_U^u)$，其中 $V$ 是
$\mathcal{D}$ 的对象，$\phi$ 是 $\mathcal{D}$ 中的态射
$\phi : U \to u(V)$，并且存在元素
$s''_a \in \mathcal{F}(V)$ 与 $f''_a \in \mathcal{O}(V)$，
它们在 $u_p\mathcal{F}(U)$ 与 $u_p\mathcal{O}(U)$ 中的像
分别等于 $s'_a$ 与 $f'_a$，且
$\sum f''_a s''_a = 0$ 在 $\mathcal{F}(V)$ 中成立。
于是得到复形
$$
\mathcal{O}|_V \xrightarrow{(f''_1, \ldots, f''_n)}
\mathcal{O}|_V^{\oplus n} \xrightarrow{(s''_1, \ldots, s''_n)}
\mathcal{F}|_V
$$
应用引理 \ref{sites-modules-lemma-flat-eq} 的另一方向，可知存在
$\mathcal{D}$ 中的覆盖 $\{V_i \to V\}$，并且对每个 $i$
都有分解
$$
\mathcal{O}|_{V_i}^{\oplus n}
\xrightarrow{B''_i}
\mathcal{O}|_{V_i}^{\oplus l_i} \xrightarrow{(t''_{i1}, \ldots, t''_{il_i})}
\mathcal{F}|_{V_i}
$$
它是 $(s''_1, \ldots, s''_n)|_{V_i}$ 的分解，且
$B_i \circ (f''_1, \ldots, f''_n)|_{V_i} = 0$.
置 $U_i = U \times_{\phi, u(V)} u(V_i)$，记
$B_i \in \text{Mat}(l_i \times n, f^{-1}\mathcal{O}(U_i))$
为 $B''_i$ 的像，并记
$t_{ij} \in f^{-1}\mathcal{F}(U_i)$ 为 $t''_{ij}$ 的像。
于是得到分解
$$
f^{-1}\mathcal{O}|_{U_i}^{\oplus n}
\xrightarrow{B_i}
f^{-1}\mathcal{O}|_{U_i}^{\oplus l_i}
\xrightarrow{(t_{i1}, \ldots, t_{il_i})}
\mathcal{F}|_{U_i}
$$
它是 $(s_1, \ldots, s_n)|_{U_i}$ 的分解，且
$B_i \circ (f_1, \ldots, f_n)|_{U_i} = 0$.
证明完成。
\end{proof}

\begin{lemma}
\label{sites-modules-lemma-stalk-flat}
设 $(\mathcal{C}, \mathcal{O})$ 为环站点，
$p$ 为 $\mathcal{C}$ 的一个点。若 $\mathcal{F}$
是平坦 $\mathcal{O}$-模，则 $\mathcal{F}_p$
是平坦 $\mathcal{O}_p$-模。
\end{lemma}

\begin{proof}
第 \ref{sites-modules-section-skyscraper} 节表明，可以把 $p$ 看成环拓扑斯态射
$$
(p, \text{id}_{\mathcal{O}_p}) :
(\Sh(pt), \mathcal{O}_p)
\longrightarrow
(\Sh(\mathcal{C}), \mathcal{O}).
$$
使得拉回函子
$p^* : \textit{Mod}(\mathcal{O}) \to \textit{Mod}(\mathcal{O}_p)$
等于茎函子。因此结论由引理 \ref{sites-modules-lemma-pullback-flat} 得到。
\end{proof}

\begin{lemma}
\label{sites-modules-lemma-check-flat-stalks}
设 $(\mathcal{C}, \mathcal{O})$ 为环站点，
$\mathcal{F}$ 为 $\mathcal{O}$-模层，并设
$\{p_i\}_{i \in I}$ 为 $\mathcal{C}$ 的一族保守点。
则 $\mathcal{F}$ 平坦，当且仅当对所有 $i \in I$，
$\mathcal{F}_{p_i}$ 是平坦 $\mathcal{O}_{p_i}$-模。
\end{lemma}

\begin{proof}
一个方向由引理 \ref{sites-modules-lemma-stalk-flat} 得到。反之，使用引理
\ref{sites-modules-lemma-tensor-product-pullback}（在 $p$ 处取茎由
$p^{-1}$ 给出）中的等式
$(\mathcal{F} \otimes_\mathcal{O} \mathcal{G})_p =
\mathcal{F}_p \otimes_{\mathcal{O}_p} \mathcal{G}_p$，
再应用引理 \ref{sites-modules-lemma-check-exactness-stalks}。
\end{proof}

\begin{lemma}
\label{sites-modules-lemma-pullback-ses}
设
$f : (\Sh(\mathcal{C}'), \mathcal{O}') \to (\Sh(\mathcal{C}), \mathcal{O})$
为环拓扑斯态射。设
$0 \to \mathcal{F} \to \mathcal{G} \to \mathcal{H} \to 0$
为 $\mathcal{O}$-模的短正合列，且 $\mathcal{H}$
是平坦 $\mathcal{O}$-模。则序列
$0 \to f^*\mathcal{F} \to f^*\mathcal{G} \to f^*\mathcal{H} \to 0$
也正合。
\end{lemma}

\begin{proof}
由于 $f^{-1}$ 正合，有 $f^{-1}\mathcal{O}$-模的短正合列
$0 \to f^{-1}\mathcal{F} \to f^{-1}\mathcal{G} \to f^{-1}\mathcal{H} \to 0$
由引理 \ref{sites-modules-lemma-pullback-flat}，
$f^{-1}\mathcal{O}$-模 $f^{-1}\mathcal{H}$ 平坦。由引理
\ref{sites-modules-lemma-flat-tor-zero}，这蕴含该序列在
$f^{-1}\mathcal{O}$ 上与 $\mathcal{O}'$ 作张量积后仍正合。由于
$f^*\mathcal{F} = f^{-1}\mathcal{F} \otimes_{f^{-1}\mathcal{O}} \mathcal{O}'$
，而 $\mathcal{G}$ 与 $\mathcal{H}$ 也同样，故得结论。
\end{proof}





\section{局部环拓扑斯}
\label{sites-modules-section-locally-ringed}


\noindent
本节的参考文献为 \cite[Expos\'e IV, Exercice 13.9]{SGA4}。

\begin{lemma}
\label{sites-modules-lemma-locally-ringed}
设 $(\mathcal{C}, \mathcal{O})$ 为环站点。下列条件等价：
\begin{enumerate}
\item 对 $\mathcal{C}$ 的每个对象 $U$ 及
$f \in \mathcal{O}(U)$，存在覆盖 $\{U_j \to U\}$，使得对每个
$j$，$f|_{U_j}$ 或 $(1 - f)|_{U_j}$ 可逆。
\item 若 $U \in \Ob(\mathcal{C})$、$n \geq 1$，且
$f_1, \ldots, f_n \in \mathcal{O}(U)$ 在 $\mathcal{O}(U)$
中生成单位理想，则存在覆盖 $\{U_j \to U\}$，使得对每个 $j$
都存在 $i$，使 $f_i|_{U_j}$ 可逆。
\item 集合层的映射
$$
(\mathcal{O} \times \mathcal{O})
\amalg
(\mathcal{O} \times \mathcal{O})
\longrightarrow
\mathcal{O} \times \mathcal{O}
$$
把第一分支中的 $(f, a)$ 映到 $(f, af)$，并把第二分支中的
$(f, b)$ 映到 $(f, b(1 - f))$；此映射是满射。
\end{enumerate}
\end{lemma}

\begin{proof}
(2) 显然蕴含 (1)。为证明 (1) 蕴含 (2)，对 $n$ 作归纳。
第一个非平凡情形是 $n = 2$（$n = 1$ 显然）。此时存在
$a_1, a_2 \in \mathcal{O}(U)$，使
$a_1f_1 + a_2f_2 = 1$。由假设，可找到覆盖
$\{U_j \to U\}$，使得对每个 $j$，
$a_1f_1|_{U_j}$ 或 $a_2f_2|_{U_j}$ 可逆。因此
$f_1|_{U_j}$ 或 $f_2|_{U_j}$ 可逆，如所需。当 $n > 2$ 时，
存在 $a_1, \ldots, a_n \in \mathcal{O}(U)$，使
$a_1f_1 + \ldots + a_nf_n = 1$。由 $n = 2$ 的情形，存在覆盖
$\{U_j \to U\}_{j \in J}$，使得对每个 $j$，
$f_n|_{U_j}$ 或
$a_1f_1 + \ldots + a_{n - 1}f_{n - 1}|_{U_j}$ 可逆。
设第一种情形恰发生于 $j \in J_n$，并置
$J' = J \setminus J_n$。由归纳假设，对每个 $j \in J'$
可找到覆盖 $\{U_{jk} \to U_j\}_{k \in K_j}$，使得对每个
$k \in K_j$ 都存在 $i \in \{1, \ldots, n - 1\}$，使
$f_i|_{U_{jk}}$ 可逆。由站点公理，态射族
$\{U_j \to U\}_{j \in J_n} \cup \{U_{jk} \to U\}_{j \in J', k \in K_j}$
是具有所需性质的覆盖。

\medskip\noindent
假设 (1)。为证 (3) 中的映射满射，设 $(f, c)$ 是
$\mathcal{O} \times \mathcal{O}$ 在 $U$ 上的截面。由假设，
存在覆盖 $\{U_j \to U\}$，使得对每个 $j$，$f$ 或 $1-f$
的限制可逆。第一种情形取
$a = c|_{U_j} (f|_{U_j})^{-1}$，第二种情形取
$b = c|_{U_j} (1 - f|_{U_j})^{-1}$。故在每个覆盖成员上，
$(f,c)$ 都属于该映射的像。反之，假设 (3) 成立。对任意
$U$ 及 $f \in \mathcal{O}(U)$，存在 $U$ 的覆盖
$\{U_j \to U\}$，使截面 $(f,1)|_{U_j}$ 属于 (3)
中的映射在 $U_j$ 上截面的像。这恰好意味着 $f$ 或 $1-f$
在 $U_j$ 上的限制可逆，故 (1) 成立。
\end{proof}

\begin{lemma}
\label{sites-modules-lemma-locally-ringed-stalk}
设 $(\mathcal{C}, \mathcal{O})$ 为环站点。考虑下列条件：
\begin{enumerate}
\item 对 $\mathcal{C}$ 的每个对象 $U$ 及
$f \in \mathcal{O}(U)$，存在覆盖 $\{U_j \to U\}$，使得对每个
$j$，$f|_{U_j}$ 或 $(1-f)|_{U_j}$ 可逆。
\item 对 $\mathcal{C}$ 的每个点 $p$，茎 $\mathcal{O}_p$
是零环或局部环。
\end{enumerate}
总有 (1) $\Rightarrow$ (2)。若 $\mathcal{C}$ 有足够多的点，
则 (1) 与 (2) 等价。
\end{lemma}

\begin{proof}
假设 (1)。设 $p$ 是由函子
$u : \mathcal{C} \to \textit{Sets}$ 给出的 $\mathcal{C}$ 的点，
并取 $f_p \in \mathcal{O}_p$。由于 $\mathcal{O}_p$
由《站点》，公式 (\ref{sites-equation-stalk}) 计算，可以用三元组
$(U,x,f)$ 表示 $f_p$，其中 $x \in U(U)$ 且
$f \in \mathcal{O}(U)$。由假设，存在覆盖
$\{U_i \to U\}$，使得对每个 $i$，$f$ 或 $1-f$
在 $U_i$ 上可逆。由于 $u$ 定义站点的点，对某个 $i$
存在 $x_i \in u(U_i)$ 映到 $x \in u(U)$。由《站点》，公式
(\ref{sites-equation-stalk}) 附近的讨论，
$(U,x,f)$ 与 $(U_i,x_i,f|_{U_i})$ 定义
$\mathcal{O}_p$ 中的同一元素。因此 $f_p$ 或 $1-f_p$ 可逆。
于是环 $\mathcal{O}_p$ 的每个元素 $a$ 都满足
$a$ 或 $1-a$ 可逆。这意味着 $\mathcal{O}_p$ 是零环或局部环；
见《代数》，引理 \ref{algebra-lemma-characterize-local-ring}。

\medskip\noindent
假设 (2)，并假设 $\mathcal{C}$ 有足够多的点。考虑集合层的映射
$$
\mathcal{O} \times \mathcal{O} \amalg \mathcal{O} \times \mathcal{O}
\longrightarrow
\mathcal{O} \times \mathcal{O}
$$
即引理 \ref{sites-modules-lemma-locally-ringed} 的 (3) 中的映射。
对任意局部环 $R$，相应映射
$(R \times R) \amalg (R \times R) \to R \times R$
是满射；例如见《代数》，引理
\ref{algebra-lemma-characterize-local-ring}。由于每个
$\mathcal{O}_p$ 是局部环或零环，该映射在茎上满射。
由 $\mathcal{C}$ 有足够多的点的假设，该映射满射，结论得证。
\end{proof}

\noindent
《模》，第 \ref{modules-section-pathology} 节指出，在环化空间
$(X,\mathcal{O}_X)$ 中可能存在一个开子空间，使结构层在其上为零。
为排除这种情形，可以要求截面 $1$ 与 $0$ 在空间 $X$
的每个茎中取不同值。在环拓扑斯与环站点的语境下，该条件是
\begin{equation}
\label{sites-modules-equation-one-is-never-zero}
\emptyset^\# \longrightarrow
\text{Equalizer}(0, 1 : * \longrightarrow \mathcal{O})
\end{equation}
是层的同构。这里 $*$ 是单元素层，而 $\emptyset^\#$ 是“空层”，
即分别为层范畴的终对象与始对象；见《站点》，例
\ref{sites-example-singleton-sheaf} 及第
\ref{sites-section-almost-cocontinuous} 节。换言之，该条件是：
只要 $U \in \Ob(\mathcal{C})$ 在层论意义下不为空，
就有 $1,0 \in \mathcal{O}(U)$ 不相等。下面陈述必需的引理。

\begin{lemma}
\label{sites-modules-lemma-ringed-stalk-not-zero}
设 $(\mathcal{C}, \mathcal{O})$ 为环站点。考虑下列命题：
\begin{enumerate}
\item (\ref{sites-modules-equation-one-is-never-zero}) 是同构；
\item 对 $\mathcal{C}$ 的每个点 $p$，茎 $\mathcal{O}_p$ 不是零环。
\end{enumerate}
总有 (1) $\Rightarrow$ (2)；若 $\mathcal{C}$ 有足够多的点，
则 (1) $\Leftrightarrow$ (2)。
\end{lemma}

\begin{proof}
从略。
\end{proof}

\noindent
引理 \ref{sites-modules-lemma-locally-ringed}、
\ref{sites-modules-lemma-locally-ringed-stalk} 及
\ref{sites-modules-lemma-ringed-stalk-not-zero} 引出下列定义。

\begin{definition}
\label{sites-modules-definition-locally-ringed}
若 (\ref{sites-modules-equation-one-is-never-zero}) 是同构，且引理
\ref{sites-modules-lemma-locally-ringed} 中的等价性质成立，则称环站点
$(\mathcal{C}, \mathcal{O})$ 为{\it 局部环站点}。
\end{definition}

\noindent
\cite[Expos\'e IV, Exercice 13.9]{SGA4} 中缺少
(\ref{sites-modules-equation-one-is-never-zero}) 为同构这一条件，因而其中的
局部环站点与局部环拓扑斯概念略有不同。由于我们的动机来自
局部环化空间的概念，故加入了这一条件（见上文解释）。

\begin{lemma}
\label{sites-modules-lemma-locally-ringed-intrinsic}
成为局部环站点是内蕴性质。更准确地：
\begin{enumerate}
\item 若 $f : \Sh(\mathcal{C}') \to \Sh(\mathcal{C})$
为拓扑斯态射，且 $(\mathcal{C}, \mathcal{O})$
为局部环站点，则 $(\mathcal{C}', f^{-1}\mathcal{O})$
为局部环站点；
\item 若
$(f, f^\sharp) : (\Sh(\mathcal{C}'), \mathcal{O}')
\to (\Sh(\mathcal{C}), \mathcal{O})$
为环拓扑斯的等价，则 $(\mathcal{C}, \mathcal{O})$
局部环，当且仅当
$(\mathcal{C}', \mathcal{O}')$
局部环。
\end{enumerate}
\end{lemma}

\begin{proof}
(2) 显然由 (1) 得到。为证明 (1)，注意 $f^{-1}$ 正合，故
$f^{-1}* = *$、$f^{-1}\emptyset^\# = \emptyset^\#$，
且 $f^{-1}$ 与积及等化子交换，并把同构与满射分别变为同构与满射。
因此 $f^{-1}$ 把同构 (\ref{sites-modules-equation-one-is-never-zero})
变为 $f^{-1}\mathcal{O}$ 的相应同构，并把引理
\ref{sites-modules-lemma-locally-ringed} 的 (3) 中的满射变为
$f^{-1}\mathcal{O}$ 的相应满射。
\end{proof}

\noindent
事实上，引理 \ref{sites-modules-lemma-locally-ringed-intrinsic} 的 (2)
是《概形》，引理 \ref{schemes-lemma-isomorphism-locally-ringed}
的对应版本。它保证下列定义有意义。

\begin{definition}
\label{sites-modules-definition-locally-ringed-topos}
若环拓扑斯 $(\Sh(\mathcal{C}), \mathcal{O})$ 所基于的环站点
$(\mathcal{C}, \mathcal{O})$ 局部环，则称该环拓扑斯
{\it 局部环}。
\end{definition}

\noindent
下面给出局部环性质的一个推论。

\begin{lemma}
\label{sites-modules-lemma-invertible-is-locally-free-rank-1}
设 $(\Sh(\mathcal{C}), \mathcal{O})$ 为环拓扑斯。
任何秩为 $1$ 的局部自由 $\mathcal{O}$-模都可逆。
若 $(\mathcal{C}, \mathcal{O})$ 局部环，则逆命题也成立
（一般情形下则不成立）。
\end{lemma}

\begin{proof}
假设 $\mathcal{L}$ 是秩为 $1$ 的局部自由模，并考虑求值映射
$$
\mathcal{L} \otimes_\mathcal{O}
\SheafHom_\mathcal{O}(\mathcal{L}, \mathcal{O})
\longrightarrow \mathcal{O}
$$
对 $\mathcal{C}$ 的任意对象 $U$，限制到使 $\mathcal{L}$
平凡化的某个覆盖的各成员上，可见此映射是同构（细节从略）。
因此由引理 \ref{sites-modules-lemma-invertible}，$\mathcal{L}$ 可逆。

\medskip\noindent
假设 $(\Sh(\mathcal{C}), \mathcal{O})$ 局部环，并设
$U$ 为 $\mathcal{C}$ 的对象。引理 \ref{sites-modules-lemma-invertible}
的证明表明，存在覆盖 $\{U_i \to U\}$，使
$\mathcal{L}|_{\mathcal{C}/U_i}$ 是某个有限自由
$\mathcal{O}_{U_i}$-模的直和因子。以 $U_i$ 替换 $U$ 后，设
$p : \mathcal{O}_U^{\oplus r} \to \mathcal{O}_U^{\oplus r}$
为投影算子，其像同构于 $\mathcal{L}|_{\mathcal{C}/U}$。
则 $p$ 对应于矩阵
$$
P = (p_{ij}) \in \text{Mat}(r \times r, \mathcal{O}(U))
$$
它是投影矩阵：$P^2 = P$。置 $A = \mathcal{O}(U)$，
于是 $P \in \text{Mat}(r \times r, A)$。由《代数》，引理
\ref{algebra-lemma-finite-projective}，$P$ 的像是
$A$ 上的有限局部自由模 $M$。故存在生成单位理想的
$f_1,\ldots,f_t \in A$，使 $M_{f_i}$ 有限自由。由引理
\ref{sites-modules-lemma-locally-ringed}，将 $U$ 替换为某个开覆盖的各成员后，
可以假设 $M$ 自由。这意味着 $\mathcal{L}|_U$ 自由
（细节从略）。当然，由于 $\mathcal{L}$ 可逆，只有当
$\mathcal{L}|_U$ 的秩为 $1$ 时才可能如此，证明完成。
\end{proof}

\noindent
下面说明局部环化空间之间的态射意味着什么。为此需要下列引理。

\begin{lemma}
\label{sites-modules-lemma-locally-ringed-morphism}
设
$(f, f^\sharp) : (\Sh(\mathcal{C}), \mathcal{O}_\mathcal{C})
\to (\Sh(\mathcal{D}), \mathcal{O}_\mathcal{D})$
为环拓扑斯态射。考虑下列条件：
\begin{enumerate}
\item 层的图
$$
\xymatrix{
f^{-1}(\mathcal{O}^*_\mathcal{D}) \ar[r]_-{f^\sharp} \ar[d] &
\mathcal{O}^*_\mathcal{C} \ar[d] \\
f^{-1}(\mathcal{O}_\mathcal{D}) \ar[r]^-{f^\sharp} &
\mathcal{O}_\mathcal{C}
}
$$
是笛卡尔图。
\item 对 $\mathcal{C}$ 的任意点 $p$，置 $q = f \circ p$，则集合的图
$$
\xymatrix{
\mathcal{O}^*_{\mathcal{D}, q} \ar[r] \ar[d] &
\mathcal{O}^*_{\mathcal{C}, p} \ar[d] \\
\mathcal{O}_{\mathcal{D}, q} \ar[r] &
\mathcal{O}_{\mathcal{C}, p}
}
$$
是笛卡尔图。
\end{enumerate}
总有 (1) $\Rightarrow$ (2)。若 $\mathcal{C}$ 有足够多的点，
则 (1) 与 (2) 等价。若
$(\Sh(\mathcal{C}), \mathcal{O}_\mathcal{C})$
与
$(\Sh(\mathcal{D}), \mathcal{O}_\mathcal{D})$
均为局部环拓扑斯，则 (2) 等价于
\begin{enumerate}
\item[(3)] 对 $\mathcal{C}$ 的任意点 $p$，置 $q = f \circ p$，
则环映射
$\mathcal{O}_{\mathcal{D}, q} \to \mathcal{O}_{\mathcal{C}, p}$
是局部环同态。
\end{enumerate}
事实上，若对一族保守点有性质 (2) 或 (3)，则 (1) 成立。
\end{lemma}

\begin{proof}
展开各定义即可证明。
\end{proof}

\begin{definition}
\label{sites-modules-definition-morphism-locally-ringed-topoi}
设 $(f, f^\sharp) : (\Sh(\mathcal{C}), \mathcal{O}_\mathcal{C})
\to (\Sh(\mathcal{D}), \mathcal{O}_\mathcal{D})$
为环拓扑斯态射。假设
$(\Sh(\mathcal{C}), \mathcal{O}_\mathcal{C})$
与
$(\Sh(\mathcal{D}), \mathcal{O}_\mathcal{D})$
均为局部环拓扑斯。若且唯若层的图
$$
\xymatrix{
f^{-1}(\mathcal{O}^*_\mathcal{D}) \ar[r]_-{f^\sharp} \ar[d] &
\mathcal{O}^*_\mathcal{C} \ar[d] \\
f^{-1}(\mathcal{O}_\mathcal{D}) \ar[r]^-{f^\sharp} &
\mathcal{O}_\mathcal{C}
}
$$
是笛卡尔图（见引理 \ref{sites-modules-lemma-locally-ringed-morphism}），
则称 $(f,f^\sharp)$ 为{\it 局部环拓扑斯态射}。
若 $(f,f^\sharp)$ 是环站点态射，并且其关联的环拓扑斯态射
是局部环拓扑斯态射，则称其为{\it 局部环站点态射}。
\end{definition}

\noindent
显然，局部环拓扑斯之间的环拓扑斯同构自动是局部环拓扑斯同构。

\begin{lemma}
\label{sites-modules-lemma-composition-morphisms-locally-ringed-topoi}
设
$(f, f^\sharp) :
(\Sh(\mathcal{C}_1), \mathcal{O}_1)
\to (\Sh(\mathcal{C}_2), \mathcal{O}_2)$ and
$(g, g^\sharp) :
(\Sh(\mathcal{C}_2), \mathcal{O}_2) \to
(\Sh(\mathcal{C}_3), \mathcal{O}_3)$
为局部环拓扑斯态射。则复合
$(g, g^\sharp) \circ (f, f^\sharp)$（见定义
\ref{sites-modules-definition-ringed-topos}）也是局部环拓扑斯态射。
\end{lemma}

\begin{proof}
从略。
\end{proof}

\begin{lemma}
\label{sites-modules-lemma-locally-ringed-intrinsic-morphism}
设 $f : \Sh(\mathcal{C}') \to \Sh(\mathcal{C})$
为拓扑斯态射。若 $\mathcal{O}$ 是 $\mathcal{C}$ 上的环层，则
$$
f^{-1}(\mathcal{O}^*) = (f^{-1}\mathcal{O})^*.
$$
特别地，若 $\mathcal{O}$ 使 $\mathcal{C}$ 成为局部环站点，
则置 $f^\sharp = \text{id}$ 后，环拓扑斯态射
$$
(f, f^\sharp) :
(\Sh(\mathcal{C}'), f^{-1}\mathcal{O})
\to
(\Sh(\mathcal{C}, \mathcal{O})
$$
是局部环拓扑斯态射。
\end{lemma}

\begin{proof}
注意，图
$$
\xymatrix{
\mathcal{O}^* \ar[rr] \ar[d]_{u \mapsto (u, u^{-1})} & &
{*} \ar[d]^{1} \\
\mathcal{O} \times \mathcal{O} \ar[rr]^-{(a, b) \mapsto ab} & &
\mathcal{O}
}
$$
是笛卡尔图。由于 $f^{-1}$ 正合，图
$$
\xymatrix{
f^{-1}(\mathcal{O}^*)
\ar[d]_{u \mapsto (u, u^{-1})} \ar[rr] & &
{*} \ar[d]^{1} \\
f^{-1}\mathcal{O} \times f^{-1}\mathcal{O} \ar[rr]^-{(a, b) \mapsto ab} & &
f^{-1}\mathcal{O}
}
$$
也是笛卡尔图，这蕴含第一个断言。对第二个断言，注意由引理
\ref{sites-modules-lemma-locally-ringed-intrinsic}，
$(\mathcal{C}', f^{-1}\mathcal{O})$ 是局部环站点，故该断言有意义。
此时第一部分蕴含该态射是局部环拓扑斯态射。
\end{proof}

\begin{lemma}
\label{sites-modules-lemma-localize-morphism-locally-ringed-topoi}
局部环站点与拓扑斯的局部化。
\begin{enumerate}
\item 设 $(\mathcal{C}, \mathcal{O})$ 为局部环站点，
$U$ 为 $\mathcal{C}$ 的对象。则局部化
$(\mathcal{C}/U, \mathcal{O}_U)$ 是局部环站点，且局部化态射
$$
(j_U, j_U^\sharp) :
(\Sh(\mathcal{C}/U), \mathcal{O}_U)
\to
(\Sh(\mathcal{C}), \mathcal{O})
$$
是局部环拓扑斯态射。
\item 设 $(\mathcal{C}, \mathcal{O})$ 为局部环站点，
$f : V \to U$ 为 $\mathcal{C}$ 中的态射。则引理
\ref{sites-modules-lemma-relocalize} 中的态射
$$
(j, j^\sharp) :
(\Sh(\mathcal{C}/V), \mathcal{O}_V)
\to
(\Sh(\mathcal{C}/U), \mathcal{O}_U)
$$
是局部环拓扑斯态射。
\item 设
$(f, f^\sharp) :
(\mathcal{C}, \mathcal{O})
\longrightarrow
(\mathcal{D}, \mathcal{O}')$
为局部环站点态射，其中 $f$ 由连续函子
$u : \mathcal{D} \to \mathcal{C}$ 给出。设 $V$ 为
$\mathcal{D}$ 的对象，并令 $U = u(V)$。则引理
\ref{sites-modules-lemma-localize-morphism-ringed-sites} 中的态射
$$
(f', (f')^\sharp) :
(\Sh(\mathcal{C}/U), \mathcal{O}_U)
\to
(\Sh(\mathcal{D}/V), \mathcal{O}'_V)
$$
是局部环站点态射。
\item 设
$(f, f^\sharp) :
(\mathcal{C}, \mathcal{O})
\longrightarrow
(\mathcal{D}, \mathcal{O}')$
为局部环站点态射，其中 $f$ 由连续函子
$u : \mathcal{D} \to \mathcal{C}$ 给出。设
$V \in \Ob(\mathcal{D})$、$U \in \Ob(\mathcal{C})$，且
$c : U \to u(V)$。则引理
\ref{sites-modules-lemma-relocalize-morphism-ringed-sites} 中的态射
$$
(f_c, (f_c)^\sharp) :
(\Sh(\mathcal{C}/U), \mathcal{O}_U)
\to
(\Sh(\mathcal{D}/V), \mathcal{O}'_V)
$$
是局部环拓扑斯态射。
\item 设 $(\Sh(\mathcal{C}), \mathcal{O})$ 为局部环拓扑斯，
$\mathcal{F}$ 为 $\mathcal{C}$ 上的层。则局部化
$(\Sh(\mathcal{C})/\mathcal{F}, \mathcal{O}_\mathcal{F})$
是局部环拓扑斯，且局部化态射
$$
(j_\mathcal{F}, j_\mathcal{F}^\sharp) :
(\Sh(\mathcal{C})/\mathcal{F}, \mathcal{O}_\mathcal{F})
\to
(\Sh(\mathcal{C}), \mathcal{O})
$$
是局部环拓扑斯态射。
\item 设 $(\Sh(\mathcal{C}), \mathcal{O})$ 为局部环拓扑斯，
$s : \mathcal{G} \to \mathcal{F}$ 为 $\mathcal{C}$ 上层的映射。
则引理 \ref{sites-modules-lemma-relocalize-ringed-topos} 中的态射
$$
(j, j^\sharp) :
(\Sh(\mathcal{C})/\mathcal{G}, \mathcal{O}_\mathcal{G})
\longrightarrow
(\Sh(\mathcal{C})/\mathcal{F}, \mathcal{O}_\mathcal{F})
$$
是局部环拓扑斯态射。
\item 设
$f :
(\Sh(\mathcal{C}), \mathcal{O})
\longrightarrow
(\Sh(\mathcal{D}), \mathcal{O}')$
为局部环拓扑斯态射，$\mathcal{G}$ 为 $\mathcal{D}$ 上的层，
并置 $\mathcal{F} = f^{-1}\mathcal{G}$。则引理
\ref{sites-modules-lemma-localize-morphism-ringed-topoi} 中的态射
$$
(f', (f')^\sharp) :
(\Sh(\mathcal{C})/\mathcal{F}, \mathcal{O}_\mathcal{F})
\longrightarrow
(\Sh(\mathcal{D})/\mathcal{G}, \mathcal{O}'_\mathcal{G})
$$
是局部环拓扑斯态射。
\item 设
$f :
(\Sh(\mathcal{C}), \mathcal{O})
\longrightarrow
(\Sh(\mathcal{D}), \mathcal{O}')$
为局部环拓扑斯态射，$\mathcal{G}$ 为 $\mathcal{D}$ 上的层，
$\mathcal{F}$ 为 $\mathcal{C}$ 上的层，并设
$s : \mathcal{F} \to f^{-1}\mathcal{G}$ 为层态射。则引理
\ref{sites-modules-lemma-relocalize-morphism-ringed-topoi} 中的态射
$$
(f_s, (f_s)^\sharp) :
(\Sh(\mathcal{C})/\mathcal{F}, \mathcal{O}_\mathcal{F})
\longrightarrow
(\Sh(\mathcal{D})/\mathcal{G}, \mathcal{O}'_\mathcal{G})
$$
是局部环拓扑斯态射。
\end{enumerate}
\end{lemma}

\begin{proof}
(1) 显然，因为 $\mathcal{O}_U$ 就是 $\mathcal{O}$ 的限制，
故可应用引理 \ref{sites-modules-lemma-locally-ringed-intrinsic} 与
\ref{sites-modules-lemma-locally-ringed-intrinsic-morphism}。(2) 也显然，因为态射
$(j,j^\sharp)$ 实际上是局部环站点的局部化，故可应用 (1)。
(3) 同样显然，因为 $(f')^\sharp$ 就是 $f^\sharp$
在拓扑斯 $\Sh(\mathcal{C})/\mathcal{F}$ 上的限制；见引理
\ref{sites-modules-lemma-localize-morphism-ringed-topoi} 的证明
（因此定义 \ref{sites-modules-definition-morphism-locally-ringed-topoi}
中态射 $f'$ 的图，就是态射 $f$ 的相应图的限制，而限制函子正合）。
结合 (2) 与 (3)，形式地得到 (4)。按引理
\ref{sites-modules-lemma-morphism-ringed-topoi-comes-from-morphism-ringed-sites}
扩大站点，并利用引理 \ref{sites-modules-lemma-localize-compare}、
\ref{sites-modules-lemma-relocalize-compare}、
\ref{sites-modules-lemma-localize-morphism-compare} 及
\ref{sites-modules-lemma-relocalize-morphism-compare} 的比较，把一切翻译成
站点及站点态射的语言，即由 (1)、(2)、(3)、(4) 分别得到
(5)、(6)、(7)、(8)。（也可以直接使用证明 (1)、(2)、(3)、(4)
时相同的论证来证明 (5)、(6)、(7)、(8)。）
\end{proof}






\section{模的下叹号函子}
\label{sites-modules-section-lower-shriek-modules}

\noindent
本节把第 \ref{sites-modules-section-exactness-lower-shriek} 节讨论的
$g_!$ 构造推广到模层的情形。

\begin{lemma}
\label{sites-modules-lemma-lower-shriek-modules}
设 $u : \mathcal{C} \to \mathcal{D}$ 为站点之间既连续又余连续的函子，
记 $g : \Sh(\mathcal{C}) \to \Sh(\mathcal{D})$
为关联的拓扑斯态射。设 $\mathcal{O}_\mathcal{D}$
为 $\mathcal{D}$ 上的环层，并置
$\mathcal{O}_\mathcal{C} = g^{-1}\mathcal{O}_\mathcal{D}$。
于是 $g$ 成为满足 $g^* = g^{-1}$ 的环拓扑斯态射。
此时存在函子
$$
g_! :
\textit{Mod}(\mathcal{O}_\mathcal{C})
\longrightarrow
\textit{Mod}(\mathcal{O}_\mathcal{D})
$$
它是 $g^*$ 的左伴随。
\end{lemma}

\begin{proof}
设 $U$ 为 $\mathcal{C}$ 的对象。对任意
$\mathcal{O}_\mathcal{D}$-模 $\mathcal{G}$，有
\begin{align*}
\Hom_{\mathcal{O}_\mathcal{C}}(j_{U!}\mathcal{O}_U, g^{-1}\mathcal{G})
& =
g^{-1}\mathcal{G}(U) \\
& =
\mathcal{G}(u(U)) \\
& =
\Hom_{\mathcal{O}_\mathcal{D}}(j_{u(U)!}\mathcal{O}_{u(U)}, \mathcal{G})
\end{align*}
因为 $g^{-1}$ 由限制描述；见《站点》，引理
\ref{sites-lemma-when-shriek}。对形如
$j_{U!}\mathcal{O}_U$ 的模的直和，当然有类似公式。
由《同调》，引理 \ref{homology-lemma-partially-defined-adjoint}
及引理 \ref{sites-modules-lemma-module-quotient-flat}，可知 $g_!$ 存在。
\end{proof}

\begin{remark}
\label{sites-modules-remark-when-shriek-equal}
警告！设 $u : \mathcal{C} \to \mathcal{D}$、$g$、
$\mathcal{O}_\mathcal{D}$ 与 $\mathcal{O}_\mathcal{C}$
如引理 \ref{sites-modules-lemma-lower-shriek-modules}。一般而言，下图{\bf 不}交换：
$$
\xymatrix{
\textit{Mod}(\mathcal{O}_\mathcal{C}) \ar[r]_{g_!} \ar[d]_{forget} &
\textit{Mod}(\mathcal{O}_\mathcal{D}) \ar[d]^{forget} \\
\textit{Ab}(\mathcal{C}) \ar[r]^{g^{Ab}_!} &
\textit{Ab}(\mathcal{D})
}
$$
（这里 $g^{Ab}_!$ 是引理 \ref{sites-modules-lemma-g-shriek-adjoint}
中的函子。）有函子变换
$$
g_!^{Ab} \circ forget \longrightarrow forget \circ g_!
$$
由引理 \ref{sites-modules-lemma-lower-shriek-modules} 的证明可知，
该变换为同构，当且仅当对 $\mathcal{C}$ 的每个对象 $U$，
$g^{Ab}_!j_{U!}\mathcal{O}_U \to g_!j_{U!}\mathcal{O}_U$
都是同构。由于
$g_!j_{U!}\mathcal{O}_U = j_{u(U)!}\mathcal{O}_{u(U)}$，
这又等价于
$$
g^{Ab}_!j_{U!}\mathcal{O}_U \longrightarrow j_{u(U)!}\mathcal{O}_{u(U)}
$$
对 $\mathcal{C}$ 的每个对象 $U$ 都是同构。注意，对这样的 $U$
有交换图
$$
\xymatrix{
\mathcal{C}/U \ar[r]_-{u'} \ar[d]_{j_U} & \mathcal{D}/u(U) \ar[d]^{j_{u(U)}} \\
\mathcal{C} \ar[r]^u & \mathcal{D}
}
$$
其中各箭头是站点的余连续函子；见《站点》，引理
\ref{sites-lemma-localize-cocontinuous}。因此
$g^{Ab}_!j_{U!} = j_{u(U)!}(g')^{Ab}_!$，其中
$g' : \Sh(\mathcal{C}/U) \to \Sh(\mathcal{D}/u(U))$
是由余连续函子 $u'$ 诱导的拓扑斯态射。于是，若对
$\mathcal{C}$ 的每个对象 $U$，典范映射
\begin{equation}
\label{sites-modules-equation-compare-on-localizations}
(g')^{Ab}_!\mathcal{O}_U \longrightarrow \mathcal{O}_{u(U)}
\end{equation}
都是同构，则 $g_! = g^{Ab}_!$。
\end{remark}

\noindent
下面两个结论的性质略有不同。

\begin{lemma}
\label{sites-modules-lemma-special-square-cocontinuous}
假设给定环拓扑斯的交换图
$$
\xymatrix{
(\Sh(\mathcal{C}'), \mathcal{O}_{\mathcal{C}'})
\ar[r]_{(g', (g')^\sharp)} \ar[d]_{(f', (f')^\sharp)} &
(\Sh(\mathcal{C}), \mathcal{O}_\mathcal{C}) \ar[d]^{(f, f^\sharp)} \\
(\Sh(\mathcal{D}'), \mathcal{O}_{\mathcal{D}'}) \ar[r]^{(g, g^\sharp)} &
(\Sh(\mathcal{D}), \mathcal{O}_\mathcal{D})
}
$$
并假设：
\begin{enumerate}
\item $f$、$f'$、$g$、$g'$ 分别对应于余连续函子
$u$、$u'$、$v$、$v'$，如《站点》，引理
\ref{sites-lemma-cocontinuous-morphism-topoi}；
\item $v \circ u' = u \circ v'$,
\item $v$ 与 $v'$ 既连续又余连续；
\item 对 $\mathcal{D}'$ 的任意对象 $V'$，由 $v$ 给出的函子
${}^{u'}_{V'}\mathcal{I} \to {}^{\ \ \ u}_{v(V')}\mathcal{I}$
是余终的；
\item $g^{-1}\mathcal{O}_{\mathcal{D}} = \mathcal{O}_{\mathcal{D}'}$
且 $(g')^{-1}\mathcal{O}_{\mathcal{C}} = \mathcal{O}_{\mathcal{C}'}$。
\end{enumerate}
则在模上有 $f'_* \circ (g')^* = g^* \circ f_*$ 及
$g'_! \circ (f')^{-1} = f^{-1} \circ g_!$。
\end{lemma}

\begin{proof}
由条件 (5)，
$(g')^*\mathcal{F} = (g')^{-1}\mathcal{F}$ 且
$g^*\mathcal{G} = g^{-1}\mathcal{G}$。因此第一个等式立即由
《站点》，引理 \ref{sites-lemma-special-square-cocontinuous}
中的相应等式得到。由引理 \ref{sites-modules-lemma-lower-shriek-modules}，
$g^*$ 与 $(g')^*$ 的左伴随函子 $g_!$ 与 $g'_!$ 存在，
故第二个等式由伴随函子的唯一性得到。
\end{proof}

\begin{lemma}
\label{sites-modules-lemma-special-square-continuous}
考虑环拓扑斯的交换图
$$
\xymatrix{
(\Sh(\mathcal{C}'), \mathcal{O}_{\mathcal{C}'})
\ar[r]_{(g', (g')^\sharp)} \ar[d]_{(f', (f')^\sharp)} &
(\Sh(\mathcal{C}), \mathcal{O}_\mathcal{C}) \ar[d]^{(f, f^\sharp)} \\
(\Sh(\mathcal{D}'), \mathcal{O}_{\mathcal{D}'}) \ar[r]^{(g, g^\sharp)} &
(\Sh(\mathcal{D}), \mathcal{O}_\mathcal{D})
}
$$
并假设有函子
$$
\xymatrix{
\mathcal{C}' \ar[r]_{v'} &
\mathcal{C} \\
\mathcal{D}' \ar[r]^v \ar[u]^{u'} &
\mathcal{D} \ar[u]_u
}
$$
使得（采用《站点》，第 \ref{sites-section-morphism-sites} 节及
\ref{sites-section-cocontinuous-morphism-topoi} 节中的记号）：
\begin{enumerate}
\item $u$ 与 $u'$ 连续，并诱导态射 $f$ 与 $f'$；
\item $v$ 与 $v'$ 余连续，并诱导态射 $g$ 与 $g'$；
\item $u \circ v = v' \circ u'$,
\item $v$ 与 $v'$ 既连续又余连续；
\item $g^{-1}\mathcal{O}_{\mathcal{D}} = \mathcal{O}_{\mathcal{D}'}$
且 $(g')^{-1}\mathcal{O}_{\mathcal{C}} = \mathcal{O}_{\mathcal{C}'}$。
\end{enumerate}
则在模上有 $f'_* \circ (g')^* = g^* \circ f_*$ 及
$g'_! \circ (f')^{-1} = f^{-1} \circ g_!$。
\end{lemma}

\begin{proof}
由条件 (5)，
$(g')^*\mathcal{F} = (g')^{-1}\mathcal{F}$ 且
$g^*\mathcal{G} = g^{-1}\mathcal{G}$。因此第一个等式立即由
《站点》，引理 \ref{sites-lemma-special-square-continuous}
中的相应等式得到。由引理 \ref{sites-modules-lemma-lower-shriek-modules}，
$g^*$ 与 $(g')^*$ 的左伴随函子 $g_!$ 与 $g'_!$ 存在，
故第二个等式由伴随函子的唯一性得到。
\end{proof}







\section{常值层}
\label{sites-modules-section-constant}

\noindent
设 $E$ 为集合，$\mathcal{C}$ 为站点。取值为 $E$ 的常值层记为
$\underline{E}$；见《站点》，例 \ref{sites-example-constant-sheaf}。
若 $E$ 是阿贝尔群、环、模等，则 $\underline{E}$
分别是阿贝尔群层、环层、模层等。

\begin{lemma}
\label{sites-modules-lemma-constant-exact}
设 $\mathcal{C}$ 为站点。若
$0 \to A \to B \to C \to 0$ 是阿贝尔群的短正合列，则
$0 \to \underline{A} \to \underline{B} \to \underline{C} \to 0$
是阿贝尔层的正合列；事实上，它作为阿贝尔预层的序列也正合。
\end{lemma}

\begin{proof}
由于层化正合，显然
$0 \to \underline{A} \to \underline{B} \to \underline{C} \to 0$
是阿贝尔层的正合列。因此
$0 \to \underline{A} \to \underline{B} \to \underline{C}$
是阿贝尔预层的正合列。为证
$\underline{B} \to \underline{C}$ 满射，选取一个集合论截面
$s : C \to B$。这诱导集合层的截面
$\underline{s} : \underline{C} \to \underline{B}$，
它是满射 $\underline{B} \to \underline{C}$ 的左逆。
\end{proof}

\begin{lemma}
\label{sites-modules-lemma-tensor-with-finitely-presented}
设 $\mathcal{C}$ 为站点，$\Lambda$ 为环，而 $M$ 与 $Q$
为 $\Lambda$-模。若 $Q$ 是有限表示的 $\Lambda$-模，则对所有
$U \in \Ob(\mathcal{C})$，有
$\underline{M \otimes_\Lambda Q}(U) = \underline{M}(U) \otimes_\Lambda Q$
。
\end{lemma}

\begin{proof}
选取表示
$\Lambda^{\oplus m} \to \Lambda^{\oplus n} \to Q \to 0$。
它给出正合列
$M^{\oplus m} \to M^{\oplus n} \to M \otimes Q \to 0$.
由引理 \ref{sites-modules-lemma-constant-exact}，得到正合列
$$
\underline{M}(U)^{\oplus m} \to
\underline{M}(U)^{\oplus n} \to
\underline{M \otimes Q}(U) \to 0
$$
这证明了引理。（注意，在 $U$ 上取截面总与有限直和交换，
但不一定与任意直和交换。）
\end{proof}

\begin{lemma}
\label{sites-modules-lemma-flat-sections}
设 $\mathcal{C}$ 为站点，$\Lambda$ 为凝聚环，
$M$ 为平坦 $\Lambda$-模。对 $U \in \Ob(\mathcal{C})$，
模 $\underline{M}(U)$ 是平坦 $\Lambda$-模。
\end{lemma}

\begin{proof}
设 $I \subset \Lambda$ 为有限生成理想。由《代数》，引理
\ref{algebra-lemma-flat}，只需证明
$\underline{M}(U) \otimes_\Lambda I \to \underline{M}(U)$
单射。由于 $\Lambda$ 凝聚，$I$ 作为 $\Lambda$-模有限表示。
由引理 \ref{sites-modules-lemma-tensor-with-finitely-presented}，
$\underline{M}(U) \otimes I = \underline{M \otimes I}$。
由于 $M$ 平坦，映射 $M \otimes I \to M$ 单射，因而
$\underline{M \otimes I} \to \underline{M}$ 单射。
\end{proof}

\begin{lemma}
\label{sites-modules-lemma-completion-flat}
设 $\mathcal{C}$ 为站点，$\Lambda$ 为 Noether 环，
$I \subset \Lambda$ 为理想。层
$\underline{\Lambda}^\wedge = \lim \underline{\Lambda/I^n}$
是平坦 $\underline{\Lambda}$-代数。此外有典范等同
$$
\underline{\Lambda}/I\underline{\Lambda} =
\underline{\Lambda}/\underline{I} =
\underline{\Lambda}^\wedge/I\underline{\Lambda}^\wedge =
\underline{\Lambda}^\wedge/\underline{I} \cdot \underline{\Lambda}^\wedge =
\underline{\Lambda}^\wedge/\underline{I}^\wedge =
\underline{\Lambda/I}
$$
其中 $\underline{I}^\wedge = \lim \underline{I/I^n}$。
\end{lemma}

\begin{proof}
为证明 $\underline{\Lambda}^\wedge$ 平坦，只需证明对每个
$U \in \Ob(\mathcal{C})$，$\underline{\Lambda}^\wedge(U)$
作为 $\Lambda$-模平坦；见引理 \ref{sites-modules-lemma-flatness-presheaves}
及 \ref{sites-modules-lemma-flatness-sheafification}。由引理
\ref{sites-modules-lemma-flat-sections}，
$$
\underline{\Lambda}^\wedge(U) = \lim \underline{\Lambda/I^n}(U)
$$
是平坦 $\Lambda/I^n$-模系统的极限。由引理
\ref{sites-modules-lemma-constant-exact}，转移映射均满射。再由
《更多代数》，引理 \ref{more-algebra-lemma-limit-flat} 即得结论。

\medskip\noindent
为证明这些等式，注意由引理
\ref{sites-modules-lemma-tensor-with-finitely-presented}，
$\underline{\Lambda}(U)/I\underline{\Lambda}(U) = \underline{\Lambda/I}(U)$
。因此
$\underline{\Lambda}/I\underline{\Lambda} =
\underline{\Lambda}/\underline{I} = \underline{\Lambda/I}$。
短正合列系统
$$
0 \to \underline{I/I^n}(U) \to \underline{\Lambda/I^n}(U) \to
\underline{\Lambda/I}(U) \to 0
$$
具有满的转移映射，故给出短正合列
$$
0 \to \lim \underline{I/I^n}(U) \to \lim \underline{\Lambda/I^n}(U) \to
\lim \underline{\Lambda/I}(U) \to 0
$$
见《同调》，引理 \ref{homology-lemma-Mittag-Leffler}。因此
$\underline{\Lambda}^\wedge/\underline{I}^\wedge =
\underline{\Lambda/I}$。由于
$$
I \underline{\Lambda}^\wedge \subset
\underline{I} \cdot \underline{\Lambda}^\wedge \subset
\underline{I}^\wedge
$$
只需证明对所有 $U$，
$I \underline{\Lambda}^\wedge(U) = \underline{I}^\wedge(U)$
。选取生成元 $I = (f_1, \ldots, f_r)$。对每个 $n$，得到短正合列
$$
0 \to K_n/(I^n)^{\oplus r} \to (\Lambda/I^n)^{\oplus r}
\xrightarrow{(f_1, \ldots, f_r)}
I/I^{n + 1} \to 0
$$
其中
$K_n =
\{(x_1, \ldots, x_r) \in \Lambda^{\oplus r} \mid \sum x_i f_i \in I^{n + 1}\}$.
由此得到短正合列
$$
0 \to \underline{K_n/(I^n)^{\oplus r}}(U) \to
\underline{(\Lambda/I^n)^{\oplus r}}(U) \to
\underline{I/I^{n + 1}}(U) \to 0
$$
计算表明 $K_n = K + (I^n)^{\oplus r}$，故转移映射
$K_{n + 1}/(I^{n + 1})^{\oplus r} \to K_n/(I^n)^{\oplus r}$
满射。因此左侧模的系统具有满的转移映射，从而更满足 ML 条件。
由此可知
$(f_1, \ldots, f_r) :
(\underline{\Lambda}^\wedge)^{\oplus r}(U) \to \underline{I}^\wedge(U)$
由《同调》，引理 \ref{homology-lemma-Mittag-Leffler} 是满射，
这正是所需结论。
\end{proof}

\begin{lemma}
\label{sites-modules-lemma-locally-constant-finite-type}
设 $\mathcal{C}$ 为站点，$\Lambda$ 为环，$M$ 为 $\Lambda$-模。
假设 $\Sh(\mathcal{C})$ 不是空拓扑斯。则
\begin{enumerate}
\item $\underline{M}$ 是有限型 $\underline{\Lambda}$-模层，
当且仅当 $M$ 是有限 $\Lambda$-模；
\item $\underline{M}$ 是有限表示的 $\underline{\Lambda}$-模层，
当且仅当 $M$ 是有限表示的 $\Lambda$-模。
\end{enumerate}
\end{lemma}

\begin{proof}
证明 (1)。若 $M$ 由 $x_1,\ldots,x_r$ 生成，则这些元素定义
$\underline{M}$ 的整体截面并生成该层，故 $\underline{M}$ 有限型。
反之，假设 $\underline{M}$ 有限型。取一个在层论意义下不为空的
对象 $U \in \mathcal{C}$（《站点》，定义
\ref{sites-definition-empty}）。因为假设
$\Sh(\mathcal{C})$ 不是空拓扑斯，这样的对象存在。于是存在覆盖
$\{U_i \to U\}$ 及有限多个截面
$s_{ij} \in \underline{M}(U_i)$，它们生成
$\underline{M}|_{U_i}$。加细覆盖后，可以假设 $s_{ij}$
来自 $M$ 的元素 $x_{ij}$。于是 $x_{ij}$ 定义
$\underline{M}$ 的整体截面，它们在 $U$ 上的限制生成
$\underline{M}$。

\medskip\noindent
假设存在 $M$ 的元素 $x_1,\ldots,x_r$，它们定义
$\underline{M}$ 的整体截面，并作为
$\underline{\Lambda}$-模层生成 $\underline{M}$。我们证明
$x_1,\ldots,x_r$ 作为 $\Lambda$-模生成 $M$。取 $x \in M$。
可找到覆盖 $\{U_i \to U\}_{i \in I}$ 及
$f_{i,j} \in \underline{\Lambda}(U_i)$，使
$x|_{U_i} = \sum f_{i,j}x_j|_{U_i}$。加细覆盖后，可以假设
$f_{i,j} \in \Lambda$。由于 $U$ 在层论意义下不为空，
至少存在一个 $i \in I$，使 $U_i$ 在层论意义下不为空。
此时映射 $M \to \underline{M}(U_i)$ 单射（细节从略）。
因此 $M$ 中有 $x = \sum f_{i,j}x_j$，如所需。

\medskip\noindent
证明 (2)。假设 $\underline{M}$ 是有限表示的
$\underline{\Lambda}$-模。由 (1)，$M$ 有限型。选取
$M$ 作为 $\Lambda$-模的一组生成元 $x_1,\ldots,x_r$。
这确定短正合列
$0 \to K \to \Lambda^{\oplus r} \to M \to 0$，
由引理 \ref{sites-modules-lemma-constant-exact}，它变为短正合列
$$
0 \to \underline{K} \to \underline{\Lambda}^{\oplus r} \to \underline{M} \to 0
$$
由引理 \ref{sites-modules-lemma-kernel-surjection-finite-onto-finite-presentation}，
$\underline{K}$ 有限型。故由 (1)，$K$ 是有限
$\Lambda$-模。因此 $M$ 是有限表示的 $\Lambda$-模。
\end{proof}




\section{局部常值层}
\label{sites-modules-section-locally-constant}

\noindent
一般定义如下。

\begin{definition}
\label{sites-modules-definition-locally-constant}
设 $\mathcal{C}$ 为站点，$\mathcal{F}$ 为集合、群、阿贝尔群、环、
固定环 $\Lambda$ 上的模等的层。
\begin{enumerate}
\item 若 $\mathcal{F}$ 作为集合、群、阿贝尔群、环、固定环
$\Lambda$ 上的模等的层，同构于第 \ref{sites-modules-section-constant} 节所述的
某个常值层 $\underline{E}$，则称 $\mathcal{F}$ 是相应类型的
{\it 常值层}。
\item 若对 $\mathcal{C}$ 的每个对象 $U$，都存在覆盖
$\{U_i \to U\}$，使 $\mathcal{F}|_{U_i}$ 为常值层，
则称 $\mathcal{F}$ {\it 局部常值}。
\item 若 $\mathcal{F}$ 是集合层或群层，且其常值取值是有限集合
或有限群，则称 $\mathcal{F}$ {\it 有限局部常值}。
\end{enumerate}
\end{definition}

\begin{lemma}
\label{sites-modules-lemma-pullback-locally-constant}
设 $f : \Sh(\mathcal{C}) \to \Sh(\mathcal{D})$ 为拓扑斯态射。
若 $\mathcal{G}$ 是 $\mathcal{D}$ 上集合、群、阿贝尔群、环、
固定环 $\Lambda$ 上的模等的局部常值层，则
$f^{-1}\mathcal{G}$ 是 $\mathcal{C}$ 上相应类型的局部常值层。
\end{lemma}

\begin{proof}
从略。
\end{proof}

\begin{lemma}
\label{sites-modules-lemma-morphism-locally-constant}
设 $\mathcal{C}$ 为具有终对象 $X$ 的站点。
\begin{enumerate}
\item 设 $\varphi : \mathcal{F} \to \mathcal{G}$ 为
$\mathcal{C}$ 上局部常值集合层的映射。若 $\mathcal{F}$
有限局部常值，则存在覆盖 $\{U_i \to X\}$，使
$\varphi|_{U_i}$ 是由集合映射诱导的常值层映射。
\item 设 $\varphi : \mathcal{F} \to \mathcal{G}$ 为
$\mathcal{C}$ 上局部常值阿贝尔群层的映射。若 $\mathcal{F}$
有限局部常值，则存在覆盖 $\{U_i \to X\}$，使
$\varphi|_{U_i}$ 是由阿贝尔群同态诱导的常值阿贝尔层映射。
\item 设 $\Lambda$ 为环，
$\varphi : \mathcal{F} \to \mathcal{G}$ 为
$\mathcal{C}$ 上局部常值 $\Lambda$-模层的映射。若
$\mathcal{F}$ 有限型，则存在覆盖 $\{U_i \to X\}$，使
$\varphi|_{U_i}$ 是由 $\Lambda$-模同态诱导的常值
$\Lambda$-模层映射。
\end{enumerate}
\end{lemma}

\begin{proof}
证明从略。
\end{proof}

\begin{lemma}
\label{sites-modules-lemma-locally-constant}
设 $\mathcal{C}$ 为站点，$\Lambda$ 为环，
$M,N$ 为 $\Lambda$-模，而 $\mathcal{F},\mathcal{G}$
为局部常值 $\Lambda$-模层。
\begin{enumerate}
\item 若 $M$ 有限表示，则
$$
\underline{\Hom_\Lambda(M, N)} =
\SheafHom_{\underline{\Lambda}}(\underline{M}, \underline{N})
$$
\item 若 $M$ 与 $N$ 都有限表示，则
$$
\underline{\text{Isom}_\Lambda(M, N)} =
\mathit{Isom}_{\underline{\Lambda}}(\underline{M}, \underline{N})
$$
\item 若 $\mathcal{F}$ 有限表示，则
$\SheafHom_{\underline{\Lambda}}(\mathcal{F}, \mathcal{G})$
是局部常值 $\Lambda$-模层。
\item 若 $\mathcal{F}$ 与 $\mathcal{G}$ 都有限表示，则
$\mathit{Isom}_{\underline{\Lambda}}(\mathcal{F}, \mathcal{G})$
是局部常值集合层。
\end{enumerate}
\end{lemma}

\begin{proof}
证明 (1)。置 $E = \Hom_\Lambda(M, N)$。须证典范映射
$$
\underline{E}
\longrightarrow
\SheafHom_{\underline{\Lambda}}(\underline{M}, \underline{N})
$$
是同构。模 $M$ 有表示
$\Lambda^{\oplus s} \to \Lambda^{\oplus t} \to M \to 0$.
于是 $E$ 位于正合列
$$
0 \to E \to \Hom_\Lambda(\Lambda^{\oplus t}, N) \to
\Hom_\Lambda(\Lambda^{\oplus s}, N)
$$
中；类似地有
$$
0 \to
\SheafHom_{\underline{\Lambda}}(\underline{M}, \underline{N}) \to
\SheafHom_{\underline{\Lambda}}(\underline{\Lambda^{\oplus t}}, \underline{N})
\to
\SheafHom_{\underline{\Lambda}}(\underline{\Lambda^{\oplus s}}, \underline{N})
$$
这把问题归约到 $M$ 为有限自由模的情形，而此时结论显然。

\medskip\noindent
证明 (3)。该问题在 $\mathcal{C}$ 上是局部的，故可假设
$\mathcal{F} = \underline{M}$ 且 $\mathcal{G} = \underline{N}$，
其中 $M,N$ 为某些 $\Lambda$-模。由引理
\ref{sites-modules-lemma-locally-constant-finite-type}，模 $M$ 有限表示。
故结论由 (1) 得到。

\medskip\noindent
(2) 与 (4) 由 (1)、(3) 及下述事实得到：
$\mathit{Isom}$ 可视为
$\SheafHom_{\underline{\Lambda}}(\mathcal{F}, \mathcal{G})$
的截面所成的子层，其中这些截面在
$\SheafHom_{\underline{\Lambda}}(\mathcal{G}, \mathcal{F})$
中有逆。
\end{proof}

\begin{lemma}
\label{sites-modules-lemma-kernel-finite-locally-constant}
设 $\mathcal{C}$ 为站点。
\begin{enumerate}
\item 有限局部常值集合层的范畴在 $\Sh(\mathcal{C})$
中对有限极限与有限余极限封闭。
\item 有限局部常值阿贝尔层的范畴是
$\textit{Ab}(\mathcal{C})$ 的弱 Serre 子范畴。
\item 设 $\Lambda$ 为 Noether 环。$\mathcal{C}$ 上有限型局部常值
$\Lambda$-模层的范畴是
$\textit{Mod}(\mathcal{C}, \Lambda)$ 的弱 Serre 子范畴。
\end{enumerate}
\end{lemma}

\begin{proof}
证明 (1)。可以在 $\mathcal{C}$ 上局部地工作。因此由引理
\ref{sites-modules-lemma-morphism-locally-constant}，可以假设给定有限集合的有限图，
而我们的层图是其关联的常值层图。此时只须在集合范畴中取极限或余极限，
再取关联的常值层。略去部分细节。

\medskip\noindent
为证明 (2) 与 (3)，使用《同调》，引理
\ref{homology-lemma-characterize-weak-serre-subcategory} 的判据。
核与余核的存在性用与上面相同的方法论证。当然，(3)
中使用 Noether 环，是为了保证有限 $\Lambda$-模之间映射的核仍是
有限 $\Lambda$-模。为证明该范畴对扩张封闭
（在 $\Lambda$-模层的情形），假设给定 $\mathcal{C}$ 上
$\Lambda$-模层的扩张
$$
0 \to \mathcal{F} \to \mathcal{E} \to \mathcal{G} \to 0
$$
其中 $\mathcal{F},\mathcal{G}$ 有限型且局部常值。在
$\mathcal{C}$ 上局部化后，可以假设 $\mathcal{F}$ 与
$\mathcal{G}$ 都是常值层，即有
$$
0 \to \underline{M} \to \mathcal{E} \to \underline{N} \to 0
$$
其中 $M,N$ 为某些 $\Lambda$-模。选取 $N$ 的生成元
$y_1,\ldots,y_m$，得到 $\Lambda$-模的短正合列
$0 \to K \to \Lambda^{\oplus m} \to N \to 0$。
进一步局部化后，可以假设 $y_j$ 提升为 $\mathcal{E}$ 的截面
$s_j$。于是 $\mathcal{E}$ 是下图所示的推出：
$$
\xymatrix{
0 \ar[r] &
\underline{K} \ar[d] \ar[r] &
\underline{\Lambda^{\oplus m}} \ar[d] \ar[r] &
\underline{N} \ar[d] \ar[r] & 0 \\
0 \ar[r] &
\underline{M} \ar[r] &
\mathcal{E} \ar[r] &
\underline{N} \ar[r] & 0
}
$$
再次应用引理 \ref{sites-modules-lemma-morphism-locally-constant}
（并使用 $\Lambda$ 为 Noether 环从而 $K$ 是有限
$\Lambda$-模这一事实），可见映射
$\underline{K} \to \underline{M}$ 局部常值，故得结论。
\end{proof}

\begin{lemma}
\label{sites-modules-lemma-tensor-product-locally-constant}
设 $\mathcal{C}$ 为站点，$\Lambda$ 为环。$\mathcal{C}$
上两个局部常值 $\Lambda$-模层的张量积仍是局部常值
$\Lambda$-模层。
\end{lemma}

\begin{proof}
从略。
\end{proof}











\section{环层的局部化}
\label{sites-modules-section-localizing-sheaves-rings}

\noindent
设 $(\mathcal{C}, \mathcal{O})$ 为环站点，
$\mathcal{S} \subset \mathcal{O}$ 为集合子预层，满足：
对所有 $U \in \Ob(\mathcal{C})$，集合
$\mathcal{S}(U) \subset \mathcal{O}(U)$ 是乘法子集；见《代数》，
定义 \ref{algebra-definition-multiplicative-subset}。此时可以考虑环预层
$$
\mathcal{S}^{-1}\mathcal{O} :
U \longmapsto \mathcal{S}(U)^{-1}\mathcal{O}(U).
$$
对 $\mathcal{C}$ 中的 $V \to U$，限制映射把截面
$f/s$（其中 $f \in \mathcal{O}(U)$、$s \in \mathcal{S}(U)$）
映到 $(f|_V)/(s|_V)$。

\begin{lemma}
\label{sites-modules-lemma-simple-invert}
在上述情形下，到层化的映射
$$
\mathcal{O} \longrightarrow (\mathcal{S}^{-1}\mathcal{O})^\#
$$
是环层同态，并具有如下泛性质：对任意环层同态
$\mathcal{O} \to \mathcal{A}$，若 $\mathcal{S}$
的每个局部截面都映到 $\mathcal{A}$ 的可逆截面，则存在唯一分解
$(\mathcal{S}^{-1}\mathcal{O})^\# \to \mathcal{A}$。
\end{lemma}

\begin{proof}
从略。
\end{proof}

\noindent
设 $(\mathcal{C}, \mathcal{O})$ 为环站点，
$\mathcal{S} \subset \mathcal{O}$ 为集合子预层，满足：
对所有 $U \in \mathcal{C}$，集合
$\mathcal{S}(U) \subset \mathcal{O}(U)$ 是乘法子集。
设 $\mathcal{F}$ 为 $\mathcal{O}$-模层。此时可以考虑
$\mathcal{S}^{-1}\mathcal{O}$-模预层
$$
\mathcal{S}^{-1}\mathcal{F} :
U \longmapsto \mathcal{S}(U)^{-1}\mathcal{F}(U).
$$
若 $V \to U$ 是 $\mathcal{C}$ 中的态射，则限制映射把截面
$t/s$（其中 $t \in \mathcal{F}(U)$、$s \in \mathcal{S}(U)$）
映到 $(t|_V)/(s|_V)$。于是 $\mathcal{S}^{-1}\mathcal{F}$
是 $\mathcal{S}^{-1}\mathcal{O}$-模预层。

\begin{lemma}
\label{sites-modules-lemma-simple-invert-module}
在上述情形下，到层化的映射
$$
\mathcal{F} \longrightarrow (\mathcal{S}^{-1}\mathcal{F})^\#
$$
具有如下泛性质：对任意 $\mathcal{O}$-模同态
$\mathcal{F} \to \mathcal{G}$，若 $\mathcal{S}$
的每个局部截面都在 $\mathcal{G}$ 上可逆地作用，则存在唯一分解
$(\mathcal{S}^{-1}\mathcal{F})^\# \to \mathcal{G}$。
此外有
$$
(\mathcal{S}^{-1}\mathcal{F})^\#
=
(\mathcal{S}^{-1}\mathcal{O})^\# \otimes_\mathcal{O} \mathcal{F}
$$
作为 $(\mathcal{S}^{-1}\mathcal{O})^\#$-模层。
\end{lemma}

\begin{proof}
从略。
\end{proof}









\section{带基点集合层}
\label{sites-modules-section-pointed}

\noindent
本节汇集一些关于带基点集合层的事实；此前我们只在阿贝尔层的情形下
提到过这些事实。

\medskip\noindent
带基点集合是一个二元组 $(S, 0)$，其中 $S$ 是集合，$0 \in S$
是 $S$ 的一个元素。带基点集合的态射
$(S, 0) \to (S', 0')$ 就是把 $0$ 映到 $0'$ 的集合映射
$S \to S'$。我们将稍微滥用记号，把“设 $S$ 为带基点集合”
理解为 $S$ 已配备一个指定元素 $0 \in S$。带基点集合层
等同于配备“基点标记” $0 : * \to \mathcal{F}$ 的集合层
$\mathcal{F}$，其中 $*$ 是终层
（《站点》，例 \ref{sites-example-singleton-sheaf}）。

\medskip\noindent
给定一个站点态射或拓扑斯态射，在带基点集合层的范畴上有直像函子
和拉回函子；见《站点》，第
\ref{sites-section-sheaves-algebraic-structures} 节。
它们由底层集合层分别取直像与拉回并赋予适当的基点标记而构造
（这里使用了终层的拉回仍是终层这一事实）。

\medskip\noindent
设 $u : \mathcal{C} \to \mathcal{D}$ 为站点之间的连续且余连续函子。
设 $g : \Sh(\mathcal{C}) \to \Sh(\mathcal{D})$ 为与 $u$ 相关联的
拓扑斯态射；见《站点》，引理
\ref{sites-lemma-cocontinuous-morphism-topoi}。则在带基点集合层上，
$g^{-1}$ 也有左伴随 $g_!$。这个函子的构造与第
\ref{sites-modules-section-exactness-lower-shriek} 节中阿贝尔层上的 $g_!$
的构造完全类似。

\medskip\noindent
类似地，若 $j : \mathcal{C}/U \to \mathcal{C}$ 如第
\ref{sites-modules-section-localize} 节所述，则带基点集合层上的函子 $j^{-1}$
有左伴随 $j_!$。

\medskip\noindent
若日后需要这些事实与构造，我们将在此精确陈述并证明相应的引理。
